% Options for packages loaded elsewhere
\PassOptionsToPackage{unicode}{hyperref}
\PassOptionsToPackage{hyphens}{url}
%
\documentclass[
  10pt,
]{report}
\usepackage{lmodern}
\usepackage{amssymb,amsmath}
\usepackage{ifxetex,ifluatex}
\ifnum 0\ifxetex 1\fi\ifluatex 1\fi=0 % if pdftex
  \usepackage[T1]{fontenc}
  \usepackage[utf8]{inputenc}
  \usepackage{textcomp} % provide euro and other symbols
\else % if luatex or xetex
  \usepackage{unicode-math}
  \defaultfontfeatures{Scale=MatchLowercase}
  \defaultfontfeatures[\rmfamily]{Ligatures=TeX,Scale=1}
  \setmainfont[]{DejaVu Serif}
  \setmonofont[]{DejaVu Sans Mono}
\fi
% Use upquote if available, for straight quotes in verbatim environments
\IfFileExists{upquote.sty}{\usepackage{upquote}}{}
\IfFileExists{microtype.sty}{% use microtype if available
  \usepackage[]{microtype}
  \UseMicrotypeSet[protrusion]{basicmath} % disable protrusion for tt fonts
}{}
\makeatletter
\@ifundefined{KOMAClassName}{% if non-KOMA class
  \IfFileExists{parskip.sty}{%
    \usepackage{parskip}
  }{% else
    \setlength{\parindent}{0pt}
    \setlength{\parskip}{6pt plus 2pt minus 1pt}}
}{% if KOMA class
  \KOMAoptions{parskip=half}}
\makeatother
\usepackage{xcolor}
\IfFileExists{xurl.sty}{\usepackage{xurl}}{} % add URL line breaks if available
\IfFileExists{bookmark.sty}{\usepackage{bookmark}}{\usepackage{hyperref}}
\hypersetup{
  pdftitle={Einstein Gaps: Overnight Mathematical Source Companion},
  hidelinks,
  pdfcreator={LaTeX via pandoc}}
\urlstyle{same} % disable monospaced font for URLs
\usepackage[margin=24mm]{geometry}
\usepackage{longtable,booktabs}
% Correct order of tables after \paragraph or \subparagraph
\usepackage{etoolbox}
\makeatletter
\patchcmd\longtable{\par}{\if@noskipsec\mbox{}\fi\par}{}{}
\makeatother
% Allow footnotes in longtable head/foot
\IfFileExists{footnotehyper.sty}{\usepackage{footnotehyper}}{\usepackage{footnote}}
\makesavenoteenv{longtable}
\setlength{\emergencystretch}{3em} % prevent overfull lines
\providecommand{\tightlist}{%
  \setlength{\itemsep}{0pt}\setlength{\parskip}{0pt}}
\setcounter{secnumdepth}{5}
\usepackage{newunicodechar}
\setmathfont{texgyredejavu-math.otf}
\AtBeginDocument{%
  \renewcommand{\backslash}{\mathord{\text{\textbackslash}}}%
  \renewcommand{\setminus}{\mathbin{\backslash}}}
\newfontfamily\sourceSymbolFont[ItalicFont=STIXGeneral,BoldFont=STIXGeneral,BoldItalicFont=STIXGeneral]{STIXGeneral}
\newcommand{\sourceUnicode}[3]{\newunicodechar{#1}{\ifmmode#3\else{\sourceSymbolFont\char"#2}\fi}}
\sourceUnicode{ℒ}{2112}{\mathscr{L}}
\sourceUnicode{ℓ}{2113}{\ell}
\sourceUnicode{ℛ}{211B}{\mathscr{R}}
\sourceUnicode{ℝ}{211D}{\mathbb{R}}
\sourceUnicode{∅}{2205}{\varnothing}
\sourceUnicode{∎}{220E}{\blacksquare}
\sourceUnicode{∖}{2216}{\setminus}
\sourceUnicode{∮}{222E}{\oint}
\sourceUnicode{≅}{2245}{\cong}
\sourceUnicode{≪}{226A}{\ll}
\sourceUnicode{≫}{226B}{\gg}
\sourceUnicode{≲}{2272}{\lesssim}
\sourceUnicode{≳}{2273}{\gtrsim}
\sourceUnicode{⋯}{22EF}{\cdots}
\sourceUnicode{✓}{2713}{\checkmark}
\sourceUnicode{𝒜}{1D49C}{\mathscr{A}}
\sourceUnicode{𝒞}{1D49E}{\mathscr{C}}
\sourceUnicode{𝒦}{1D4A6}{\mathscr{K}}
\sourceUnicode{𝒩}{1D4A9}{\mathscr{N}}
\sourceUnicode{𝒬}{1D4AC}{\mathscr{Q}}
\sourceUnicode{𝒯}{1D4AF}{\mathscr{T}}
\sourceUnicode{𝓑}{1D4D1}{\symbfscr{B}}
\sourceUnicode{𝓡}{1D4E1}{\symbfscr{R}}
\sourceUnicode{𝔟}{1D51F}{\mathfrak{b}}

\title{Einstein Gaps: Overnight Mathematical Source Companion}
\author{}
\date{08 October 2026, 09:58:28 CST}

\begin{document}
\maketitle

{
\setcounter{tocdepth}{2}
\tableofcontents
}
\hypertarget{reading-this-source-companion}{%
\chapter{Reading this source
companion}\label{reading-this-source-companion}}

This file accompanies \textbf{Einstein Gaps: Four Public Problems,
Overnight Research Synthesis, 7--8 October 2026}. Read the synthesis
first. The present volume retains the eighteen mathematical working
notes, then the Step chronology and actual collaboration messages from
the overnight interval. It is self-contained: no server-local input
files are needed to read or compile it.

These are historical working notes, not eighteen independently
established theorems. Later corrections control earlier statements. In
particular:

\begin{itemize}
\tightlist
\item
  A01 Step 55 replaces the negative formal-dimension interpretation in
  Steps 53--54; finite jets are not local solutions or closed
  counterexamples.
\item
  A02 Step 58 consolidates the conditional assumptions of Steps 53--57.
\item
  A02 Step 59 \textbf{Part D} overrides Parts A--C. Its general-KE upper
  construction survives; its zero-gap decomposition is for selected
  sequences under the cited analytic hypotheses, not all sequences.
\item
  A02 Step 61 \textbf{Section 5} overrides the earlier neighbourhood and
  remainder statements in that note.
\item
  A02 Step 63's correction section overrides Step 62's Corollary 62.D,
  its monotonicity range, and the earlier collar-twist assertion.
\item
  A02 Step 64 explains why arbitrary almost-complex structures, or
  unweighted normalization of a spinor with zeros, do not close the
  desired lower bound.
\item
  A01's Step 53 Ricci corrector is an algebraic advance; extending a
  full insertion theorem to a general critical bulk still requires its
  remainder argument. A01 communicated this distinction in message 834.
\end{itemize}

Source-local notation is retained. In several A02 notes \texttt{alpha}
denotes \(\inf P\), whereas earlier project papers use it for a
normalized defect. The synthesis uses \(p_* = \inf P\) throughout to
remove this ambiguity.

The equations and mathematical prose below are preserved through
Markdown-to- LaTeX format conversion. Code snippets, numerical
observations, hypotheses, and withdrawn claims keep their original
status. Runtime databases, full model prompts, credentials and unrelated
Session material are not exported.

Snapshot: 2026-10-08 09:58:28 CST (UTC+8).

\hypertarget{s1003a01-step-48-problems-1-and-4-supply-analysis-the-ak-weylbochner-identity-and-the-near-ke-reduction}{%
\chapter{S1003A01: Step 48 --- Problems 1 and 4: supply analysis, the AK
Weyl--Bochner identity, and the near-KE
reduction}\label{s1003a01-step-48-problems-1-and-4-supply-analysis-the-ak-weylbochner-identity-and-the-near-ke-reduction}}

Source:
\texttt{step48\_problems1\_4\_goldberg\_bochner\_and\_supply.md}.

Session S1003A01, Agent2. Status labels: {[}P{]} proved here (hand
derivation, not yet independently verified or symbolically checked),
{[}K{]} known/cited, {[}O{]} open. Sources read:
\texttt{Skills/geometric-analysis/references/einstein-problems-20261007.md};
\texttt{five\_research\_problems.tex} §§1--4 (lines 120--700);
\texttt{web\_research\_collected.tex} Ch.1 §N1 (lines 130--330) and Ch.6
(2296--2893); Knowledge: Kim arXiv:2311.16729 (full), AAD
arXiv:math/0109098 §2.3--2.5, Kirchberg math/0605611 abstract, AAD
math/9911197 abstract; A02 notes step42, step44. Literature request
s48\_t11\_tool1 (Armstrong ansatz, local λ\textless0 strictly AK
Einstein existence, 2015--2026 progress, LeBrun 2005) was still running
at step close --- read its result before relying on any novelty
statement below.

\hypertarget{known-integrability-results-that-fix-the-open-configuration-k}{%
\subsection{1. Known integrability results that fix the open
configuration
{[}K{]}}\label{known-integrability-results-that-fix-the-open-configuration-k}}

For closed AK Einstein (M,g,J,ω), \textbar ω\textbar²=2, s\textless0
const, q=\textbar∇ω\textbar², s\emph{=s+q: - Sekigawa: s≥0 ⇒ Kähler. AAD
(Math. Ann. 2002): Ric J-invariant + ω eigenform of W⁺ (η:=W⁺(ω)\^{}⊥=0)
⇒ Kähler (compact). Kirchberg: Einstein + \textbar W⁺\textbar{} const ⇒
Kähler. Drăghici/Kim: c₁·{[}ω{]}≥0 ⇒ Kähler. Kim Cor.7: π₁=1 +
degenerate W⁺ spectrum ⇒ KE. - LeBrun's δW⁺=0 identity ∫
sW⁺(ω,ω)=8∫(\textbar W⁺\textbar²−½\textbar η\textbar²) is, on Einstein
metrics, IDENTICAL to the Chern--Weil c₁² identity (Kim Cor.3):
substituting W⁺(ω,ω)=s}/2−s/6 gives
∫(ss*/8−s²/24)=∫(2\textbar W⁺\textbar²−\textbar η\textbar²) both ways.
So it carries no information beyond the five-problems AK balance (eq.
AKbalance). - In the negative case c₁·{[}ω{]}\textless0 always:
4πc₁·ω=V(s+q̄/2) and AK balance + Cauchy--Schwarz give
q̄≤(2/3)\textbar s\textbar. Hence the Drăghici/Kim criterion is vacuous
for s\textless0. - N nowhere zero ⇒ 5χ+6τ=0 (Armstrong/LeBrun, AAD Lemma
2); with Hitchin--Thorpe this forces flatness, so for s\textless0 the
Kähler-nullity set \{∇ω=0\} is nonempty (weighted count
5χ+6τ≥χ\textgreater0).

\hypertarget{ak-weylbochner-identity-p}{%
\subsection{2. AK Weyl--Bochner identity
{[}P{]}}\label{ak-weylbochner-identity-p}}

Frame e₀=ω/√2, e₁=φ/√2, e₂=Jφ/√2 (AAD gauge: ∇Ω=a⊗φ−Ja⊗Jφ,
∇φ=−a⊗Ω+b⊗Jφ). Then ∇e₀=A₁e₁+A₂e₂, A₁=a, A₂=−Ja,
\textbar A\_i\textbar²=\textbar a\textbar²=q/4, A₁⊥A₂. Write W⁺ entries
W₀₀=ν=s/6+q/4, W₀ᵢ=ηᵢ, block −ν/2·I+V,
V={[}{[}v₁₁,v₁₂{]},{[}v₁₂,−v₁₁{]}{]}, \textbar η\textbar²=η₁²+η₂²,
\textbar v\textbar²=v₁₁²+v₁₂² (entry norms; five-problems norms are
2\textbar η\textbar², 2\textbar v\textbar²). (i) Divergence identity:
from ∇\emph{∇Ω=2ρ}'\,'+(q/2)Ω, η=−½ div¹(A) (canonical-connection
divergence of the torsion), i.e.~−div a−⟨Ja,b⟩=2η₁, div(Ja)−⟨a,b⟩=2η₂.
(ii) Derdziński Weitzenböck (δW⁺=0):
∫(2\textbar∇W⁺\textbar²+s\textbar W⁺\textbar²−36detW⁺)=0, with
36detW⁺−s\textbar W⁺\textbar²=(9/4)qν²+(8\textbar s\textbar−9q)\textbar v\textbar²+(9q/2−\textbar s\textbar)\textbar η\textbar²+36C₃,
C₃=v₁₁(η₁²−η₂²)+2v₁₂η₁η₂. (iii) Splitting ∇W⁺=∇¹W⁺+T(A)·W⁺ and
integrating cross terms by parts with (i): ∫{[}
(3/2)\textbar dν\textbar²+2\textbar∇¹η\textbar²+2\textbar∇¹v\textbar²+(9/8)qν²+(11/2)q\textbar v\textbar²+((13/4)q−(3/2)\textbar s\textbar)\textbar η\textbar²−4\textbar s\textbar\textbar v\textbar²−26C₃−12⟨dν,η·A⟩+8Y
{]}=0, Y=⟨∇¹v₁₁,η₁A₁−η₂A₂⟩+⟨∇¹v₁₂,η₂A₁+η₁A₂⟩. (iv) Eliminating
4\textbar s\textbar∫\textbar v\textbar² with AK balance
∫(2\textbar η\textbar²+4\textbar v\textbar²+3q²/16)=(\textbar s\textbar/8)∫q
leaves ∫{[}
(3/2)\textbar dν\textbar²+2\textbar∇¹η\textbar²+2\textbar∇¹v\textbar²−(3/32)s²q+(3/32)\textbar s\textbar q²+(9/128)q³+(11/2)q\textbar v\textbar²+((13/4)q+½\textbar s\textbar)\textbar η\textbar²−26C₃−12⟨dν,η·A⟩+8Y
{]}=0. Reading: in the case η=V=0 the identity immediately forces q=0
(short Bochner proof of the Einstein eigenform case). In general the
zeroth-order term −(3/32)s²q is the obstruction: integral formulas alone
cannot beat it (no Poincaré inequality in the right direction). The only
plausible source of a compensating term is the AAD Bianchi component
(38), which fixes the canonical divergence of V as ≈−(κ/8)∇\_\{J·\}Ω
plus η, dq terms, so
\textbar∇¹v\textbar²≥\textbar δ¹V\textbar²\textasciitilde(s²/64)q·const;
whether this, combined with the line-bundle Weitzenböck on K⁻² (∂ vs ∂̄
parts of ∇¹V) and the η-equation (37), beats 3s²q/32 is the decisive
computation {[}O{]}. Coefficients above must be symbolically verified
before use.

\hypertarget{near-ke-and-rigidity-reduction-p-modulo-cited-inputs}{%
\subsection{3. Near-KE and rigidity reduction {[}P, modulo cited
inputs{]}}\label{near-ke-and-rigidity-reduction-p-modulo-cited-inputs}}

Lemma KE. On a compact KE surface with s\textless0, every self-dual
(complex orientation) harmonic 2-form of constant length √2 is ±ω₀.
Proof: ω=cω₀+Reθ, θ∈H⁰(K); constant length forces \textbar θ\textbar{}
const; i∂∂̄log\textbar θ\textbar²=ρ=λω₀≠0 where θ≠0, so θ=0.
Consequences: (a) if every Einstein metric on M is KE for some complex
structure (with the AK orientation equal to the complex one), Goldberg
holds on M. (b) On manifolds underlying compact ball quotients: LeBrun's
Einstein rigidity (every Einstein metric is CH up to diffeo/scale) +
Lemma KE handles the complex orientation; for the reverse orientation
W⁺(M̄)=W⁻(CH)=0 gives q=−2s/3 const\textgreater0, contradicting nonempty
Nijenhuis zeros since 5χ+6τ(M̄)=3χ≠0. So Goldberg holds on every
ball-quotient manifold {[}likely folklore --- check survey{]}. (c) Near
a negative KE metric, AK Einstein pairs converge smoothly to (g₀,±ω₀);
Goldberg near KE reduces to Koiso's statement that Einstein metrics near
negative KE are KE (needs the obstructed-Kuranishi case checked). The
linearization is degenerate along ω₀+tReθ (length violated only at order
t²), so any near-KE quantitative Goldberg estimate is
Łojasiewicz/square-root type, not linear.

\hypertarget{problem-1-actual-field-supply-p}{%
\subsection{4. Problem 1: actual-field supply
{[}P{]}}\label{problem-1-actual-field-supply-p}}

\begin{itemize}
\tightlist
\item
  Reverse orientation: M̄=X̄\#kCP². If b⁻(X)≥1 both summands have b⁺≥1, so
  ordinary SW of M̄ vanish; Bauer--Furuta of kCP² is trivial. Fake
  projective planes (b⁻=0) give only a b⁺=1 chamber input at k=1,2.
  Hence the Chapter-6 two-orientation cross-pairing has no supply on
  blowups of general type.
\item
  Same orientation: the 2\^{}k classes a\_ε=K+Σε\_iE\_i give only
  max/average information (LeBrun; N1); near k/A=1/3 the average forces
  ‖E\_i⁺‖→0 and ‖K⁻‖→0. FD with each class (budget A−k, bounded density
  B≲6+36/A+k/A since c₂≤5c₁²+36 for general type) reproduces exactly the
  N1 window: on an Einstein metric κ\_β=D\_β/V, and D\_β≥0 ∀β is
  equivalent to LeBrun's k≤A/3 (extremal β≈8, x=V/(2π²A)≈1.78 at
  k/A=1/3); FD's c(B)(β−2)²/β² gives the ε-extension and nothing far
  below 1/3.
\item
  The exact-Einstein entrance: LeBrun's loss is the Cauchy--Schwarz
  ratio t=√6\textbar W⁺\textbar/\textbar s\textbar{} (extremal t=1/8;
  Kähler t=1/2). If an Einstein metric with the blowup SW classes had
  t≥1/2 in the weighted sense,
  ∫(s²/24+2\textbar W⁺\textbar²)≥4π²(c₁⁺)²≥4π²A would exclude every k≥1.
  δW⁺=0 alone does not force large t (constant-\textbar W⁺\textbar{}
  Einstein at t=1/8 only needs
  \textbar∇W⁺\textbar²≥\textbar s\textbar³/1024). Missing information:
  an Einstein/SW mechanism controlling the
  \textbar W⁺\textbar:\textbar s\textbar{} ratio {[}O{]}.
\end{itemize}

\hypertarget{collaboration}{%
\subsection{5. Collaboration}\label{collaboration}}

Reply \#825 to A02 \#824: the S⁴/Z₂ absorption cell is the ι-quotient of
a two-ended connected-sum neck with σ=diag(1,−1,−1,−1); shrinking cost
16π² = 8π² per absorbed class (standard), so savings can come only from
the CH germ; ι maps left- to right-invariant Berger sections, so the
cell lies outside the SU(2)-diagonal class (T² reduction needed). A02
\#826 (Step 53, unanswered at close): Anderson case-(I) energy equality,
bubble-movement theorem 53.B forcing R⁺≤0 on bulk components, Theorem
53.D/E, and a question whether the Chapter-5 x²log(1/x) kernel
coefficient stays negative on non-CH KE backgrounds for all W⁻.

\hypertarget{s1003a01-step-49-the-logarithmic-burns-saving-on-an-arbitrary-negative-kuxe4hlereinstein-background-the-w-cost-is-a-removable-corrector-artifact-and-the-ch-coefficient-9-is-universal}{%
\chapter{S1003A01: Step 49 --- The logarithmic Burns saving on an
arbitrary negative Kähler--Einstein background: the W⁻ cost is a
removable corrector artifact, and the CH coefficient −9 is
universal}\label{s1003a01-step-49-the-logarithmic-burns-saving-on-an-arbitrary-negative-kuxe4hlereinstein-background-the-w-cost-is-a-removable-corrector-artifact-and-the-ch-coefficient-9-is-universal}}

Source:
\texttt{step49\_kernel\_log\_coefficient\_general\_background.md}.

Session S1003A01, Agent2, Execute. Script
\texttt{workspace/scripts/step49\_kernel\_log\_coefficient.py} (usage in
docstring); data
\texttt{workspace/data/step49/summary\_\textless{}BG\textgreater{}\_NS*\_A\textless{}deg\textgreater{}\_K*.json},
\texttt{quadform\_*.npz}, logs \texttt{log\_*\_A4.txt},
\texttt{log\_H2H2\_A6.txt}. Status labels: {[}P{]} proved here; {[}N{]}
numerically exact computation (Cauchy/Laurent extraction, errors ≤
2e-11) inside a stated corrector class; {[}C{]} conditional / outline.

\hypertarget{question-a02-826-8274-spirit-58}{%
\subsection{0. Question (A02 \#826, \#827(4); Spirit
\#58)}\label{question-a02-826-8274-spirit-58}}

At a negative KE point p, R⁺\_p = diag(s/4,0,0) has kernel ω\^{}⊥, so
the first-order Burns response 8λ²⟨R⁺\_p,UΦU\^{}T⟩ (A02 Thm 53.F′)
vanishes for a perpendicular insertion (Burns Kähler form ζ ⊥ ω). On
complex-hyperbolic (CH) backgrounds Ch.5 of
\texttt{web\_research\_collected.tex} gives the next term P/8π² = 1 −
9ε² log(1/ε) + O\_R(ε²) (Ric = −6g), i.e.~1 − ¼x²log(1/x) + O(x²) with x
= Area(E)/π in Ric = −g. Does the logarithmic saving survive on a non-CH
KE background, for arbitrary W⁻\_p? A W⁻ jet is invisible to Bianchi-IX
radial functions; the computation must use the full curvature jet and
non-invariant angular modes.

\hypertarget{exact-reduction-to-a-corner-taylor-coefficient-p}{%
\subsection{1. Exact reduction to a corner Taylor coefficient
{[}P{]}}\label{exact-reduction-to-a-corner-taylor-coefficient-p}}

Cartesian normal coordinates x of g\_bg; neck metric g\_ε = g\_bg + εk,
k = k\_B + k\_0 + k\_2 + \ldots, with k\_B = σ₁²+σ₂² (σ\_m =
(q̄dq)\emph{m/\textbar q\textbar², left-invariant). Exact fact: Ch.5's
Burns core a = (√(t²+ε), √(t²+ε), t) is \textbf{exactly} δ + εk\_B
(a\_m² − t² = ε). Correctors k\_0 (components homogeneous of degree 0)
and k\_2 (degree 2) are free. For Einstein g\_bg, P(g\_ε) =
ε²∫\textbar Ric°′}\{g\_bg\}(k)\textbar²dV + O(ε³) in √ε ≪ r ≪ 1. With z
= ε/t², y = t², the density is (ε²/t)F(z,y) and the log coefficient is
the z²y² Taylor coefficient F(0,0): (a) total t-derivatives cannot alter
it: t d/dt(z\^{}a y\^{}b) = (2b−2a)z\^{}a y\^{}b vanishes at a = b = 2;
(b) F(0,0) = C/(8π²), C = ∫\emph{\{S³\} f}\{−4\}dΩ, f\_\{−4\} the degree
−4 homogeneous part of \textbar Ric°′(k)\textbar²\_\{g\_bg\}√det g\_bg;
then P/8π² ⊃ (F/2)ε²log(1/ε). Only the background jet to degree 4 enters
(Rm\_p quadratically via \textbar x₁\textbar², x₁ = L\_flat k\_0 +
L\_\{Rm\}k\_B; Rm⊗Rm and ∇²Rm linearly via 2⟨b₁,x₂⟩, b₁ = L\_flat k\_B;
norm/volume corrections against \textbar b₁\textbar²). k\_2 enters only
through ∫⟨b₁,L\_flat k\_2⟩; Burns is Bach- and scalar-flat, so L*\_flat
b₁ = 0 and the term is a radius-independent flux difference, i.e.~zero
(numerically ≤ 1e-11).

\hypertarget{representation-theoretic-structure-p}{%
\subsection{2. Representation-theoretic structure
{[}P{]}}\label{representation-theoretic-structure-p}}

Burns is U(2)\emph{ζ-invariant, U(2)\emph{ζ = (U(1)}+×SU(2)}−)/Z₂
(SU(2)\_− trivial on Λ⁺, SO(3) on Λ⁻). The optimal corner coefficient is
a U(2)\_ζ-invariant quadratic polynomial in Rm\_p plus a linear
functional of ∇²Rm\_p. - W⁻\emph{p is the spin-2 representation of
SU(2)}−: only quadratic invariant \textbar W⁻\emph{p\textbar²; no linear
invariant; no invariant cross terms with W⁺, s, Ric°. - ∇²W⁻ has no
linear invariant (symmetric ∇² ∈ (0,0)⊕(1,1) has SU(2)}− spin ≤ 1); ∇²W⁺
contributes only ΔW⁺(ζ,ζ), fixed on Einstein metrics by the Weitzenböck
formula and zero for KE-type W⁺. - Hence for every Einstein background
whose W⁺\_p is of KE type with ζ in the −s/12 eigenspace (every point of
every negative KE surface): F(p) = F\_CH + κ \textbar W⁻\_p\textbar²/s²,
with one universal constant κ (depending on the admitted corrector
class). - The Bianchi-IX/SU(2)-diagonal class cannot see W⁻:
SU(2)\_L-invariant curvature at the centre has W⁻\_p = 0 and Ric°\emph{p
= 0 (SU(2)\emph{L = SU(2)}−), and the squash modes c are
SU(2)}−-invariant, so they do not couple linearly to W⁻ (this corrects
A02 \#826's expectation). W⁻ effects can only be corrected by
non-invariant angular modes.

\hypertarget{computation-n}{%
\subsection{3. Computation {[}N{]}}\label{computation-n}}

Exact backgrounds: CH, a = (sinh t, sinh t cosh t, sinh t);
H²(−6)×H²(−6) in product normal coordinates; both Ric = −6g (residual ≤
5e-15), complex structure J₂ (CH fibre σ₂), Burns ζ = J₃ (perpendicular,
same Λ⁺). H²×H² has \textbar W⁻\textbar² = \textbar W⁺\textbar² = s²/24.
Laurent coefficients by Cauchy averages over λ = 0.3e\^{}\{iθ\}, x = λω
(holomorphic JAX derivatives), Gauss--Legendre×trapezoid on S³, exact
quadratic form over the tangent tensors listed.

\begin{longtable}[]{@{}llllll@{}}
\toprule
\begin{minipage}[b]{0.14\columnwidth}\raggedright
background\strut
\end{minipage} & \begin{minipage}[b]{0.14\columnwidth}\raggedright
F at c = 0\strut
\end{minipage} & \begin{minipage}[b]{0.14\columnwidth}\raggedright
F at c*\strut
\end{minipage} & \begin{minipage}[b]{0.14\columnwidth}\raggedright
opt. SU(2)-diag. c\strut
\end{minipage} & \begin{minipage}[b]{0.14\columnwidth}\raggedright
opt. all E\_ab·(angular monomials ≤ 2) (215 dims)\strut
\end{minipage} & \begin{minipage}[b]{0.14\columnwidth}\raggedright
opt. all E\_ab·(angular monomials ≤ 4) (465 dims)\strut
\end{minipage}\tabularnewline
\midrule
\endhead
\begin{minipage}[t]{0.14\columnwidth}\raggedright
CH\strut
\end{minipage} & \begin{minipage}[t]{0.14\columnwidth}\raggedright
14.75 (=59/4 ✓)\strut
\end{minipage} & \begin{minipage}[t]{0.14\columnwidth}\raggedright
−18.000 ✓\strut
\end{minipage} & \begin{minipage}[t]{0.14\columnwidth}\raggedright
−18.000\strut
\end{minipage} & \begin{minipage}[t]{0.14\columnwidth}\raggedright
−18.000\strut
\end{minipage} & \begin{minipage}[t]{0.14\columnwidth}\raggedright
−18.000\strut
\end{minipage}\tabularnewline
\begin{minipage}[t]{0.14\columnwidth}\raggedright
H²×H²\strut
\end{minipage} & \begin{minipage}[t]{0.14\columnwidth}\raggedright
20.35\strut
\end{minipage} & \begin{minipage}[t]{0.14\columnwidth}\raggedright
−12.40\strut
\end{minipage} & \begin{minipage}[t]{0.14\columnwidth}\raggedright
−12.40\strut
\end{minipage} & \begin{minipage}[t]{0.14\columnwidth}\raggedright
−15.900\strut
\end{minipage} & \begin{minipage}[t]{0.14\columnwidth}\raggedright
\textbf{−18.000} (\strut
\end{minipage}\tabularnewline
\bottomrule
\end{longtable}

Checks: CH reproduces Ch.5's F(0,0;c) at c = 0 and c* to 1e-12;
enlarging the corrector class never changes the CH value; k\_2 couplings
vanish; imaginary parts ≤ 4e-13. Angular degree 6 (1305 dims) on H²×H²
is running (task s49\_t16\_tool1); a larger class can only lower F.

Reading. In the Ch.5 diagonal ansatz the W⁻ background \emph{costs} 5.6
units of F; angular-degree-2 non-invariant correctors recover 3.5 of
them and angular-degree-4 correctors recover all of them exactly. Thus
within the degree-4 class κ = 0 to 2e-11: \textbf{on every negative KE
background, at every point, the optimal perpendicular Burns insertion
has the CH logarithmic coefficient} P/8π² = 1 − ¼ x² log(1/x) + O(x²), x
= Area(E)/π (Ric = −g), independently of W⁻\_p.~The exact cancellation
strongly suggests a structural identity (the W⁻-part of the
Burns--background interaction is entirely in the range of the flat
linearized trace-free Ricci operator on degree-0 tensors, with matching
quadratic jet terms), not yet proved. If degree-6 correctors lowered F
below −18, W⁻ would \emph{help}; either way the saving persists.

\hypertarget{consequences-c-modulo-the-rigorous-general-background-version-of-ch.5s-remainder-estimate-5}{%
\subsection{4. Consequences {[}C: modulo the rigorous general-background
version of Ch.5's remainder estimate,
§5{]}}\label{consequences-c-modulo-the-rigorous-general-background-version-of-ch.5s-remainder-estimate-5}}

\begin{enumerate}
\def\labelenumi{\arabic{enumi}.}
\tightlist
\item
  p\_*(X\#CP̄²) \textless{} 8π² --- indeed P/8π² ≤ 1 − ¼x²log(1/x) + C\_X
  x² --- for \textbf{every} compact negative KE surface X (every smooth
  canonical model of general type with K ample), at an arbitrary
  insertion point, not only for ball quotients or ramified CH-containing
  towers. The only background input is KE-type W⁺\_p and ζ ∈ ker R⁺\_p.
\item
  A02's attainment chain (Thm 53.B/D, 54.A): a KE bulk component cannot
  carry a Burns-type bubble anywhere (moving the bubble to any point
  gives strict decrease, contradicting energy equality). This extends
  ``no exact CH region'' to ``no KE bulk component at all'', i.e.~the
  Einstein-bulk alternative (c) of Thm 53.E is excluded when the bulk is
  KE for the orientation selected by the bubble. The non-KE Einstein
  bulk and non-Einstein bulk cases remain (they bring W⁺ shapes other
  than KE type and Ric°\_p invariants \textbar Ric°\_p\textbar²,
  \textbar Ric°\_p(ζ)\textbar²).
\item
  Replication: since the coefficient is the same at every point, fixed-ε
  cells can be placed on arbitrary unramified covers of any negative KE
  surface; the strict sub-four linear coefficient of Ch.5 (5.6.1)
  extends from compact CH surfaces to every fixed negative KE base (with
  base-dependent δ\_X, by the same order of limits).
\item
  The CH ball is not special for the logarithmic mechanism: the saving
  is a property of the R⁺-kernel orientation alone (to this order).
\end{enumerate}

\hypertarget{what-remains}{%
\subsection{5. What remains}\label{what-remains}}

\begin{itemize}
\tightlist
\item
  Degree-6 convergence (running) and a proof of κ = 0 (identify the
  explicit corrector; representation theory restricts it to the SU(2)\_−
  spin-2 isotypic component).
\item
  Rigorous general-background construction: Ch.5 analytic factorization
  with g\_ε = cutoff interpolation of g\_bg + εk (exact Burns core = δ +
  εk\_B makes the inner region Burns + O(r²)); prove the analytic
  (z,y)-representation of the angular-integrated density and the
  O\_R(ε²) remainder; order-ε terms: only the invariant R⁺(ζ,ζ) can
  appear (zero here), k\_0-linear ones vanish by Burns criticality.
\item
  Non-KE W⁺ shapes with R⁺ ≤ 0, ζ ∈ ker R⁺ (W⁺ = (\textbar s\textbar/12,
  w, −\textbar s\textbar/12−w)): needs an Einstein 4-jet framework; ∇²Rm
  enters only via ΔW⁺(ζ,ζ).
\item
  Non-Einstein bulk minimizers (Problem 3 setting): add Ric°\_p
  invariants and the ε²⟨Ric°\_bg, Ric°″⟩ term.
\end{itemize}

\hypertarget{s1003a01-step-50-the-w-sector-of-the-burnsbackground-interaction-is-an-exact-completed-square-and-the-logarithmic-sub-8ux3c0uxb2-one-blow-up-construction-holds-on-every-negative-kuxe4hlereinstein-surface}{%
\chapter{S1003A01: Step 50: the W⁻ sector of the Burns--background
interaction is an exact completed square, and the logarithmic sub-8π²
one-blow-up construction holds on every negative Kähler--Einstein
surface}\label{s1003a01-step-50-the-w-sector-of-the-burnsbackground-interaction-is-an-exact-completed-square-and-the-logarithmic-sub-8ux3c0uxb2-one-blow-up-construction-holds-on-every-negative-kuxe4hlereinstein-surface}}

Source:
\texttt{step50\_wminus\_completed\_square\_and\_general\_KE\_construction.md}.

Session S1003A01, Agent2, Execute. Supersedes the interpretation in
§§3--5 of
\texttt{step49\_kernel\_log\_coefficient\_general\_background.md}. The
Step-49 numbers there remain valid.

Status labels: - {[}E{]} exact symbolic (rational-function) identity; -
{[}P{]} proof given here; - {[}N{]} floating-point computation, recorded
as evidence only; - {[}C{]} proof outline whose remaining analytic
details are routine but not yet written out line by line.

Scripts: - \texttt{workspace/scripts/step50\_wminus\_compatibility.py}:
CGB-form W⁻ residual over corrector classes; -
\texttt{step50\_natural\_corrector.py}: sparse search for the corrector;
- \texttt{step50\_check\_corrector\_ric0.py}: check in the independent
Step-49 Ric°/Laurent quadratic form; -
\texttt{step50\_exact\_identity.py}: exact sympy identity.

Data are in \texttt{workspace/data/step50/}:
\texttt{wminus\_compat\_A\{0,2,4,6\}.json},
\texttt{natural\_corrector.json},
\texttt{corrector\_check\_ric0form.json}, \texttt{exact\_identity.json}.

\hypertarget{setting-and-notation}{%
\subsection{0. Setting and notation}\label{setting-and-notation}}

\begin{itemize}
\tightlist
\item
  (X, g\_bg) is a compact negative Kähler--Einstein (KE) surface,
  normalized Ric = −6g as in Ch.5. Its Kähler form is ω and p ∈ X is
  arbitrary.
\item
  x are g\_bg-normal coordinates at p, with r = \textbar x\textbar.
\item
  ζ ∈ Λ⁺\_p is a unit 2-form with ζ ⊥ ω, the perpendicular kernel
  orientation, ker R⁺\_p = ω\^{}⊥.
\item
  Burns tail: k\_B = σ₁²+σ₂² = (v₀v₀ᵀ+v₁v₁ᵀ)/\textbar x\textbar⁴. Flat
  Burns is \textbf{exactly} δ+εk\_B, and it is scalar-flat and ASD for
  the orientation in which ζ is self-dual (W⁺ = 0, s = 0).
\item
  Normal-coordinate jet: g\_bg = δ + h₂ + h₃ + h₄ + \ldots, with h₂ =
  −(1/3)R\_\{ikjl\}x\textsuperscript{kx}l. Split h₂ = h\_KE + h\_W,
  where h\_W is the W⁻\_p part.
\item
  Neck metric: g\_bg + ε(k\_B + k₀), where k₀ is a corrector whose
  components are homogeneous of degree 0.
\item
  The log coefficient F is the z²y² corner coefficient (z = ε/r², y =
  r²) of the S³-averaged density per d log r, normalized as in Ch.5. The
  relevant expansion is P/8π² ⊃ (F/2)ε²log(1/ε), and Ch.5 gives F = −18
  on CH.
\end{itemize}

\hypertarget{cherngaussbonnet-form-of-the-corner-coefficient-p}{%
\subsection{1. Chern--Gauss--Bonnet form of the corner coefficient
{[}P{]}}\label{cherngaussbonnet-form-of-the-corner-coefficient-p}}

In dimension 4, pointwise: \textbar Ric°\textbar² dV =
(4\textbar W⁺\textbar² + s²/12) dV − 16π²e − 24π²l. Here
\textbar W⁺\textbar² is the Frobenius norm of W⁺ as an operator on Λ²,
so that ∫\_\{CP²\}\textbar W⁺\textbar² = 12π². The forms e and l are the
Euler and \textbf{signature} densities, normalized by ∫e = χ and ∫l = τ
(in terms of the first Pontryagin form, 24π²l = 8π²p₁). This correction
is from the TTCS verifier.

\textbf{Lemma 1.1.} - (a) Compute both densities of g\_bg + ε(k\_B+k₀),
relative to g\_bg, on the annulus. After S³-averaging they are analytic
functions of (y,z), the S³ average killing odd powers of r. They differ
by r∂\_r of an analytic function G(y,z): the averaged Chern--Weil
transgression on the sphere of radius r. - (b) At fixed ε, r∂\_r =
2y∂\_y − 2z∂\_z multiplies y\^{}b z\^{}a by (2b−2a). It therefore kills
every diagonal coefficient, so F computed from \textbar Ric°\textbar²
equals F computed from f := 4\textbar W⁺\textbar² + s²/12. - (c) In the
f-form only Rm\_p enters F. Burns is exactly ASD and scalar-flat, so
every flat-Burns term in W⁺ and s vanishes identically. Degree counting
then leaves exactly three contributions to the degree −4 density at
order ε²: - (i) \textbar W⁺₁\textbar² and s₁², where W⁺₁ and s₁ are the
order-ε parts of degree −2; they are linear in h₂ ⊗ k\_B and in k₀; -
(ii) ⟨W⁺\emph{bg(p), W⁺₂⟩ and s\_bg s₂, where W⁺₂ and s₂ are the
order-ε² parts of degree −4; they are linear in h₂ or in k₀; - (iii)
volume and norm factors multiplying background constants; - (iv)
first-order volume and norm corrections multiplied by ⟨𝓡₀(p),
𝓡}\{1,−2\}⟩. This term is linear in (h₂, k₀); the verifier noted it was
missing from the first version, and it does not affect the W⁻ square.
The jets h₃, h₄, \ldots{} contribute only to degrees ≥ −2 at order ε².
(In the Ric°-form the 4-jet does appear, through ⟨b₁,x₂⟩. That apparent
dependence is exact-form bookkeeping.)

\textbf{Corollary 1.2 (W⁻ enters only as a nonnegative square).} Assume
Rm\_p is of KE type with ζ ⊥ ω, so that Ric°\emph{p = 0 and W⁺\emph{p =
(s/12)diag(2,−1,−1) in the ω-adapted basis. Let U(2)\emph{ζ = U(1)}+ ×
SU(2)}−; here SU(2)}− fixes ζ and ω and acts on W⁻ by spin 2. Then: -
Every term linear in W⁻ has zero S³ average: its other ingredients
(k\_B, W⁺\emph{bg, s\_bg and the SU(2)}−-equivariant correctors) are
SU(2)\_−-invariant or equivariant, and spin 2 has no invariant. - The
only W⁻-quadratic terms are \textbar W⁺₁\^{}W\textbar² and (s₁\^{}W)².
In particular W⁺₂ carries no W⁻⊗W⁻ terms at degree −4: Rm²-terms land in
degree −2. - Hence for every corrector class 𝒦 that is invariant under
SO(4):

F\_X(p) = F\_CH + κ\_𝒦(W⁻\_p),

κ\_𝒦(W⁻) = min\_\{k₀∈𝒦\} (1/8π²)∫\_\{S³\} {[} 4\textbar W⁺'`(h\_W,k\_B)
+ W⁺'(k₀)\textbar² + (s'`(h\_W,k\_B) + s'(k₀))²/12 {]} dΩ ≥ 0.

Here F\_CH is the CH optimum over the same class, primes denote flat
derivatives, and the two arguments of W⁺'\,' are the two perturbation
directions.

So W⁻ can never improve the coefficient; it can only cost. The cost
vanishes exactly when the W⁻ source (S\_W, σ\_W) := (W⁺'`, s'')(h\_W,
k\_B) lies in the range of the flat linearized (W⁺, s) operator on
degree-0 tensors.

Numerical confirmation of the reduction ({[}N{]},
\texttt{wminus\_compat\_A*.json}, all values in F units): \textbar{}
corrector class \textbar{} κ from (W⁺,s)-form, Step 50 \textbar{}
F\_H²×H² − F\_CH from Ric°/Laurent form, Step 49 \textbar{}
\textbar---\textbar---\textbar---\textbar{} \textbar{} none /
SU(2)\_−-invariant modes \textbar{} 5.6 \textbar{} 20.35 − 14.75 = −12.4
+ 18 = 5.6 \textbar{} \textbar{} angular degree ≤ 2 \textbar{} 2.1
\textbar{} −15.9 + 18 = 2.1 \textbar{} \textbar{} angular degree ≤ 4
\textbar{} 3·10⁻²⁹ \textbar{} 0 (2·10⁻¹¹) \textbar{} \textbar{} angular
degree ≤ 6 \textbar{} 8·10⁻²⁹ \textbar{} 0 (task s49\_t16\_tool1:
−18.000000000017) \textbar{}

Two further checks: - The flat corrector quadratic blocks agree in the
two forms to 2·10⁻¹⁵ (relative). This is the exact-form statement (b) at
the quadratic level. - The coupling of the W⁻ source to the
SU(2)\_−-invariant modes is 2·10⁻¹³, as the spin argument requires.

The degree-6 Ric°-form run stops at −18 rather than going below it, as
Corollary 1.2 predicts (κ ≥ 0).

\hypertarget{the-explicit-corrector-and-the-exact-completed-square-ep}{%
\subsection{2. The explicit corrector and the exact completed square
{[}E{]}+{[}P{]}}\label{the-explicit-corrector-and-the-exact-completed-square-ep}}

\textbf{Theorem 2.1 (first-order W⁻ compatibility of the Burns tail).}
For every algebraic W⁻ (pure anti-self-dual Weyl jet h\_W in the
ζ-orientation), the degree-0 tensor

k₀\^{}W = (3/2)·sym(h\_W k\_B) − (1/4)·tr(h\_W k\_B)·δ, with sym(M) =
(M+Mᵀ)/2 (matrix products in normal coordinates),

solves

W⁺'\,'\_flat(h\_W, k\_B) + W⁺'\_flat(k₀\^{}W) = 0, s'\,'\_flat(h\_W,
k\_B) + s'\_flat(k₀\^{}W) = 0.

Here tr(h\_W k\_B) = −v₂ᵀh\_W v₂/\textbar x\textbar⁴, because h\_W x = 0
by the Gauss lemma and tr h\_W = 0.

Proof. - For W⁻ equal to the W⁻-jet of H²(−6)×H²(−6), all entries are
rational functions. \texttt{step50\_exact\_identity.py} verifies all ten
components (five of W⁺ in the J-basis and s) to be identically 0 in
sympy rational arithmetic, using the exact bilinear curvature formula R
= ½∂²g + Σ\_e(Γ\_\{ebc\}Γ\_\{ead\} − Γ\_\{ebd\}Γ\_\{eac\}) with
first-order frame corrections. - The same run checks exactly that h\_W
is pure W⁻ (W⁺`(h\_W) = s'(h\_W) = Ric°`(h\_W) = 0), that W⁺'(k\_B) =
s'(k\_B) = 0, and that the source without corrector is nonzero. -
General W⁻: both sides of the identity are SU(2)\_−-equivariant and
linear in W⁻, because k₀\^{}W is built naturally from h\_W, k\_B and δ.
The set of W⁻ for which the identity holds is therefore an invariant
linear subspace of the irreducible spin-2 representation. It contains a
nonzero element, so it is everything. ∎

Interpretation. - The trace term only corrects s, since W⁺ is
conformally invariant. - The coefficient 3/2 is the factor by which the
``metric product'' sym(h\_W k\_B) must be amplified for the composite
metric δ + h\_W + εk\_B to stay ASD and scalar-flat to first order in
(W⁻ × ε). - Geometrically, the perpendicular Burns tail can be coupled
to an arbitrary background W⁻ field at first order with no W⁺ or scalar
response. W⁻ is the curvature the Burns bubble itself carries, so this
compatibility has a natural ASD-deformation flavour; the exact identity
is what is proved.

\textbf{Completed square in the independent Ric°/Laurent form} ({[}N{]},
\texttt{corrector\_check\_ric0form.json}, Step-49 quadratic forms,
angular degree ≤ 4): - F\_H²×H²(c* + t·k₀\^{}W) = −18 + 5.6(1−t)²:
values −12.4, −16.6, −18.000000000019, −16.6 at t = 0, ½, 1, 3/2. - The
cross term with the Ch.5 c*-ansatz is −5.6 and the quadratic term is
+5.6. - On CH the same tensor (built from the H²×H² h\_W) has zero cross
term (5·10⁻¹⁵) and costs +5.6t², as it must, since CH has no W⁻ to
cancel.

This is exactly the completed-square structure seen in the saved
matrices (Spirit \#59): identical corrector blocks, an extra 5.6 in the
constant term, a negligible pairing with the CH source, and a Schur
complement recovering 5.6.

\textbf{Corollary 2.2 {[}P, given Ch.5's symbolic F\_CH(c*) = −18{]}.}
At every point of every negative KE surface, the corrector k₀ =
k₀\^{}\{c\emph{\} + k₀\^{}W gives corner coefficient F = −18. Here
k₀\^{}\{c}\} is the degree-0 part of the Ch.5 neck: v₁v₁ᵀ/u + 2c₁v₀v₀ᵀ/u
+ 2c₂v₁v₁ᵀ/u + 2c₃v₂v₂ᵀ/u with c* = (9/4, 3/4, 1/2), written in the
ω,ζ-adapted frame at p.~No CH region, symmetry, or Bianchi-IX structure
of the background is used.

When the corrector class is a linear, SO(4)-stable space 𝒦, isotypic
decomposition gives inf\_𝒦 F\_X = inf\_𝒦 F\_CH + κ\_𝒦 with κ\_𝒦 ≥ 0.
This is \textbf{not} zero for every such class: the degree-≤2 class
gives κ = 2.1. It is zero whenever the W⁻ source lies in L(𝒦). That
happens for the degree-≤4 class, for every larger class, and for
unrestricted smooth degree-0 correctors. For those classes W⁻ is
invisible to the optimal logarithmic coefficient.

\hypertarget{the-construction-on-an-arbitrary-negative-ke-surface-pc}{%
\subsection{3. The construction on an arbitrary negative KE surface
{[}P/C{]}}\label{the-construction-on-an-arbitrary-negative-ke-surface-pc}}

\textbf{Theorem 3.1.} Let (X, g\_bg) be a compact negative KE surface
(Ric = −6g), p ∈ X, ζ ⊥ ω. Then for every sufficiently small T
\textgreater{} 0 there exist ε\_T \textgreater{} 0 and smooth metrics
g\_ε on X\#CP̄² (0 \textless{} ε \textless{} ε\_T) with the following
properties: - g\_ε = g\_bg outside the normal ball B\_\{2T\}(p); - g\_ε
is exactly a flat-scale Burns metric on r ≤ √ε; - P(g\_ε)/8π² = 1 −
9ε²log(1/ε) + O\_\{X,p,T\}(ε²).

In Ric = −g units, with x = Area(E)/π: P/8π² = 1 − ¼x²log(1/x) + O(x²).

Construction (additive form of Ch.5 (5.1.1)). Let χ\_in(r) = χ(r/√ε) and
χ\_out(r) = 1 − χ(r/T), with χ as in Ch.5. Set - r ≤ √ε: g\_ε = δ +
εk\_B (Burns, smooth across the bolt); - otherwise g\_ε =
(1−χ\_in)(δ+εk\_B) + χ\_in·{[}g\_bg + χ\_out·ε(k\_B + k₀){]}, with k₀ =
k₀\^{}\{c*\} + k₀\^{}W(W⁻\_p) from Corollary 2.2.

Proof structure. Three facts make the general case work like Ch.5,
though the background has no symmetry.

\begin{enumerate}
\def\labelenumi{(\Alph{enumi})}
\tightlist
\item
  Analytic two-scale representation {[}C{]}.
\end{enumerate}

\begin{itemize}
\tightlist
\item
  KE metrics are real-analytic in normal coordinates (DeTurck--Kazdan,
  together with analyticity of exp).
\item
  On the neck 2√ε ≤ r ≤ T, write g\_ε(x) = G(ω; r, z) = δ + Σ\_\{n≥2\}
  rⁿG\_n(ω) + zK\_B(ω) + z r²k₀(ω), which is analytic in (r, z) on
  {[}0,T{]}×{[}0,1/4{]}.
\item
  The \textbar Ric°\textbar² density per d log r dΩ, S³-averaged, is
  ψ(r, z). It is even in r, because (ω, r) ↦ (−ω, −r) fixes x and z. So
  ψ = ψ(y, z) is analytic.
\item
  Exact boundary data: ψ(y, z) = O(z²), since Ric°(g\_bg) = 0; and ψ(0,
  z) = ψ\_B(z), the flat-Burns density.
\item
  Hence ψ − ψ\_B = z²y·φ(y, z) with φ analytic. Writing φ = φ(0, z) +
  y·φ₃(y, z) gives the density z²y·φ(0, z) + ε²φ₃(y, z).
\item
  The second piece gives ε²φ₃(0,0)·log(T/2√ε) + O\_T(ε²) by the
  mean-value bound of Ch.5 (5.3.2), with φ₃(0,0) = F/2-normalized,
  i.e.~the corner coefficient.
\item
  The first piece is (ε/2)∫\_\{ε/T²\}\^{}\{1/4\} φ(0, z) dz = ε·ℓ\_N +
  O(ε²/T²). This is a pure order-ε bucket, linear in (h₂, k₀).
\end{itemize}

\begin{enumerate}
\def\labelenumi{(\Alph{enumi})}
\setcounter{enumi}{1}
\tightlist
\item
  Core and transitions.
\end{enumerate}

\begin{itemize}
\tightlist
\item
  r ≤ √ε: exact Burns.
\item
  √ε ≤ r ≤ 2√ε: in ξ = x/√ε the metric is ε·(Burns + η), where η =
  ε(h₂(ξ)+k₀) + O(ε\^{}\{3/2\}) in C\^{}k on 1 ≤ \textbar ξ\textbar{} ≤
  2.
\item
  Burns is P-critical: it is Bach-flat and scalar-flat, so ∇P = 0. The
  P-change is therefore ε·(flux bilinear in Burns and the asymptotic
  form of η₁) + O(ε²).
\item
  The ε\^{}\{3/2\} term is odd under the isometry −I ∈ U(2) of Burns, so
  its linear contribution vanishes.
\item
  Outer transition T ≤ r ≤ 2T: g\_ε − g\_bg = O\_T(ε) in C², with
  Ric°(g\_bg) = 0, so the cost is O\_T(ε²). There is no first-order term
  because the background is Einstein.
\end{itemize}

\begin{enumerate}
\def\labelenumi{(\Alph{enumi})}
\setcounter{enumi}{2}
\tightlist
\item
  Vanishing of the order-ε bucket {[}P{]}.
\end{enumerate}

\begin{itemize}
\tightlist
\item
  The total order-ε coefficient ℓ = ℓ\_core + ℓ\_N is a
  U(2)\_ζ-invariant linear functional of (Rm\_p, k₀).
\item
  It is independent of the inner blending and of the precise neck form:
  compactly supported changes of η₁ pair to zero with the P-critical
  Burns metric.

  \begin{itemize}
  \tightlist
  \item
    The same holds for decaying changes (for example additive versus
    Ch.5's multiplicative radii, which differ by
    ε·O(\textbar ξ\textbar⁻²) in scaled variables). The boundary flux
    against the Burns tail is
    O(\textbar ξ\textbar⁻⁷)·\textbar ξ\textbar³ → 0.
  \item
    So ℓ depends only on the non-decaying asymptotics: the growth ε
    h₂(ξ) and the constants ε k₀.
  \end{itemize}
\item
  Direct flux evaluation (TTCS verifier,
  \texttt{trace/ttcs/results/step50\_ttcs01\_verifier1\_integration.md}
  §4; supersedes the CH-comparison argument of the first version):

  \begin{itemize}
  \tightlist
  \item
    Burns has s = W⁺ = 0, so for variations vanishing near the bolt
    DP\_\{g\_B\}(q) = 𝓑\_R(q), the outer boundary flux of the first
    variation.
  \item
    Bounded degree-0 correctors k₀ give 𝓑\_R = O(R⁻²) → 0. So the bucket
    is independent of \textbf{every} degree-0 corrector, and of the
    blending.
  \item
    The growing jet h₂ = −(1/3)R\_\{ikjl\}x\textsuperscript{kx}l is
    evaluated against the leading Burns Ricci tensor b = (2/r⁴)(δ − 2θ⊗θ
    − 2J\_ζθ⊗J\_ζθ), with θ = x/r. The S³ average uses
    ⨍R(J\_ζθ,θ,J\_ζθ,θ) = s/12 + ½W⁺(ζ,ζ). The result is 𝓑\_∞(h₂) =
    −8π²(W⁺(ζ,ζ) + s/12) = −8π²R⁺\_p(ζ,ζ), so ℓ = −R⁺\_p(ζ,ζ) (in P/8π²
    and Burns-ε units).
  \item
    At a KE point with ζ ⊥ ω, W⁺(ζ,ζ) = −s/12, so ℓ = 0.
  \item
    This agrees in form with A02's first-order Burns term −8λ²⟨R⁺\_p,
    UΦUᵀ⟩, and with Ch.5 (5.A), which has no O(ε) term on CH.
  \item
    Status: the derivation is from the verifier; this session has not
    yet rechecked the constant independently. Only the vanishing in the
    kernel orientation is used.
  \end{itemize}
\end{itemize}

Combining (A)--(C) with F = −18 from Corollary 2.2 gives the theorem. ∎

What still has to be written in full: the analytic-representation
details of (A), the divisibility statements, and the precise flux
formula for ℓ. These follow Ch.5 §§5.3--5.5 verbatim, with ``radial
ODE'' replaced by ``S³-averaged analytic density''; no estimate beyond
Ch.5's mean-value bound is needed. Constants depend on (X, p, T). No
uniformity over different KE surfaces is claimed or needed.

\hypertarget{closed-upper-bounds-via-the-yamabe-shortcut-p}{%
\subsection{4. Closed upper bounds via the Yamabe shortcut
{[}P{]}}\label{closed-upper-bounds-via-the-yamabe-shortcut-p}}

On a closed four-manifold, conformal invariance of ∫\textbar W⁺\textbar²
gives P(ĝ) − P(g) = {[}S(ĝ) − S(g){]}/12 for conformal metrics. If Y(M)
\textless{} 0, the constant-scalar-curvature (Yamabe) representative ĝ
of {[}g{]} minimizes S in {[}g{]}. Hence P(ĝ) ≤ P(g) and L\_s(ĝ) = 0, so

E(ĝ) = 4P(ĝ)/(P(ĝ) + S(ĝ)/4) ≤ 4P(g)/(P(g) + 8π²A), A = c₁²(X),

whenever the actual Seiberg--Witten supply gives S(ĝ) ≥ 32π²A. This is
LeBrun's estimate on X\#kCP̄² with X minimal of general type (b⁺ ≥ 2, or
the stated b⁺ = 1 chamber), and on the inherited connected-sum families
under their stated supply. No local J\_6 expansion or conformal
normalization of the new cells is needed.

\textbf{Corollary 4.1.} Let X₀ be any fixed compact negative KE surface
and X\_j → X₀ unramified covers of degree d\_j → ∞, with A\_j = d\_jA₀.
Fix T and a ball B\_\{2T\}(p₀). Theorem 3.1 gives a fixed cell with P =
8π²(1 − δ\_\{X₀\}), δ\_\{X₀\} \textgreater{} 0, which is then frozen.
Replicate it in the lifted balls, together with the inherited shrinking
handles. Then for the balanced families (Q\_j = k\_j + 4ℓ\_j, h\_j → 0,
θ\_j = k\_j/Q\_j → θ):

limsup α(M\_j)/h\_j ≤ 4(1 − δ\_\{X₀\}θ).

For θ = 2/9 this is 4 − (8/9)δ\_\{X₀\} \textless{} 4. For pure blow-ups
(θ = 1) it is 4(1 − δ\_\{X₀\}).

Order of limits: first T, then a fixed ε, then j → ∞. This extends Ch.5
(5.6.1) from compact CH surfaces to covers of every negative KE base. It
is an upper construction only; it does not extend the all-metric lower
window. A single blow-up already gives α(X\#CP̄²) ≤ 4P/(P+8π²A)
\textless{} 4/(1+A) under the supply hypothesis.

\hypertarget{consequences-for-the-peer-programme-a02-thm-53.bd-54.a-p-given-theorem-3.1}{%
\subsection{5. Consequences for the peer programme (A02 Thm 53.B/D,
54.A) {[}P given Theorem
3.1{]}}\label{consequences-for-the-peer-programme-a02-thm-53.bd-54.a-p-given-theorem-3.1}}

\begin{itemize}
\tightlist
\item
  p\_\emph{(X\#CP̄²) \textless{} 8π² = p\_}(X) + 8π² for every negative
  KE surface X.
\item
  Consider an Anderson case-(I) energy-equality configuration whose bulk
  is (X, g\_KE) and which carries one Burns-type bubble.

  \begin{itemize}
  \tightlist
  \item
    For the bubble orientation at the bubble point, an aligned insertion
    pays +8λ²\textbar s\textbar/4 at first order.
  \item
    A perpendicular insertion is first-order neutral and strictly saves
    at log order.
  \item
    Bubble movement (A02 53.B) together with Theorem 3.1 then
    contradicts minimality.
  \item
    So such a configuration cannot occur.
  \end{itemize}
\item
  This excludes a KE bulk carrying a Burns bubble. It does \textbf{not}
  exclude non-KE Einstein bulk, prove attainment, or give
  p\_\emph{(M\_k) \textless{} p\_}(M\_\{k−m\}) + 8π²m at non-KE bulk
  minimizers. Those need F for non-KE R⁺ ≤ 0 shapes and for Ric°\_p ≠ 0
  (§6).
\end{itemize}

\hypertarget{what-the-cgb-form-makes-computable-next}{%
\subsection{6. What the CGB form makes computable
next}\label{what-the-cgb-form-makes-computable-next}}

The degree count of Lemma 1.1(c) holds for any background, Einstein or
not. For a non-Einstein background there is, however, an additional
first-order bulk variation ε·DP\_\{g\_bg\}{[}χ(k\_B + k₀){]}. It reduces
to the pointwise ℓ = −R⁺\emph{p(ζ,ζ) only when the bulk is P-critical,
as it is for A02's minimizers; otherwise it needs separate control.
Hence F(Rm\_p, ζ) is a U(2)\emph{ζ-invariant quadratic form in Rm\_p
alone, independent of higher jets. The invariants are: - W⁻: only
\textbar W⁻\textbar², with κ = 0 (Theorem 2.1); - Ric° ∈ Λ⁺⊗Λ⁻:
\textbar Ric°(ζ)\textbar² and \textbar Ric°}\{ζ\^{}⊥\}\textbar², each
entering as a nonnegative minimized square, with no cross terms with s
or W⁺ (SU(2)}− spin 1); - (s, W⁺) under U(1)\emph{+: s², s·W⁺}\{ζζ\},
W⁺\emph{\{ζζ\}², \textbar W⁺}\{ζ⊥\}\textbar² (weight 1) and
\textbar W⁺\_\{⊥⊥,0\}\textbar² (weight 2).

So 5 + 2 coefficients determine the log coefficient at every point of
every background. They can be computed as the a²ε² coefficient of the
f-density of δ + a h₂ + ε(k\_B + a k₀). That coefficient is homogeneous
of degree −4, so no Laurent extraction is needed, and the computation
fits directly into \texttt{step50\_wminus\_compatibility.py}. For
non-Einstein bulk there is also an order-ε flux term ℓ ∝ R⁺\_p(ζ,ζ),
which A02's R⁺ ≤ 0 selects to vanish on ker R⁺. This is the input to
A02's general-k comparison and to the Problem-3 attainment chain.

\hypertarget{s1003a01-step-52-the-kernel-direction-burns-insertion-theorem-on-an-arbitrary-einstein-background}{%
\chapter{S1003A01: Step 52 --- The kernel-direction Burns insertion
theorem on an arbitrary Einstein
background}\label{s1003a01-step-52-the-kernel-direction-burns-insertion-theorem-on-an-arbitrary-einstein-background}}

Source: \texttt{step52\_kernel\_insertion\_theorem.md}.

Session S1003A01, Agent2, Step 52. Status labels: \textbf{{[}E{]}} exact
symbolic/rational computation; \textbf{{[}P{]}} written proof;
\textbf{{[}C{]}} conditional or relies on an inherited result named
explicitly; \textbf{{[}N{]}} floating evidence only.

This note supersedes the open items of Step 51 (item (A) of its handoff)
and the conditional status of Step 50 Theorem 3.1. The inherited inputs
are: the Step-50 exact W⁻ identity
(\texttt{workspace/scripts/step50\_exact\_identity.py},
\texttt{data/step50/exact\_identity.json}), the CGB pointwise identity,
and \textbf{one} statement of Ch. 5 of
\texttt{web\_research\_collected.tex} (the CH perpendicular construction
has no order-ε term, eq. (5.A)), used only to fix the order-ε flux
constant up to a nonzero factor (§3.3). The Ch. 5 value −18 is
\textbf{no longer an input}: it is reproduced exactly below.

\hypertarget{setting-and-notation-1}{%
\subsection{1. Setting and notation}\label{setting-and-notation-1}}

(X,g) Einstein, Ric = (s/4)g, p ∈ X, normal coordinates x on
B\_\{4T\}(p). Λ⁺\_p is identified with span(J₁,J₂,J₃) (quaternionic
structures, quaternion convention of the Step-49--52 scripts), the Burns
orientation ζ = J₃; v\_m = J\_m x, v₀ = x, u = \textbar x\textbar², P\_m
= v\_m v\_mᵀ/u.

R⁺ = W⁺ + s/12 on Λ⁺ (eigenvalues sum to s/4); μ\_ζ = R⁺(ζ,ζ); T =
trace-free part of R⁺ restricted to ζ\^{}⊥ ∩ Λ⁺ (a 2×2 symmetric matrix
in an orthonormal basis), \textbar T\textbar² = (μ₁−μ₂)²/2 when ζ is an
eigenvector.

Flat Burns: δ + εk\_B, k\_B = (v₁v₁ᵀ + v₂v₂ᵀ)/u² (indices m = 1,2 ⊥ ζ),
exactly ASD and scalar-flat, exceptional sphere of area πε, Kähler form
ζ (self-dual in X's orientation, so the blow-up is X\#CP̄²).

Chiral curvature 𝒞 = (2W⁺, s/√12); f(g) = \textbar 𝒞\textbar² dV\_g.
Pointwise CGB: \textbar Ric°\textbar² dV = f − 16π² e − 24π² l (e Euler,
l signature density). Hence for closed M, P(g) = ∫\_M f(g) −
8π²(2χ+3τ)(M).

\hypertarget{the-exact-local-coefficient-e}{%
\subsection{2. The exact local coefficient
{[}E{]}}\label{the-exact-local-coefficient-e}}

Scripts: \texttt{workspace/scripts/step52\_exact\_su2minus\_fast.py}
(exact rational quadratic form by polarization, truncated bivariate
(a,ε) arithmetic in Fractions, 0.2 s per evaluation) and
\texttt{step52\_analyze\_Q.py}. Data:
\texttt{workspace/data/step52a/exact\_su2minus\_Q.json},
\texttt{exact\_su2minus\_analysis.json}.

\textbf{Reduction.} If W⁻\_p = 0 and Ric°\emph{p = 0, every ingredient
(the 2-jet h₂ = −(1/3)R}\{ikjl\}x\textsuperscript{kx}l, k\_B, and
correctors built from x and the v\_m) is SU(2)₋-equivariant; SU(2)₋
commutes with the J\_m and acts transitively on S³. For SU(2)₋-invariant
correctors the a²ε² CGB density is therefore \textbf{constant on S³}, so
F = density(x₀)/4 exactly. The corrector functional is convex (its
Hessian is the flat (W⁺,s)-linearization norm), so its infimum over any
SU(2)₋-stable class of degree-0 correctors is attained on the invariant
ones: the 10 tensors sym(v\_i v\_jᵀ)/u, 0 ≤ i ≤ j ≤ 3. This is the full
answer in the (s,W⁺) sector, for an \textbf{arbitrary} W⁺ (U(1)₊
rotation allows w₂₃ = 0 without loss).

\textbf{Checks inside the run.} The a-linear W⁺ and s of the jet equal
the prescribed (w, s) exactly; CH k\_B-only gives F = 12 exactly (Step
49 floating 11.9999999999); the {[}aε{]} coefficient and the pure-Burns
ε² coefficient vanish identically for all 120 polarization vectors.

\textbf{Theorem 2.1 {[}E{]}.} Let ζ be an eigenvector of W⁺\_p (w₁₃ =
w₂₃ = 0). Over all degree-0 correctors in the (s,W⁺) sector,

F\_opt(Rm\_p, ζ) = −μ\_ζ(μ\_ζ − s/8) − \textbar T\textbar²,

exactly (output \texttt{Fopt},
\texttt{Fopt\_minus\_Phi\_w13=0\ =\ −2w₁₂²}, which is precisely the
off-diagonal part of −(μ₁−μ₂)²/2 in a non-eigenbasis of ζ\^{}⊥). This is
A02's Theorem 57.Φ, now an exact theorem in this sector, obtained by an
independent normal-coordinate computation. Special values: U(2) locus
reproduces (s−24λ)(s+12λ)/288 (Step 51) exactly; W⁺ = 0 gives +s²/288
\textgreater{} 0.

\textbf{Theorem 2.2 (explicit kernel corrector) {[}E{]}.} On the kernel
locus μ\_ζ = 0, with A = −s/12, the family

k₀\^{}\{ker\} = (7A/4)(P₁+P₂) + (A/2)P₃ + (1/3) Σ\_\{i,j≤2\} T\_ij (J\_i
x)⊗(J\_j x)/u

(Spirit \#61) satisfies F(k₀\^{}\{ker\}) = −\textbar T\textbar²
identically in (s, T), and the gradient of F at k₀\^{}\{ker\} vanishes
in all 10 directions: it is an exact minimizer, not only a competitor.
At s = −24, T = diag(3,−3) its coefficients are (4.5, 2.5, 1), the
degree-zero part of the Ch. 5 neck c\_* = (9/4,3/4,1/2) (the extra +1 on
the Kähler fibre is the C² = 1 + t² + \ldots{} factor of the CH radius).
Hence the Ch. 5 value −18 = −\textbar T\textbar² is reproduced exactly.

\textbf{W⁻ and Ric° sectors.} The integrated F is SU(2)₋-invariant on
(jet ⊕ corrector) space. The (s,W⁺)-jets and k₀\^{}\{ker\} are
SU(2)₋-invariant; W⁻ and its corrector k₀\^{}W = (3/2)sym(h\_W k\_B) −
(1/4)tr(h\_W k\_B)δ lie in the spin-2 isotypic part. Cross terms vanish
by irreducibility, and the W⁻ self-term vanishes pointwise by the
Step-50 exact identity. So **for every Einstein 2-jet with ζ ∈ ker
R⁺\_p,**

F(k₀\^{}\{ker\} + k₀\^{}W) = −\textbar T\textbar² {[}E + representation
theory{]}.

(Ric° = 0 for Einstein. For non-Einstein jets the Ric° self-term is ≥ 0
and floating-zero, Step 51 {[}N{]}; a closed-form Ric° corrector is
being sought in \texttt{scripts/step52\_ric0\_corrector.py}, see §7.)

\textbf{Mixed block {[}E{]}.} If w₁₃ ≠ 0, F contains 2w₁₃(−9c₂ + 15c₆ +
7w₁₃)/9, where c₂ multiplies sym(x⊗J₂x)/u, a null direction of the
corrector Hessian (it is L\_\{(log r)J₂x\}δ, a rotation of the Burns
orientation by an angle growing like log r). The log coefficient is then
unbounded below. This happens whenever the mixed block is nonzero;
combined with μ\_ζ = 0 it forces R⁺ indefinite; there a direction with μ
\textgreater{} 0 already gives a first-order saving, and rotating ζ
along the neck exploits the gradient 2R⁺(ζ,ζ'). On R⁺ ≤ 0, μ\_ζ = 0
forces R⁺ζ = 0, so this does not occur.

\hypertarget{the-two-scale-construction-and-its-remainder-p}{%
\subsection{3. The two-scale construction and its remainder
{[}P{]}}\label{the-two-scale-construction-and-its-remainder-p}}

\hypertarget{the-metric}{%
\subsubsection{3.1 The metric}\label{the-metric}}

Fix T ≤ inj(p)/4 small, cutoffs χ\_in(ρ) (0 on ρ ≤ 1, 1 on ρ ≥ 2),
χ\_out(ρ) (1 on ρ ≤ 1, 0 on ρ ≥ 2), and k = k\_B + k₀, k₀ =
k₀\^{}\{ker\} + k₀\^{}W (degree 0, smooth on R⁴∖0). Define

\begin{itemize}
\tightlist
\item
  r ≤ T: g\_ε = δ + εk\_B + χ\_in(r/√ε)·{[}(g − δ) + εk₀{]};
\item
  r ≥ T: g\_ε = g + ε χ\_out(r/T)·k; g\_ε = g outside B\_\{2T\}.
\end{itemize}

The formulas agree on 2√ε ≤ r ≤ T (both equal g + εk). On r ≤ √ε, g\_ε
is exactly flat Burns, which extends smoothly over the exceptional
sphere (blow-up of (C², J₃) at 0, self-intersection −1 in X's
orientation). Positivity holds for ε small since k\_B ≥ 0, g − δ =
O(r²), εk₀ = O(ε). So g\_ε is a smooth metric on M = X\#CP̄², identical
to g outside B\_\{2T\}(p).

Since 2χ+3τ drops by one, P(g\_ε) − P(g) = 8π² + Δ, Δ = ∫\emph{\{B̃\}
f(g\_ε) − ∫}\{B\_\{2T\}\} f(g).

\hypertarget{region-estimates}{%
\subsubsection{3.2 Region estimates}\label{region-estimates}}

Scale invariance: in polar coordinates
∫\emph{\{r₁\textless r\textless r₂\} f(g) = ∫ (dr/r) ∫}\{S³\} f(ĝ\_r)(θ)
dθ with ĝ\_r = r⁻²Φ\_r\^{}* g (Φ\_r(x) = rx).

\textbf{Core, r ≤ 2√ε.} ε⁻¹Φ\_\{√ε\}\^{}* g\_ε = Burns₁ +
χ(\textbar x̂\textbar)·ε·m\_ε(x̂) with m\_ε bounded in C\^{}k on 1 ≤
\textbar x̂\textbar{} ≤ 2; on \textbar x̂\textbar{} ≤ 1 it is exactly
Burns (f = 0). Since 𝒞(Burns) ≡ 0 and f is quadratic in 𝒞, f = O(ε²)
there; the core contributes O(ε²), as does ∫\emph{\{B}\{2√ε\}\} f(g).

\textbf{Neck, 2√ε ≤ r ≤ T.} With z = ε/r², ĝ\_\{r\} = δ + G(r,x) + z
k\_B(x) + z r² k₀(x), G(r,x) = g(rx) − δ = r²h₂ + r³h₃ + \ldots{}
(Cartesian components; TTCS repair of an erroneous r⁻² factor), smooth
in (r, x) up to r = 0. Put Ψ(r,z) = (8π²)⁻¹∫\emph{\{S³\} f(ĝ}\{r,z\}); Ψ
is smooth on {[}0,T{]}×{[}0,1/4{]}. (a) Ψ(0,z) = 0 (flat Burns). (b) Ψ
is even in r: the antipodal map x ↦ −x has det +1, preserves S³ with its
orientation, fixes k\_B and k₀ (even tensors) and maps h\_m to (−1)\^{}m
h\_m, i.e.~r to −r. (c) 𝒞(ĝ) = r²𝒞₂(z) + r³𝒞₃(z) + \ldots{} (no r¹
term), so Ψ = O(r⁴) uniformly in z. Let R = Ψ(r,z) − Ψ(r,0) −
z∂\emph{zΨ(r,0). R vanishes to order 4 in r and order 2 in z, so by
Hadamard R = r⁴z²ρ(r,z) with ρ smooth, ρ even in r, and ρ(0,0) =
(8π²)⁻¹∫}\{S³\}{[}a²ε²{]}f = F (§2: only h₂ and k₀ enter, since the r⁴
coefficient of \textbar 𝒞\textbar² is \textbar 𝒞₂\textbar²). With
\textbar ρ − F\textbar{} ≤ C(r² + z): ∫\_\{2√ε\}\^{}\{T\} R dr/r = ε² ∫
ρ(r, ε/r²) dr/r = (F/2) ε² log(1/ε) + O\_T(ε²), using ∫ r dr ≤ T²/2 and
∫ ε r⁻³ dr ≤ 1/8 as in Ch. 5. The term Ψ(r,0) is the background and
cancels.

\textbf{Outer transition, T ≤ r ≤ 2T.} f(g + εχk) = f(g) + ε
Df\_g{[}χk{]} + O\_T(ε²).

\hypertarget{the-order-ux3b5-term-is-a-flux-and-vanishes-on-the-kernel}{%
\subsubsection{3.3 The order-ε term is a flux, and vanishes on the
kernel}\label{the-order-ux3b5-term-is-a-flux-and-vanishes-on-the-kernel}}

The z-linear neck term and the outer linear term together equal
ε∫\emph{\{r≥2√ε\} Df\_g{[}χ\_out k{]} dV (the z-coefficient at fixed r
is the ε-derivative of f(g + εk)). Since g is Einstein,
D(\textbar Ric°\textbar² dV)\emph{g = 0 pointwise, and the variations of
the Chern--Weil forms are exact: Df\_g{[}h{]} = dΘ\_g(h), Θ linear and
local in h. Because χ\_out k = 0 near r = 2T, ε∫}\{r≥2√ε\} Df\_g{[}χk{]}
= −Λ(2√ε), Λ(r) := ε∮}\{S\_r\} Θ\_g(k). Θ is scale invariant and
vanishes at the flat metric (CW forms are quadratic in curvature), so
Λ(r) = ε{[}θ₂ + θ₃ r + O(r²){]} + ε·O(r²) (the k₀ part starts at r²),
and θ₃ = 0 by the antipodal parity. Thus the order-ε contribution is
−εθ₂ + O(ε²), where θ₂ is bilinear in (Rm\_p, k\_B): a U(2)\_ζ-invariant
linear function of Rm\_p.~The only such invariants are s and W⁺(ζ,ζ)
(W⁻, Ric° and T carry nontrivial SU(2)₋ or U(1)₊ weight), so θ₂ = αs +
βW⁺(ζ,ζ).

\textbf{Direct computation (TTCS step52 verifier, replaces the
calibration below) {[}P{]}.} With f − \textbar Ric°\textbar²dV =
8π²p₁(Λ⁺) and Θ\_g(h) = −2tr(Ȧ⁺{[}h{]}∧F⁺): the flat Burns family is
Kähler for J\_ζ, its Λ⁺-connection variation is Ȧ⁺{[}k\_B{]} = L\_ζ·r⁻²η
(η = r⁻²ι\_xΩ, Ω = √2ζ, ∫\_\{S³\}η∧dη = 4π²), and the L\_ζ-component of
F⁺ pairs through μ\_ζΩ, giving θ₂ = 8π²μ\_ζ. Hence the order-ε term is
exactly −8π²εR⁺\_p(ζ,ζ), i.e.~ℓ = −μ\_ζ in P/8π² units (agreeing with
Step 50). The Ch. 5 calibration that follows is retained only as a
consistency remark: the verifier notes that its normal-coordinate
equivalence (singular linear part a multiple of k\_B modulo gauge and
degree ≥ 0 terms) was not established from the supplied material.

Calibration {[}C: Ch. 5 (5.A), superseded{]}. Ch. 5's CH construction
(Ric = −6g, ζ ⊥ ω, a kernel direction since R⁺\_CH = (−6,0,0)) has in
normal coordinates the same two-scale form with k\_B at leading order
(its extra terms are degree ≥ 0 and enter Λ only at O(εr²)), and its
theorem has no order-ε term. Hence α·(−24) + β·2 = 0, i.e.~θ₂ =
β·R⁺(ζ,ζ) (β = 12α). The Kähler period law of Step 8 (aligned Kähler
collar, P = 2(2π − λt)² + L/4, t the exceptional period, λ the Einstein
constant) has a nonzero first-order term in t whenever λ ≠ 0, and the
aligned direction has R⁺(ω,ω) = s/4 ≠ 0; so β ≠ 0, with cost −εθ₂ ∝
−R⁺(ζ,ζ) (positive for negative KE, μ\_ω = s/4 \textless{} 0, and a
saving for positive KE, as in Step 8). This agrees with the TTCS-derived
flux ℓ = −R⁺\_p(ζ,ζ) of Step 50; only its vanishing on ker R⁺ is used
below.

\hypertarget{theorem}{%
\subsubsection{3.4 Theorem}\label{theorem}}

\textbf{Theorem 52.A (kernel insertion) {[}P, with §2 exact and the Ch.
5 order-ε calibration{]}.} Let (X,g) be an Einstein four-manifold, p ∈
X, ζ ∈ Λ⁺\_p a unit vector with R⁺\_pζ = 0, and T the trace-free part of
R⁺\emph{p on ζ\^{}⊥. The metrics g\_ε of §3.1 (with k₀ = k₀\^{}\{ker\} +
k₀\^{}W) are smooth on X\#CP̄², equal g outside B}\{2T\}(p), and

P(g\_ε) = P(g) + 8π² − 4π²\textbar T\textbar² ε² log(1/ε) + O(ε²),

with O(ε²) depending on g\textbar{}\emph{\{B}\{4T\}(p)\}, ζ, T and
uniform for (p,ζ) in compact sets. If X is closed, P(g) = 0 and, for T ≠
0 and ε small, P(g\_ε) \textless{} 8π².

Remarks. (i) Every negative or positive KE point, ζ ⊥ ω: T = diag(s/8,
−s/8), \textbar T\textbar² = s²/32 (= 18 at s = −24, giving −72π²ε²
log(1/ε) = 8π²·(−9ε² log(1/ε)), Ch. 5). (ii) No smallness of W⁻ and no
CH collar is required. (iii) The fixed-configuration use needs only one
such (p,ζ) and one fixed small ε; uniformity over moving insertion data
is not needed for that contradiction (Spirit \#61). (iv) This is an
upper construction only.

\textbf{Corollary 52.B.} For an Einstein point p, a Burns insertion
lowers P below P(g)+8π² at leading order iff either R⁺\_p has a positive
eigenvalue (first order, ζ = that eigenvector) or R⁺\_p has a kernel
vector with T ≠ 0 (ε² log order) --- within this two-scale ansatz with
fixed degree-0 correctors and at these two orders. Full expansion for an
eigenvector ζ: (P(g\_ε) − P(g) − 8π²)/8π² = −μ\_ζε + ½{[}−μ\_ζ(μ\_ζ −
s/8) − \textbar T\textbar²{]}ε²log(1/ε) + O(ε²). No decrease at these
orders exactly in the cases (i) R⁺\_p \textless{} 0; (ii) s \textless{}
0, spec R⁺\_p = (0, s/8, s/8) (W⁺ = diag(−s/12, s/24, s/24), Derdziński
degenerate type); (iii) R⁺\_p = 0 (s = 0, W⁺ = 0).

\hypertarget{no-open-degenerate-kernel-stratum-on-an-einstein-metric-p}{%
\subsection{4. No open degenerate-kernel stratum on an Einstein metric
{[}P{]}}\label{no-open-degenerate-kernel-stratum-on-an-einstein-metric-p}}

\textbf{Lemma 52.C.} Let g be Einstein with s \textless{} 0 on a
connected open set U on which W⁺ has eigenvalues (a, −a/2, −a/2), a =
−s/12. Then U = ∅.

Proof. The simple eigenvalue gives a smooth (up to sign, locally) unit
self-dual form u with W⁺ = (3a/2)(u⊗u) − (a/2)·Id\_\{Λ⁺\}
(\textbar u\textbar² normalized so u⊗u is the orthogonal projector), a
constant. Einstein metrics satisfy δW⁺ = 0 (second Bianchi + Ric
parallel). The identity term is parallel, so ∇\^{}i(u\_\{ij\}u\_\{kl\})
= 0, i.e.~(δu)\emph{j u}\{kl\} + u\_\{ij\}∇\^{}i u\_\{kl\} = 0. Contract
with u\^{}\{kl\}: \textbar u\textbar² constant gives
u\^{}\{kl\}∇u\_\{kl\} = 0, hence δu = 0. Then u\_\{ij\}∇\^{}iu\_\{kl\} =
0 for all j,k,l; u is nondegenerate (u² = −c·Id), so multiplying by
u\^{}j\_m gives ∇\emph{m u}\{kl\} = 0. A parallel self-dual form makes g
Kähler with W⁺(u) = (s/6)u, while here W⁺(u) = a u = −(s/12)u; equality
forces s = 0, contradiction. ∎

Consequently, for s \textless{} 0, the contact set Z = \{p:
λ\_max(W⁺\_p) = −s/12, λ\_mid = λ\_min\} of an Einstein metric with R⁺ ≤
0 is closed with empty interior (for s = 0, R⁺ ≤ 0 forces R⁺ ≡ 0, a
separate branch; s \textgreater{} 0 is impossible). The maximum
principle does not remove isolated contact: at p ∈ Z, with e the top
eigenvector extended radially parallel, φ = ⟨W⁺e,e⟩ has Δφ(p) = ⟨ΔW⁺
e,e⟩ = (s/2)a − 6a² + 2\textbar W⁺\textbar² = −9a² \textless{} 0
(Einstein formula ΔW⁺ = (s/2)W⁺ − 6W⁺∘W⁺ + 2\textbar W⁺\textbar²Id,
checked on KE), consistent with a maximum. The lemma does not apply to
non-Einstein P-critical bulk, and does not by itself produce T ≠ 0
kernel points nearby.

\hypertarget{consequences-for-the-degeneration-chain-with-the-correct-sign-statement}{%
\subsection{5. Consequences for the degeneration chain, with the correct
sign
statement}\label{consequences-for-the-degeneration-chain-with-the-correct-sign-statement}}

In A02's energy-equality configuration (Theorem 53.B), bubble movement
plus Theorem 52.A excludes Einstein bulk points with a positive R⁺
eigenvalue and kernel points with T ≠ 0. Any surviving Einstein bulk
satisfies R⁺ ≤ 0 everywhere, R⁺ \textless{} 0 off a closed set Z with
empty interior, and λ\_max(W⁺) ≤ \textbar s\textbar/12, so
\textbar W⁺\textbar² ≤ s²/24 with equality nowhere (equality needs
eigenvalues (a,a,−2a), a two-dimensional kernel with T ≠ 0).

Correction to the ``negative on the whole cone'' reading: F ≤ 0 is the
conclusion \textbf{on the kernel}. In the interior R⁺ \textless{} 0 the
log coefficient can be positive (W⁺ = 0: F = s²/288 \textgreater{} 0;
generally −μ\_ζ(μ\_ζ − s/8) \textgreater{} 0 when s/8 \textless{} μ\_ζ
\textless{} 0); it is subordinate to the positive first-order cost −εθ₂
∝ ε\textbar μ\_ζ\textbar. The strictly negative cone (e.g.~real
hyperbolic metrics, W⁺ = 0) is a genuine survivor and is a global
problem, not a remainder term. LeBrun's Weyl estimate with Gauss--Bonnet
does not exclude it when c₁² = A (1/4 + 2t − 5t² ≥ 0 on {[}0,1/2{]}, t =
√6\textbar W⁺\textbar/\textbar s\textbar; Step 51).

\hypertarget{relation-to-the-four-public-problems}{%
\subsection{6. Relation to the four public
problems}\label{relation-to-the-four-public-problems}}

\begin{itemize}
\tightlist
\item
  Problem 2 (optimal P on blow-ups): p\_*(X\#CP̄²) \textless{} 8π² for
  every Einstein X having a kernel point with T ≠ 0 --- all KE surfaces
  of either sign --- strengthening Step 50 from KE-collars to arbitrary
  Einstein kernel points, with an explicit closed-form corrector.
  Replication on covers and the Yamabe shortcut (Step 50 §4) give the
  strict sub-four slope on covers of any fixed such base. No lower bound
  is implied.
\item
  Problem 3 (attainment / non-Bach-flat minimizers): the theorem
  supplies the local strict-decrease input in the fixed energy-equality
  configuration; what remains open is the R⁺ \textless{} 0 bulk branch
  and isolated contact.
\item
  Problem 1: the surviving pointwise constraint t ≤ 1/2 is exactly the
  missing t-distribution information of Step 48; excluding Einstein bulk
  with t ≤ 1/2 on a manifold with an SW basic class and c₁² = A needs a
  differential/global input.
\end{itemize}

\hypertarget{open-items}{%
\subsection{7. Open items}\label{open-items}}

\begin{enumerate}
\def\labelenumi{\arabic{enumi}.}
\tightlist
\item
  Closed-form Ric° corrector (non-Einstein P-critical bulk).
  \texttt{step52\_ric0\_corrector.py}
  (data/step52/ric0\_corrector.json): the 17-tensor natural span
  contains a floating solution (least-squares residual 1e-14, rank 12 of
  17, so non-unique), but naive rationalization (limit\_denominator)
  gave wrong rationals and the symbolic identity check failed. Next:
  exact rational solve with the 5 free parameters fixed (e.g.~eliminate
  dependent columns), then re-verify. Averaging-closed corrector classes
  are required in the §2 averaging statement (TTCS repair).
\item
  Non-Einstein bulk: the order-ε term acquires the bulk first variation
  DP\_g{[}χk{]} unless g is P-critical, and Lemma 52.C uses δW⁺ = 0.
\item
  The ε² analysis at isolated contact points of Z and whether bubbles
  can be moved to nearby T ≠ 0 kernel points (kernel points are only on
  Z when R⁺ ≤ 0, so this reduces to the R⁺ \textless{} 0 branch near Z).
\item
  The R⁺ \textless{} 0 global branch (Problem 1 link).
\end{enumerate}

\hypertarget{s1003a01-step-53-the-trace-free-ricci-compatibility-of-the-perpendicular-burns-tail-is-conformaldiffeomorphism-covariance-plus-harmonicity}{%
\chapter{S1003A01: Step 53 --- The trace-free-Ricci compatibility of the
perpendicular Burns tail is conformal--diffeomorphism covariance plus
harmonicity}\label{s1003a01-step-53-the-trace-free-ricci-compatibility-of-the-perpendicular-burns-tail-is-conformaldiffeomorphism-covariance-plus-harmonicity}}

Source: \texttt{step53\_ric0\_geometric\_corrector.md}.

Session S1003A01, Agent2. Inputs: Spirit \#62 (candidate formula), Step
50 (W⁻ corrector), Step 52 (Theorem 52.A,
\texttt{step52\_kernel\_insertion\_theorem.md}). Script:
\texttt{workspace/scripts/step53\_ric0\_geometric\_corrector.py}; output
\texttt{workspace/data/step53/ric0\_geometric\_corrector.json}, log
\texttt{ric0\_geo\_log.txt}.

\hypertarget{setting}{%
\subsection{Setting}\label{setting}}

Flat chart ℝ⁴∖0, ζ = J = J₃, Burns g\_ε = δ + εk\_B, k\_B =
(v₀v₀ᵀ+v₁v₁ᵀ)/u², v\_m = J\_m x, u = \textbar x\textbar². The Burns
metric is exactly ASD and scalar-flat: W⁺(g\_ε) = 0 and s(g\_ε) = 0 for
every ε. For E trace-free symmetric, q = xᵀEx, E\_H = (E + JᵀEJ)/2, q\_H
= xᵀE\_Hx. The pure-Ric° normal 2-jet is h\_E = −(1/6)(E u + q δ − Ex xᵀ
− x (Ex)ᵀ) (Kulkarni--Nomizu ½E∧δ, h\_E x = 0).

\hypertarget{lemma-53.1-geometric-ric-corrector}{%
\subsection{Lemma 53.1 (geometric Ric°
corrector)}\label{lemma-53.1-geometric-ric-corrector}}

With f = −q/4, V = −uEx/12 + qx/6, f\_ε = f − ε q\_H/(8u), and k₀\^{}E =
L\_V k\_B − (q/2) k\_B − (q\_H/(4u)) δ, the following hold identically
on ℝ⁴∖0: W⁺'\,'\_δ(h\_E,k\_B) + W⁺'\_δ(k₀\^{}E) = 0, s'\,'\_δ(h\_E,k\_B)
+ s'\_δ(k₀\^{}E) = 0. k₀\^{}E is homogeneous of degree 0, so it belongs
to the same corrector class as the Step-50 W⁻ corrector k₀\^{}W.

Proof. (i) Exactly h\_E = 2fδ + L\_Vδ (checked symbolically). (ii)
Consider G\_t = φ\_t\^{}*(e\^{}\{2tf\_ε\}g\_ε) with φ\_t the flow of V.
W⁺ is conformally covariant and natural, so W⁺(G\_t) = 0 for all t. In
dimension four s(e\^{}\{2u\}g) = e\^{}\{−2u\}(s − 6Δu −
6\textbar du\textbar²), hence d/dt\textbar₀ s(G\_t) = −6Δ\_\{g\_ε\}f\_ε
(s(g\_ε) = 0). (iii) f\_ε is g\_ε-harmonic to first order in ε: Δ\_δ f =
−½ tr E = 0, and the ε-coefficient −k\_B\^{}\{ij\}∂\emph{\{ij\}f − (div
k\_B − ½d tr k\_B)·df + Δ\_δ(−q\_H/(8u)) vanishes identically (checked
exactly; the same expression for the uncorrected f is nonzero). (iv)
Therefore DW⁺}\{g\_ε\}{[}L\_Vg\_ε + 2f\_εg\_ε{]} = 0 and
Ds\_\{g\_ε\}{[}L\_Vg\_ε + 2f\_εg\_ε{]} = O(ε²). The ε⁰ coefficient is
the statement W⁺`(h\_E) = s'(h\_E) = 0 (pure Ric° jet). The ε¹
coefficient is D²(W⁺,s)\_δ{[}k\_B,h\_E{]} + D(W⁺,s)\_δ{[}L\_Vk\_B +
2fk\_B + 2f₁δ{]} = 0 with f₁ = −q\_H/(8u), which is the lemma. Vanishing
of W⁺ is frame independent, so the orthonormal-frame expansion used by
the curvature engine vanishes order by order. ∎

Independent exact check (curvature engine of
\texttt{step52\_ric0\_corrector.py}, two unrelated E, all 7 components
(6 of W⁺, 1 of s) at 5 random rational points each): all zero; the
uncorrected source is nonzero. Control: dropping the q\_H term leaves
all W⁺ components zero and only the scalar component nonzero (12/1285,
\ldots), exactly as conformal covariance predicts.

\hypertarget{meaning}{%
\subsection{Meaning}\label{meaning}}

\begin{itemize}
\tightlist
\item
  The Ric° coefficient of the logarithmic Burns corner term vanishes
  because the Ric° jet is pure gauge-plus-conformal at the 2-jet level
  (h\_E = L\_Vδ + 2fδ), and Burns admits the conformal continuation
  f\_ε. The ε²log(1/ε) cross term between Ric°\_p and Burns is a
  completed square with zero minimum, just as for W⁻. The floating
  least-squares solutions of Step 52 (rank 12 in a 17-tensor span) were
  non-unique because of this gauge/conformal redundancy; the
  rationalized coefficients were spurious.
\item
  Together with Step 50 (W⁻) and Step 52 (exact (s,W⁺) sector), the
  formal corner coefficient at any point of any metric, in the CGB form
  4\textbar W⁺\textbar²+s²/12, is F = −μ\_ζ(μ\_ζ − s/8) −
  \textbar T\textbar² (ζ an R⁺-eigenvector), with no W⁻ or Ric°
  contribution. The Gauss--Bonnet/signature reduction P =
  2∫(2\textbar W⁺\textbar² + s²/24) − const holds for every metric, so
  the corner computation itself needs no Einstein hypothesis.
\item
  What still differs for non-Einstein (e.g.~P-critical) bulk is not the
  corner coefficient but (a) the order-ε bulk variation DP\_g{[}χk{]},
  absent only when g is P-critical; (b) whether higher jets still enter
  only through exact forms (proved via Einstein in Step 50). A02's Part
  A (P-critical remainder) is the place where (a)--(b) are handled;
  Lemma 52.C (no open degenerate stratum) still requires δW⁺ = 0 and
  does not transfer.
\end{itemize}

\hypertarget{s1003a01-step-53-problem-4-negative-scalar-goldberg-formal-jet-structure-and-the-closed-canonical-form-ux3b2}{%
\chapter{S1003A01: Step 53 --- Problem 4 (negative-scalar Goldberg):
formal-jet structure and the closed canonical form
β}\label{s1003a01-step-53-problem-4-negative-scalar-goldberg-formal-jet-structure-and-the-closed-canonical-form-ux3b2}}

Source: \texttt{step53\_goldberg\_jets\_and\_canonical\_form.md}.

Session S1003A01, Agent2. Labels: {[}P{]} proved here (by hand), {[}N{]}
numerical evidence, {[}K{]} cited, {[}O{]} open.

\hypertarget{setting-and-conventions}{%
\subsection{0. Setting and conventions}\label{setting-and-conventions}}

(M⁴,g,J,ω) closed, ω(X,Y)=g(JX,Y), dω=0, Ric=λg, λ\textless0, s=4λ.
Orientation of J; ω self-dual harmonic, \textbar ω\textbar²=2 (form
norm). q=\textbar∇ω\textbar², s\emph{=s+q. W⁺ in the splitting ℝω/√2 ⊕
ω\^{}⊥ (five\_research\_problems.tex §4, eq. (eta),(canonicalcurv)): W⁺
= {[}{[}ν, η*/√2{]},{[}η/√2, −ν/2 I + V{]}{]}, ν = s/6 + q/4. Canonical
(Blair/Chern) connection 𝒜 on K⁻¹: iF\_𝒜 = η + ((s+s})/8) ω + ((s−s*)/8)
ω̂, ω̂ ASD with \textbar ω̂\textbar²=2 on \{q\textgreater0\}. Integrability
bridge {[}K, AAD{]}: η ≡ 0 ⇒ Kähler on closed M.

\hypertarget{local-status-k}{%
\subsection{1. Local status {[}K{]}}\label{local-status-k}}

AAD, ``Local rigidity of certain classes of almost Kähler 4-manifolds''
(Knowledge object Local-rigidity-\ldots c0d51aa8ffb7, intro lines
33--38): - Nurowski--Przanowski and Tod give \emph{local} strictly AK
Einstein metrics, all Ricci-flat; in all of them W⁻=0 and η=0. -
{[}AAD/Armstrong, ref. {[}6{]} there{]} every strictly AK Einstein
4-manifold with ω an eigenform of W⁺ (η=0) is of NPT type, hence
Ricci-flat and never compact. - ``The case of negative scalar curvature
is still wide open.'' Consequence: for λ\textless0 the existence of
\emph{local} strictly AK Einstein metrics is itself not supplied by
these sources (literature request s53\_t7\_tool1 pending). If no local
λ\textless0 examples exist, a finite-order prolongation obstruction
would prove local, hence global, Goldberg; if they exist formally, the
problem is genuinely global. Deciding which is the case is therefore a
decisive question, not a side computation.

\hypertarget{formal-system-and-corrected-symbol-interpretation-p}{%
\subsection{2. Formal system and corrected symbol interpretation
{[}P{]}}\label{formal-system-and-corrected-symbol-interpretation-p}}

Unknown jets in coordinates with g=δ+O(\textbar x\textbar²),
ω(0)=e⁰¹+e²³: g (10 comps, degree 2..D), ω (6 comps, degree 1..D).
Equations: E1 ωᵀg⁻¹ω=g, E2 dω=0, E3 Ric=λg. The Step-53 count omitted an
independent scalar Ricci-symbol relation. In Darboux gauge, a compatible
variation h satisfies hJ+Jh=0 and
σ\_Ric(ξ)h(Jξ,Jξ)−(1/2)\textbar ξ\textbar² tr σ\_Ric(ξ)h=0. Its
prolongations add n\_\{N−4\}; the corrected top-degree free count is
N²+8N−5 for N≥4, so the earlier negative count at N=15 and estimate near
18 do not predict a Goldberg obstruction. Nor can the full nonlinear
Jacobian rank detect a genuine lower-jet condition reliably: one must
hold a lower jet fixed and project the next-order source onto the symbol
cokernel after removing known identities. Exact details and the
intrinsic AAD interpretation are in
\texttt{step55\_darboux\_ricci\_compatibility.md}.

\hypertarget{numerical-jet-engine-n}{%
\subsection{3. Numerical jet engine
{[}N{]}}\label{numerical-jet-engine-n}}

\texttt{workspace/scripts/step53\_ak\_einstein\_jets.py} uses JAX
truncated multivariate polynomials and all-at-once Gauss--Newton. The
exact H²(−1)×H²(−1) Kähler--Einstein jet has residual 0 at D=4. D=4 and
independently optimized D=5 truncations exist numerically for q₀=0, RF
(λ=0,q₀=0.3), and NEG (λ=−1,q₀=0.3), with residual squares about 10⁻²⁷.
At D=5, ranks are 1035 for q₀=0 and 1036 for RF/NEG. These independently
optimized jets are not compatible formal germs. The q₀=0 random jets are
also not Kähler germs: their pointwise q is zero but the reported η and
W⁺ are generic, so only the explicit H²×H² jet is an exact KE control.
D6/D7 files from the cancelled high-memory task are incomplete and have
no mathematical meaning. A first fixed D4→D5 extension of the numerical
random jets gave nonzero cokernel projections for every branch,
including the controls; it therefore establishes only that these random
D4 truncations do not lie on the sampled degree-five branches. Before
further prolongation the code must use exact KE and exact Tod sources
and quotient the cokernel by algebraic compatibility, closure, Bianchi,
and the Darboux Ricci relation. Re-running the unconstrained degree
ladder is not warranted.

\hypertarget{the-closed-canonical-form-ux3b2-p}{%
\subsection{4. The closed canonical form β
{[}P{]}}\label{the-closed-canonical-form-ux3b2-p}}

Einstein gives ρ=(s/4)ω for the Riemannian Ricci form, closed. Put β :=
iF\_𝒜 − (s/4)ω. Then dβ=0, {[}β{]}=2πc₁ − (s/4){[}ω{]}, and β⁺ = η +
(q/8)ω, β⁻ = −(q/8)ω̂, β − η = (q/8)(ω − ω̂) = ± a∧Ja (decomposable,
\textbar·\textbar² = q²/16). Since η ∈ {[}{[}Λ\^{}\{2,0\}{]}{]} and a∧Ja
is of type (1,1), β∧β = \textbar η\textbar² vol pointwise. Hence
∫\textbar η\textbar² = {[}β{]}² --- this is exactly the AK balance
(AKbalance) again, now as a cup-square identity. Hodge comparison with
the harmonic representative H of {[}β{]} (H⁺ = 2πc₁\^{}\{+,h\} − (s/4)ω,
H⁻ = 2πc₁\^{}\{−,h\}): ‖β‖² ≥ ‖H‖² and ‖β⁺‖²−‖β⁻‖² = ‖H⁺‖²−‖H⁻‖² give ∫
q²/16 ≥ 2‖H⁻‖² = 8π² ‖c₁\^{}\{−,h\}‖²\_\{L²\}. So the anti-self-dual
harmonic part of c₁ is a lower bound for nonintegrability (vanishes for
KE, consistent). It does not force q=0; it shows that a strictly AK
Einstein metric whose harmonic c₁ is not self-dual must pay q in L². The
upper side still needs a new compatibility relation.

\hypertarget{open-o}{%
\subsection{5. Open {[}O{]}}\label{open-o}}

\begin{itemize}
\tightlist
\item
  Local λ\textless0,q\textgreater0 existence remains unresolved. The
  next meaningful local calculation is an intrinsic prolongation of the
  AAD Ricci/Bianchi system, or a fixed-lower-jet calculation based on
  exact KE and exact Tod controls with all universal identities removed.
\item
  Globally, the checked Einstein AAD equations couple q, η, W₃⁺ and the
  Nijenhuis derivative; the exact identities are recorded in
  \texttt{step55\_darboux\_ricci\_compatibility.md}. The missing input
  is canonical-connection/actual-field coercivity for the transverse
  terms, not another scalar counting argument.
\end{itemize}

\hypertarget{s1003a01-step-55-darboux-gauge-ricci-compatibility-for-negative-scalar-goldberg}{%
\chapter{S1003A01: Step 55 --- Darboux-gauge Ricci compatibility for
negative-scalar
Goldberg}\label{s1003a01-step-55-darboux-gauge-ricci-compatibility-for-negative-scalar-goldberg}}

Source: \texttt{step55\_darboux\_ricci\_compatibility.md}.

Session S1003A01, Agent3. This note replaces the Step-54 interpretation
that a negative formal-dimension count by degree 15 by itself predicts a
new Goldberg obstruction. It records a universal Ricci-symbol relation,
derives its geometric form, and places it inside the existing
almost-Kähler identities.

\hypertarget{set-up-and-the-symbol-relation}{%
\subsection{1. Set-up and the symbol
relation}\label{set-up-and-the-symbol-relation}}

Work at a point of an almost-Kähler four-manifold in Darboux
coordinates, so the symplectic form is fixed and the compatible metric
varies through symmetric tensors satisfying \[
hJ+Jh=0.
\] In particular, \(\operatorname{tr}h=0\). For a real covector \(\xi\),
the flat Ricci symbol is, up to the common sign convention, \[
\sigma_{\mathrm{Ric}}(\xi)h
=\frac12\{\xi\otimes h\xi+h\xi\otimes\xi-|\xi|^2h\}.
\] Since \(h(J\xi,J\xi)=-h(\xi,\xi)\), direct contraction gives \[
\boxed{
(\sigma_{\mathrm{Ric}}(\xi)h)(J\xi,J\xi)
-\frac12|\xi|^2\operatorname{tr}(\sigma_{\mathrm{Ric}}(\xi)h)=0.}
\tag{1.1}
\] This is independent of the contracted-Bianchi symbol. Its polynomial
prolongations contribute \(n_{N-4}\) relations at Taylor degree \(N\),
which changes the Step-54 free count to \[
f_{\mathrm{corr}}(N)=N^2+8N-5\qquad(N\ge4).
\] Hence the earlier negative count at \(N=15\) is explained by an
omitted universal symbol relation. It gave no evidence of a
negative-scalar-specific condition and no finite bound for nonlinear
completion.

\hypertarget{coordinate-free-fourth-order-operator}{%
\subsection{2. Coordinate-free fourth-order
operator}\label{coordinate-free-fourth-order-operator}}

For a symmetric two-tensor \(S\), define the fourth-order principal
contraction \[
\mathcal C_J(S)
:=J_i{}^aJ_j{}^b\nabla^i\nabla^jS_{ab}
-\frac12L(\operatorname{tr}S),
\qquad L:=\nabla^i\nabla_i.
\tag{2.1}
\] Freezing \(J\) and replacing \(\nabla\) by \(\xi\), its symbol is
exactly the scalar in (1.1). In AAD's nonnegative convention
\(\Delta=-L\), the second term is \(+\frac12\Delta\operatorname{tr}S\).
Thus for a compatible metric the fourth derivatives cancel in
\(\mathcal C_J(\operatorname{Ric})\); what remains is a lower-jet
relation involving derivatives of \(J\), curvature, and the lower-order
nonlinearities of Ricci.

The first term has a useful exact divergence expansion. Put \[
D^a:=J_i{}^a\nabla^i=-J^a{}_i\nabla^i,
\qquad \nabla^iJ_i{}^a=0,
\tag{2.2}
\] where the last equality is the almost-Kähler identity
\(\delta\omega=0\). Then \[
\begin{aligned}
J_i{}^aJ_j{}^b\nabla^i\nabla^jS_{ab}
={}&\nabla^i\!\left(J_i{}^aJ_j{}^b\nabla^jS_{ab}\right)
-J_i{}^a(\nabla^iJ_j{}^b)\nabla^jS_{ab} \\
={}&\nabla^i\!\left(J_i{}^aJ_j{}^b\nabla^jS_{ab}\right)
-\langle\nabla J,\nabla S\rangle_{J},
\end{aligned}
\tag{2.3}
\] with the displayed contraction defining the final shorthand.
Therefore on a closed manifold \[
\int_M J_i{}^aJ_j{}^b\nabla^i\nabla^jS_{ab}\,d\mu
=-\int_M J_i{}^a(\nabla^iJ_j{}^b)\nabla^jS_{ab}\,d\mu.
\tag{2.4}
\] For an Einstein tensor \(S=\operatorname{Ric}=\lambda g\), both sides
vanish identically because \(\nabla S=0\) and
\(\operatorname{tr}S=4\lambda\) is constant. Consequently the unreduced
contraction (2.1), evaluated only after imposing Einstein, is a
differential consequence of \(\nabla\operatorname{Ric}=0\) and by itself
cannot distinguish \(\lambda<0\) from the Kähler-Einstein and Ricci-flat
Tod controls. This does \textbf{not} prove that the lower-order equation
obtained after eliminating fourth derivatives is redundant among a
truncated lower-jet system; such a resolved differential consequence can
still be an integrability relation and may contain sign-sensitive terms.

This does not make the Darboux symbol relation useless. It identifies a
real Spencer compatibility that must be included before any symbol
count, and it shows exactly where information can enter after the
Einstein equation: through the \emph{resolved lower-order expression}
for the same identity in the variables \(J,\nabla J\), not through a
bare Ricci projection. Such a resolved expression is one component of
the Ricci identities for \(\nabla\omega\) and the differential Bianchi
identities.

\hypertarget{identification-inside-the-aad-curvature-system}{%
\subsection{3. Identification inside the AAD curvature
system}\label{identification-inside-the-aad-curvature-system}}

Apostolov--Armstrong--Drăghici use \[
\nabla\Omega=a\otimes\phi-Ja\otimes J\phi,
\qquad |\nabla\Omega|^2=4|a|^2,
\tag{3.1}
\] and derive the Ricci identities \[
da-Ja\wedge b=-R(J\phi),\qquad
d(Ja)+a\wedge b=-R(\phi),
\tag{3.2}
\] together with \[
db=a\wedge Ja-R(\Omega).
\tag{3.3}
\] These equations determine the self-dual curvature and Ricci tensor
from the lower jets of \(J\). Their differentiated system is the
intrinsic form of the Darboux-gauge compatibility mechanism.

For a general almost-Kähler metric their differential Bianchi equation
is (AAD Lemma 4, equation (33)) \[
\begin{aligned}
0={}&-\tfrac14(d^J(\kappa-s))_Z\Omega
+\tfrac16(d(\kappa-s)\wedge Z^\flat)^+
+(\delta\rho_*'')_Z\Omega \\
&+\tfrac\kappa4\nabla_{JZ}\Omega
-\nabla_{\rho_0(Z)}\Omega
+\nabla_{\rho_*''(Z)}\Omega
+\nabla_{JZ}\rho_*''
+2(\delta W_3^+)_Z+A_Z^+.
\end{aligned}
\tag{3.4}
\] On an Einstein almost-Kähler metric, \(s\) is constant, \(\rho_0=0\),
\(A=0\), while \[
\kappa-s=\frac32q,
\qquad q=|\nabla\Omega|^2,
\qquad \rho_*''\equiv\eta .
\tag{3.5}
\] Thus the surviving equation is \[
\boxed{
0=-\frac38(d^Jq)_Z\Omega
+\frac14(dq\wedge Z^\flat)^+
+(\delta\eta)_Z\Omega
+\frac\kappa4\nabla_{JZ}\Omega
+\nabla_{\eta(Z)}\Omega
+\nabla_{JZ}\eta
+2(\delta W_3^+)_Z .}
\tag{3.6}
\] This is the substantive lower-jet compatibility absent from the
scalar symbol projection. It retains the transverse canonical curvature
\(\eta\), the derivative of \(W_3^+\), and their coupling to torsion. It
is not a Yang--Mills equation for the canonical connection and does not
permit treating \(\eta\) as freely prescribed forcing.

AAD's general Weitzenböck formula (their equation (47)) specializes
under Einstein to \[
\begin{aligned}
(\nabla^*\nabla(\nabla\Omega))_X={}&-\frac{3s}{4}\nabla_X\Omega
+\left[\frac13d(\kappa-s)_X
-2\langle(\delta W_3^+)_{JX},\Omega\rangle\right]\Omega \\
&+2(\nabla_X\eta)''+2\nabla_{J\eta(X)}\Omega,
\end{aligned}
\tag{3.7}
\] where the other terms vanish because \(\operatorname{Ric}_0=0\) and
\(A=0\). Using (3.5), \[
\boxed{
(\nabla^*\nabla(\nabla\Omega))_X
+\frac{3s}{4}\nabla_X\Omega
=\left[\frac12dq_X
-2\langle(\delta W_3^+)_{JX},\Omega\rangle\right]\Omega
+2(\nabla_X\eta)''+2\nabla_{J\eta(X)}\Omega .}
\tag{3.8}
\] The sign \(3s/4<0\) explains why a naive integration does not force
torsion to vanish, but it is not evidence for local freedom. Equations
(3.6) and (3.8) constrain \(dq,\nabla\eta,\delta W_3^+\), and
\(\nabla\Omega\) together. The \(\Omega\)-component of (3.6) removes the
explicit \(dq\) and \(\delta W_3^+\) forcing in (3.8): \[
\frac12dq_X-2\langle(\delta W_3^+)_{JX},\Omega\rangle
=2(\delta\eta)_{JX}+\langle\eta,\nabla_X\Omega\rangle,
\tag{3.9}
\] so the Helmholtz equation can be written with \(\eta,\nabla\eta\) as
its only explicit sources. For Einstein metrics, the third Gray
condition is \(\eta=0\) alone; \(W_3^+=0\) is not additionally required.
Then AAD's closed Helmholtz equation follows, and their local
classification---not the rough-Laplacian equation alone---forces any
strict solution in this branch into the Ricci-flat NPT/Tod class. Hence
it supplies no strictly almost-Kähler negative-Einstein germ.

\hypertarget{consequences-for-the-jet-computation}{%
\subsection{4. Consequences for the jet
computation}\label{consequences-for-the-jet-computation}}

\begin{enumerate}
\def\labelenumi{\arabic{enumi}.}
\tightlist
\item
  The original all-at-once degree-5 computations are usable only as
  existence of independently optimized truncated jets. They do
  \textbf{not} form compatible formal germs. The preserved D5 runs have
  residual squares between \(1.1\times10^{-27}\) and
  \(5.7\times10^{-27}\); the full Jacobian ranks are 1035 for KE and
  1036 for RF/NEG. As at degree 4, the KE one-row difference is expected
  from the singular linearization of the imposed scalar condition
  \(q(0)=0\). Moreover the D5 \texttt{KE} jets are not the exact
  Kähler-Einstein control: their computed \(q\) is numerical zero, but
  their reported \(|\eta|=0.156,0.298\) and generic \(W^+\) show that
  imposing only the pointwise scalar equation \(q(0)=0\) does not impose
  a Kähler germ. Only the explicit exact \(H^2\times H^2\) jet is a KE
  control.
\item
  D6/D7 empty or incomplete outputs resulted from resource termination
  and carry no mathematical content.
\item
  The first fixed-lower-jet test (files
  \texttt{data/step54\_prolong/prolong\_D4\_*\_z0\_to\_D5.json})
  produced nonzero projection norms for every branch: \(1.567\) for the
  mislabeled \texttt{KE} D4 jet, \(0.669\) for a generic RF jet, and
  \(0.741\) for a NEG jet, with exact affine-linearity errors
  \(4.4\)--\(5.8\times10^{-14}\). Because even the controls fail and the
  source jets were independently optimized truncations rather than exact
  extendible germs, these numbers are \textbf{not} obstructions. They
  demonstrate that the present least-squares jets do not lie on a
  degree-five extendible branch and that sector-by-sector residual norms
  are basis-dependent (the total cokernel vector mixes E1/E2/E3 rows).
  The fixed-lower-jet method becomes meaningful only with an exact KE
  and exact Tod source and with its cokernel quotiented by (1.1),
  contracted Bianchi, closure, and algebraic compatibility identities.
\item
  A projection along (1.1) should vanish as a consequence of
  (3.2)--(3.4); it is a control on the implementation, not a
  negative-scalar obstruction.
\item
  The next local computation should use an intrinsic first-order
  prolongation of (3.2)--(3.3), (3.6), and (3.8), or Darboux variables
  with the universal identities explicitly removed. It should seek a
  residual relation not generated by these equations and compare an
  exact KE germ and an exact Tod germ against negative jets at several
  \(q/|\lambda|\). Re-running the unconstrained degree ladder is not
  justified.
\end{enumerate}

\hypertarget{global-implication}{%
\subsection{5. Global implication}\label{global-implication}}

The local scalar candidate has collapsed to a known compatibility
identity rather than an integrability theorem. Therefore the
highest-value Goldberg route is the coupled canonical-connection system:
\[
\eta=-\tfrac12\operatorname{div}^1(A),
\quad\text{AAD (3.6),}
\quad\text{AAD (3.8),}
\quad
\beta=iF_{\mathcal A}-\tfrac s4\omega,
\quad
\int q^2/16\ge8\pi^2\|c_1^{-,h}\|^2.
\] A successful rigidity proof must obtain a global coercive relation
that controls the transverse terms in (3.6)--(3.8), possibly from the
actual canonical/Seiberg--Witten field, rather than discarding them.
Conversely, a local counterexample construction must solve this coupled
system with \(\eta\ne0\) for \(s<0\), not merely the scalar Ricci
compatibility.

\hypertarget{evidence-status}{%
\subsection{6. Evidence status}\label{evidence-status}}

\begin{itemize}
\tightlist
\item
  Equations (1.1), (2.3), (3.6), and (3.8) are direct
  algebraic/specialization derivations from the displayed definitions
  and the cited AAD equations.
\item
  The identification of the entire nonlinear Darboux remainder with the
  AAD system is structural; a component-by-component coordinate
  expansion has not been performed.
\item
  No negative-scalar strictly almost-Kähler Einstein germ, compact
  example, or integrability proof is obtained here.
\end{itemize}

\hypertarget{global-identities-extracted-from-the-coupled-equations}{%
\subsection{7. Global identities extracted from the coupled
equations}\label{global-identities-extracted-from-the-coupled-equations}}

An independent coefficient check
(\texttt{trace/ttcs/results/step55\_ttcs01\_verifier1\_integration.md})
confirms (3.6), (3.8), the identification
\(\eta=\rho_*''=W^+(\Omega)_{\Omega^\perp}\), and the form of (3.9).
Pairing (3.8) with \(T=\nabla\Omega\) and integrating gives \[
\boxed{
\int_M|\nabla^2\Omega|^2+\frac{3s}{4}\int_Mq
=4\int_M|\eta|^2.}
\tag{7.1}
\] The \(\Omega\)-valued forcing is orthogonal to \(T\), and the
\(J\eta\)-directional term contracts a skew tensor with the symmetric
tensor \(\langle T_X,T_Y\rangle\), hence vanishes. AAD's identity \[
\nabla^*\nabla\Omega=2\eta+\frac q2\Omega
\] then supplies the right-hand side. For \(s<0\), (7.1) reads \[
\int|\nabla^2\Omega|^2=\frac{3|s|}{4}\int q+4\int|\eta|^2,
\tag{7.2}
\] so the Helmholtz equation measures, but does not suppress, torsion.

The canonical closed form \(\beta=iF_{\mathcal A}-(s/4)\Omega\)
satisfies \(\beta\wedge\beta=|\eta|^2d\mu\), so (7.1) also gives the
topological identity \[
\boxed{
\int|\nabla^2\Omega|^2+\frac{3s}{4}\int q
=4[\beta]^2,
\qquad [\beta]=2\pi c_1-\frac{s}{4}[\Omega].}
\tag{7.3}
\] This connects the full second derivative of the almost-Kähler field
to the same canonical class appearing in the Hodge inequality; it
remains non-coercive in the negative-scalar regime.

The transverse Bianchi component provides the additional exact equation
\[
0=\frac14(dq\wedge Z^\flat)''+\frac\kappa4T_{JZ}
+2(\delta W_3^+)_Z''+\nabla_{\eta(Z)}\Omega+(\nabla_{JZ}\eta)'',
\tag{7.4}
\] which is the correct place to seek coercivity. It retains the
derivative of \(W_3^+\), rather than replacing it by a scalar norm. In
particular, pairing (7.4) with \(T_{JZ}\), integrating, and using the
established almost-Kähler balance yields \[
\boxed{
\int_M\sum_i\langle(\delta W_3^+)_{e_i},T_{Je_i}\rangle
=2\int_M|W_3^+|^2.}
\tag{7.5}
\] Thus the transverse Weyl component has a signed global cross-pairing
with the Nijenhuis derivative. The missing step for Goldberg is a
canonical-connection estimate controlling the left side of (7.5),
together with the \(\eta\) terms in (7.4), strongly enough to contradict
the negative-scalar balance unless \(T=0\). A pointwise or \(L^2\) bound
that discards these coupled derivatives will reproduce the already
non-coercive identities.

\hypertarget{s1003a02-step-53-andersons-idealized-minimizers-as-the-entrance-attainment-transfer-under-smooth-bubbling-first-order-bubble-exclusion-reduction-of-the-non-bach-flat-optimum-problem}{%
\chapter{S1003A02: Step 53 --- Anderson's idealized minimizers as the
entrance; attainment transfer under smooth bubbling; first-order bubble
exclusion; reduction of the non-Bach-flat optimum
problem}\label{s1003a02-step-53-andersons-idealized-minimizers-as-the-entrance-attainment-transfer-under-smooth-bubbling-first-order-bubble-exclusion-reduction-of-the-non-bach-flat-optimum-problem}}

Source:
\texttt{step53\_anderson\_entrance\_bubble\_exclusion\_attainment.md}.

Session S1003A02, Agent2, Execute. Conventions:
\(P(g)=\int|\operatorname{Ric}^\circ|^2\), operator norms on
\(\Lambda^2\); \(\mathbf R^\pm=W^\pm+\tfrac s{12}I\) on \(\Lambda^\pm\);
\(\mathcal K(g)=\int|\mathrm{Rm}_{\rm op}|^2=8\pi^2\chi+P\) on closed
\(M\) (five\_research\_problems.tex eq. (fullP)). Sources read this
step: the 2026-10-07 supplements
(\texttt{Skills/geometric-analysis/research/supplements/20261007-einstein-gaps/five\_research\_problems.tex}
§§1--3, \texttt{web\_research\_collected.tex} Ch. 5 and Ch. 7),
Anderson, \emph{Canonical metrics on 3-manifolds and 4-manifolds}, Asian
J. Math. 10 (2006) (Knowledge object
\texttt{Canonical-Metrics-on-3-Manifolds-and-4-Manifolds-\/-8d02c53d4380-\/-0b5d0f7e0ac4},
Thm 4.1, Cor 4.4, Lemma 4.5, Prop 4.6, (4.32)--(4.33), Cor 4.8, Thm 4.10
and its proof, Remark 4.12, (4.65)), Tian--Viaclovsky, \emph{Bach-flat
ALE metrics} Thm 1.1/1.3 (Knowledge object
\texttt{Bach-flat-asymptotically-locally-Euclidean-metrics-\/-3c3dd988289c}),
and our Theorem F
(\texttt{step20\_burns\_response\_and\_neck\_quantization.md} §1).

This note supersedes the Step-52 framing in which ``step 0'' (an
own-scale regularization of \emph{arbitrary} minimizing sequences) was
the open entrance. For existence and structure questions one only needs
\textbf{some} minimizing sequence with controlled structure, and
Anderson's Theorem 4.10 is an existing theorem providing one. The
Step-52 residual analysis remains correct as a statement about arbitrary
sequences; it is no longer on the critical path.

\begin{center}\rule{0.5\linewidth}{0.5pt}\end{center}

\hypertarget{the-entrance-andersons-perturbed-minimizers}{%
\subsection{1. The entrance: Anderson's perturbed
minimizers}\label{the-entrance-andersons-perturbed-minimizers}}

Anderson perturbs exactly our functional: \[
\mathcal Z^{2p}_\varepsilon(g)=\int\Big(|\operatorname{Ric}|^2-\tfrac14s^2\Big)+\varepsilon\int(1+|\operatorname{Ric}|^2)^p
=P(g)+\varepsilon\int(1+|\operatorname{Ric}|^2)^p ,\qquad p>1 \text{ close to }1,
\] on unit-volume metrics ((4.32)). A bound on
\(\mathcal Z^{2p}_\varepsilon\) bounds
\(\int|\operatorname{Ric}|^{2p}\), so his Theorem 4.1 (convergence under
\(L^{2p}\) Ricci bounds) produces minimizing pairs
\((\Omega_\varepsilon,g_\varepsilon)\) (Cor 4.4/4.8); local minimizers
are \(L^{2,4}\) (Lemma 4.5) and \(C^\infty\) (Prop 4.6, via the
decomposition \(h=\delta^*X+\phi g+k\) and the elliptic Euler--Lagrange
system (4.33)). Since \(P\) is scale invariant and every fixed smooth
\(g\) has \(\varepsilon\int(1+|\operatorname{Ric}_g|^2)^p\to0\), these
are minimizing for \(P\) at cost \(o(1)\): this is the Sacks--Uhlenbeck
mechanism, and it replaces the flow-based entrance of Steps 50--52. The
sixth-order objection of Step 19 does not apply: the perturbation is
zeroth order in curvature, and the Euler--Lagrange equation stays fourth
order with uniform elliptic estimates on each summand as
\(\varepsilon\to0\) once \(C^\alpha\) harmonic control holds ((4.49),
proved by contradiction from the \(L^q\)/\(L^2\) curvature radii).

\textbf{Theorem 4.10 (Anderson), specialized to \(P\).} For every closed
oriented 4-manifold \(M\) there is a minimizing sequence \(g_i\) for
\(P\) (unit volume) converging in pointed Gromov--Hausdorff sense to one
of:

\begin{itemize}
\tightlist
\item
  \textbf{(I)} a compact, possibly reducible orbifold \((V,g_0)\),
  \(\operatorname{vol}V=1\), \(g_0\) smooth on \(V_0\) and \(C^0\)
  across the singular points (which may have \(\Gamma=\{e\}\)), together
  with ALE spaces \((N_k,g^k_\infty)\) associated with the singular
  points, with \textbf{energy equality}
  \(\mathcal F(g_0)+\sum_k\mathcal F(N_k)=\inf\mathcal F\);
\item
  \textbf{(II)} a maximal orbifold domain \(\Omega\subset\subset M\)
  (possibly empty) with the inequality \(\le\) in place of equality and
  \(F\)-structure collapse outside a compact set.
\end{itemize}

In both cases \(\nabla\mathcal F(g_0)=0\) and \(g_0\) minimizes
\(\mathcal F\) among compact perturbations. In case (I), \(M\) is the
generalized connected sum of \(V_0\) with the ALE spaces along spherical
space forms ((4.65)).

\textbf{Lemma 53.0 (the energy equality is the \(P\)-ledger).} In case
(I), \[
p_*(M)=P(V,g_0)+\sum_kP(N_k).
\] \emph{Proof.} For an orbifold \(V\) and a (possibly orbifold) ALE
space \(N\) with group \(\Gamma\) at infinity, Chern--Gauss--Bonnet
reads \(\int|\mathrm{Rm}_{\rm op}|^2=8\pi^2\chi_{\rm orb}+P\), with
\(\chi_{\rm orb}(V)=\chi(V)-\sum_j(1-|\Gamma_j|^{-1})\) and
\(\chi_{\rm orb}(N)=\chi(N)-|\Gamma|^{-1}-\sum(\text{interior defects})\).
Gluing along \(S^3/\Gamma\) (Euler characteristic 0) gives
\(\chi(M)=\chi(V)-\#\{\text{sing}\}+\sum\chi(N_k)\) at one level, and
the \(|\Gamma|^{-1}\) terms cancel pairwise; by induction on the depth
of the bubble tree the Euler terms sum to \(8\pi^2\chi(M)\). Subtracting
from Anderson's equality for \(\mathcal R^2=\mathcal K\) gives the
claim. \(\square\)

The same telescoping is the Theorem-G characteristic ledger at the level
of Euler terms; the signature/eta ledger (Step 52, Prop 52.10) is the
companion statement for \(T_\pm\).

\textbf{What is used from the source, and where its proof is delicate.}
(a) The local-minimality of the Cor-4.4 limit pair among compact
perturbations is asserted from a weak \(L^{2,2p}\) limit; a complete
proof needs strong convergence (uniform convexity of
\(\int(1+|\operatorname{Ric}|^2)^p\) in \(\partial^2g\), as in
Sacks--Uhlenbeck). (b) The case-(I) energy equality (4.43) is derived by
citation to the orbifold compactness/no-neck theory
({[}8{]},{[}9{]},{[}11{]} there); Remark 4.12 lists the corresponding
case-(II) equality as a conjecture. (c) Removability of \(\Gamma=\{e\}\)
points and smoothness of \(g_0\) in orbifold charts use finite-energy
removable-singularity theory for critical metrics (Tian--Viaclovsky). We
use Theorem 4.10 as stated, with these dependencies recorded; every
result below is stated as conditional on it.

\begin{center}\rule{0.5\linewidth}{0.5pt}\end{center}

\hypertarget{reassembly-and-gluing}{%
\subsection{2. Reassembly and gluing}\label{reassembly-and-gluing}}

\textbf{Lemma 53.R (reassembly).} Let a case-(I) configuration be given,
and suppose each ALE space \(N_k\) has an end of order \(\tau_k>0\). For
any choice of gluing scales \(\lambda_k\to0\) (each tree glued
recursively, scales separated) the glued metrics \(g_\lambda\) live on
\(M\) and satisfy \(P(g_\lambda)\to P(g_0)+\sum_kP(N_k)=p_*(M)\). The
same holds if any smooth closed component of \(V\) carries an arbitrary
fixed smooth metric \(h\) in place of \(g_0\), with \(P(g_0)\) replaced
by \(P(h)\) for that component.

\emph{Proof sketch.} Topology: (4.65). Energy: on a cut annulus
\(\{r\sim r_0\}\) at a gluing point the glued metric is \(g_0+O(r^2)\)
on the bulk side and \(g_N^{\lambda}=\delta+O(\lambda^\tau r^{-\tau})\)
on the bubble side; with a cutoff the curvature error is
\(O(1+\lambda^\tau r_0^{-\tau-2})\) on volume \(O(r_0^4)\), so the cut
cost is \(O(r_0^4+\lambda^{2\tau}r_0^{-2\tau})\to0\) for
\(\lambda\ll r_0\to0\); outside the annuli the metrics are the
unperturbed pieces up to errors that vanish with \(r_0,\lambda/r_0\).
\(\square\)

\textbf{Theorem 53.F′ (first-order response of an anti-self-dual
scalar-flat AE bubble; general bulk).} Let \((N,g_N)\) be a complete AE
(\(\Gamma=\{e\}\)) manifold with \(W^+_{g_N}\equiv0\),
\(s_{g_N}\equiv0\), and an AE chart \(g_N=\delta+h_2+O(r^{-2-\delta})\)
(with derivatives) where \(h_2\) is homogeneous of degree \(-2\). Define
its chiral flux \[
\Phi^{++}(N)=-\int P_+\mathcal R^{\rm lin}_\delta(\chi h_2)P_+\,dx\quad(\text{a symmetric endomorphism of }\Lambda^+_0),
\] which is independent of the cutoff \(\chi\) because
\(P_+\mathcal R^{\rm lin}(h_2)P_+=0\) and \(s^{\rm lin}(h_2)=0\)
pointwise (the degree \(-4\) parts of \(W^+_{g_N}=0\), \(s_{g_N}=0\)).
Let \((V,g)\) be closed, \(p\in V\), and let \(U\in SO(4)\) identify
\(T_pV\) with \(\mathbb R^4\); \(U\) acts on \(\Lambda^+\) through
\(SO(3)\). Gluing \(N\) at \(p\) at scale \(\lambda\)
(\(g_N\mapsto\lambda^2g_N\)) with cutoff radius \(r_0=\lambda^{2/5}\)
gives, on \(V\#\widehat N\), \[
\boxed{\;P(\hat g_\lambda)=P(g)+P(N)-8\lambda^2\big\langle\mathbf R^+_p,\;U\Phi^{++}(N)U^{T}\big\rangle+O(\lambda^{12/5}).\;}
\] For Burns (\(\lambda^2=x_0\)) \(\Phi^{++}=\pi^2w_1w_1^T\) and this is
Theorem F.

\emph{Proof.} Identical to Theorem F (§1 of the Step-20 note) with two
observations. (i) The exact topological reduction
\(P(\hat g)-P(g)=4\Delta\!\int|W^+|^2+\tfrac1{12}\Delta S-8\pi^2(T_+(\widehat N)-4)\),
and the same identity on \(N\) itself, give
\(P(\hat g)-P(g)-P(N)=4\Delta W^++\tfrac1{12}\Delta S-4\int_N|W^+_N|^2-\tfrac1{12}\int_Ns_N^2\),
and the last two terms vanish. (ii) In the core and intermediate regions
the mixed finite difference argument uses only
\(W^+(\lambda^2g_N)=s(\lambda^2g_N)=0\) identically and the decay
\(\lambda^2r^{-2}\) of the tail; the cut annulus produces
\(-2\lambda^2\langle W^+_p,U\Phi^{++}U^T\rangle\) and
\(-2\lambda^2s_pF_s\) with \(F_s=4\operatorname{tr}\Phi^{++}\), which
assemble to \(-8\lambda^2\langle\mathbf R^+_p,U\Phi^{++}U^T\rangle\).
\(\square\)

By von Neumann's trace inequality, with eigenvalues ordered
decreasingly, \[
\max_{U}\big\langle\mathbf R^+_p,U\Phi^{++}U^T\big\rangle=\sum_{i=1}^3\lambda_i(\mathbf R^+_p)\,\mu_i(\Phi^{++}(N))=: \mathfrak r_N(p).
\] For Burns, \(\mathfrak r_N(p)=\pi^2\lambda_{\max}(\mathbf R^+_p)\).
If \(W^+_N\not\equiv0\) the cross term
\(2\langle W^+_g,W^+(\lambda^2g_N)\rangle\) in the core is of order
\(\lambda^2\log(1/\lambda)\) whenever the \(r^{-4}\) tail of \(W^+_N\)
has nonzero spherical mean; the response of non-ASD bubbles must be
computed separately.

\begin{center}\rule{0.5\linewidth}{0.5pt}\end{center}

\hypertarget{bubble-exclusion-and-attainment-transfer}{%
\subsection{3. Bubble exclusion and attainment
transfer}\label{bubble-exclusion-and-attainment-transfer}}

Throughout this section a case-(I) configuration is fixed, and the
hypotheses of Lemma 53.R hold.

\textbf{Theorem 53.B (first-order bubble exclusion).} Suppose the
configuration contains a bubble tree at a smooth point whose outermost
bubble \(N\) satisfies the hypotheses of Theorem 53.F′. Then for every
component \(V_a\) of \(V\) and every regular point \(p\in V_a\), \[
\mathfrak r_N(p)=\sum_i\lambda_i(\mathbf R^+_{g_0}(p))\,\mu_i(\Phi^{++}(N))\le0 .
\] If \(N\) is Burns, or more generally if
\(\widehat N\cong\overline{\mathbb{CP}}^2\) and \(N\) is ASD scalar-flat
with rank-one flux, then \(\mathbf R^+_{g_0}\le0\) on every component of
the bulk.

\emph{Proof.} Let \(M'\) be the reassembled configuration without that
tree; \(M'\) is connected and \(M\cong M'\#\widehat T\), \(\widehat T\)
the compactified tree. The connected sum with \(\widehat T\) at any
point and through any orientation-preserving frame is well defined up to
diffeomorphism (Palais), so gluing the tree at \(p\) with frame \(U\)
still gives \(M\). Glue all other trees by Lemma 53.R at scales tending
to 0 first; then glue the tree at \(p\) with optimal \(U\) and a small
fixed \(\lambda\). By Lemma 53.0, Lemma 53.R and Theorem 53.F′,
\(P(\hat g)=p_*(M)-8\lambda^2\mathfrak r_N(p)+O(\lambda^{12/5})+o(1)\).
If \(\mathfrak r_N(p)>0\), fixing \(\lambda\) small and then sending the
other scales to zero gives \(P(\hat g)<p_*(M)\). \(\square\)

The point is that a minimizing configuration cannot leave first-order
interaction energy unused: a bubble that has concentrated somewhere may
be glued back anywhere, and the energy equality makes the configuration
value the competitor's zeroth-order value.

\textbf{Corollary 53.C (curvature of the bulk when
\(\overline{\mathbb{CP}}^2\)-bubbling occurs).} Suppose a Burns-type
bubble occurs. Since \(g_0\) is \(P\)-critical,
\(\operatorname{tr}\nabla P=-\tfrac12\Delta s\) forces \(s_0\) constant
on each component. Then on every component:

\begin{enumerate}
\def\labelenumi{\arabic{enumi}.}
\tightlist
\item
  \(s_0\le0\) (since \(\operatorname{tr}\mathbf R^+=s/4\));
\item
  if \(s_0=0\), then \(W^+\equiv0\): the component is scalar-flat
  anti-self-dual;
\item
  if \(s_0<0\), then \(\lambda_{\max}(W^+)\le|s_0|/12\), hence
  \(|W^+|^2\le s_0^2/24\) pointwise, with equality exactly at points
  where \(W^+\) has the Kähler spectrum \((|s|/12,|s|/12,-|s|/6)\);
  consequently
  \(P(g_0|_{V_a})\le\tfrac14\int_{V_a}s_0^2-8\pi^2T_+^{\rm orb}(V_a)\);
\item
  no component contains an open set on which \(g_0\) is complex
  hyperbolic. (On such a ball the Chapter-5 cell of the 2026-10-07
  collection --- exact CH outside a fixed collar,
  \(P/(8\pi^2)=1-9\epsilon^2\log(1/\epsilon)+O_R(\epsilon^2)\) ---
  replaces the bubble. Every AE bubble with
  \(\widehat N\cong\overline{\mathbb{CP}}^2\) costs \(P(N)\ge8\pi^2\) by
  the AE cost formula
  \(P_N=4\int|W^+|^2+\tfrac1{12}\int s^2+8\pi^2b_2+32\pi^2b_1\). Hence
  the replacement lowers the total below \(p_*(M)\).)
\end{enumerate}

The Kähler--Einstein bulk is exactly the borderline of item 3:
\(\mathbf R^+=\operatorname{diag}(s/4,0,0)\le0\) with a 2-sphere of
kernel directions. This is consistent with the fixed-complex-structure
recovery sequence (KE canonical model plus Burns bubbles,
\(\inf_{\rm K\ddot ahler}P=8\pi^2k\)). Item 4 says that, for ball
quotients, the unrestricted minimizing configuration is never this one.

\textbf{Theorem 53.D (attainment transfer under smooth bubbling).}
Suppose the configuration has no orbifold points with
\(\Gamma\ne\{e\}\). Then every component \(V_a\) of the bulk is a smooth
closed 4-manifold, \(g_0|_{V_a}\) is smooth, and \[
P(g_0|_{V_a})=p_*(V_a)\quad\text{for every }a,\qquad p_*(M)=\sum_ap_*(V_a)+\sum_kP(N_k).
\] Thus each bulk component attains its own infimum. In particular
\(p_*(M)\) is attained on \(M\) iff some configuration is trivial (one
component \(V\cong M\), no bubbles).

\emph{Proof.} Smoothness at \(\Gamma=\{e\}\) points: removable
singularities. Lemma 53.R with arbitrary smooth metrics \(h_a\) on the
components gives \(p_*(M)\le\sum_aP(h_a)+\sum_kP(N_k)\), hence, by Lemma
53.0, \(\sum_aP(g_0|_{V_a})\le\sum_ap_*(V_a)\). The reverse holds
termwise. \(\square\)

\begin{center}\rule{0.5\linewidth}{0.5pt}\end{center}

\hypertarget{consequence-for-the-non-bach-flat-global-optimum-problem-3-and-for-problem-2}{%
\subsection{4. Consequence for the non-Bach-flat global optimum (Problem
3) and for Problem
2}\label{consequence-for-the-non-bach-flat-global-optimum-problem-3-and-for-problem-2}}

\textbf{Lemma 53.Y (Yamabe sign of bubble compactifications).} If \(N\)
is AE with \(s\equiv0\) and an end of order \(\tau>1\), then its
one-point compactification \(\widehat N\) has \(Y(\widehat N)>0\).
(Multi-ended: compactify each end.)

\emph{Proof.} In inverted coordinates \(\hat g=|y|^4g\) extends to a
\(C^{1,\tau-1}\) metric with \(s_{\hat g}\in L^q\), \(q>2\). The
function \(u=|y|^{-2}(1+o(1))\) satisfies \(g=u^2\hat g\), so
\(L_{\hat g}u=0\) away from the point and
\(L_{\hat g}u=24\pi^2\delta_0\) as a distribution. If the first
eigenfunction \(\phi>0\) of \(L_{\hat g}=-6\Delta+s\) had eigenvalue
\(\lambda_1\le0\), then
\(24\pi^2\phi(0)=\langle L u,\phi\rangle=\lambda_1\int u\phi\le0\), a
contradiction. Hence \(Y([\hat g])>0\); smooth approximation in
\(W^{2,q}\) keeps \(\lambda_1>0\). \(\square\)

\textbf{Theorem 53.E (reduction of Problem 3).} Let \(M\) be closed with
\(Y(M)<0\) (e.g.~any \(M\) with a monopole class of positive square,
such as \(X\#k\overline{\mathbb{CP}}^2\), \(X\) minimal of general
type). Suppose \(P\) on \(M\) has a case-(I) configuration without
\(\Gamma\ne\{e\}\) points whose bubbles have ends of order \(\tau>1\).
Then:

\begin{enumerate}
\def\labelenumi{\arabic{enumi}.}
\tightlist
\item
  some bulk component \(V_a\) has \(Y(V_a)<0\); in fact
  \(Y(M)^2\le\sum_{a:\,Y(V_a)<0}Y(V_a)^2\);
\item
  on every such component \(g_0\) is a smooth global minimizer of \(P\)
  which is not scalar-flat;
\item
  if such a \(V_a\) carries no Einstein metric, \(g_0|_{V_a}\) is a
  non-Bach-flat global optimum of the total-curvature functional on
  \(V_a\) (five\_research\_problems.tex Prop. nonbach);
\item
  otherwise \(g_0|_{V_a}\) is Einstein.
\end{enumerate}

\emph{Proof.} \(M\) is the connected sum of the closed pieces \(V_a\),
\(\widehat N_k\) along the gluing graph, together with one
\(S^1\times S^3\) per independent loop. Pieces of the second and third
kind have positive Yamabe invariant (Lemma 53.Y, and \(S^1\times S^3\)).
By Kobayashi's connected-sum inequalities, if every \(V_a\) had
\(Y\ge0\) then \(Y(M)\ge0\), and the quantitative form gives item 1.
Item 2 is Theorem 53.D together with the fact that \(Y<0\) excludes
scalar-flat metrics. Items 3--4 follow from Prop. nonbach and
\(p_*=0\iff\) Einstein attainment. \(\square\)

So \textbf{smooth bubbling is no obstruction to Problem 3}. It moves
attainment to a bulk manifold that inherits negative Yamabe invariant.
The genuine obstructions to a smooth non-Bach-flat optimum are only:

\begin{itemize}
\item
  \begin{enumerate}
  \def\labelenumi{(\alph{enumi})}
  \tightlist
  \item
    collapse (case II);
  \end{enumerate}
\item
  \begin{enumerate}
  \def\labelenumi{(\alph{enumi})}
  \setcounter{enumi}{1}
  \tightlist
  \item
    orbifold points with nontrivial \(\Gamma\), i.e.~Kronheimer-type
    absorption (Step 52);
  \end{enumerate}
\item
  \begin{enumerate}
  \def\labelenumi{(\alph{enumi})}
  \setcounter{enumi}{2}
  \tightlist
  \item
    Einstein bulk components.
  \end{enumerate}
\end{itemize}

For \(M=X\#k\overline{\mathbb{CP}}^2\) the third alternative is itself
constrained:

\begin{itemize}
\tightlist
\item
  An Einstein bulk \(V_a\) has \(Y(V_a)<0\), and \(Y(V_a)\le Y(M)\) if
  it is the only negative component.
\item
  It must satisfy Hitchin--Thorpe and, where it carries a monopole
  class, LeBrun's and N1's obstructions. Which summand carries the basic
  classes needs a separate Seiberg--Witten connected-sum argument; that
  argument is not supplied here.
\item
  If a Burns bubble is present, it must satisfy
  \(\lambda_{\max}(W^+)\le|s|/12\) (Cor 53.C).
\item
  If \(X\) is a ball quotient and \(V_a\cong X\), Cor 53.C(4) excludes
  it whenever a \(\overline{\mathbb{CP}}^2\)-bubble is present.
\end{itemize}

\textbf{Consequence for Problem 2 (exact value of
\(p_*(X\#k\overline{\mathbb{CP}}^2)\)).} Combine Theorem 53.D with
Theorem 53.B. In any noncollapsed, \(\Gamma\)-free degenerate minimizing
configuration, \(p_*(M)=\sum_ap_*(V_a)+\sum_kP(N_k)\), where the bulk
minimizers lie in the cone \(\mathbf R^+\le0\) as soon as a
\(\overline{\mathbb{CP}}^2\)-bubble is present. The strict inequality
that rules such a configuration out is therefore
\(p_*(V\#\overline{\mathbb{CP}}^2)<p_*(V)+8\pi^2\) at a minimizer of
\(V\). It holds at first order whenever the minimizer has
\(\lambda_{\max}(\mathbf R^+)>0\) somewhere (Theorem F). It holds at
second order on CH regions (Chapter 5). It is open exactly for
minimizers with \(\mathbf R^+\le0\) and no CH region. That borderline
needs a general-background second-order (kernel-direction) response
formula, uniform in the background. This is the precise location where
the ``strictly below every degeneration cost implies smooth attainment''
mechanism must be completed.

\begin{center}\rule{0.5\linewidth}{0.5pt}\end{center}

\hypertarget{status}{%
\subsection{5. Status}\label{status}}

\begin{itemize}
\tightlist
\item
  \textbf{Proved here, conditional only on Anderson Thm 4.10 as stated
  plus standard gluing/removability:} Lemma 53.0, Theorem 53.F′, Theorem
  53.B, Cor 53.C, Theorem 53.D, Theorem 53.E. Lemma 53.Y uses order
  \(\tau>1\). Tian--Viaclovsky give \(\tau<2\) arbitrary for ASD
  scalar-flat or harmonic-curvature ALE; for general Bach-flat
  scalar-flat ends only order 0 is in their Thm 1.1, and higher order
  must be imported from later decay results.
\item
  \textbf{Withdrawn framing:} ``step 0 is the open entrance'' (Step 52
  §5) --- not needed for existence/structure.
\item
  \textbf{Open and decisive:}

  \begin{enumerate}
  \def\labelenumi{\arabic{enumi}.}
  \tightlist
  \item
    exclusion of collapse for \(Y(M)<0\);
  \item
    exclusion or pricing of nontrivial-\(\Gamma\) points (absorption);
  \item
    the second-order kernel-direction response in a general background
    with \(\mathbf R^+\le0\);
  \item
    responses of non-ASD bubbles.
  \end{enumerate}
\end{itemize}

\begin{center}\rule{0.5\linewidth}{0.5pt}\end{center}

\hypertarget{corrections-after-verification-step-54}{%
\subsection{6. Corrections after verification (Step
54)}\label{corrections-after-verification-step-54}}

The Step-53 verifier
(\texttt{trace/ttcs/results/step53\_ttcs01\_verifier1\_integration.md})
confirmed the Euler telescope, the half-flat response, one-attachment
exclusion, and attainment transfer, with the following corrections. They
are binding over §§1--4 above.

\begin{itemize}
\tightlist
\item
  Lemma 53.0: count cone branches at reducible vertices (truncated-piece
  proof).
\item
  53.B: one-attachment trees only. The flux must be positive; the
  positive-mass theorem gives
  \(\operatorname{tr}\Phi^{++}=3\pi^2m_{\rm ADM}>0\).
\item
  53.C: KE kernel directions form \(S^1\), not \(S^2\). Item 4 needs the
  CH-cell theorem with its collar-radius scope.
\item
  53.F′: the \(O(\lambda^{12/5})\) remainder needs the half-flat
  cancellation, and the homogeneous \(r^{-2}\) expansion is a
  hypothesis. The \(\lambda^2\log\) remark for non-ASD bubbles is
  deleted (zero spherical mean).
\item
  53.D: attainment ⇔ existence of a trivial generalized configuration.
\item
  53.Y: any order \(\tau>0\) with derivative control suffices.
\item
  Prop. nonbach is replaced by: Bach-flat \(+\) \(P\)-critical \(+\)
  constant \(s\neq0\) ⇒ Einstein.
\item
  Anderson's Prop 4.6 regularity and its uniformity under rescaling are
  soft points of the source.
\end{itemize}

The k=1 ball-quotient theorem built on these is in
\texttt{step54\_k1\_ball\_quotient\_attainment\_theorem.md}.

\hypertarget{s1003a02-step-54-the-k-1-ball-quotient-theorem-smooth-bubbling-is-topologically-excluded-and-the-minimizer-is-einstein-or-a-non-bach-flat-global-optimum}{%
\chapter{S1003A02: Step 54 --- The k = 1 ball-quotient theorem: smooth
bubbling is topologically excluded, and the minimizer is Einstein or a
non-Bach-flat global
optimum}\label{s1003a02-step-54-the-k-1-ball-quotient-theorem-smooth-bubbling-is-topologically-excluded-and-the-minimizer-is-einstein-or-a-non-bach-flat-global-optimum}}

Source: \texttt{step54\_k1\_ball\_quotient\_attainment\_theorem.md}.

Session S1003A02, Agent2, Execute. Builds on
\texttt{step53\_anderson\_entrance\_bubble\_exclusion\_attainment.md}
(as corrected by the Step-53 verifier,
\texttt{trace/ttcs/results/step53\_ttcs01\_verifier1\_integration.md},
whose corrections are listed in §0) and on the supplement
\texttt{web\_research\_collected.tex} Ch. 5 (analytic CH-collar cell).
Checks:
\texttt{workspace/scripts/step54\_charge\_and\_entropy\_checks.py} →
\texttt{workspace/data/step54\_charge\_and\_entropy\_checks.json}.

Conventions. \(P(g)=\int|\operatorname{Ric}^\circ|^2\),
\(p_*(M)=\inf P\), \(T_+=2\chi+3\tau\). For closed \(V\),
\(P=4\int|W^+|^2+\frac1{12}\int s^2-8\pi^2T_+(V)\).

\begin{center}\rule{0.5\linewidth}{0.5pt}\end{center}

\hypertarget{step-53-corrections-now-in-force}{%
\subsection{0. Step-53 corrections now in
force}\label{step-53-corrections-now-in-force}}

\begin{enumerate}
\def\labelenumi{\arabic{enumi}.}
\tightlist
\item
  Euler bookkeeping counts \textbf{cone branches} at reducible vertices;
  the truncated-piece proof of \(p_*(M)=P(V,g_0)+\sum_kP(N_k)\) covers
  all tree depths, reducible vertices and multi-ended necks (verifier
  §2). Schwarzschild check \(P=32\pi^2\) agrees with A01's neck
  baseline.
\item
  Bubble relocation (53.B) needs a one-attachment tree with
  \(M\cong M'\#\widehat T\); a two-ended neck is not movable.
\item
  Rank-one flux gives \(\mathbf R^+\le0\) only for \textbf{positive}
  flux; positivity follows from the positive-mass theorem
  (\(\operatorname{tr}\Phi^{++}=3\pi^2m_{\rm ADM}\)), and then exclusion
  gives \(s\le0\) for every nonflat ASD bubble.
\item
  Negative KE: \(\mathbf R^+=\operatorname{diag}(s/4,0,0)\), kernel of
  dimension 2, \textbf{unit kernel directions an \(S^1\)}.
\item
  53.F′ remainder \(O(\lambda^{12/5})\) is correct but needs the
  half-flat cancellation \(A(k)=O(r^{-6})\); the homogeneous \(r^{-2}\)
  expansion is a hypothesis; the claimed \(\lambda^2\log\) non-ASD
  interaction is deleted.
\item
  Lemma 53.Y holds for any order \(\tau>0\) with derivative-controlled
  decay (\(2<q<4/(2-\tau)\)).
\item
  Attainment ⇔ existence of a trivial energy-equality \emph{generalized
  configuration}, not necessarily Anderson's limit.
\item
  Bach-flat + \(P\)-critical + \(s\equiv\)const\(\neq0\) ⇒ Einstein
  (since \(\nabla\!\int s^2=-2s\operatorname{Ric}^\circ\) at constant
  \(s\)); ``Prop. nonbach'' is replaced by this one-line proof.
\item
  Anderson's printed proof has three soft spots: local minimality of
  weak limits; Prop. 4.6 regularity (the ``linearly independent
  summands'' inference is invalid) and its uniformity under rescaling,
  where the perturbation becomes
  \(\varepsilon\ell^{4-4p}\int(\ell^4+|\operatorname{Ric}|^2)^p\); the
  case-(I) equality (4.43) by citation. Theorem 4.10 is a published
  statement; we use it with these dependencies named.
\end{enumerate}

\begin{center}\rule{0.5\linewidth}{0.5pt}\end{center}

\hypertarget{topological-cost-of-smooth-bubbles-in-both-orientations}{%
\subsection{1. Topological cost of smooth bubbles in both
orientations}\label{topological-cost-of-smooth-bubbles-in-both-orientations}}

\textbf{Lemma 54.1 (two-orientation bubble cost).} Let \(N\) be a smooth
complete AE manifold with \(e\ge1\) ends, all with trivial group, with
derivative-controlled decay of some order \(\tau>0\) (so the
Chern--Gauss--Bonnet and signature boundary terms converge to the flat
values). Let \(\widehat N\) be the closed manifold obtained by adding
one point per end. Then \[
P(N)=4\!\int_N|W^+|^2+\tfrac1{12}\!\int_Ns^2+8\pi^2c(N)=4\!\int_N|W^-|^2+\tfrac1{12}\!\int_Ns^2+8\pi^2c(\bar N),
\] \[
c(N)=4(e-1)+4b_1+b^- -5b^+,\qquad c(\bar N)=4(e-1)+4b_1+b^+-5b^-,
\] Betti numbers of \(\widehat N\). Hence
\(P(N)\ge8\pi^2\max\{c(N),c(\bar N),0\}\).

\emph{Proof.} \(\chi(N)=\chi(\widehat N)-e\), Gauss--Bonnet integrand
total \(e_{GB}(N)=\chi(N)-e\), \(\tau(N)=\tau(\widehat N)\),
\(\chi(\widehat N)=2-2b_1+b_2\). Insert into
\(P=4\int|W^\pm|^2+\frac1{12}\int s^2-8\pi^2(2e_{GB}\pm3\tau)\).
\(\square\)

Checks (script): Burns \(c=1\); reversed Burns \(c(\bar N)=1\);
Schwarzschild \(c=4\) (\(=32\pi^2\)); three-ended neck 8;
\(S^1\times S^3\) cap 4; homotopy-sphere bubble 0; \(S^2\times S^2\) cap
\(-4\) in both orientations (no bound). \textbf{{[}Corrected after
verification.{]}}
\(c,\bar c\le0\iff D:=4(e-1)+4b_1\le\min\{5b^+-b^-,5b^--b^+\}\). This
does \textbf{not} force \(e=1\) or \(b_1=0\). Example: a two-ended
scalar-flat AE metric on \(S^2\times S^2\) minus two points (Green
functions of a PSC metric) has \(c=\bar c=0\); so does
\((S^1\times S^3)\#(S^2\times S^2)\) minus a point. The usable form is:
\[
b^+(\widehat N)=0\ \Longrightarrow\ P(N)\ge8\pi^2\,[4(e-1)+4b_1+b_2].\tag{1}
\] So \(b^+\) must be eliminated first, by Seiberg--Witten (Lemma 54.3).

\textbf{Lemma 54.1′ (region form; no bubble equation needed).} Let
\(g_i\) converge to a case-(I) configuration, and let \(p\in V\) be a
concentration point with trivial group. Assume the bulk satisfies
small-energy regularity on punctured balls around \(p\): since \(g_0\)
is \(P\)-critical on \(V\setminus\{p\}\) with finite energy and a
Sobolev bound, \(|\mathrm{Rm}_{g_0}|(x)=o(d(x,p)^{-2})\) with
derivatives. For small \(r_0\), choose by Fubini a good radius
\(r\in(r_0/2,r_0)\). Let \(U_i\subset M\) be the region bounded by the
sphere \(S_r(p)\) transplanted to \((M,g_i)\), and \(\widehat U\) its
cap (with \(M\cong\widehat U\#R\)). Then \[
P(g_i;U_i)\ \ge\ 8\pi^2\max\{c(\widehat U),c(\overline{\widehat U}),0\}-o_{r_0}(1)-o_i(1),
\] with \(c\) as above and \(e=1\). For a region with \(e\) boundary
spheres at bulk points (a neck joining bulk pieces) the same holds with
\(4(e-1)\) included.

\emph{Proof.} Chern--Gauss--Bonnet and the APS-signature formula with
boundary on \(U_i\). The boundary integrands (second fundamental form
and tangential curvature, Chern--Simons form) converge, for a.e. good
radius, under the \(C^{1,\alpha}\cap W^{2,q}\) bulk convergence on the
annulus. By the punctured small-energy estimate they converge to the
round-sphere values as \(r_0\to0\). These are the values used in Lemma
54.1. \(\square\)

So once \(b^+(\widehat U)=0\) is known, the bound (1) for the region
needs \textbf{no} bubble equation, decay, or scalar-flatness. What is
needed is only bulk criticality, which Anderson's statement supplies at
unit scale where the perturbation \(\varepsilon\to0\), together with
lower semicontinuity. Scalar-flatness and decay of bubbles
(\(\mathrm{H_{dec}}\)) enter only through Lemma 53.Y (PSC of
\(\widehat N\)). That is needed only in the indefinite subcase with
\(b^+(R)=0\); when \(b^+(R)\ge1\) the standard connected-sum vanishing
theorem suffices.

\textbf{Lemma 54.2 (summand structure).} In a case-(I) configuration
with only trivial groups, each capped bubble \(\widehat N_k\) is a
connected summand of \(M\): \(M\cong\widehat N_k\#R_k\) for a closed
\(R_k\). Globally
\(M\cong\#_aV_a\#\#_k\widehat N_k\#b_1(\mathcal G)(S^1\times S^3)\),
\(\mathcal G\) the gluing graph. (An \(e\)-ended bubble attached to
\(e\) boundary spheres of the rest is \(\widehat N_k\) summed to the
rest, with the extra attachments producing self-sums,
i.e.~\(S^1\times S^3\) summands.) Each \(\widehat N_k\) carries a metric
of positive scalar curvature if \(N_k\) is scalar-flat (Lemma 53.Y,
corrected form, each end compactified separately).

\textbf{Lemma 54.3 (PSC summands with \(b^+\ge1\) kill SW).} Let
\(M=Z\#R\) be closed with \(b^+(M)\ge2\), where \(Z\) admits a PSC
metric and \(b^+(Z)\ge1\). Then every Seiberg--Witten invariant of \(M\)
vanishes.

\emph{Proof.} If \(b^+(R)\ge1\) this is the standard connected-sum
vanishing theorem. In general stretch the neck: \(g_T=\) PSC metric on
\(Z\setminus B\) with round cylindrical end
\(\cup\,S^3\times[-T,T]\cup\) any metric on \(R\setminus B\). Choose the
self-dual perturbation \(\eta\) supported in \(Z\setminus B\), small and
generic. Monopoles for \((g_T,\eta)\) with \(T\to\infty\) converge,
after gauge, on compact subsets to finite-energy monopoles on both
cylindrical-end pieces, whose cylindrical limits are the (unique,
nondegenerate) flat reducible of the round \(S^3\). On the \(Z\)-side,
irreducibles are excluded by the Weitzenböck formula (\(s>0\), \(\eta\)
small); reducibles require
\(2\pi\,\mathcal H^+(c_1(L)|_Z)=\mathcal H^+(\eta)\) in the \(L^2\)
harmonic self-dual space, which is \(\cong H^+(Z)\neq0\), a
codimension-\(b^+(Z)\) condition avoided by generic \(\eta\). So the
\(Z\)-side has no limit and the moduli space of \(M\) is empty for
\(T\gg1\); as \(b^+(M)\ge2\) the invariant is metric-independent and
vanishes. \(\square\) (The compactness step is the standard
neck-stretching theorem for monopoles across a PSC \(S^3\) neck; to be
cited precisely, e.g.~Kronheimer--Mrowka/Morgan--Szabó--Taubes; not
re-proved here.)

\begin{center}\rule{0.5\linewidth}{0.5pt}\end{center}

\hypertarget{the-theorem}{%
\subsection{2. The theorem}\label{the-theorem}}

Let \(X\) be a compact complex-hyperbolic surface (smooth ball quotient
\(\Gamma\backslash\mathbb{CH}^2\)) other than a fake projective plane.
Since \(\chi(\mathcal O_X)=c_2/3\), \(p_g=c_2/3-1+q\). For \(c_2\ge6\)
this gives \(p_g\ge1\); for \(c_2=3\) it gives \(p_g=q\), and \(q\ge1\)
(e.g.~Cartwright--Steger) gives \(p_g\ge1\). Thus \(b^+(X)=2p_g+1\ge3\)
for every ball quotient except the fake projective planes
(\(c_2=3,q=0\)), which are exactly those with \(b^+=1\). Put
\(M=X\#\overline{\mathbb{CP}}^2\), \(A=c_1^2(X)=3c_2(X)\).

\textbf{Unconditional input (upper bound).} By the analytic CH-collar
cell (supplement Ch. 5: a single blow-up cell equal to the CH metric
outside a fixed collar,
\(P/(8\pi^2)=1-9\epsilon^2\log(1/\epsilon)+O_R(\epsilon^2)\)), inserted
into any embedded CH ball of \(X\) of radius below the injectivity
radius, \[
p_*(M)<8\pi^2 .\tag{U}
\] (The certified perpendicular cell, \(P/(8\pi^2)\le0.99566273277\),
gives the same where its collar fits.)

\textbf{Theorem 54.A.} Let \(X,M\) be as above. Consider a minimizing
configuration for \(P\) on \(M\) of Anderson's case (I) in which every
cone point has trivial group (hypothesis \(\mathrm{H_{sm}}\)), and
assume:

\begin{itemize}
\tightlist
\item
  \(\mathrm{H_{lsc}}\): \(p_*(M)\ge\sum_aP(V_a,g_0)+\sum_kP(N_k)\)
  (lower semicontinuity on separated scales);
\item
  \(\mathrm{H_{dec}}\): each bubble \(N_k\) is smooth, complete, AE with
  derivative-controlled decay of some order \(\tau>0\), and scalar-flat;
\item
  for part (iii) also \(\mathrm{H_{eq}}\): Anderson's case-(I) energy
  equality, and \(\mathrm{H_{rem}}\): smooth removability at the
  trivial-group cone points.
\end{itemize}

Then:

\begin{enumerate}
\def\labelenumi{(\roman{enumi})}
\item
  \textbf{No smooth bubble carries homology or connects components.} The
  bulk \(V\) is connected; every bubble is one-ended; every
  \(\widehat N_k\) is a homotopy 4-sphere. In particular \(E\) is
  carried by the bulk.
\item
  \textbf{The bulk is a homeomorphic copy of \(M\).} With
  \(\Sigma=\#_k\widehat N_k\) (a homotopy sphere), \(M\cong V\#\Sigma\),
  so \(V\) is homeomorphic to \(M\); each bubble has
  \(P(N_k)=4\int|W^+|^2=4\int|W^-|^2<8\pi^2\).
\item
  \textbf{Attainment and the dichotomy.} Under
  \(\mathrm{H_{eq}},\mathrm{H_{rem}}\): \(g_0\) is a smooth metric on
  \(V\) with \(P(g_0)=p_*(V)\), and \(p_*(M)=p_*(V)+\sum_kP(N_k)\). If
  every \(\widehat N_k\) is diffeomorphic to \(S^4\) (automatic if the
  smooth 4-dimensional Poincaré conjecture holds), there are \textbf{no
  bubbles}: \(V\cong M\) and \(g_0\) attains \(p_*(M)\). In all cases
  \(s_{g_0}\) is a negative constant, and \textbf{either \(g_0\) is
  Einstein, or \(g_0\) is a smooth non-Bach-flat global minimizer of
  \(P\)} (equivalently of \(\int|\mathrm{Rm}|^2\)) on \(V\).
\end{enumerate}

\emph{Proof (corrected order).} (i) Step 1: by Lemma 54.2,
\(M\cong\widehat N_k\#R_k\). If \(b^+(\widehat N_k)\ge1\), then either
\(b^+(R_k)\ge1\), and ordinary connected-sum vanishing applies; or
\(b^+(R_k)=0\), and Lemma 54.3 applies with \(\widehat N_k\) PSC
(\(\mathrm{H_{dec}}\)). Both contradict \(SW_M\ne0\). Here
\(b^+(M)\ge3\), and the blow-up formula applies to the canonical class
of \(X\). So \(b^+(\widehat N_k)=0\), with no energy used. Step 2: by
(U), \(\mathrm{H_{lsc}}\) and (1),
\(\sum_k[4(e_k-1)+4b_1+b_2](\widehat N_k)<1\), so \(e_k=1\) and
\(b_1=b_2=0\). Step 3: one-ended bubbles cannot bridge bulk components,
so \(V\) is connected and \(M\cong V\#\Sigma\). Graph loops are also
excluded by \(\pi_1\), since \(\Gamma\) is freely indecomposable and not
\(\mathbb Z\).

\(\pi_1\): \(\pi_1(M)=\Gamma=\pi_1(V)*\pi_1(\widehat N_1)*\cdots\).
\(\Gamma\) is torsion-free and a \(PD_4\) group (\(X\) is aspherical),
hence one-ended and freely indecomposable (Stallings). If some
\(\pi_1(\widehat N_k)\neq1\), then \(\pi_1(\widehat N_k)\cong\Gamma\),
and Hopf's theorem (\(H_2(Y;\mathbb Q)\to H_2(\pi_1Y;\mathbb Q)\) is
onto) gives \(b_2(\widehat N_k)\ge b_2(\Gamma)=b_2(X)\ge1\),
contradiction. So each \(\widehat N_k\) is a simply connected rational
homology sphere; \(H_2\) is torsion-free for simply connected closed
4-manifolds, so \(\widehat N_k\) is an integral homology sphere, hence a
homotopy sphere, homeomorphic to \(S^4\) (Freedman).

\begin{enumerate}
\def\labelenumi{(\roman{enumi})}
\setcounter{enumi}{1}
\item
  \(\Sigma\) is homeomorphic to \(S^4\), so
  \(V\approx V\#\Sigma\cong M\). For \(e=1\), \(b_1=b_2=0\), Lemma 54.1
  gives \(P(N_k)=4\int|W^\pm|^2\) (scalar-flat), and the sum is below
  \(8\pi^2\).
\item
  Theorem 53.D (as corrected, §0.1/0.7) gives \(P(g_0)=p_*(V)\) and
  \(p_*(M)=p_*(V)+\sum P(N_k)\). If all \(\widehat N_k\cong S^4\) then
  \(V\cong M\), \(p_*(V)=p_*(M)\), so \(\sum_kP(N_k)=0\); a scalar-flat
  AE bubble with \(\Gamma=1\) and \(P=0\) is Ricci-flat AE, hence flat
  by Bishop--Gromov --- no bubble. \(P\)-criticality gives
  \(\Delta s=0\), so \(s\) is constant. SW: \(SW_V\neq0\), because
  \(SW_{V\#\Sigma}=SW_V\) for a PSC homotopy sphere \(\Sigma\)
  (neck-stretching as in 54.3, now with \(b^+(\Sigma)=0\): the
  \(\Sigma\)-side contributes the unique unobstructed reducible) and
  \(SW_M\ne0\). Then \(s>0\) is impossible, and \(s\equiv0\) is
  impossible since \(V\) carries a monopole class \(c\) of square
  \(A-1>0\) (LeBrun: \(\int s^2\ge32\pi^2c^2\)). So \(s<0\) constant,
  and by §0.8 the metric is Einstein or non-Bach-flat. It is a global
  minimizer because it attains \(p_*(V)\). \(\square\)
\end{enumerate}

\textbf{Remarks on what has been gained.}

\begin{enumerate}
\def\labelenumi{\arabic{enumi}.}
\setcounter{enumi}{-1}
\item
  \emph{Hypothesis accounting after Lemma 54.1′.} (i)--(ii) hold under
  \(\mathrm{H_{sm}}\), \(\mathrm{H_{lsc}}\) and bulk punctured
  small-energy regularity, except in one subcase. That subcase is a ball
  region with indefinite \(\widehat U\) carrying all of \(b^+(M)\), so
  \(b^+(R)=0\), with \(b^-(\widehat U)\le5b^+\); there
  \(\mathrm{H_{dec}}\) is used for PSC. In that subcase the bulk has
  \(b^+(V)=0\), and \(\widehat U\) carries the entire positive part of
  the form of \(X\).
\item
  The exclusion of smooth bubbling uses only lower semicontinuity, decay
  and topology. It needs neither the second-order kernel coefficient,
  nor energy equality, nor the moving-bubble argument. It holds
  uniformly for every non-fake-plane ball quotient at once. The
  mechanism: the analytic CH cell puts \(p_*(M)\) strictly below the
  cheapest topologically nontrivial smooth bubble. Lemma 54.1 shows the
  cheapest such bubble costs \(8\pi^2\) --- semidefinite forms by the
  two-orientation charge, indefinite forms by PSC + SW, \(\pi_1\) by
  Hopf.
\item
  The Kähler recovery picture (KE ball quotient \(+\) Burns bubble, cost
  \(\to8\pi^2\)) and the neck-pinch picture
  \(X\sqcup(\overline{\mathbb{CP}}^2,g_{FS})\), both Einstein with
  \(P=0\), are excluded. The first is excluded by (U). The second is
  excluded because the connecting neck has \(e=2\) and costs
  \(\ge32\pi^2\). The verifier's Schwarzschild value and A01's baseline
  confirm this number.
\item
  Problem-level consequence. Under the gates,
  \(p_*(X\#\overline{\mathbb{CP}}^2)\) is attained (modulo exotic
  homotopy-sphere bubbles) by a smooth constant-negative-scalar
  \(P\)-critical global minimizer. Either it is Einstein, answering
  Problem 1 at \(k=1\) affirmatively in a range LeBrun's
  \(k\ge A/3=c_2\) obstruction leaves open. Or it is non-Bach-flat,
  answering Problem 3. Its value satisfies
  \(p_*(M)<8\pi^2(1-9\epsilon^2\log(1/\epsilon)+O(\epsilon^2))\) for
  every admissible \(\epsilon\), so the minimizer is strictly cheaper
  than every CH-plus-Burns configuration.
\item
  Entropy does not decide the dichotomy (script, \texttt{BCG\_Bishop}).
  For an Einstein \(g\) with \(\operatorname{Ric}=-6g\) on \(M\), the
  BCG volume-entropy inequality with Bishop's bound \(h\le3\sqrt2\)
  against \(h(\mathbb{CH}^2)=4\) gives only
  \(\operatorname{Vol}(g)\ge(64/81)\operatorname{Vol}(X)\). LeBrun's
  scalar estimate \(\int s^2\ge32\pi^2A\) already gives the stronger
  \(\operatorname{Vol}(g)\ge\operatorname{Vol}(X)\). Topology gives only
  the upper bounds
  \(\operatorname{Vol}(g)/\operatorname{Vol}(X)\le3-1/c_2\) (from
  \(T_+\)) and \(\le2+2/c_2\) (from \(\chi\)), so the window is
  nonempty. A side remark (direction corrected after verification): for
  Einstein \(g\) on \(M\),
  \(\int|W^+|^2-\int s^2/24=2\pi^2(A-1)-\int s^2/16\le-2\pi^2\). So any
  Einstein metric on \(M\) has \(W^+\)-energy strictly \textbf{below}
  the Kähler-type value \(\int s^2/24\). Equality would force
  \(\operatorname{Vol}=(1-\frac1{3c_2})\operatorname{Vol}(X)\), which
  LeBrun's estimate excludes. The real-hyperbolic Sambusetti mechanism
  does not transfer: complex-hyperbolic entropy \(4\) is not extremal
  among \(\operatorname{Ric}\ge-6\) metrics (\(3\sqrt2\)).
\end{enumerate}

\textbf{Extension to general \(k\) (same hypotheses,
\(M_k=X\#k\overline{\mathbb{CP}}^2\)).} Disjoint analytic cells give
\(p_*(M_k)<8\pi^2k\) for every fixed \(k\). The same proof shows that in
a smooth-point configuration \[
\sum_k\big[4(e_k-1)+4b_1(\widehat N_k)+\beta(\widehat N_k)\big]<k,\qquad
\beta=\max\{b^--5b^+,\,b^+-5b^-,0\},
\] with \(b^+(\widehat N_k)=0\) for every bubble (Lemma 54.3,
\(b^+(M_k)\ge3\)). Hence \(\sum_kb_2(\widehat N_k)<k\): \textbf{not all
exceptional classes can escape into smooth bubbles}, and the bulk is
connected (all \(e_k=1\)) as soon as \(p_*(M_k)<32\pi^2\). Converting
this into attainment for \(k\ge2\) needs strict subadditivity
\(p_*(M_k)<p_*(M_{k-m})+8\pi^2m\) at bulk minimizers, which is the
second-order kernel problem of Step 53 §4.

\begin{center}\rule{0.5\linewidth}{0.5pt}\end{center}

\hypertarget{the-remaining-gates-stated-as-mathematics}{%
\subsection{3. The remaining gates, stated as
mathematics}\label{the-remaining-gates-stated-as-mathematics}}

\textbf{(G1) Collapse (case II).} Scalar curvature alone cannot exclude
it: finite-volume complex-hyperbolic cusps collapse with bounded
curvature and constant \(s<0\). The relevant questions are topological
and charge-theoretic.

\begin{itemize}
\tightlist
\item
  \emph{Proved here.} A region \(C\) carrying a free circle action has
  \(TC\cong\pi^*TB\oplus\underline{\mathbb R}\), with \(B\) an
  orientable 3-manifold, so \(TC\) is trivial and \(w_2(C)=0\). Every
  class in the image of \(H_2(C)\to H_2(M)\) therefore has even square
  (Wu), and \(E\) (\(E^2=-1\)) is not carried by such a region.
\item
  \emph{Open.} (a) The same statement for positive-rank \(F\)-structures
  with finite isotropy (Seifert-type pieces). (b) The case where \(E\)
  is split between \(\Omega\) and \(C\): the Mayer--Vietoris
  intersection form across the graph-manifold interface
  \(\partial\Omega\). (c) Pricing the case-(II) energy loss with the
  Hopf-adiabatic charges of Theorem G. Anderson's case-(II) inequality
  is only \(\le\), but \(\mathrm{H_{lsc}}\) is all that (i) needs, so
  (i)'s argument applies verbatim to the bubbles of \(\Omega\).
\end{itemize}

\textbf{(G2) Orbifold points.} Here \(E\) can be absorbed only together
with an \(X\)-class: roots \(a+mE\) with \(a^2=m^2-2\), \(a\in H_2(X)\).
Kronheimer bubbles cost zero, so the bound must come from the orbifold
bulk. The lattice-free part of the argument above carries over. If the
orbifold bulk is good (a global quotient \(\tilde V/G\)), an orbifold SW
package gives a LeBrun bound that is useless here (\(1-A/3<0\)). Pricing
must instead come from the CH-germ interaction (A01's
\(S^4/\mathbb Z_2\) baseline \(8\pi^2\) per class). Note that absorption
of \(E\) needs \(m\) odd, hence \(a^2\) odd; for \(m=1\), \(a^2=-1\). So
the intersection form of \(X\) must be odd with a class of square
\(-1\). The parity lemma (52.9(v)) makes such absorbing orbifolds
non-Kähler. \textbf{{[}Corrected:{]}} \(a^2=m^2-2\) alone does not force
\(m\) odd. Even \(m\) (e.g.~\(a+2E\), \(a^2=2\)) is lattice-allowed.
``Absorbing \(E\)'' needs a precise definition and an integral interface
computation before any parity statement; the inference to a
square-\(-1\) class is withdrawn.

\textbf{(G3) Analytic gates.} \(\mathrm{H_{lsc}}\) is standard for weak
\(W^{2,2}\) convergence on compact regions at separated scales.
\(\mathrm{H_{dec}}\) needs (a) that bubbles of Anderson's perturbed
sequence are critical for \(P\) itself; the rescaled perturbation
coefficient \(\varepsilon\ell^{4-4p}\) must tend to zero, which is the
point the strategist and verifier flag. It also needs (b) a decay
theorem for scalar-flat Bach-flat AE ends with trivial group. With
\(s=0\) established, (b) needs only some \(\tau>0\) with derivative
control; Bando--Kasue--Nakajima-type arguments for \(L^2\)-curvature
critical ALE limits give order 0, so an improvement is needed.
\(\mathrm{H_{eq}}\) (Anderson (4.43)) and \(\mathrm{H_{rem}}\) are used
only in (iii).

\textbf{Gate economy.} (i)--(ii) need
\(\mathrm{H_{sm}},\mathrm{H_{lsc}},\mathrm{H_{dec}}\) only. The analytic
point that must be proved with our own tools is \(\mathrm{H_{dec}}\)(a):
bubble-scale criticality, i.e.~the vanishing of the rescaled
perturbation. A route: choose \(p=p(\varepsilon)\downarrow1\) with
\(\varepsilon\,\ell_{\min}(\varepsilon)^{4-4p}\to0\). Here
\(\ell_{\min}\) is the smallest \(L^6\)-curvature radius. A lower bound
\(\ell_{\min}\ge\exp(-C/\varepsilon)\) would suffice with
\(p-1=o(\varepsilon)\), since then
\(\varepsilon\ell^{4-4p}\le\varepsilon e^{4C(p-1)/\varepsilon}\to0\).
Such a lower bound plausibly follows from the energy bound
\(\int|\operatorname{Ric}|^{2p}\le C/\varepsilon\) via the \(L^{2p}\)
\(\varepsilon\)-regularity at scale \(\ell\):
\(\ell^{4p-4}\int_{B_\ell}|\operatorname{Ric}|^{2p}\ge c\) forces
\(\ell^{4-4p}\le C/(c\varepsilon)\), so
\(\varepsilon\ell^{4-4p}\le C/c\) is merely bounded. The bubble equation
then carries a bounded zeroth-order perturbation
\(\kappa\,\nabla\!\int|\operatorname{Ric}|^{2p}\), and \(\kappa\to0\)
requires a sharper choice. This is the precise open analytic estimate,
recorded rather than claimed.

\textbf{{[}Superseded in part by §5(c): the compact-replacement argument
gives bubble \(P\)-criticality directly. The verifier's estimate
\(\varepsilon_iQ_p(g_i)\le I_{\varepsilon_i}-p_*+o(1)\to0\) for genuine
smooth near-minimizers makes \(\kappa_i\to0\) at any scale with a Ricci
lower bound. The scaling identity below remains a valid alternative
under \(\mathrm{H_{eq}}\).{]}}

\textbf{Proposition 54.4 (bubble-scale criticality from a scaling
identity).} Fix \(p>1\) as in Anderson. Let
\(\hat g_i=\ell_i^{-2}g_{\varepsilon_i}\) be a bubble rescaling
(\(\ell_i\to0\)) converging to a nonflat complete limit \(N\), and put
\(\kappa_i=\varepsilon_i\ell_i^{4-4p}\). Then:

\begin{enumerate}
\def\labelenumi{(\alph{enumi})}
\item
  \textbf{\(\kappa_i\) is bounded} whenever
  \(\operatorname{Ric}_N\not\equiv0\). Indeed
  \(\kappa_i\int_{B_R}|\operatorname{Ric}_{\hat g_i}|^{2p}\le\varepsilon_i\int_M|\operatorname{Ric}_{g_{\varepsilon_i}}|^{2p}\le\mathcal Z_{\varepsilon_i}\le C\),
  and the left integral has a positive liminf by lower semicontinuity.
\item
  \textbf{Scaling (Pohozaev) identity.} Test the Euler--Lagrange
  equation of \(\hat g_i\) with the conformal variation
  \(\phi_R\hat g_i\). Here \(\phi_R=1\) on \(B_R\), \(\phi_R=0\) off
  \(B_{2R}\), \(|\nabla^2\phi_R|\lesssim R^{-2}\). The perturbed
  functional is \(P+\kappa_i\int(\ell_i^4+|\operatorname{Ric}|^2)^p+\) a
  volume multiplier \(\mu_i\ell_i^4\operatorname{Vol}\), with
  \(\mu_i=O(1)\). Since \(P\) is scale invariant, its conformal
  variation is
  \(\delta P(\phi)=-2\int\langle\operatorname{Ric}^\circ,\nabla^2\phi\rangle\).
  The perturbation's constant-scaling derivative is
  \((2-2p)\int|\operatorname{Ric}|^{2p}\) up to \(O(\ell_i^4)\). This
  gives \[
  (2p-2)\,\kappa_i\!\int_{B_R}\!|\operatorname{Ric}_{\hat g_i}|^{2p}
  \ \lesssim\ \Big(\int_{A_R}|\operatorname{Ric}^\circ_{\hat g_i}|^2\Big)^{1/2}
  +\kappa_i^{1/(2p)}R^{-(2p-2)/p}+O(\ell_i^4R^4),
  \] using Hölder on the annulus \(A_R=B_{2R}\setminus B_R\) together
  with \(\kappa_i\int|\operatorname{Ric}|^{2p}\le C\).
\item
  With case-(I) no-neck energy (\(\mathrm{H_{eq}}\): the \(L^2\) Ricci
  energy in the neck annuli between the bubble scale and the parent
  scale tends to zero), let \(i\to\infty\) and then \(R\to\infty\). This
  gives \(\lim\kappa_i\cdot\int|\operatorname{Ric}_N|^{2p}=0\). Hence
  \(\kappa_i\to0\) or \(\operatorname{Ric}_N\equiv0\). In both cases the
  limit equation is that of \(P\): the perturbation gradient tends to
  zero in distributions, since
  \(\kappa_i\int|\operatorname{Ric}|^{2p-1}|\nabla^2\psi|\le C\kappa_i^{1/(2p)}\to0\).
  By the Step-51/52 trace argument the limit is then scalar-flat
  (\(\Delta s=0\), \(s\in L^2\)).
\end{enumerate}

So the strategist's caution is answered. The rescaled perturbation is
automatically bounded, and the scaling identity forces it to vanish on
every bubble that is not Ricci-flat. The remaining analytic input is the
passage to the limit in the quadratic curvature terms of the
\(P\)-equation. This needs strong local \(L^2\) curvature convergence,
which follows from \(L^q\) curvature control with \(q>4\) and the Prop.
52.1 bootstrap. Together with decay, it is what is left of
\(\mathrm{H_{dec}}\). Note that 54.4 uses \(\mathrm{H_{eq}}\), whereas
Theorem 54.A(i) outside the indefinite subcase does not.

\begin{center}\rule{0.5\linewidth}{0.5pt}\end{center}

\hypertarget{answer-to-a01-825-local-ke-rigidity-near-g_gamma}{%
\subsection{\texorpdfstring{4. Answer to A01 \#825 (local KE rigidity
near
\(g_\Gamma\))}{4. Answer to A01 \#825 (local KE rigidity near g\_\textbackslash Gamma)}}\label{answer-to-a01-825-local-ke-rigidity-near-g_gamma}}

On \(\mathbb{CH}^2\) (\(\operatorname{Ric}=-6g\)) the curvature action
\(\mathring R\) on trace-free symmetric 2-tensors has eigenvalues \(-2\)
(J-invariant, dim 3) and \(+4\) (J-anti, dim 6) (Step 44,
\texttt{step44\_perturbative\_ke.json}). The eigenvalues alone do not
give the gap. The TT Codazzi--Bochner identity
\(\|\nabla h\|^2\ge\langle\mathring Rh,h\rangle+6\|h\|^2\) does, as an
\textbf{integrated} (not pointwise) estimate:
\(\langle(\nabla^*\nabla-2\mathring R)h,h\rangle\ge2\|h\|^2\) on TT
tensors. Hence there are no infinitesimal Einstein deformations, and by
Koiso's local rigidity theorem every Einstein metric sufficiently
\(C^{2,\alpha}\)-close (or \(H^s\), \(s>4\), in Bianchi gauge) to
\(g_\Gamma\) is isometric to it up to scale. The gap is integrated and
\(\Gamma\)-independent, and Lemma M (Step 45--46) supplies
quotient-uniform \(L^\infty\)/Hölder inverses, so the closeness
threshold is uniform over finite covers in the uniformly-local Hölder
norm. \(H^2\)-closeness is not enough (\(H^2\not\subset L^\infty\)).
Theorem \(A_C\) is a statement about the canonical monopole branch and
is not needed for this rigidity. For a general negative KE surface (not
a ball quotient), Koiso shows that infinitesimal Einstein deformations
are infinitesimal complex deformations. Einstein metrics near a KE
metric are KE for a nearby complex structure when the Kuranishi space is
smooth. In general this is the remaining condition.

\begin{center}\rule{0.5\linewidth}{0.5pt}\end{center}

\hypertarget{verification-outcome-binding-tracettcsresultsstep54_ttcs01_verifier1_integration.md}{%
\subsection{\texorpdfstring{5. Verification outcome (binding;
\texttt{trace/ttcs/results/step54\_ttcs01\_verifier1\_integration.md})}{5. Verification outcome (binding; trace/ttcs/results/step54\_ttcs01\_verifier1\_integration.md)}}\label{verification-outcome-binding-tracettcsresultsstep54_ttcs01_verifier1_integration.md}}

The verifier confirmed the following:

\begin{itemize}
\tightlist
\item
  the charge formula, with the Schwarzschild value \(32\pi^2\);
\item
  the summand decomposition (2) for arbitrary graphs;
\item
  Lemma 54.3, including the case \(b^+(R)=0\), via the one-sided PSC
  neck argument (an exact citation is still to be registered);
\item
  the \(\pi_1\)/Hopf/Freedman chain;
\item
  part (iii): \(SW_{V\#\Sigma}=\pm SW_V\), \(s<0\) constant, the
  Einstein/non-Bach-flat dichotomy, and the no-bubble conclusion for
  standard caps;
\item
  the BCG arithmetic and the Koiso gap \(2\) (integrated form).
\end{itemize}

It also confirmed that (U) holds for \textbf{every} compact ball
quotient, fake planes included: Ch. 5 builds the cell for every fixed
\(R>0\). Corrections applied above: the order of proof (SW first),
Remark 4's direction, orbifold parity, and the Koiso wording. Additional
results recorded:

\textbf{(a) Degree-one strengthening, general \(k\).} The blowdown
\(M_k\to X\) is a classifying map. By free indecomposability, exactly
one vertex cap \(Y\) carries \(\Gamma\), and \(Y\to X\) has degree one,
so \(Q_X\hookrightarrow Q_Y\) and \(b^+(Y)\ge b^+(X)\ge3\). Since every
bubble has \(b^+=0\), that vertex is bulk. Hence \textbf{every bubble
cap is simply connected} for every \(k\), and the \(H_2\)-image of the
bubbles lies in \(\operatorname{span}\{E_1,\dots,E_k\}\). With
\(p_*(M_k)<8\pi^2k\): \[
\sum_j[4(e_j-1)+4b_1(\widehat N_j)+b_2(\widehat N_j)]<k ,
\] so the total exceptional rank in smooth bubbles is \(<k\), and all
bubbles are one-ended for \(k\le4\).

\textbf{(b) Hypothesis ledger.}

\begin{longtable}[]{@{}ll@{}}
\toprule
\begin{minipage}[b]{0.47\columnwidth}\raggedright
Conclusion\strut
\end{minipage} & \begin{minipage}[b]{0.47\columnwidth}\raggedright
Hypotheses\strut
\end{minipage}\tabularnewline
\midrule
\endhead
\begin{minipage}[t]{0.47\columnwidth}\raggedright
(U)\strut
\end{minipage} & \begin{minipage}[t]{0.47\columnwidth}\raggedright
Ch. 5 cell only\strut
\end{minipage}\tabularnewline
\begin{minipage}[t]{0.47\columnwidth}\raggedright
\(b^+(\widehat N_k)=0\)\strut
\end{minipage} & \begin{minipage}[t]{0.47\columnwidth}\raggedright
\(\mathrm{H_{sm}}\), summand topology; PSC (from \(\mathrm{H_{dec}}\))
only if \(b^+(R_k)=0\)\strut
\end{minipage}\tabularnewline
\begin{minipage}[t]{0.47\columnwidth}\raggedright
one-ended, homologically trivial, simply connected bubbles; connected
bulk \(V\approx M\)\strut
\end{minipage} & \begin{minipage}[t]{0.47\columnwidth}\raggedright
add (U), \(\mathrm{H_{lsc}}\)\strut
\end{minipage}\tabularnewline
\begin{minipage}[t]{0.47\columnwidth}\raggedright
\(P(g_0)=p_*(V)\)\strut
\end{minipage} & \begin{minipage}[t]{0.47\columnwidth}\raggedright
add \(\mathrm{H_{rem}}\), \(\mathrm{H_{eq}}\)\strut
\end{minipage}\tabularnewline
\begin{minipage}[t]{0.47\columnwidth}\raggedright
\(s<0\), Einstein/non-Bach-flat\strut
\end{minipage} & \begin{minipage}[t]{0.47\columnwidth}\raggedright
add the PSC-homotopy-sphere SW gluing\strut
\end{minipage}\tabularnewline
\begin{minipage}[t]{0.47\columnwidth}\raggedright
attainment on \(M\), no bubbles\strut
\end{minipage} & \begin{minipage}[t]{0.47\columnwidth}\raggedright
all caps \(\cong S^4\)\strut
\end{minipage}\tabularnewline
\bottomrule
\end{longtable}

\(\mathrm{H_{eq}}\) becomes redundant if vanishing-error reassembly is
proved (reassembly gives \(\le\), and \(\mathrm{H_{lsc}}\) gives
\(\ge\)).

\textbf{(c) Bubble criticality by compact replacement.} Suppose the
actual smooth \(P\)-minimizing sequence converges smoothly on compact
sets at a bubble scale. A compactly supported \(P\)-lowering variation
of the bubble, transplanted into \(g_i\), would give \(P<p_*\). So the
bubble locally minimizes \(P\), hence is \(P\)-critical, has
\(s\equiv0\), and is Bach-flat. This needs neither \(\mathrm{H_{eq}}\)
nor \(\kappa\to0\). The genuine analytic residue of \(\mathrm{H_{dec}}\)
is therefore: smooth (or strong enough) compact convergence at bubble
scales, plus derivative-controlled AE decay of some order \(\tau>0\) for
scalar-flat Bach-flat AE ends with trivial group.

\textbf{(d) What remains geometrically.} Theorem 54.A excludes all
homologically or topologically charged smooth bubbling on
\(X\#\overline{\mathbb{CP}}^2\). The geometries still open are:

\begin{enumerate}
\def\labelenumi{\arabic{enumi}.}
\tightlist
\item
  collapse (case II);
\item
  nontrivial-group orbifold points (absorption);
\item
  nonflat bubbles whose caps are exotic homotopy spheres, which would
  also disprove smooth 4D Poincaré.
\end{enumerate}

\hypertarget{s1003a02-step-55-the-k-1-orbifold-gate-two-orientation-charges-with-ux3b7-invariants-the-degree-one-lattice-constraint-and-what-survives}{%
\chapter{S1003A02: Step 55 --- The k = 1 orbifold gate: two-orientation
charges with η-invariants, the degree-one lattice constraint, and what
survives}\label{s1003a02-step-55-the-k-1-orbifold-gate-two-orientation-charges-with-ux3b7-invariants-the-degree-one-lattice-constraint-and-what-survives}}

Source:
\texttt{step55\_k1\_orbifold\_gate\_charges\_lattices\_collapse.md}.

\begin{quote}
\textbf{Step-56 binding corrections:} see
\texttt{step56\_orbifold\_attainment\_zero\_energy\_bound\_Tpoint\_interaction.md}
§0. The changes are: - Prop. 55.3 is restricted to simply connected
regions. - Prop. 55.6 is reproved via Kähler + Siu. - The §6 collapse
ledger is withdrawn. - Theorem 55.7 is reduced to a coarse trichotomy
plus the additive budget. Session S1003A02, Agent2, Execute. Continues
\texttt{step54\_k1\_ball\_quotient\_attainment\_theorem.md} (§5 binding)
and the strategist's note \#70. Scripts and data:
\end{quote}

\begin{itemize}
\tightlist
\item
  \texttt{workspace/scripts/step55\_orbifold\_bubble\_charges.py} →
  \texttt{workspace/data/step55\_orbifold\_bubble\_charges.json} (exact
  Dedekind-sum η-invariants, charges, convention checks, numerical \(P\)
  of LeBrun's \(\mathcal O(-n)\) metrics);
\item
  \texttt{workspace/scripts/step55\_Tchain\_embeddings.py} →
  \texttt{workspace/data/step55\_Tchain\_embeddings.json} (primitive
  embeddings of T-singularity chain lattices into \(-I_{r+1}\), all
  \(n\le7\)).
\end{itemize}

Setting throughout: \(X\) a compact ball quotient, \(\Gamma=\pi_1X\)
(torsion-free, Gromov hyperbolic, one-ended, freely indecomposable),
\(M=X\#\overline{\mathbb{CP}}^2\), \(A=c_1^2(X)\), \(T_+=2\chi+3\tau\),
\(T_+(M)=A-1\), and (U): \(p_*(M)<8\pi^2\). All statements about
minimizing configurations are inside Anderson's case (I) with the
analytic ledger of step54 §5(b) (\(\mathrm H_{\rm lsc}\), bubble
criticality by compact replacement, and derivative-controlled ALE decay
of the bubbles, \(\mathrm H_{\rm dec}\), now with nontrivial group).

\textbf{Outcome in one paragraph.} The strategist's degree-one reduction
is correct after one added Seiberg--Witten input, and it does kill the
Kronheimer hyperkähler channel at \(k=1\). It does \emph{not} close the
orbifold gate. Pricing every remaining bubble with the exact
two-orientation charge (η included) gives: every orbifold bubble
carrying \(E\) or \(2E\) through a \textbf{cyclic} end costs
\(\ge8\pi^2\) and is excluded by (U) (LeBrun's \(\mathcal O(-4)\) costs
exactly \(8\pi^2\)); but charge-free or negatively charged channels
survive: classes \(mE\), \(m\ge3\), through lens ends (first:
\(L(9,4)\), \(c_+=-23/9\)); \(2E\) through binary-dihedral ends
\(D_5,D_9\); \(E\) itself through the binary-icosahedral end
(\(c_+=-7\)); and rational-ball bubbles. Among rational balls, the
zero-charge ones are exactly the T-singularity Milnor fibres (Ricci-flat
Kähler ALE, \(P=0\)), and the degree-one map forces a sharp lattice
condition: the T-chain must embed \textbf{primitively} in a
negative-definite unimodular lattice of rank \(r+1\). This fails for the
strategist's simplest case \(B_2=\mathrm{EH}/\mathbb Z_2\) (\(n=2\)) and
exactly for \(n\equiv2,4,6\pmod 8\); it holds otherwise (first:
\(1/9(1,2)\), chain \([5,2]\)). The zero-charge bulk is provably
\textbf{non-Kähler}, and a conservation identity shows each needs
\(C_+(V)<A-w(N)\) with \(w(N)\) the bubble's self-dual Weyl charge. The
gate is therefore reduced to one analytic question --- an unrestricted
lower bound for \(P\) on non-Kähler orbifolds rationally equivalent to
\(X\) or \(M\) --- plus a finite explicit list for each \(|\Gamma|\).

\begin{center}\rule{0.5\linewidth}{0.5pt}\end{center}

\hypertarget{charges-of-ale-bubbles-with-nontrivial-group}{%
\subsection{1. Charges of ALE bubbles with nontrivial
group}\label{charges-of-ale-bubbles-with-nontrivial-group}}

Let \(N\) be smooth, complete, ALE with one end
\(\mathbb R^4/\Gamma_\infty\), oriented by \(M\), with
derivative-controlled decay of some order \(\tau_0>0\). For a cyclic end
the generator acts by \((\zeta,\zeta^q)\), \(\zeta=e^{2\pi i/n}\), in
the orientation of \(N\). Put \(\eta(n,q):=-4\,s(q,n)\), where \(s\) is
the Dedekind sum (APS signature-defect convention). Chern--Gauss--Bonnet
and APS give \[
P(N)=4\!\int_N|W^\pm|^2+\tfrac1{12}\!\int_Ns^2+8\pi^2c_\pm(N),\qquad
c_\pm(N)=\frac{2}{|\Gamma_\infty|}-2\chi(N)\mp3\tau(N)\mp3\eta_{\Gamma_\infty}. \tag{1.1}
\] Hence \(P(N)\ge8\pi^2\max\{c_+,c_-,0\}\).

\textbf{Convention checks (all exact, script).}

\begin{longtable}[]{@{}lll@{}}
\toprule
bubble & \(c_+\) & \(c_-\)\tabularnewline
\midrule
\endhead
Burns \(\mathcal O(-1)\) & 1 & −5\tabularnewline
Eguchi--Hanson & 0 & −6\tabularnewline
\(A_3\) hyperkähler & 0 & −15\tabularnewline
\(A_6\) hyperkähler & 0 & −192/7\tabularnewline
\(\mathrm{EH}/\mathbb Z_2\), complex orientation & 0 & −3\tabularnewline
\bottomrule
\end{longtable}

\textbf{New exact value (checked against the explicit metric).} LeBrun's
scalar-flat Kähler ALE metric on \(\mathcal O(-n)\), \[
g=\frac{dr^2}{V}+r^2(\sigma_1^2+\sigma_2^2)+r^2V\sigma_3^2,\qquad V=1+\frac{n-2}{r^2}+\frac{1-n}{r^4},
\] has \(s\equiv0\) (symbolic) and \(W^+=0\), so (1.1) predicts \[
P(\mathcal O(-n),g_{\rm LeB})=8\pi^2\frac{(n-2)^2}{n}.
\] Direct numerical integration of \(\int|\operatorname{Ric}|^2\), using
the orthonormal Bianchi-IX Ricci formulas of
\texttt{step48\_bianchi\_indicial.py} in the \(r\)-coordinate, agrees
for \(n=1,\dots,8\) to \(10^{-14}\): the values are
\(1,0,1/3,1,9/5,8/3,25/7,9/2\) in units of \(8\pi^2\). This fixes the η
sign: the \(\mathcal O(-n)\) end has \(q=1\) and
\(\eta(n,1)=-(n-1)(n-2)/(3n)\). The \(A_{n-1}\) end has \(q=-1\) and
\(\eta=+(n-1)(n-2)/(3n)\).

\textbf{Conservation (ledger) identity.} If a bubble region \(N\) with
end \(S^3/\Gamma\) is glued into an orbifold point of the bulk \(V\),
then \(\chi_{\rm orb}(V)=\chi(V^\circ)+1/|\Gamma|\) and
\(\tau_{\rm orb}(V)=\tau(V^\circ)-\eta_N\). Hence \[
T_+^{\rm orb}(V)=T_+(M)+c_+(N),\tag{1.2}
\] and, with the bubble scalar-flat, \[
\frac{P(V)+P(N)}{8\pi^2}=C_+(V)+w(N)-T_+(M),\qquad w(N):=\frac1{2\pi^2}\int_N|W^+|^2=\frac{P(N)}{8\pi^2}-c_+(N).\tag{1.3}
\] Here \(C_+(V)=\frac1{2\pi^2}\int|W^+|^2+\frac1{96\pi^2}\int s^2\) is
the orbifold integral. A whole subtree hanging below an orbifold point
telescopes to one effective smooth region \(N_{\rm eff}\): internal
links cancel exactly as in Theorem G, and \(4\int|W^+|^2\ge0\) on every
piece. So (1.1)--(1.3) apply to \(N_{\rm eff}\).

(1.3) says what a bubble does. It removes \(w(N)\) units of self-dual
Weyl energy from the bulk's budget, and the bulk then needs
\(C_+(V)<A-w(N)\) for the configuration to beat (U). The bulk always has
\(C_+(V)\ge T_+^{\rm orb}(V)\); (1.2) makes this the same necessary
condition \(\max\{c_+,c_-,0\}<1\).

\textbf{Two-piece duality.} If a closed \(Z\) splits as \(N\cup_LW\)
along a spherical space form and both pieces are read as ALE with
opposite ends, then \(c_+(N)+c_+(W)=4/|\Gamma|-T_+(Z)\). For
\(Z=\overline{\mathbb{CP}}^2=C_2\cup(-B_2)\): \(1+(-3)=-2=1-3\).

\begin{center}\rule{0.5\linewidth}{0.5pt}\end{center}

\hypertarget{degree-one-reduction-for-orbifold-bubbles}{%
\subsection{2. Degree-one reduction for orbifold
bubbles}\label{degree-one-reduction-for-orbifold-bubbles}}

\textbf{Lemma 55.1.} In a case-(I) configuration on \(M\), let \(N\) be
the effective smooth region below any orbifold point \(p\in V\) of the
bulk, with link \(L=S^3/\Gamma_p\). Assume:

(SW\(_L\)) if \(M=A\cup_LB\) with \(L\) a spherical space form and
\(b^+(A),b^+(B)\ge1\), then \(SW_M=0\). The one-sided version also
holds: if \(b^+(B)\ge1\), \(b^+(A)=0\), and
\(\widehat B=B\cup\mathrm{cone}\) carries a PSC orbifold metric, then
\(SW_M=0\).

Assume also that \(X\) is not a fake projective plane (so \(b^+(M)\ge3\)
and \(SW_M\ne0\)). Then \(\pi_1(N)\to\Gamma\) is trivial,
\(H_2(N;\mathbb Q)\to H_2(M;\mathbb Q)\) is injective with image in
\(\ker(c_*:H_2(M)\to H_2(X))=\mathbb ZE\), and \(b_2(N)\le1\). At most
one orbifold point of \(V\) has \(b_2(N_{\rm eff})=1\); all the others
are rational homology balls.

\emph{Proof.} Van Kampen: \(\Gamma=\pi_1(V^\circ)*_{\Gamma_p}\pi_1(N)\)
as a pushout. \(\Gamma_p\) is finite and \(\Gamma\) is torsion-free, so
\(\Gamma_p\) dies and
\(\Gamma=\pi_1(V^\circ)/\langle\!\langle\Gamma_p\rangle\!\rangle*\pi_1(N)/\langle\!\langle\Gamma_p\rangle\!\rangle\).
Free indecomposability makes one factor trivial. Suppose the bulk factor
were trivial. Then the classifying map \(c:M\to X\) is null on
\(V^\circ\), factors through \(M/V^\circ=\widehat N\), and has degree
one. Then \(Q_X\) embeds in \(Q_N\), so \(b^+(N)\ge3\) while
\(b^+(V^\circ)\) is whatever remains. If \(b^+(V^\circ)\ge1\),
(SW\(_L\)) gives a contradiction. If \(b^+(V^\circ)=0\), \(N\) is
scalar-flat ALE (bubble criticality, step54 §5(c)), its Green-function
compactification \(\widehat N\) is a PSC orbifold (Lemma 53.Y, orbifold
point at infinity), and the one-sided (SW\(_L\)) again contradicts
\(SW_M\ne0\). So the bubble factor is trivial and \(c|_N\) lifts to
\(\widetilde X\) (contractible). Hence \(H_2(N)\to H_2(X)\) is zero.
Mayer--Vietoris over \(\mathbb Q\) (with \(H_{1,2}(L;\mathbb Q)=0\))
gives
\(H_2(M;\mathbb Q)=H_2(V^\circ;\mathbb Q)\oplus\bigoplus H_2(N_j;\mathbb Q)\)
orthogonally, so the \(N_j\) contribute orthogonal subspaces of the line
\(\mathbb QE\). \(\square\)

(SW\(_L\)) is the standard vanishing theorem for splittings along PSC
rational homology spheres when both sides have \(b^+\ge1\). The
one-sided orbifold-cap form is the analogue of Lemma 54.3. Both are to
be cited precisely; they are recorded as hypotheses, not re-proved.

\textbf{Corollary 55.2 (Kronheimer is a \(k\ge2\) phenomenon).} A
hyperkähler ALE bubble has \(H_2\) spanned by \((-2)\)-classes, while a
\(b_2=1\) bubble region at \(k=1\) carries only \(\sigma\mapsto mE\)
with \(\sigma^2=-m^2\neq-2\). So no Kronheimer (\(\Gamma\subset SU(2)\),
hyperkähler) bubble occurs at \(k=1\). For general \(k\) the same proof
puts every bubble lattice inside
\(\operatorname{span}\{E_i\}\cong-I_k\): this is exactly the Step-52
root-lattice absorption (\(A_1\oplus A_1\subset\mathbb Z^2\), \ldots).

\begin{center}\rule{0.5\linewidth}{0.5pt}\end{center}

\hypertarget{pricing-the-b_21-orbifold-bubbles}{%
\subsection{\texorpdfstring{3. Pricing the \(b_2=1\) orbifold
bubbles}{3. Pricing the b\_2=1 orbifold bubbles}}\label{pricing-the-b_21-orbifold-bubbles}}

Here \(\chi(N)=2\), \(\tau(N)=-1\), and the class is \(mE\). If \(N\) is
simply connected, \(|H_1(\partial N)|=m^2\). For a cyclic end this means
\(n=m^2\), and the linking form \(\langle\pm1/m^2\rangle\) forces \(q\)
to be a unit square mod \(m^2\) (\(q=1\) is the \(\mathcal O(-m^2)\)
end). Then \[
c_+=\frac{2}{m^2}-1-3\eta(m^2,q),\qquad c_-=\frac{2}{m^2}-7+3\eta(m^2,q).
\]

\textbf{Proposition 55.3.} (a) For \(m=1\) (trivial group: Burns type)
and \(m=2\) (the only admissible \(q\) is \(1\), the \(\mathcal O(-4)\)
end), \(c_+=1\). So every such bubble costs \(\ge8\pi^2\), and by (U)
and \(P(V)\ge0\) it cannot occur. LeBrun's \(\mathcal O(-4)\) metric
realizes the bound: \(P=8\pi^2\) exactly. \emph{The class \(E\) or
\(2E\) cannot escape into an orbifold bubble with a cyclic end.}

\begin{enumerate}
\def\labelenumi{(\alph{enumi})}
\setcounter{enumi}{1}
\item
  For \(m\ge3\), charge does not exclude: the admissible \(q\) with
  \(\max\{c_\pm\}<1\) are listed in the JSON. The first ones are
  \(m=3\), \(q\in\{4,7\}\) (\(c_+=-23/9\), \(c_-=-5\)); \(m=4\), \(q=9\)
  (\(c_+=1/4\)); \(m=5\), \(q\in\{6,11,16,21\}\). These lens spaces
  occur as boundaries of \(\langle-m^2\rangle\) traces (e.g.~lens-space
  surgeries on \(T(2,5)\)), so the topology is not empty.
\item
  Non-cyclic ends. With \(\Gamma\subset SU(2)\) in the \emph{reversed}
  orientation, \(c_+=4/|\Gamma|+r-3\ge1\) for every \(D_r,E_r\), so all
  of those are excluded. With \(\Gamma\subset SU(2)\) in \(M\)'s
  orientation (η from Kronheimer's \(P=0\):
  \(\eta=(2/|\Gamma|+r-2)/3\)): \(c_+=1-r\), \(c_-=4/|\Gamma|-9+r\).
  Charge excludes \(D_r\) for \(r\ge10\). It does not exclude
  \(D_4,\dots,D_9,E_6,E_7,E_8\). Simple connectivity and the linking
  form (\(H_1\) cyclic of order \(m^2\) with form
  \(\langle-1/m^2\rangle\)) leave \(m=2\): \(D_5,D_9\)
  (\(H_1=\mathbb Z_4\), \(r\equiv1\bmod4\)), and \(m=1\): \(E_8\). The
  last is a bubble carrying \(E\) itself with Poincaré-sphere end,
  topologically the \(-1\)-trace of the left-handed trefoil, with
  \(c_+=-7\) and \(c_-=-29/30\).
\end{enumerate}

Interpretation through (1.3). Each surviving case has \(c_+<1\) and
pushes charge into the bulk: \(T_+^{\rm orb}(V)=A-1+c_+\) (for \(E_8\),
\(A-8\)). To beat (U) the bulk needs \(C_+(V)<A-w(N)\) with
\(w(N)\ge\max\{0,-c_+,c_--c_+\}\) (for \(E_8\), \(w\ge7\)). An orbifold
LeBrun estimate (\(C_+(V)\ge2A/3\), assuming an orbifold SW package on
\(V\) that retains the canonical class) would exclude a channel whenever
\(A\le3w(N)\). So it removes the strongly negative channels for small
\(X\) only.

\begin{center}\rule{0.5\linewidth}{0.5pt}\end{center}

\hypertarget{rational-ball-bubbles-and-the-t-singularity-channel}{%
\subsection{4. Rational-ball bubbles and the T-singularity
channel}\label{rational-ball-bubbles-and-the-t-singularity-channel}}

Here \(\chi=1\) and \(\tau=0\), \(n=|H_1(\partial N)|=m^2\), and
\(c_\pm=2/n-2\mp3\eta\). Exact computation for \(m\le10\) and all \(q\):
about half the lens ends have \(\max c_\pm\ge1\) and are excluded. For
\(m\ge3\) the maximum over \(q\) grows like \(m^2\).

\textbf{Proposition 55.4 (zero charge = T-singularity Milnor fibres).}
For every T-singularity \(\tfrac1{n^2}(1,na-1)\), \(\gcd(n,a)=1\),
\(n\le10\), the end has \(c_+=0\) in the complex orientation and
\(c_-=-4(1-1/n^2)\) (exact, script). The Milnor fibre \(B_{n,a}\)
carries a Ricci-flat Kähler ALE metric, the free
\(\mathbb Z_n\)-quotient of the \(A_{n-1}\) hyperkähler metric
(Şuvaina--Wright type; citation to register). In its complex orientation
\(W^+=0\), so \(P=0\) attains the charge bound. In the reversed
orientation \(w=4(1-1/n^2)\). Example:
\(\mathrm{EH}/\mathbb Z_2=D(T^*\mathbb{RP}^2)\) has charges \((0,-3)\).

So a zero-charge rational-ball bubble is precisely the Riemannian shadow
of a \(\mathbb Q\)-Gorenstein smoothing. Write the bulk as
\(V=(Y\setminus C)\cup\mathrm{cone}\), where \(C\) is the resolution
chain. Then \(M=Y_C\) is the rational blowdown, and (1.2) gives
\(T_+^{\rm orb}(V)=A-1\): no charge moves.

\textbf{Theorem 55.5 (lattice condition for the zero-charge channel at
\(k=1\)).} Suppose a case-(I) configuration on \(M\) has a \(B_{n,a}\)
bubble in its complex orientation. Let \(Y=V^\circ\cup C\), with \(C\)
the minimal-resolution plumbing of length \(r\). Then:

\begin{enumerate}
\def\labelenumi{(\roman{enumi})}
\item
  \(\pi_1(Y)=\Gamma\), and \(Y\to X\) has degree one.
\item
  Every sphere of \(C\) lies in \(K=\ker c_*\subset H_2(Y)\), a
  negative-definite unimodular lattice of rank \(r+1\); for \(r\le6\),
  \(K\cong-I_{r+1}\).
\item
  The chain lattice \(C_{n,a}\) embeds \textbf{primitively} in \(K\).
\end{enumerate}

Consequently: \textbf{a primitive embedding into some negative-definite
unimodular lattice of rank \(r+1\) exists iff \(1-na\) is a unit square
mod \(n^2\), iff \(n\) is odd or \(8\mid n\).} The zero-charge T-channel
at \(k=1\) is therefore impossible for \(n\equiv2,4,6\pmod 8\) and
lattice-allowed otherwise. Exhaustive search into \(-I_{r+1}\)
(\texttt{step55\_Tchain\_embeddings.json}) confirms both directions for
all \(n\le7\). Discriminant values were checked for all 79
T-singularities with \(n\le16\)
(\texttt{step55\_Tchain\_discriminants.json}): \(\beta\equiv1-na\), and
the unit-square test matches the parity law.

\emph{Proof.} (i) \(\pi_1(M)\) is the pushout of
\(\pi_1(V^\circ)\leftarrow\pi_1(L)\to\pi_1(B)=\mathbb Z_n\), and the
right map is onto. The image of the finite group \(\pi_1(L)\) in the
torsion-free \(\Gamma\) is trivial, so
\(\Gamma=\pi_1(V^\circ)/\langle\!\langle\pi_1L\rangle\!\rangle*\pi_1(B)/\langle\!\langle\pi_1L\rangle\!\rangle\)
and the second factor is trivial. Since \(C\) is simply connected,
\(\pi_1(Y)=\pi_1(V^\circ)/\langle\!\langle\pi_1L\rangle\!\rangle=\Gamma\).
Both classifying maps are null on \(C\) and on \(B\), so both factor
through \(V^\circ/\partial V^\circ\); hence
\(\deg(Y\to X)=\deg(M\to X)=1\).

\begin{enumerate}
\def\labelenumi{(\roman{enumi})}
\setcounter{enumi}{1}
\item
  A degree-one map splits \(H_2(Y)=K\oplus c^!H_2(X)\) orthogonally with
  \(Q_Y|_{c^!H_2X}\cong Q_X\), so \(K\) is unimodular. Its rank is
  \(b_2(M)+r-b_2(X)=r+1\), and its signature
  \(\tau(M)-r-\tau(X)=-(r+1)\). Spheres map to zero in \(H_2(X)\)
  because \(\pi_2(X)=0\). Unimodular definite lattices of rank \(\le7\)
  are diagonal.
\item
  This is the step that uses torsion. The sequence of \((Y,V^\circ)\)
  gives
  \(0\to\operatorname{coker}(H_2(Y)\to C^*)\to H_1(V^\circ)\to H_1(Y)\to0\).
  The cokernel \(G_c\), of order \(o\), equals
  \(\operatorname{coker}(K\to C^*)\) because \(Q_X\perp C\), and it is
  the image of \(H_1(L)=\mathbb Z_{n^2}\). Mayer--Vietoris for
  \(M=V^\circ\cup B\), with \(H_1(B)=\mathbb Z_n\) and
  \(H_1(L)\to H_1(B)\) onto, gives
  \(|\mathrm{tors}\,H_1(M)|=|\mathrm{tors}\,H_1(Y)|\cdot\gcd(o,n)\). But
  \(\pi_1(M)=\pi_1(Y)=\Gamma\), so \(\gcd(o,n)=1\). Since \(o\mid n^2\),
  \(o=1\): \(K\to C^*\) is onto, i.e.~\(C\) is primitive. For the
  congruence: primitivity in unimodular \(K\) makes
  \(C^\perp=\langle w\rangle\) with \(w^2=-n^2\)
  (\(\det C^\perp=\det C\)), and the gluing needs an anti-isometry
  \(\operatorname{disc}(C)\cong-\operatorname{disc}\langle-n^2\rangle=\langle1/n^2\rangle\).
  The chain's discriminant form is \(\langle\beta/n^2\rangle\) with
  \(\beta\equiv1-na\) (computed), so the condition is
  \(1-na\equiv u^2\pmod{n^2}\). Conversely, an anti-isometry produces a
  unimodular overlattice, so the condition is also sufficient. Odd
  squares are \(\equiv1\pmod8\), while \(1-na\equiv3,5,7\pmod8\) when
  \(v_2(n)\in\{1,2\}\); for \(v_2(n)\ge3\), \(1-na\equiv1\pmod8\). For
  odd primes, \(1-na\equiv1\pmod p\) is a square by Hensel. For
  \(r+1\ge8\) the overlattice need not be \(-I_{r+1}\), but any
  unimodular \(K\) is admissible in (ii). \(\square\)
\end{enumerate}

\textbf{The strategist's \(S^4\)-free example is excluded; the first
real one is \(n=3\).} For \(n=2\) (\(B_2=\mathrm{EH}/\mathbb Z_2\)) the
only embedding of \(\langle-4\rangle\) in \(-I_2\) is \(2E'\), which is
not primitive. Equivalently, the standard
\(B_2\subset\overline{\mathbb{CP}}^2\) (the complement of the conic)
sits in the \emph{reversed} orientation: charge \(-3\), bulk
\(V=X\#\overline{\mathbb P}(1,1,4)\) with \(T_+^{\rm orb}=A-4\) and
\(w=3\). This is a negatively charged channel. For \(n=3\), \(a=1\)
(chain \([5,2]\)), \(v_1=2e_1+e_2\), \(v_2=-e_2-e_3\) is primitive in
\(-I_3\). It is realized smoothly in \(\#3\overline{\mathbb{CP}}^2\) by
the spheres \(2E_1+E_2\) (conic \(\#\) line) and \(E_2+E_3\) (tube),
meeting once. If the rational blowdown
\(Z=(\#3\overline{\mathbb{CP}}^2)_{C_{9,2}}\) is diffeomorphic to
\(\overline{\mathbb{CP}}^2\) (it is a closed negative-definite \(b_2=1\)
manifold with \(H_1=0\); simple connectivity of the complement is still
to be checked), then \(M=X\#Z\) carries the zero-charge configuration.
The bulk is the orbifold \(V=X\#\mathcal O_{9,2}\), rationally \(M\),
with one \(\tfrac19(1,2)\) point and \(T_+^{\rm orb}=A-1\); the bubble
is the Ricci-flat \(B_{3,1}\) at cost \(0\).

\textbf{Proposition 55.6 (the zero-charge bulk is non-Kähler).} Suppose
the orbifold \(V\) of the zero-charge channel were a complex orbifold,
with the T-point a complex quotient singularity. Then its minimal
resolution \(Y\) is a compact complex surface containing the holomorphic
chain \(C\). Moreover \(\pi_1(Y)=\Gamma\), and \(Y\) has a degree-one
map to \(X\) whose kernel lattice is \(K\) of rank \(r+1\). \(Y\) has
\(b^+\ge3\) and nonzero SW invariants (rational blowdown formula from
\(SW_M\ne0\)), so it is of general type. Its minimal model \(X'\)
satisfies \(b_2(Y)-b_2(X')=r+1\) (the blow-down classes span \(K\)).
Every curve whose class lies in \(K=\ker c_*\) is contracted to points
by \(Y\to X'\), so \(C\) lies in the exceptional locus of \(r+1\)
blow-ups. \emph{Counting:} in a sequence of \(r+1\) point blow-ups each
new exceptional curve starts with square \(-1\), and each later centre
lies on at most two earlier exceptional curves. Hence
\(\sum_{\text{all }r+1\text{ components}}(-D_j^2)\le(r+1)+2r=3r+1\). A
T-chain (\(d=1\)) has \(\sum b_i=3r+1\): it is generated from \([4]\) by
\([b_1,\dots,b_r]\mapsto[2,b_1,\dots,b_r+1]\) or
\([b_1+1,\dots,b_r,2]\), each step adding \(3\). The remaining
\((r+1)\)-st component contributes \(\ge1\), so \(3r+2\le3r+1\), a
contradiction. \(\square\)

So in the Kähler category this channel cannot reach \(M\). Instead, a KE
orbifold with T-points has \(\mathbb Q\)-Gorenstein smoothings
(Biquard--Rollin type, citation to register), which are minimal surfaces
with two SW basic classes rather than the four of \(M\). The \(b_2=1\)
channels of §3 are not covered by this counting argument.

\begin{center}\rule{0.5\linewidth}{0.5pt}\end{center}

\hypertarget{the-reduced-k1-theorem-orbifold-part}{%
\subsection{\texorpdfstring{5. The reduced \(k=1\) theorem (orbifold
part)}{5. The reduced k=1 theorem (orbifold part)}}\label{the-reduced-k1-theorem-orbifold-part}}

\textbf{Theorem 55.7.} Let \(X\) be a ball quotient other than a fake
projective plane, and \(M=X\#\overline{\mathbb{CP}}^2\). In a case-(I)
minimizing configuration, assume \(\mathrm H_{\rm lsc}\),
\(\mathrm H_{\rm dec}\) (any group), (SW\(_L\)) and bubble criticality.
Then exactly one of the following holds:

\begin{enumerate}
\def\labelenumi{(\Alph{enumi})}
\item
  every cone point has trivial group, and Theorem 54.A applies
  (attainment / Einstein-or-non-Bach-flat, modulo exotic \(S^4\) caps
  and \(\mathrm H_{\rm eq},\mathrm H_{\rm rem}\));
\item
  there is one orbifold point whose bubble region carries \(mE\), with
  \((m,\Gamma)\) in the explicit charge-allowed list of Prop. 55.3 (so
  \(m\ge3\) for cyclic ends; \(m=2\) only with \(D_5,D_9\) ends; \(m=1\)
  only with the \(E_8\) end);
\item
  all orbifold points are rational-ball bubbles with \(\max c_\pm<1\).
  Each is a T-type zero-charge ball with primitive chain embedding (Thm
  55.5; \(n\) odd or \(8\mid n\)), a negatively charged ball
  (e.g.~reversed T-type), or a non-T ball from the JSON list.
\end{enumerate}

In (B) and (C) the bulk is a non-Kähler \(P\)-critical orbifold with \[
p_*(M)\ \ge\ 8\pi^2\big[C_+(V)+\textstyle\sum w(N_j)-(A-1)\big],
\] and, if every bubble in the configuration has \(P=0\) (the Ricci-flat
T-balls), reassembly gives \(p_*(M)=\alpha_{\rm orb}(V)=P(V,g_0)\),
attained on the orbifold.

So the orbifold gate at \(k=1\) is equivalent to the following
\textbf{orbifold threshold problem}: show
\(\alpha_{\rm orb}(V)\ge8\pi^2\bigl(1-\max\{c_+,c_-,0\}\bigr)\) (more
precisely \(C_+(V)+w\ge A\)) for the explicit non-Kähler orbifold
degenerations \(V\) of \(M\) listed above, or construct one that beats
the analytic CH cell. That construction would be a genuine orbifold
absorption of \(E\) (or \(3E\), \ldots) and a new nonattainment
mechanism.

What this does \textbf{not} do: it gives no unrestricted lower bound for
\(P\) on orbifolds. LeBrun's \(2A/3\) is too weak except for \(A\le3w\).
The Kähler orbifold bound \(C_+\ge T_+^{\rm orb}\) is sharp (KE
orbifolds) but irrelevant, since the bulks are non-Kähler.

\begin{center}\rule{0.5\linewidth}{0.5pt}\end{center}

\hypertarget{collapse-gate-what-the-hyperbolicity-of-gamma-gives}{%
\subsection{\texorpdfstring{6. Collapse gate: what the hyperbolicity of
\(\Gamma\)
gives}{6. Collapse gate: what the hyperbolicity of \textbackslash Gamma gives}}\label{collapse-gate-what-the-hyperbolicity-of-gamma-gives}}

\textbf{Lemma 55.8 (cyclic-image lemma).} Let \(f:\Sigma\to M\) be a map
of a closed oriented surface whose composite \(\pi_1\)-image in
\(\Gamma\) is cyclic (e.g.~\(\Sigma=S^2\) or \(T^2\), since abelian
subgroups of the torsion-free hyperbolic group \(\Gamma\) are cyclic).
Then \(f_*[\Sigma]\in\operatorname{span}\{E_i\}\). More generally, every
class in the image of \(H_2(C)\to H_2(M)\) lies in
\(\operatorname{span}E_i\) for any region \(C\) whose \(\pi_1\)-image in
\(\Gamma\) is cyclic.

\emph{Proof.} \(X=K(\Gamma,1)\), so \(c\circ f\) is determined up to
homotopy by \(\pi_1\). A cyclic image factors through
\(K(\mathbb Z,1)=S^1\), so \((c\circ f)_*[\Sigma]=0\), and
\(\ker c_*=\operatorname{span}E_i\). (The bounded-cohomology route, in
which \(H^2_b\to H^2\) is onto for hyperbolic \(\Gamma\) and the
\(\ell^1\)-seminorm is a norm on \(H_2(\Gamma;\mathbb R)\), gives the
same for amenable images, which are cyclic here. It would add genuinely
new information only through Gromov's multiplicity-\(\le2\) vanishing,
which \(F\)-structure covers do not satisfy in general.) \(\square\)

\textbf{What this settles and what it does not.}

\begin{itemize}
\tightlist
\item
  \emph{Settled.} Combined with the step-54 free-circle spin lemma: a
  collapsing region with a free circle action, or with cyclic
  \(\pi_1\)-image, cannot carry \(E\) (\(E^2\) odd), and in the cyclic
  case carries no class of \(X\).
\item
  \emph{Not settled --- the obstruction, stated sharply.} A collapsed
  region \(C\) carrying a negative class changes the ledger in the
  \textbf{wrong} direction for an exclusion. With \(\chi(C)=0\), (1.2)
  becomes \(T_+(\Omega)=A-1-3\tau(C)\) up to boundary η-terms. If
  \(\tau(C)=-1\), then \(P(\Omega)=8\pi^2[C_+(\Omega)-(A+2)]\) while
  \(P(C)\to0\) under bounded-curvature collapse. Charge accounting
  therefore forces no positive cost on the configuration. The whole
  burden moves to whether a complete finite-volume \(P\)-critical
  \(\Omega\) with collapsing ends can have \(C_+(\Omega)<A+3\) (to beat
  (U)), a regime where no estimate is available. Collapse is thus the
  one degeneration in which \(E\) can disappear from the topological
  ledger at zero bubble energy. It must be excluded topologically: show
  that no region with a positive-rank \(F\)-structure (finite isotropy
  allowed, no fixed points) contains a class of odd square. The
  supporting facts are that fixed-point-free \(S^1\)-actions on closed
  4-manifolds have \(\tau=0\) (\(G\)-signature) and that free ones are
  spin. The missing pieces are the Seifert/finite-isotropy and
  torus-isotropy cases and the boundary η-terms of graph-manifold
  interfaces.
\end{itemize}

\begin{center}\rule{0.5\linewidth}{0.5pt}\end{center}

\hypertarget{remarks}{%
\subsection{7. Remarks}\label{remarks}}

\begin{enumerate}
\def\labelenumi{\arabic{enumi}.}
\tightlist
\item
  \emph{Einstein metrics on \(M\) sit far from the Yamabe-optimal
  classes} (strategist's remark, verified). Theorem Y at \(P=0\) gives
  \(8\pi^2k\le4\Delta V\), with \(\Delta V=(Y[g]^2-32\pi^2A)/16\). Hence
  \(Y[g_E]^2\ge32\pi^2(A+k)\), while \(\inf_{[g]}Y[g]^2=32\pi^2A\)
  (LeBrun). For \(k=1\): \(Y[g_E]^2\ge32\pi^2(A+1)\).
\item
  \emph{Correction to the strategist's \#70.} ``Every bubble has
  \(H_2\subset\langle E\rangle\)'' needs (SW\(_L\)); without it a bubble
  could carry \(\Gamma\). ``Kronheimer excluded outright at \(k=1\)'' is
  right. ``What survives with \(\Gamma\ne1\) is a finite list'' is right
  for each fixed \(|\Gamma|\), but the list is infinite in \(m\) and
  contains non-cyclic \(D_5,D_9,E_8\) ends. The \(\mathcal O(-4)\)
  candidate is excluded (charge exactly 1). The rational-ball case does
  survive, but not as \(\mathrm{EH}/\mathbb Z_2\) in the zero-charge
  orientation; its first zero-charge member is \(B_{3,1}\).
\item
  \emph{Why charge alone cannot finish.} Conservation (1.3) shows every
  channel has total Kähler-type baseline exactly that of \(M\). The
  discriminating information must come from a non-Kähler orbifold
  estimate or from the CH-germ interaction, exactly as the analytic cell
  beats the Kähler baseline for smooth bubbles.
\end{enumerate}

\hypertarget{s1003a02-step-56-orbifold-category-attainment-at-k-1-the-zero-energy-upper-bound-and-the-local-interaction-at-a-wahl-point}{%
\chapter{S1003A02: Step 56 --- Orbifold-category attainment at k = 1,
the zero-energy upper bound, and the local interaction at a Wahl
point}\label{s1003a02-step-56-orbifold-category-attainment-at-k-1-the-zero-energy-upper-bound-and-the-local-interaction-at-a-wahl-point}}

Source:
\texttt{step56\_orbifold\_attainment\_zero\_energy\_bound\_Tpoint\_interaction.md}.

Session S1003A02, Agent2, Execute. This note continues
\texttt{step55\_k1\_orbifold\_gate\_charges\_lattices\_collapse.md} and
makes the Step-55 verifier corrections binding (§0). It then: - states
the orbifold-category theorem that the conservation law makes stable
(§1); - proves the unconditional zero-energy upper bound (§2); - fixes
the algebraic structure of curvature at a Wahl point (§3); - computes
the leading local interaction between a P-critical bulk and a Ricci-flat
Wahl-type Milnor-fibre bubble (§4).

The result of §4 is negative for the cheap exclusion route. The
\(\lambda^4\log(1/\lambda)\) interaction vanishes after the spherical
average, so no pointwise Biquard--Ozuch-type obstruction appears at the
logarithmic order.

Scripts and data: -
\texttt{workspace/scripts/step56\_Tpoint\_local\_structure.py} →
\texttt{workspace/data/step56\_Tpoint\_local\_structure.json}
(Γ-invariant curvature, GH quadrupole phase test); -
\texttt{workspace/scripts/step56\_Tpoint\_log\_interaction.py} →
\texttt{workspace/data/step56\_Tpoint\_log\_interaction.json} (the
\(\lambda^4\log\) coefficient as a quadratic form on invariant
curvature); -
\texttt{workspace/scripts/step56\_log\_interaction\_controls.py} →
\texttt{workspace/data/step56\_log\_interaction\_controls.json} (generic
curvature and wrong-tail controls).

Sources registered by the Step-55 request (partial; read as listed): -
\textbackslash cite\{suvaina2011ale\}: ALE Ricci-flat Kähler metrics on
every \(\mathbb Q\)-Gorenstein smoothing fibre; -
\textbackslash cite\{rasdeaconu2013ale\}: class-T classification and the
explicit \((xy=P(z^n))/\mu_n\) model; -
\textbackslash cite\{lisca2007lens\}: lens spaces bounding rational
balls; - \textbackslash cite\{fintushel1997rational\}: rational blowdown
and its SW formula; - \textbackslash cite\{agolhyperbolic\}: SW
vanishing for splittings along an \(L\)-space with \(b^+\ge1\) on both
sides. This is the two-sided (SW\(_L\)) input. Lens spaces and spherical
space forms are \(L\)-spaces, so the two-sided form is now citable; -
Biquard--Rollin smoothing of class-T cscK orbifolds (Knowledge object
\texttt{Smoothing-singular-extremal-Kähler-surfaces-and-minimal-Lagrangians-\/-76bcf086d1fb}).

\begin{center}\rule{0.5\linewidth}{0.5pt}\end{center}

\hypertarget{corrections-to-step-55-now-in-force-verifier-tracettcsresultsstep55_ttcs01_verifier1_integration.md}{%
\subsection{\texorpdfstring{0. Corrections to Step 55 now in force
(verifier
\texttt{trace/ttcs/results/step55\_ttcs01\_verifier1\_integration.md})}{0. Corrections to Step 55 now in force (verifier trace/ttcs/results/step55\_ttcs01\_verifier1\_integration.md)}}\label{corrections-to-step-55-now-in-force-verifier-tracettcsresultsstep55_ttcs01_verifier1_integration.md}}

\begin{enumerate}
\def\labelenumi{\arabic{enumi}.}
\tightlist
\item
  \textbf{Retained as proved:} the charge formula (1.1),
  \(P(\mathcal O(-n),g_{\rm LeB})=8\pi^2(n-2)^2/n\) (exact, since
  \(W^+=0\)), the ledger (1.2) and the conservation identity (1.3) (sums
  over attachment regions), and Theorem 55.5. In 55.5 the \(v_2(n)=1\)
  case is obstructed mod 4, not mod 8. The Wahl-ball charges
  \((0,-4(1-1/n^2))\) hold for \textbf{all} \(n\), by the
  Hirzebruch--Jung formula with \(\sum(b_i-3)=1\) and \(q+q'=n^2-2\).
\item
  \textbf{Lemma 55.1} gives
  \(\pi_1(N)/\langle\!\langle\pi_1L\rangle\!\rangle=1\), not
  \(\pi_1(N)=1\). Prop. 55.3 is valid \textbf{only for simply connected}
  rank-one regions. Counterexample to the unrestricted cyclic-end
  exclusion:
  \(N=\bar B_{t,a}\natural(\overline{\mathbb{CP}}^2\setminus B^4)\),
  with charges \((-3+4/t^2,-5)\); for \(t=2\) the end is \(L(4,3)\) and
  the charges are \((-2,-5)\). In general \(|H_1(L)|=m^2t^2\).
\item
  \textbf{Additive budget:} under (U) and \(\mathrm H_{\rm lsc}\),
  \(\sum_j\max\{c_+(N_j),c_-(N_j),0\}<1\), and the exact threshold is
  \(C_+(V)+\sum_jw(N_j)<A\).
\item
  ``Zero charge ⇔ Wahl ball'' is \textbf{not} proved. Charge does not
  determine the filling, and the rational-ball table is hypothetical
  charge arithmetic.
\item
  \textbf{Prop. 55.6, repaired proof} (replaces the SW/general-type
  argument):

  \begin{itemize}
  \tightlist
  \item
    \(\pi_1(Y)=\Gamma\), so \(b_1(Y)\) is even and \(Y\) is Kähler;
  \item
    the degree-one map \(Y\to X\) is homotopic to a harmonic map, which
    is \(\pm\)holomorphic by Siu rigidity;
  \item
    so the map is birational and \(Y\) is \(X\) blown up \(r+1\) times;
  \item
    since \(X\) has no rational curves, the chain is exceptional, and
    the counting argument \(3r+2\le3r+1\) applies. Scope: only the
    complex-orientation zero-charge Wahl channel with a complex-orbifold
    bulk.
  \end{itemize}
\item
  \textbf{Collapse ledger (§6 of Step 55) withdrawn.} For a
  bounded-curvature collapsing region \(C\) with volume \(\to0\),
  \(P(g_i)/8\pi^2=C_+(\Omega_i)-(A-1)+o(1)\). The threshold remains
  \(C_+(\Omega_i)<A\), because boundary terms compensate \(\tau(C)\).
  What survives is that collapsed regions cost nothing energetically and
  must be excluded topologically. The free-circle spin lemma applies;
  the cyclic-image lemma only puts classes in \(\mathbb ZE\) and does
  \textbf{not} exclude \(E\) (a punctured \(\overline{\mathbb{CP}}^2\)
  is simply connected).
\item
  \textbf{Theorem 55.7 becomes a coarse trichotomy.} In a case-(I)
  configuration, under (SW\(_L\)) and its one-sided form, exactly one
  holds:

  \begin{itemize}
  \item
    \begin{enumerate}
    \def\labelenumii{(\Alph{enumii})}
    \tightlist
    \item
      no cone point with nontrivial group;
    \end{enumerate}
  \item
    \begin{enumerate}
    \def\labelenumii{(\Alph{enumii})}
    \setcounter{enumii}{1}
    \tightlist
    \item
      exactly one nontrivial-group attachment region with \(b_2=1\)
      (image in \(\mathbb QE\)), all others rational balls;
    \end{enumerate}
  \item
    \begin{enumerate}
    \def\labelenumii{(\Alph{enumii})}
    \setcounter{enumii}{2}
    \tightlist
    \item
      nontrivial-group points exist and all their regions are rational
      balls. In every case the additive budget of item 3 applies.
    \end{enumerate}
  \end{itemize}
\item
  \textbf{The \([5,2]\) model has simply connected complement}
  (auxiliary-line argument). So
  \(Z=(\#3\overline{\mathbb{CP}}^2)_{C_{9,2}}\) is simply connected with
  form \(\langle-1\rangle\), hence \textbf{homeomorphic} to
  \(\overline{\mathbb{CP}}^2\). A diffeomorphism is still unproved; §2
  is therefore stated for the smooth manifold \(X\#Z\).
\item
  Inherited typo in Step 54: reversing Burns exchanges \((1,-5)\).
\end{enumerate}

\begin{center}\rule{0.5\linewidth}{0.5pt}\end{center}

\hypertarget{the-orbifold-category-theorem}{%
\subsection{1. The orbifold-category
theorem}\label{the-orbifold-category-theorem}}

\textbf{Conservation law, reread.} Identity (1.3) of Step 55 says that
bubbling never changes the Kähler-type baseline \(T_+(M)=A-1\). It only
moves self-dual Weyl charge between bulk and bubbles. So any exclusion
of the orbifold channels must use a lower bound on a non-Kähler bulk.
For orbifolds this is the same open-range inequality \(C_+\ge A\) that
is open on \(M\) itself; LeBrun's \(2A/3\) fails there for the same
reason. Accordingly the orbifold channels are not treated as a lemma to
be proved. They are carried into the theorem statement.

\textbf{Theorem 56.A (attainment in the orbifold category; \(k=1\)).}
Let \(X\) be a compact ball quotient other than a fake projective plane,
and \(M=X\#\overline{\mathbb{CP}}^2\). Consider a case-(I) Anderson
configuration for \(p_*(M)\).

Assume: - the step-54 ledger \(\mathrm H_{\rm lsc}\) and
\(\mathrm H_{\rm dec}\) (any group); - bubble criticality by compact
replacement; - (SW\(_L\)): the two-sided form is
\textbackslash cite\{agolhyperbolic\}; the one-sided PSC-cap form
remains a hypothesis; - smooth orbifold regularity of the bulk \(g_0\);
- \textbf{(N\(_0\))}: every bubble region at a point with nontrivial
group has \(P=0\), i.e.~is Ricci-flat ALE.

Then:

\begin{enumerate}
\def\labelenumi{(\roman{enumi})}
\item
  \(p_*(M)=P(V,g_0)=\alpha_{\rm orb}(V)\), where \(V\) is the orbifold
  bulk. Either \(V=M\) (no orbifold points), or every orbifold point is
  a cyclic point whose bubble is a Ricci-flat ALE rational ball. In the
  latter case the configuration is case (C) of §0.7, with
  \(c_+(N_j)\le0\).
\item
  \(g_0\) is \(P\)-critical on \(V\), with constant scalar curvature
  \(s_0<0\).
\item
  Exactly one of the following holds:
\end{enumerate}

\begin{itemize}
\tightlist
\item
  \(g_0\) is an Einstein orbifold metric. Then \(p_*(M)=0\), so
  \(\alpha(M)=0\); if \(V\ne M\) the infimum is attained only if \(M\)
  carries an Einstein metric.
\item
  \(g_0\) is not Bach-flat. This is Problem 3 in the orbifold category.
\end{itemize}

\begin{enumerate}
\def\labelenumi{(\roman{enumi})}
\setcounter{enumi}{3}
\tightlist
\item
  If \(V\ne M\) and some orbifold point is a Wahl point
  \(\tfrac1{n^2}(1,na-1)\) with its Milnor fibre in the complex
  orientation, then:
\end{enumerate}

\begin{itemize}
\tightlist
\item
  \(n\) is odd or \(8\mid n\) (Thm 55.5);
\item
  \(V\) is not a complex orbifold with that complex quotient singularity
  (Prop 55.6, repaired);
\item
  the minimal resolution \(Y\) has \(\pi_1=\Gamma\), and \(Y\to X\) has
  degree one with a unimodular kernel lattice containing the chain
  primitively.
\end{itemize}

\emph{Proof.} - (i) Conservation (1.3) with \(P(N_j)=0\) gives
\(p_*(M)=P(V,g_0)\) under \(\mathrm H_{\rm lsc}\). - Zero-energy
reassembly (verifier §10, made precise in §2 below) gives
\(p_*(M)\le P(V,g)\) for every smooth orbifold metric \(g\) on \(V\).
Hence \(P(V,g_0)=\alpha_{\rm orb}(V)\). - Ricci-flat ALE spaces have
\(c_+=P/8\pi^2-w=-w\le0\). A Ricci-flat ALE with \(b_2=1\) in case (B)
would carry a class of square \(-m^2\) on an Einstein ALE space with
\(P=0\). We do not exclude this in general. (N\(_0\)) together with case
(B) is therefore kept in the statement as part of case (i)'s
alternative, unless the end is simply connected cyclic with \(m\le2\),
where \(c_+=1>0\) is impossible. - (ii) Criticality on \(V\) is the
minimizing property of \(g_0\) in \(\alpha_{\rm orb}(V)\). The trace of
the Euler--Lagrange equation is \(\Delta s=0\), so \(s\) is constant. -
If \(s_0\ge0\), reassembly produces metrics on \(M\) with
\(\int s_-^2\to0\), because the bubbles are scalar-flat and the gluing
error vanishes. This contradicts LeBrun's estimate
\(\int s^2\ge\int s_-^2\ge32\pi^2c_1^2(X)\), which holds for every
metric on \(M\) since \(SW_M\ne0\) and the basic class is
\(\pm K_X\pm E\). - (iii) The proof of Step 53 applies verbatim on
orbifolds (local, away from finitely many points, plus removability at
the cone points for a smooth orbifold metric): Bach-flat +
\(P\)-critical + constant \(s\ne0\) ⇒ Einstein. If \(g_0\) is Einstein,
reassembly gives \(P\to0\). - (iv) Step 55 with the repairs of §0.
\(\square\)

What 56.A does \textbf{not} assume away: - the case-(B) channels with
non-simply-connected regions or noncyclic ends; - non-Ricci-flat
rational-ball bubbles of positive cost satisfying the additive budget; -
exotic-homotopy-sphere bubbles; - collapse. These remain hypotheses
through (N\(_0\)) and case (I).

\textbf{Corollary 56.B (orbifold Einstein metrics give zero gap).} If
some compact orbifold \(V\) with only Wahl points, whose Milnor-fibre
smoothing is diffeomorphic to \(M\), carries an Einstein orbifold
metric, then \(\alpha(M)=0\). This is Prop. 56.U applied at
\(P(V,g)=0\). - If the Wahl points are complex quotient points of a
complex orbifold structure on \(V\), the metric cannot be Kähler
(56.A(iv)). - For comparison, Theorem Y at \(P=0\) shows that a smooth
Einstein metric on \(M\) would need \(Y^2\ge32\pi^2(A+1)\). No orbifold
analogue is claimed.

This is the \(k=1\) analogue of the Step-52 Kronheimer absorption,
restricted by the lattice law (\(n\) odd or \(8\mid n\)). Absorption
through case-(B) regions is not covered by this corollary.

\begin{center}\rule{0.5\linewidth}{0.5pt}\end{center}

\hypertarget{the-unconditional-zero-energy-upper-bound}{%
\subsection{2. The unconditional zero-energy upper
bound}\label{the-unconditional-zero-energy-upper-bound}}

\textbf{Proposition 56.U.} Let \(V\) be a compact oriented 4-orbifold
with isolated cyclic points \(p_j\in V\),
\(\Gamma_j=\tfrac1{n_j^2}(1,n_ja_j-1)\). Let \(M'\) be the smooth
manifold obtained by replacing a neighbourhood of each \(p_j\) by the
Milnor fibre \(B_{n_j,a_j}\), glued in the complex orientation of the
local chart. Then for every smooth orbifold metric \(g\) on \(V\), \[
p_*(M')\le P(V,g).
\]

\emph{Proof.} 1. Each \(B_{n,a}\) carries a Ricci-flat Kähler ALE metric
\(g_B\) asymptotic to \(\mathbb C^2/\Gamma\) at rate \(O(r^{-4})\) with
all derivatives \textbackslash cite\{suvaina2011ale,rasdeaconu2013ale\}.
In GH form it is the free \(\mathbb Z_n\) quotient of the \(A_{n-1}\)
metric with centres on a \(\mathbb Z_n\)-orbit. 2. In an orbifold chart,
\(g=g_{\rm flat}+O(r^2)\), and the \(O(r^2)\) term is a smooth
\(\Gamma\)-invariant tensor. 3. Glue \(\lambda^2g_B\) to \(g\) on the
annulus \(\rho\le r\le2\rho\) by a cutoff. The interpolation error
tensor is \(O(\rho^2)+O((\lambda/\rho)^4)\) with \(j\)-th derivatives
scaled by \(\rho^{-j}\). The annular volume is \(O(\rho^4)\), so the
annular \(P\)-error is \(O(\rho^4)+O((\lambda/\rho)^8)\). 4. Inside
\(r<\rho\), \(P(\lambda^2g_B)=0\). The deleted orbifold region carries
\(\int_{r<\rho}|{\rm Ric}^\circ_g|^2=O(\rho^4)\). 5. With
\(\rho=\lambda^{1/2}\) the total error is \(O(\lambda^2)\to0\).
\(\square\)

So every orbifold degeneration through Wahl points is a competitor, and
any construction of a cheap orbifold cell bounds the smooth problem from
above. For \(V=X\#\mathcal O_{9,2}\) with \(\mathcal O_{9,2}=Z\) with
\(B_{3,1}\) collapsed (\(Z\) from §0.8), \(M'=X\#Z\). This is
homeomorphic to \(M\); a diffeomorphism follows if
\(Z\cong\overline{\mathbb{CP}}^2\).

Because the gluing in 56.U is \(O(\lambda^2)\) and not more precise, the
first-order interaction between \(g\) and the bubble is not seen. §4
computes it.

\begin{center}\rule{0.5\linewidth}{0.5pt}\end{center}

\hypertarget{curvature-at-a-wahl-point-is-forced-into-a-4-parameter-type}{%
\subsection{3. Curvature at a Wahl point is forced into a 4-parameter
type}\label{curvature-at-a-wahl-point-is-forced-into-a-4-parameter-type}}

Let \(\Gamma=\tfrac1{N}(1,q)\) with \(N=n^2\), \(q=na-1\), \(n\ge3\). In
complex coordinates of the chart: -
\(\Lambda^+=\mathbb R\omega\oplus\mathbb C\,dz_1\wedge dz_2\), with
characters \(1,\zeta^{1+q}=e^{2\pi ia/n}\); -
\(\Lambda^-=\mathbb R\bar\omega\oplus\mathbb C\,dz_1\wedge d\bar z_2\),
with characters \(1,\zeta^{1-q}\).

None of \(\zeta^{1\pm q},\zeta^{2(1\pm q)},\zeta^{2q}\) is trivial for
\(n\ge3\) (script, all Wahl points with \(n\le9\)). Hence the
Γ-invariant algebraic curvature tensors form a 4-dimensional space: \[
R_p=\Bigl(\tfrac{s}{12}+w_+D\Bigr)\oplus\Bigl(\tfrac{s}{12}+w_-D\Bigr),\qquad B=b\,\hat\omega\otimes\hat{\bar\omega},
\qquad D=\operatorname{diag}(2,-1,-1).
\] Here \(W^\pm\) are degenerate (type D) with the distinguished
eigenvector \(\omega\) (resp. \(\bar\omega\)), and
\(\mathrm{Ric}^\circ_p\propto b\,(\omega\circ\bar\omega)\), which is
\(J\)-invariant with eigenvalues \((b',b',-b',-b')\) on \(z_1,z_2\).

For \(n=2\) (\(\tfrac14(1,1)\)) the space is 12-dimensional. It is
irrelevant here, since \(n=2\) is excluded by the lattice law.

So at every Wahl point the curvature of \textbf{any} smooth orbifold
metric is pointwise of ambi-Kähler algebraic type, with the Kähler forms
of both orientations fixed by Γ. For a \(P\)-critical bulk this is the
only local datum the bubble can couple to at second order.

\textbf{Bubble moduli.} For the \(\mathbb Z_n\)-symmetric GH
configuration of \(n\ge3\) centres on a circle of radius \(R_c\) around
the \(\omega\)-axis, the leading tail is the quadrupole
\(\delta V=-(n^2R_c^2/4)\,P_2(\cos\theta)/r^3\) on the cover. It is
\textbf{independent of the phase} of the configuration (script:
\(n=3,\dots,8\); dependent only for \(n=2\)). The axial offset is pure
gauge. So the only modulus that affects the leading interaction is the
scale \(\lambda\) (\(\lambda^4\propto R_c^2\)). The relocation and frame
freedom used in Theorem 53.B is absent: the frame is fixed by Γ, and the
normaliser torus acts trivially on invariant curvature.

\begin{center}\rule{0.5\linewidth}{0.5pt}\end{center}

\hypertarget{the-lambda4log1lambda-interaction-vanishes}{%
\subsection{\texorpdfstring{4. The \(\lambda^4\log(1/\lambda)\)
interaction
vanishes}{4. The \textbackslash lambda\^{}4\textbackslash log(1/\textbackslash lambda) interaction vanishes}}\label{the-lambda4log1lambda-interaction-vanishes}}

\textbf{Setup.} On the universal cover \(\mathbb C^2\) of the chart: -
the bulk is \(g_0=\delta+q\) in normal coordinates,
\(q_{ij}=-\tfrac13R_{ikjl}y^ky^l\); - the bubble tail is \(t\),
homogeneous of degree \(-4\), obtained from the GH quadrupole in the
Hopf/moment coordinates of the triholomorphic circle
\((e^{it}y_1,e^{-it}y_2)\). The axis is the \(\omega\)-direction,
matching §3.

Script checks: - the flat GH structure reproduces \(\delta\) exactly; -
normal coordinates reproduce \(\mathrm{Ric}_p\); - the tail is
linearized Ricci-flat (\(|{\rm Ric}'(t)|\sim10^{-7}\), round-off level).
The opposite orientation of \(A_2\) gives \(|{\rm Ric}'|\sim60\), so the
sign is fixed.

In the overlap \(\lambda\ll r\ll1\) the glued metric is
\(\delta+q+\lambda^4t+\dots\). Since \({\rm Ric}'(t)=0\), the first term
of \(P\) mixing bulk and bubble is the \(e^2m\) coefficient of
\(|{\rm Ric}^\circ|^2d\mu\) for \(\delta+eq+mt\). It is homogeneous of
degree \(-4\), so \[
P(\text{glued})=P(V,g_0)+\lambda^4\log(1/\lambda)\,\frac{c_n}{|\Gamma|}\int_{S^3}c(\hat y)\,d\sigma+O(\lambda^4).
\] The log coefficient is gauge invariant: a degree \(-3\)
Lie-derivative change of \(t\) changes the density by a divergence with
no log.

\textbf{Result.} - The pointwise coefficient \(c(\hat y)\) is nonzero
and \(T^2\)-invariant. For example, at \((s,w_+,w_-,b)=(1,.5,-.7,.9)\),
\(c=-28.98\) at \(\chi=0.6\), the same for three torus angles. - Its
\(S^3\) integral, computed as a quadratic form on the 4-parameter
invariant curvature (48-point Gauss--Legendre in \(\chi\), Richardson
extrapolation in both perturbation parameters), is \[
  \int_{S^3}c\,d\sigma=0\qquad(\text{all 10 matrix entries }\lesssim2\times10^{-8},\ \text{against pointwise values }O(10)).
  \] - With \(k=2\times10^{-3}\) and no extrapolation, the entries were
\(O(1)\). That was \(m^3\) contamination, now removed.

Control runs (generic non-invariant \(R_p\); wrong-sign tail) are
recorded in §4.1. The vanishing persists in all of them. §4.1 also gives
the structural reason (Lemma 56.L).

\textbf{Consequence.} At the logarithmic order a Ricci-flat Wahl-ball
bubble does not see the bulk curvature at the singular point. There is
no analogue of the Burns kernel-direction saving \(-9x_0^2\log(1/x_0)\),
and no analogue of a Biquard--Ozuch condition
\(\mathbf R^+(\omega,\omega)=0\) at this order. The first interaction is
at order \(\lambda^4\) without a log. Its coefficient is a renormalized
local quadratic form \(\mathcal Q_\Gamma(R_p)\), the
\(P\)-polarizability of \(B_{n,a}\). It is obtained by minimizing the
second variation of \(P\) on the bubble over perturbations with
prescribed growth \(R_p\,r^2\), after subtracting the bulk's own
\(\int|{\rm Ric}^\circ_p|^2\) volume term (the log term vanishes, as
shown).

Minimality of the configuration in 56.A forces
\(\mathcal Q_\Gamma(R_p)\ge0\) at each Wahl point. Bulk nonlocal
responses enter only at \(O(\lambda^8)\), because the bubble's source in
the bulk is \(O(\lambda^4)\) and \(g_0\) is critical. Deciding whether
\(\mathcal Q_\Gamma\) is indefinite on the 4-dimensional space requires
solving a fourth-order \(\mathbb Z_n\times U(1)\)-symmetric problem on
the GH space. If it is indefinite, it gives a pointwise constraint at
every Wahl point of an orbifold minimizer. If it is positive
semidefinite, the local mechanism cannot exclude the T-channel, and the
decision returns to the global threshold or to a construction.

\hypertarget{controls-and-why-the-vanishing-is-structural}{%
\subsubsection{4.1 Controls, and why the vanishing is
structural}\label{controls-and-why-the-vanishing-is-structural}}

The control runs (\texttt{step56\_log\_interaction\_controls.json};
24×12×12 product quadrature over \(S^3\); each pair is (signed integral,
integral of the absolute value)):

\begin{longtable}[]{@{}llll@{}}
\toprule
\begin{minipage}[b]{0.22\columnwidth}\raggedright
bulk \(R_p\)\strut
\end{minipage} & \begin{minipage}[b]{0.22\columnwidth}\raggedright
tail\strut
\end{minipage} & \begin{minipage}[b]{0.22\columnwidth}\raggedright
\(\int_{S^3}c\)\strut
\end{minipage} & \begin{minipage}[b]{0.22\columnwidth}\raggedright
\(\int_{S^3}|c|\)\strut
\end{minipage}\tabularnewline
\midrule
\endhead
\begin{minipage}[t]{0.22\columnwidth}\raggedright
generic (random Bianchi tensor, not Γ-invariant)\strut
\end{minipage} & \begin{minipage}[t]{0.22\columnwidth}\raggedright
Ricci-flat (correct)\strut
\end{minipage} & \begin{minipage}[t]{0.22\columnwidth}\raggedright
\(5.5\times10^{-8}\)\strut
\end{minipage} & \begin{minipage}[t]{0.22\columnwidth}\raggedright
2484\strut
\end{minipage}\tabularnewline
\begin{minipage}[t]{0.22\columnwidth}\raggedright
invariant\strut
\end{minipage} & \begin{minipage}[t]{0.22\columnwidth}\raggedright
wrong-sign \(A_2\) (not linearized Ricci-flat)\strut
\end{minipage} & \begin{minipage}[t]{0.22\columnwidth}\raggedright
\(-1.4\times10^{-7}\)\strut
\end{minipage} & \begin{minipage}[t]{0.22\columnwidth}\raggedright
3086\strut
\end{minipage}\tabularnewline
\begin{minipage}[t]{0.22\columnwidth}\raggedright
generic\strut
\end{minipage} & \begin{minipage}[t]{0.22\columnwidth}\raggedright
wrong-sign\strut
\end{minipage} & \begin{minipage}[t]{0.22\columnwidth}\raggedright
\(-9.2\times10^{-7}\)\strut
\end{minipage} & \begin{minipage}[t]{0.22\columnwidth}\raggedright
6082\strut
\end{minipage}\tabularnewline
\begin{minipage}[t]{0.22\columnwidth}\raggedright
invariant\strut
\end{minipage} & \begin{minipage}[t]{0.22\columnwidth}\raggedright
correct\strut
\end{minipage} & \begin{minipage}[t]{0.22\columnwidth}\raggedright
\(1.7\times10^{-8}\)\strut
\end{minipage} & \begin{minipage}[t]{0.22\columnwidth}\raggedright
467\strut
\end{minipage}\tabularnewline
\bottomrule
\end{longtable}

So the vanishing holds for arbitrary curvature and does not use
Ricci-flatness of the tail. Its proof is as follows.

\textbf{Lemma 56.L (no logarithmic linear interaction).} - Let \(g_0\)
be any smooth metric near \(p\), and let \(t\) be a symmetric 2-tensor
homogeneous of degree \(-4\) on the cone. - The part of the
\(P\)-density linear in \(t\) is
\(\langle\nabla P(g_0),t\rangle\,d\mu+d\iota(B(g_0;t))\), with \(B\)
bilinear in jets. - The divergence term contributes, on the annulus
\(\{\lambda<r<1\}\), only sphere fluxes. Each homogeneous piece of
degree \(d\) gives \(\rho^{d+3}\), and a degree \(-3\) piece gives a
constant; neither produces \(\log\rho\). - Hence the
\(\lambda^4\log(1/\lambda)\) coefficient equals
\(\int_{S^3/\Gamma}\langle[\nabla P(g_0)]_0,t\rangle\), where
\([\cdot]_0\) is the degree-0 Taylor coefficient at \(p\).

Two consequences: (a) If \(g_0\) is \(P\)-critical near \(p\), then
\(\nabla P(g_0)\equiv0\) and the coefficient is \(0\) for \textbf{every}
tail. (b) For the truncated metric \(\delta+q\) used in the script,
\([\nabla P]_0=E_2(R_p,R_p)\) is a constant tensor, and the coefficient
is \(\langle E_2,\int_{S^3}t\rangle\). For any tail built equivariantly
from a quadrupole \(Q\in\mathrm{Sym}^2_0\Lambda^+\cong(2,0)\), the
average \(\int_{S^3}t_{ij}\) lies in
\(\mathrm{Sym}^2T^*=(0,0)\oplus(1,1)\), which contains no \((2,0)\). So
the average is \(0\). This explains the controls, including the
wrong-sign tail.

The numerics therefore confirm Lemma 56.L: the script measures exactly
the term the lemma isolates. For the actual problem, (a) is the
operative statement and needs no computation. \textbf{At a Wahl point of
a \(P\)-critical orbifold bulk there is no \(\lambda^4\log\)
interaction.}

\textbf{What the first interaction is.} With the log term gone, the
\(\lambda^4\) coefficient has two local parts: 1. \textbf{Bulk flux.}
The constant (degree \(-3\)) part of \(B(g_0;t)\) pairs the tail's
quadrupole with the degree-1 jets of \(g_0\)'s curvature
(\(\nabla\mathrm{Rm}_p\), Γ-invariant part). This term is linear in the
bubble. 2. \textbf{Bubble polarizability.} This is the renormalized
minimum of the second variation of \(P\) on \(B_{n,a}\) along
perturbations growing like \(R_p\,r^2\). It is quadratic in \(R_p\) and
nonlinear in the bubble.

The nonlocal bulk response is \(O(\lambda^8)\). Minimality in 56.A gives
\(\mathcal Q_\Gamma(\nabla\mathrm{Rm}_p,R_p)\ge0\). The sign structure
of this form decides whether any pointwise mechanism exists. Computing
it needs a \(\mathbb Z_n\times U(1)\)-symmetric fourth-order problem on
the GH space. This is a finite, well-posed computation, but it was not
done in this Step.

\textbf{Judgement.} The strategist's hoped-for analogue of
Biquard--Ozuch is not available at the order where Burns and EH
analogues live: - For Burns, the mass tail \(r^{-2}\) gives an
order-\(\lambda^2\) term linear in \(R^+\). - For Ricci-flat bubbles,
criticality kills the log term identically.

So the zero-charge Wahl channel at \(k=1\) cannot be closed by a cheap
local sign. It must be decided by the polarizability form, by a global
non-Kähler estimate (§1), or by a construction (§2, 56.U).

\hypertarget{s1003a02-step-57-the-kernel-direction-second-order-burns-coefficient-in-an-arbitrary-background}{%
\chapter{S1003A02: Step 57 --- The kernel-direction second-order Burns
coefficient in an arbitrary
background}\label{s1003a02-step-57-the-kernel-direction-second-order-burns-coefficient-in-an-arbitrary-background}}

Source:
\texttt{step57\_kernel\_burns\_coefficient\_general\_background.md}.

Agent2, Execute. Status: \textbf{formal coefficient computed and
cross-checked (Step 57); D = 6 fit, density-independence check and eight
independent product-background tests completed in Step 58 (§6), all
agreeing with 57.Φ to 10⁻¹⁴; independently reproduced for KE backgrounds
by A01 (\#830). Rigorous remainder (two-scale factorization in a
nonhomogeneous background) not yet written.}

Scripts: \texttt{workspace/scripts/step57\_kernel\_log\_coefficient.py}
(core: jet-based \textbar Ric°\textbar² density, Burns tail, S³
quadrature, K bases), \texttt{step57\_run\_CH.py} (CH tests),
\texttt{step57\_general\_form.py} (20×20 curvature form; args D nu\_pts
nphi {[}ell{]}), \texttt{step57\_fit\_invariants.py} (U(2)\_ζ-invariant
fit), \texttt{step57\_norm\_check.py}, \texttt{step57\_checks.py}. Data:
\texttt{workspace/data/step57\_general\_form\_D2.json} (q6×16, includes
the ℓ(C) and N(R) computation),
\texttt{step57\_general\_form\_D2\_q8\_24.json},
\texttt{...\_D2\_q10\_32.json}, \texttt{...\_D4\_q8\_24.json},
\texttt{...\_D4\_q10\_32.json} and \texttt{*\_invfit.json}.

\hypertarget{the-formal-problem}{%
\subsection{1. The formal problem}\label{the-formal-problem}}

Near p write the background as g0 = δ + νB + ν²C (+ odd degree-3 part),
the Burns metric exactly as δ + εA with A = ∂∂̄ log\textbar z\textbar²
for the constant complex structure J\_ζ (normalised so
∂∂̄\textbar z\textbar² = δ; then g\_B is exactly Burns, linear in ε, ε =
Area(E)/π, and its Kähler form ζ ∈ Λ⁺ with the complex orientation), and
a free degree-0 correction εν K. With the two-parameter homogeneity
r\textsuperscript{\{-4\}(ε/r²)}a (μr)\^{}m of the density, the only
scale-invariant (log) terms at second order are a = 2, m = 4. Hence

P(g̃)/8π² = P(g0)/8π² + 1 − ε R⁺\_p(ζ,ζ) + (Φ/2) ε² log(1/ε) + O(ε²)
(formal),

Φ = (8π²)\^{}\{-1\} ∫\_\{S³\} {[}ε²ν² coefficient of
\textbar Ric°\textbar² dμ at r = 1{]}, minimised over K. Φ is the sum of
c\_AABB(B,B) + c\_AAC(C) + c\_ABK(B,K) + c\_KK(K); c\_KK =
\textbar Ric°\emph{lin(K)\textbar² is background independent (its
restriction to the SU(2)}-‑invariant diagonal squash reproduces
16\textbar d\textbar² − 5(Σd)² of supplement Ch. 5).

\textbf{Lemma 57.1 (no ∇Rm, no ∇²Rm).} (i) ∇Rm cannot enter (odd parity
under −1 ∈ U(1)\_ζ; ∇Rm² is m = 6). (ii) The C-linear term is ∂\_c of
the ε² part of F(δ+εA+cC) = ⟨∇P(δ+εA),C⟩ + div V. The Burns family is
P-critical for all ε, so the term is a divergence of a homogeneous
degree −3 field and has zero log coefficient. Numerically ‖ℓ‖ = 5·10⁻¹⁵
on all 350 quartic jets. Hence \textbf{Φ depends only on Rm\_p}, and no
P-criticality of g0 is needed for the log coefficient.

\hypertarget{checks}{%
\subsection{2. Checks}\label{checks}}

\begin{itemize}
\tightlist
\item
  Burns δ+εA is scalar flat (\textbar s\textbar{} \textasciitilde{}
  10⁻¹⁵) and not Ricci flat; CH in potential −log(1−\textbar x\textbar²)
  has Ric = −6g exactly.
\item
  CH perpendicular (J\_CH = J2 ⊥ J\_ζ in Λ⁺): Φ = −18.000000000000 with
  the 10-dim SU(2)\emph{- invariant K, with general polynomial K of
  degree D = 2 (100 fields) and D = 4 (350). The Ch. 5 constant-mode
  minimiser is therefore optimal among all degree-0 corrections. Optimal
  invariant K: dr²: −2, σ1²: 1/2, σ2²: 5/2, σ3²: −1 (gauge-equivalent to
  c}*).
\item
  CH parallel (ζ = ω\_CH): Φ = −18 as well (first order nonzero there).
\item
  Gauge independence: CH evaluated with normal-coordinate B instead of
  Bergman coordinates gives −18 again.
\item
  Quadrature converged (q6×16, q8×24, q10×32 identical to 10⁻¹⁴).
\end{itemize}

\hypertarget{result-the-invariant-form}{%
\subsection{3. Result: the invariant
form}\label{result-the-invariant-form}}

U(2)\emph{ζ = U(1)}+ × SU(2)\emph{- decomposition: s, w = W⁺(ζ,ζ),
W⁺}\{ζ⊥\} (charge 2), W⁺\_⊥⊥ = trace-free part of W⁺ on ζ\^{}⊥ (charge
4), W⁻ (spin 2), Ric°\emph{ζ (J\_ζ-Hermitian part, spin 1), Ric°}⊥
(J\_ζ-anti-Hermitian). Norms: operator Frobenius on Λ² (session
convention, \textbar W⁺\textbar² = s²/24 for Kähler); Ric° block B with
\textbar Ric°\textbar²\_tensor = 4\textbar B\textbar²\_op.

Fits (relative residual ≤ 6·10⁻¹⁵):

\begin{longtable}[]{@{}llll@{}}
\toprule
coefficient of & K = 0 & D = 2 & D = 4 (q8×24 and q10×32)\tabularnewline
\midrule
\endhead
s² & 1/144 & 1/288 & 1/288\tabularnewline
s w & −5/24 & −1/24 & −1/24\tabularnewline
w² & 2 & −1 & −1\tabularnewline
& W⁺\_⊥⊥ & ² & −5/9\tabularnewline
& W⁻ & ² & 7/30\tabularnewline
& B\_ζ & ² & 5/27\tabularnewline
& B\_⊥ & ² & 10/27\tabularnewline
& W⁺\_\{ζ⊥\} & ² & 14/9\tabularnewline
\bottomrule
\end{longtable}

The W⁺\emph{\{ζ⊥\} direction has a b-component in ker Q (the constant-K
frame rotation): when ζ is not an R⁺-eigenvector the first-order term
can be lowered by rotating the frame, so this entry is not a log
coefficient; at a first-order-stationary frame W⁺}\{ζ⊥\} = 0.

\textbf{Theorem 57.Φ (formal kernel coefficient).} Let ζ be an
eigenvector of R⁺\_p = W⁺\emph{p + s\_p/12 with eigenvalue μ1, other
eigenvalues μ2, μ3. Then Φ = −(w + s/12)(w − s/24) −
\textbar W⁺}⊥⊥\textbar² = −μ1(μ1 − μ2 − μ3)/2 − (μ2 − μ3)²/2. In the
kernel direction μ1 = 0: \textbf{Φ = −(μ2 − μ3)²/2 ≤ 0}, independent of
W⁻, Ric°, ∇Rm, ∇²Rm.

{[}Superseded by §7 item 1: the argument below is invalid as stated; the
chiral-square argument replaces it.{]} Why W⁻ and Ric° drop out
(independent argument): for every metric P(g̃) − P(g) − 8π² =
Δ(4∫\textbar W⁺\textbar² + ∫s²/12) (Chern--Weil). If W⁺\_p = 0, s\_p = 0
the right side is ≥ 0 up to O(ε²) for any local ASD scalar-flat
representative of the jet, so the W⁻/Ric° coefficients are ≥ 0; D = 4
attains 0. Consistency: scalar-flat Kähler blow-ups have exactly
topological cost. Kähler parallel case: Φ = −s²/32, while cscK blow-ups
have analytic P in the class parameter. Both are consistent, because the
−s²/32 saving uses non-Kähler K.

CH: μ = (0, s/4, 0), s = −24, Φ = −18. ✓

\hypertarget{consequences-formal-rigorous-version-needs-5}{%
\subsection{4. Consequences (formal; rigorous version needs
§5)}\label{consequences-formal-rigorous-version-needs-5}}

\begin{enumerate}
\def\labelenumi{\arabic{enumi}.}
\tightlist
\item
  \textbf{Every Kähler--Einstein background saves.} R⁺ = diag(s/4, 0,
  0). The kernel circle ζ ⊥ ω gives Φ = −s²/32 \textless{} 0 at
  \textbf{every} point, whatever W⁻ is. So p\_*(X\#kCP̄²) \textless{}
  8π²k for every compact KE surface (K ample), using k disjoint cells.
  For canonical models with ADE points, add Kronheimer resolutions at
  zero cost. This is the category-wide strict inequality on all minimal
  general-type surfaces asked for in Problem 2. The scale-invariant
  saving per cell is (s²/64) ε² log(1/ε), i.e.~λ²ε²log(1/ε)/4 for Ric =
  λg. The earlier D = 2 estimate, which suggested a W⁻ penalty and a
  c₁²/c₂ \textgreater{} 21/37 condition, is withdrawn.
\item
  \textbf{Minimizer rigidity (with Theorem 53.B, conditional on the
  Anderson ledger).} A P-critical bulk carrying a Burns-type bubble must
  satisfy, at every point: R⁺ ≤ 0, and if λ\_max(R⁺) = 0 then the other
  two eigenvalues are equal. In particular no bulk point can be
  Kähler--Einstein-like (eigenvalues (s/4, 0, 0) with s ≠ 0). This
  strengthens Cor. 53.C (``no exact CH region'') to all KE and
  Kähler-type degenerate R⁺. It is not an exclusion of Burns bubbling:
  R⁺ \textless{} 0 everywhere (λ\_max(W⁺) \textless{}
  \textbar s\textbar/12) is untouched, because there the first order
  already costs.
\item
  Strategist \#72's hypothesis ``negative on the whole cone R⁺ ≤ 0'' is
  answered: second order is relevant only on the boundary of the cone,
  where it is ≤ 0 and vanishes exactly at μ2 = μ3. The interior is
  decided at first order (cost −εμ1 \textgreater{} 0). W⁻ enters neither
  linearly (SU(2)\_- invariance) nor quadratically.
\end{enumerate}

\hypertarget{remaining-work-on-this-line}{%
\subsection{5. Remaining work on this
line}\label{remaining-work-on-this-line}}

\begin{itemize}
\tightlist
\item
  D = 6 run (\texttt{s57\_t21\_tool2},
  \texttt{step57\_general\_form\_D6\_q8\_24.json}): confirm s², sw, w²,
  W⁺\_⊥⊥ unchanged. The W⁻/Ric° zeros are protected below by §3.
\item
  Code check prepared but not yet run:
  \texttt{density\_Fplus\_from\_jet} (4\textbar W⁺\textbar² + s²/12)
  must give identical M, b, Meff (it differs from \textbar Ric°\textbar²
  by a local divergence).
\item
  Rigorous theorem: extend the Ch. 5 two-scale factorization (z = ε/t²,
  t, ω on S³) to an analytic background with angular dependence. a = 1
  terms are divergences or average to zero, a = 2, m \textless{} 4 terms
  are matched to the core, and the first-order inner flux is exactly
  −εR⁺(ζ,ζ). This requires an explicit K (extract the D = 4 minimiser)
  and the claim O\_R(ε²) uniformly. Theorem F's O(ε\^{}\{6/5\}) is not
  adequate. Coordinate with A01, whose Ch. 5 machinery is the template.
\item
  Batch gluing and moving backgrounds, for Problem 2's matching law on
  balanced families.
\end{itemize}

\hypertarget{step-58-completion-of-the-computation}{%
\subsection{6. Step 58 completion of the
computation}\label{step-58-completion-of-the-computation}}

\textbf{(a) D = 6 invariant fit}
(\texttt{step57\_general\_form\_D6\_q8\_24\_invfit} printed by
\texttt{step57\_fit\_invariants.py}): s² 1/288, sw −1/24, w² −1,
\textbar W⁺\emph{⊥⊥\textbar² −1, W⁻ 0 (7·10⁻¹⁷), Ric°\emph{ζ 0 (10⁻¹⁵),
Ric°}⊥ 0 (3·10⁻¹⁶), residual 5·10⁻¹⁵. Identical to D = 4. The W⁺}\{ζ⊥\}
entry moves with D (−0.185, −0.081, −0.353 at D = 2, 4, 6), confirming
that it is a kernel-leak (frame-rotation) artefact and not a
coefficient.

\textbf{(b) Density independence.} With \texttt{STEP57\_DENSITY=fplus}
the driver integrates (4\textbar W⁺\textbar² + s²/12)dμ instead of
\textbar Ric°\textbar²dμ
(\texttt{step57\_general\_form\_D\{2,4\}\_q8\_24\_fplus.json}). Raw M
agrees to 5.6·10⁻¹⁶, M\_eff to 1.7·10⁻¹⁴ (D = 2) and 7.9·10⁻¹⁵ (D = 4),
the leak to 10⁻¹⁵, Φ\_CH = −18 in both. This is the numerical form of
the fact that the two densities differ by Chern--Weil densities, whose
ε²ν² parts are divergences of degree −3 homogeneous fields (zero flux
difference between spheres), so they carry no log coefficient. It also
checks the session's operator-norm convention for W⁺ against the tensor
computation of Ric°.

\textbf{(c) Independent product backgrounds}
(\texttt{workspace/scripts/step58\_product\_tests.py}, batch
\texttt{step58\_product\_batch.sh}, data
\texttt{workspace/data/step58\_products/}). Background Σ₁(K₁)×Σ₂(K₂) in
\emph{conformal} coordinates (not normal coordinates), with nonzero
quartic jet C, Ric° ≠ 0 when K₁ ≠ K₂, W⁻ ≠ 0. Products are Kähler with
R⁺ = diag(s/4,0,0), s = 2(K₁+K₂); `par' puts ζ on the product Kähler
form (μ₁ = s/4), `perp' puts ζ in the kernel (μ = (0, s/4, 0)).
Prediction −s²/32 in both layouts.

\begin{longtable}[]{@{}lllll@{}}
\toprule
layout & (K₁,K₂) & D = 4 & D = 6 & prediction\tabularnewline
\midrule
\endhead
par & (−1,−1) KE & −0.3781 & −0.5000000000000 & −1/2\tabularnewline
par & (−1,−3) & −1.5124 & −2.0000000000000 & −2\tabularnewline
par & (−1,1) scalar-flat & 0 & 0 & 0\tabularnewline
perp & (−1,−1) KE & −0.4314 & −0.5000000000000 & −1/2\tabularnewline
perp & (−1,−3) & −1.7257 & −2.0000000000000 & −2\tabularnewline
perp & (−1,1) scalar-flat & 0 & 0 & 0\tabularnewline
perp & (−2,1/2) & −0.2427 & −0.2812500000000 & −9/32\tabularnewline
perp & (1,2) positive & −0.9707 & −1.1250000000000 & −9/8\tabularnewline
\bottomrule
\end{longtable}

All D = 6 errors ≤ 3·10⁻¹⁵, kernel leak ≤ 10⁻¹⁵. Lessons: (i) Ric° (up
to \textbar K₁−K₂\textbar{} = 2.5 here) and W⁻ are genuinely invisible,
now with a background that is not in the normal-coordinate gauge of the
fit; (ii) the scalar-flat Kähler product H²×S², which has large Ric°,
gives exactly 0, the predicted topological cost; (iii)
\textbf{convergence in corrector space depends on the gauge of the
background jet}: in normal coordinates D = 4 already converged, in
conformal coordinates D = 6 is needed (the D = 4 values lie above the
truth, as a minimization over a smaller space must). The verifier
question `do degree-0 correctors exhaust the log coefficient' is
therefore also a question about this gauge freedom; see the Step-58 TTCS
record.

\textbf{(d) Independent reproduction by A01 (\#830,
\texttt{Sessions/S1003A01/workspace/notes/step50\_wminus\_completed\_square\_and\_general\_KE\_construction.md}).}
A01 reduces F to the CGB form 4\textbar W⁺\textbar² + s²/12, shows W⁻
enters only through a completed square with the explicit corrector
k₀\^{}W = (3/2)sym(h\_W k\_B) − (1/4)tr(h\_W k\_B)δ, h\_W =
−(1/3)W⁻\_\{ikjl\}x\textsuperscript{kx}l, k\_B = σ₁²+σ₂² (exact sympy),
and obtains F\_KE = −18 at s = −24 on every negative KE surface in the
kernel orientation; in Ric = −g units P/8π² = 1 − (1/4)x²log(1/x). Our Φ
= −s²/32 gives (Φ/2) = −s²/64, which is −9 at s = −24 and −1/4 at s =
−4. Agreement in both normalizations. A01's statement that the Ric°
coefficients are ``nonnegative minimized squares'' is consistent with
our finding that their minimized value is exactly 0.

\textbf{(e) Two readings.} 1. \emph{Unification with Theorem H.}
Real-hyperbolic bulks have R⁺ = (s/12)I \textless{} 0 strictly, so every
Burns cell pays −εμ₁ = ε\textbar s\textbar/12 \textgreater{} 0 at first
order: the cost 8π²k is approached from above, which is exactly Theorem
H's regime (conditional on its Yamabe hypothesis). KE and CH bulks have
a kernel circle with anisotropic ζ\^{}⊥, so they save at second order,
so 8π²k is strictly beaten (p\_* \textless{} 8π²k); the insertion family
itself is then not minimizing and says nothing about attainment of
p\_\emph{. The two known regimes are the interior and the boundary of
the cone R⁺ ≤ 0. That `Burns bubbling is not excluded where R⁺
\textless{} 0' is the correct description of the hyperbolic phenomenon,
not a gap. 2. }Strict pointwise Weyl bound at minimizers carrying Burns
bubbles* (conditional on 53.B and the ledger, and on 57.Φ becoming
rigorous). R⁺ ≤ 0 with tr R⁺ = s/4 \textless{} 0 means the eigenvalues
of W⁺ are ≤ c = \textbar s\textbar/12; the maximum of
\textbar W⁺\textbar² on \{λ\_i ≤ c, Σλ\_i = 0\} is attained only at (c,
c, −2c), i.e.~KE type, value s²/24. 57.Φ excludes exactly that type
(double top eigenvalue ⇒ kernel anisotropic). Hence \textbar W⁺\textbar²
\textless{} s²/24 \textbf{at every point}, and the allowed boundary type
is W⁺ = (\textbar s\textbar/12)(1, −1/2, −1/2), \textbar W⁺\textbar² =
s²/96. Integrated: C₊(V) \textless{} ∫s²/(32π²). This is strict but not
quantitative; it does not contradict LeBrun's direction and gives no
global estimate by itself. (If s = 0, R⁺ ≤ 0 forces W⁺ = 0 and the
strict inequality is false; the statement is for s \textless{} 0.)

\hypertarget{step-58-verification-tracettcsresultsstep58_ttcs01_verifier1_integration.md-and-exhaustive-minimization}{%
\subsection{\texorpdfstring{7. Step-58 verification
(\texttt{trace/ttcs/results/step58\_ttcs01\_verifier1\_integration.md})
and exhaustive
minimization}{7. Step-58 verification (trace/ttcs/results/step58\_ttcs01\_verifier1\_integration.md) and exhaustive minimization}}\label{step-58-verification-tracettcsresultsstep58_ttcs01_verifier1_integration.md-and-exhaustive-minimization}}

\textbf{Confirmed.} The invariant algebra (fit ⇔ eigenvalue form), the
normalization (Φ/2)ε²log(1/ε) (log annulus from r \textasciitilde{} √ε,
so ∫dr/r = ½ log(1/ε)), CH = −18 ⇒ −9ε²log, all product predictions,
Lemma 57.1's mechanism, the strict bound \textbar W⁺\textbar²
\textless{} s²/24 (proof: t\_i = c − w\_i ≥ 0, Σt\_i = 3c,
\textbar W⁺\textbar² = Σt\_i² − 3c² ≤ 6c², equality only at (c,c,−2c)),
and the KE/hyperbolic table. Real-hyperbolic points have Φ = s²/288
\textgreater{} 0, irrelevant behind the first-order cost.

\textbf{Repairs adopted.} 1. \emph{W⁻/Ric° nonnegativity.} The pointwise
Chern--Weil argument of §3 was invalid (vanishing of W⁺, s at p does not
make F₊(g̃) − F₊(g) ≥ 0). Replacement (chiral square): with 𝒞(g) = (2W⁺,
s/√12), F₊ = ∫\textbar 𝒞\textbar². Burns has 𝒞 ≡ 0, and D𝒞\_δ{[}B{]} = 0
for B in the W⁻ ⊕ Ric° sector, so 𝒞(g\_\{ε,ν\}) = εν(𝒜\_B +
D𝒞\_δ{[}K{]}) + \ldots{} and the ε²ν² coefficient of the F₊ density is
\textbar 𝒜\_B + D𝒞\_δ{[}K{]}\textbar² ≥ 0. Zero is attained iff
D𝒞\_δ{[}K{]} = −𝒜\_B is solvable; A01's explicit k₀\^{}W solves it for
W⁻ exactly; for Ric° the solvability is supported numerically (D = 4, 6
zeros) but no closed-form corrector is yet written.
U(2)\emph{ζ-equivariance removes cross terms between sectors. 2.
\emph{Lemma 57.1 wording.} The ε²C part of the first-variation field is
homogeneous of degree −3, its radial flux is r-independent, so its
divergence has zero spherical average (no log). Simpler: W⁺ = s = 0 on
Burns, so the first variation of the F₊ density vanishes pointwise. ∇²Rm
and the curvature-quadratic part of the quartic jet may still enter the
regular O(ε²) term. 3. \emph{Scope of the coefficient.} Φ is the
coefficient of the matched ansatz (Burns core + degree-0 corrector). If
E\_P(g)(p) ≠ 0, singular tails ε²χr⁻⁴T on ε\^{}α \textless{} r
\textless{} R pair linearly with E\_P at order ε²log(1/ε); that is a
lowering of the background, not a bubble effect. At P-critical bulks
(KE; minimizer bulks) this freedom has no linear contribution.
Higher-homogeneity terms such as εRm²H₂ can sit in the same corner; they
vanish by Burns criticality, not by power counting. 4. \emph{Gauge and
D.} D is angular polynomial complexity of degree-0 tensors. A cubic
coordinate change X₃ shifts K ↦ K + ℒ}\{X₃\}A (degree 0, higher angular
complexity), which explains why conformal-coordinate products need D = 6
while normal coordinates converge at D = 4. 5. \emph{KE scope.}
Ricci-flat KE (s = 0) has Φ = 0. For minimal general-type surfaces with
(−2)-curves, use the KE canonical orbifold, Burns cells at regular
points, Kronheimer resolutions at ADE points with scales chosen after
the Burns scales. Strictness then needs only the KE-cell remainder
o(ε²log) and zero-energy resolution gluing. 6. \emph{Open boundary
case.} At R⁺ = (0, s/8, s/8) both displayed terms vanish; the next
object is the regular O(ε²) coefficient.

\textbf{Exhaustive minimization in the (s, W⁺) sector (new,
\texttt{workspace/scripts/step58\_invariant\_exhaustive.py} →
\texttt{workspace/data/step58\_invariant\_exhaustive.json}).} For
curvature with only s and W⁺, the normal-coordinate jet and the Burns
tail are SU(2)\emph{−-invariant; averaging K over SU(2)}− keeps b(K) and
does not increase Q(K) (convexity), and SU(2)\emph{− acts simply
transitively on S³, so the 10-dimensional invariant space Sym²⟨dr, σ₁,
σ₂, σ₃⟩ contains the minimizer over \textbf{all} degree-0 correctors.
Result (q8×24, residual 2.4·10⁻¹⁵): s² 1/288, sw −1/24, w² −1,
\textbar W⁺}⊥⊥\textbar² −1, CH = −18.0000000000003; the W⁺\_\{ζ⊥\}
directions leak into ker Q (1.414), the frame-rotation null direction.
So in the (s, W⁺) sector 57.Φ is the exact optimum over all angular
modes, not a truncation artefact. Together with item 1, the formal
coefficient is now established in every sector except for a closed-form
Ric° corrector.

\hypertarget{s1003a02-consolidated-hypothesis-ledger-and-theorem-chain-steps-5357}{%
\chapter{S1003A02: Consolidated hypothesis ledger and theorem chain:
Steps
53--57}\label{s1003a02-consolidated-hypothesis-ledger-and-theorem-chain-steps-5357}}

Source: \texttt{step58\_consolidated\_hypothesis\_ledger.md}.

Session S1003A02, Agent2, Execute. This note replaces the layered
correction sections (§0 of Steps 54, 55, 56, 57) with a single
reference. The authoritative per-step derivations remain in their
original notes; this note tracks what is proved, conditional, and open
at the end of Step 58.

Conventions: \(P(g)=\int|{\rm Ric}^\circ|^2\), \(p_*(M)=\inf P\),
\(T_+=2\chi+3\tau\), \(A=c_1^2(X)\) for the minimal model,
\(\mathbf R^\pm=W^\pm+\frac s{12}I\). Eigenvalues of \(W^+\) in the
\(\Lambda^+\) convention sum to 0; operator norm.

\begin{center}\rule{0.5\linewidth}{0.5pt}\end{center}

\hypertarget{i.-hypotheses}{%
\subsection{I. Hypotheses}\label{i.-hypotheses}}

All refer to a sequence of unit-volume metrics \((M,g_i)\) with
\(P(g_i)\to p_*(M)\), on \(M=X\#k\overline{\mathbb{CP}}^2\).

\textbf{(I)} Anderson case (I): \(\varepsilon\to0,p\to1\) subsequence of
minimizers of
\(Z^{2p}_\varepsilon=P+\varepsilon\int(1+|{\rm Ric}|^2)^p\) converges
(orbifold Cheeger--Gromov) to a multi-component configuration
\((V_a,g_0)\cup(N_k)\) with

\begin{itemize}
\tightlist
\item
  \(V_a\): compact orbifold bulk components, together covering \(M\) up
  to bubbles and collapsed regions;
\item
  \(N_k\): ALE bubble regions at separated smaller scales.
\end{itemize}

Reference: Anderson, \emph{Canonical Metrics on 3- and 4-Manifolds}, Thm
4.10 (2006). The existence of the sequence follows from standard direct
method; the fourth-order Euler--Lagrange system and its elliptic
regularity are Anderson's.

\textbf{\(\mathrm H_{\rm sm}\)} (smooth cone points): every cone point
of \(V\) has trivial group. (The alternative is orbifold absorption,
analyzed in §III.B.)

\textbf{\(\mathrm H_{\rm lsc}\)} (lower semicontinuity):
\(p_*(M)\ge\sum_aP(V_a,g_0)+\sum_kP(N_k)\). Standard for weak
\(W^{2,2}\) convergence on compact sets at separated scales.

\textbf{\(\mathrm H_{\rm dec}\)} (decay): each bubble \(N_k\) is smooth,
complete, AE with derivative-controlled decay of some order \(\tau>0\),
and scalar-flat. Used for PSC of the one-point compactification (Lemma
53.Y with \(\tau>0\)) and for the charge formula.

\textbf{\(\mathrm H_{\rm eq}\)} (energy equality):
\(P(g_i)=\sum_aP(V_a,g_{i,a})+\sum_kP(N_k)+o(1)\). Anderson (4.43);
standard in bubble-tree theory but not proved for the fourth-order
system without extra input.

\textbf{\(\mathrm H_{\rm rem}\)} (removability): \(g_0\) extends
smoothly across the cone points (trivial group case). Used only for
attainment.

\textbf{(SW\(_L\))} (two-sided): SW vanishes for splittings along an
\(L\)-space with \(b^+\ge1\) on both sides. Citable:
\textbackslash cite\{agolhyperbolic\} with Lin's proof; lens spaces and
spherical space forms are \(L\)-spaces. One-sided form (PSC orbifold cap
+ \(b^+\ge1\) other side): remains a hypothesis.

\textbf{(N\(_0\))} (Ricci-flat orbifold bubbles): every bubble at a
point with nontrivial group is Ricci-flat ALE, i.e. \(P=0\). Eliminates
non-Ricci-flat orbifold channels and exotic-end survivors.

\textbf{(U)} (upper bound): \(p_*(M)<8\pi^2k\) (for each \(k\ge1\)).
Proved unconditionally: the analytic CH-collar cell gives
\(P/(8\pi^2)=1-9\varepsilon^2\log(1/\varepsilon)+O(\varepsilon^2)\) for
every embedded ball below the injectivity radius. For \(k=1\) the
authoritative certificate gives \(P/(8\pi^2)\le0.995662732771115\).

\begin{center}\rule{0.5\linewidth}{0.5pt}\end{center}

\hypertarget{ii.-proved-results-unconditional-or-with-stated-universal-input}{%
\subsection{II. Proved results (unconditional or with stated universal
input)}\label{ii.-proved-results-unconditional-or-with-stated-universal-input}}

\textbf{Theorem K\(_\infty\) (Kähler infimum).} For any minimal compact
general-type surface \(X\) with \(b^+(X)\ge3\) and \(k\ge1\), \[
\inf_{\text{fixed-complex Kähler metrics on }X\#k\overline{\mathbb{CP}}^2} P = 8\pi^2k,
\] not attained for \(k>0\). Uses Hodge index on total-transform
exceptional classes (strict lower bound), and successive Arezzo--Pacard
cscK blow-ups on the canonical orbifold (recovery). Scope: fixed complex
structure only.

\textbf{Theorem 57.Φ (formal kernel coefficient).} \emph{Formal status:
the \(\varepsilon^2\log(1/\varepsilon)\) coefficient is computed and
cross-checked numerically but the remainder estimate is not yet proved.}
Let \(\zeta\) be an eigenvector of \(\mathbf R^+=W^++s/12\) at \(p\)
with eigenvalue \(\mu_1\), other eigenvalues \(\mu_2,\mu_3\). Then the
coefficient \(\Phi\) in
\(P(\tilde g)/8\pi^2=P(g)/8\pi^2+1-\varepsilon\mu_1+(\Phi/2)\varepsilon^2\log(1/\varepsilon)+O(\varepsilon^2)\)
is \[
\Phi=-\mu_1(\mu_1-\mu_2-\mu_3)/2-(\mu_2-\mu_3)^2/2.
\] In the kernel (\(\mu_1=0\)): \(\Phi=-(\mu_2-\mu_3)^2/2\le0\). The
coefficient depends only on \({\rm Rm}_p\) through \(s\) and \(W^+\); it
is independent of \(W^-\), \({\rm Ric}^\circ\), \(\nabla{\rm Rm}\), and
\(\nabla^2{\rm Rm}\).

Status of the coefficient (after the Step-58 verifier,
\texttt{trace/ttcs/results/step58\_ttcs01\_verifier1\_integration.md}):
in the \((s,W^+)\) sector it is the exact optimum over all degree-0
correctors (SU(2)\(_-\) averaging reduces to the 10-dim invariant space;
\texttt{step58\_invariant\_exhaustive.json}); in the
\(W^-\oplus{\rm Ric}^\circ\) sector the chiral square gives \(\ge0\),
A01's closed-form \(k_0^W\) attains 0 for \(W^-\), and Ric° zero is
numerical (D=4,6) pending a closed-form corrector. It is the coefficient
of the matched ansatz; at non-\(P\)-critical backgrounds singular tails
pairing with \(E_P(p)\) are a separate (background-lowering) effect.
Verified: (i) reproduces \(-18\) on CH with invariant, D=2,4 and general
D=4 correctors; (ii) the Chern--Weil density check gives identical M and
M\_eff; (iii) eight product-background tests with Ric°≠0, W⁻≠0 agree at
D=6 to 10⁻¹⁵; (iv) independently confirmed for all KE backgrounds by A01
(\#830) via the CGB reduction and explicit W⁻ corrector.

\textbf{Formal consequence:} for every compact KE surface (canonical
bundle ample), \(\Phi=-s^2/32<0\) at every point in the kernel
direction. Hence \(p_*(X\#k\overline{\mathbb{CP}}^2)<8\pi^2k\).

\textbf{Proposition 56.U (zero-energy upper bound).} For a compact
4-orbifold \(V\) with Wahl points, Milnor-fibre smoothing \(M'\):
\(p_*(M')\le P(V,g)\) for every smooth orbifold metric. Unconditional.

\textbf{Corollary 56.B.} An Einstein orbifold metric on such \(V\)
forces \(\alpha(M')=0\). Non-Kähler if the bulk would be complex with
those quotient singularities (Prop 55.6 repaired). Lattice: \(n\) odd or
\(8\mid n\).

\begin{center}\rule{0.5\linewidth}{0.5pt}\end{center}

\hypertarget{iii.-conditional-results}{%
\subsection{III. Conditional results}\label{iii.-conditional-results}}

\hypertarget{a.-smooth-bubble-exclusion-and-attainment-k-1-ball-quotient}{%
\subsubsection{A. Smooth-bubble exclusion and attainment (k = 1, ball
quotient)}\label{a.-smooth-bubble-exclusion-and-attainment-k-1-ball-quotient}}

\textbf{Theorem 54.A.} Let \(X\) be a compact ball quotient with
\(b^+(X)\ge3\), \(M=X\#\overline{\mathbb{CP}}^2\). Under \textbf{(I)},
\(\mathrm H_{\rm sm}\), \(\mathrm H_{\rm lsc}\),
\(\mathrm H_{\rm dec}\):

\begin{enumerate}
\def\labelenumi{(\roman{enumi})}
\item
  Every bubble is one-ended, simply connected, has \(b_2=0\). \(E\) is
  in the bulk.
\item
  \(V\) is homeomorphic to \(M\); each \(P(N_k)<8\pi^2\).
\end{enumerate}

Under also \(\mathrm H_{\rm eq}\), \(\mathrm H_{\rm rem}\), and assuming
caps \(\cong S^4\):

\begin{enumerate}
\def\labelenumi{(\roman{enumi})}
\setcounter{enumi}{2}
\tightlist
\item
  \(g_0\) attains \(p_*(M)\) on a smooth \(V\cong M\), with \(s<0\)
  constant. Either Einstein or non-Bach-flat.
\end{enumerate}

\hypertarget{b.-orbifold-category-attainment-k-1-ball-quotient}{%
\subsubsection{B. Orbifold-category attainment (k = 1, ball
quotient)}\label{b.-orbifold-category-attainment-k-1-ball-quotient}}

\textbf{Theorem 56.A.} Under \textbf{(I)}, \(\mathrm H_{\rm lsc}\),
\(\mathrm H_{\rm dec}\) (any group), bubble criticality by compact
replacement, \textbf{(SW\(_L\))}, smooth orbifold regularity,
\textbf{(N\(_0\))}:

\begin{enumerate}
\def\labelenumi{(\roman{enumi})}
\item
  \(p_*(M)=P(V,g_0)=\alpha_{\rm orb}(V)\). Orbifold points are cyclic
  with Ricci-flat rational-ball bubbles.
\item
  \(s_0<0\) constant. Proof: LeBrun's \(\int s_-^2\ge32\pi^2c_1^2(X)\)
  on \(M\) via reassembly.
\item
  Einstein (\(\Rightarrow\alpha(M)=0\)) or non-Bach-flat.
\item
  Wahl points: lattice law, non-complex, degree-one kernel.
\end{enumerate}

\hypertarget{c.-minimizer-rigidity-formal-conditional-on-the-ledger-and-53.b}{%
\subsubsection{C. Minimizer rigidity (formal; conditional on the ledger
and
53.B)}\label{c.-minimizer-rigidity-formal-conditional-on-the-ledger-and-53.b}}

At a P-critical constant-negative-scalar bulk carrying a Burns-type
bubble: \(\mathbf R^+\le0\), and
\(\lambda_{\max}(\mathbf R^+)=0\Rightarrow\mu_2=\mu_3\). Hence no
KE-type point (eigenvalues \((s/4,0,0)\) with \(s\ne0\)), and
\(|W^+|^2<s^2/24\) strictly at every point.

\hypertarget{d.-general-k-any-ball-quotient}{%
\subsubsection{D. General-k (any ball
quotient)}\label{d.-general-k-any-ball-quotient}}

Under (I), \(\mathrm H_{\rm sm}\), \(\mathrm H_{\rm lsc}\),
\(\mathrm H_{\rm dec}\): all bubble caps are simply connected, \(H_2\)
of bubbles lies in \(\mathrm{span}\{E_1,\ldots,E_k\}\), total
exceptional rank in smooth bubbles is \(<k\). One-ended for \(k\le4\).

\begin{center}\rule{0.5\linewidth}{0.5pt}\end{center}

\hypertarget{iv.-open-gates}{%
\subsection{IV. Open gates}\label{iv.-open-gates}}

\begin{enumerate}
\def\labelenumi{\arabic{enumi}.}
\item
  \textbf{Analytic gates.} \(\mathrm H_{\rm dec}\)(a) (bubble-scale
  criticality via compact convergence) is the genuine residue. Prop 54.4
  shows it follows from \(\mathrm H_{\rm eq}\) and scaling.
  \(\mathrm H_{\rm rem}\) is separate.
\item
  \textbf{Orbifold absorption.} The threshold is
  \(C_+(V)+\sum w(N_j)\ge A\), which is the open-range inequality.
  (N\(_0\)) removes non-Ricci-flat channels by fiat; its removal
  requires pricing non-Ricci-flat orbifold bubbles.
\item
  \textbf{Collapse.} Anderson case (II): bounded curvature, volume
  \(\to0\). Free-circle regions are spin and cannot carry the odd class
  \(E\); ball-quotient \(\pi_1\) is freely indecomposable, fibre images
  are cyclic. The missing pieces: finite-isotropy control,
  graph-manifold interface splitting, and adiabatic \(\eta\)-analysis.
\item
  \textbf{Exotic homotopy spheres.} Would also disprove smooth 4D
  Poincaré.
\item
  \textbf{The rigorous remainder for 57.Φ.} Needs the two-scale
  factorization (Ch. 5 type) in an analytic background with angular
  dependence. The KE case (convergent jets, averaged CGB form) is the
  shortest path. Division of labour with A01 is active.
\item
  \textbf{The second-order matching law for balanced families} (Problem
  2). Per-cell saving \((\lambda^2\varepsilon^2\log(1/\varepsilon))/4\)
  for \(\mathrm{Ric}=\lambda g\); optimize \(\varepsilon\) at fixed
  density; batch-gluing interactions. Deferred since Step 50.
\end{enumerate}

\begin{center}\rule{0.5\linewidth}{0.5pt}\end{center}

\hypertarget{iv.-zero-gap-without-nearby-einstein-metrics-spherical-and-hyperbolic-orbifold-resolutions-step-58}{%
\subsection{IV′. Zero gap without nearby Einstein metrics: spherical and
hyperbolic orbifold resolutions (Step
58)}\label{iv.-zero-gap-without-nearby-einstein-metrics-spherical-and-hyperbolic-orbifold-resolutions-step-58}}

\textbf{Proposition 58.O.} Let \((V,g_o)\) be a compact spherical or
hyperbolic 4-orbifold whose singular points are isolated, of type
\(\mathbb R^4/\Gamma_j\) with \(\Gamma_j\subset SO(4)\) acting freely on
\(S^3\) and conjugate in \(O(4)\) to the end group of some Ricci-flat
ALE space \(N_j\) (e.g.~\(\Gamma=\{\pm1\}\), \(N=\) Eguchi--Hanson; ADE
groups, Kronheimer; class-T cyclic groups, Şuvaina). Let \(M'\) be the
smooth manifold obtained by replacing each singular neighbourhood by
\(N_j\) (or \(\overline{N_j}\) if the conjugation reverses orientation;
Ricci-flatness is orientation-free). Then \(\alpha(M')=0\). If at least
one singular point is \(\mathbb R^4/\mathbb Z_2\), then no sequence of
Einstein metrics on \(M'\) converges to \((V,g_o)\) in the
Gromov--Hausdorff sense \textbackslash cite\{ozuch2021integrability\}
(Thm 0.9). Hence on \(M'\) the zero gap is realized only by degenerating
approximations: \(\alpha(M')=0\) is attained iff \(M'\) carries an
Einstein metric, and any such Einstein metric is GH-far from \(g_o\).

\emph{Proof.} \(g_o\) is Einstein, so \(P(g_o)=0\). Glue
\(\lambda^2g_{N_j}\) at each point exactly as in Prop. 56.U / Cor. 52.8:
annular error \(O(\rho^4)+O((\lambda/\rho)^{2\tau})\), \(\tau=4\) for
these ALE spaces, interior \(P(\lambda^2g_N)=0\), deleted orbifold
region \(O(\rho^4)\). With \(\rho=\lambda^{1/2}\) the total is
\(O(\lambda^2)\to0\). The second statement is Ozuch's theorem; the
attainment statement is the equality case of \(P\ge0\). \(\square\)

Examples: \(S^4/\{\pm1\}\) (two \(\mathbb Z_2\) points) resolves to
\(\overline{\mathbb{CP}}^2\#\overline{\mathbb{CP}}^2\) (Thm 52.7(a));
hence \(\alpha(\mathbb{CP}^2\#\mathbb{CP}^2)=0\), attainment is
equivalent to Einstein existence on \(\mathbb{CP}^2\#\mathbb{CP}^2\)
(open), and the S⁴/Z₂ degenerations are not limits of Einstein metrics.
Hyperbolic 4-orbifolds with isolated \(\mathbb Z_2\) points give, in the
same way, smooth manifolds with \(\alpha=0\) and the same
non-approximability. This is the cleanest present instance of the
separation between \(\alpha=0\), approximation, and existence.

\textbf{Target 58.Q (quantitative Biquard--Ozuch obstruction in \(P\)).}
At an Einstein (hence \(P\)-critical) orbifold point with an EH bubble
of scale \(\lambda\), Lemma 56.L kills the \(\lambda^4\log(1/\lambda)\)
term. The leading cost is a finite renormalized polarizability \[
P=\lambda^4\,\mathcal Q_{EH}(\mathbf R_p)+o(\lambda^4),\qquad
\mathcal Q_{EH}(\mathbf R)=\inf_{k}\int_{EH}\big|D{\rm Ric}^\circ_{g_{EH}}(\chi B_{\mathbf R}+k)\big|^2,
\] \(k\) decaying, \(B_{\mathbf R}\) the quadratic background jet.
Expected structure: \(\mathcal Q_{EH}\) is a nonnegative quadratic form
on the curvature at \(p\) whose kernel is the vanishing locus of
Biquard's obstruction (pairing of \(\mathbf R^\pm_p\) with the three
\(L^2\) deformations of EH), positive at spherical/hyperbolic points. If
true, it gives the near-extremal statement \(P\ge c\lambda^4\) along
EH-degenerations of nonintegrable Einstein orbifolds, i.e.~a
quantitative form of Ozuch's theorem in the gap functional. EH is
\(U(2)\)-invariant, so \(\mathcal Q_{EH}\) reduces to cohomogeneity-one
ODE systems by \(U(2)\)-type of the background jet. Not yet computed.

\begin{center}\rule{0.5\linewidth}{0.5pt}\end{center}

\hypertarget{v.-what-is-not-proved}{%
\subsection{V. What is not proved}\label{v.-what-is-not-proved}}

\begin{itemize}
\tightlist
\item
  Any unrestricted positive gap in the classical window \(0<k<A/3\).
\item
  A universal \(k^2/A\) lower bound.
\item
  A closed family realizing \(P/k\to0\) or the formal \(k^2/A\) scale.
\item
  An unconditional Anderson decomposition or smooth minimizer on
  \(X\#\overline{\mathbb{CP}}^2\).
\item
  An unconditional O10 finite-block theorem.
\item
  Any topology-only or field-free per-class tariff.
\end{itemize}

\hypertarget{s1003a02-step-59-a-rigorous-second-order-burns-expansion-the-strict-sub-threshold-bound-on-every-ke-surface-and-the-structure-of-the-zero-gap-set}{%
\chapter{S1003A02: Step 59 --- A rigorous second-order Burns expansion,
the strict sub-threshold bound on every KE surface, and the structure of
the zero-gap
set}\label{s1003a02-step-59-a-rigorous-second-order-burns-expansion-the-strict-sub-threshold-bound-on-every-ke-surface-and-the-structure-of-the-zero-gap-set}}

Source:
\texttt{step59\_rigorous\_burns\_expansion\_and\_zero\_gap\_structure.md}.

Agent2, Execute.

Status of the parts: - Part A (Theorem 59.R and Corollary 59.KE): proved
here. The (s, W⁺) coefficient is exact rational arithmetic. For the W⁻
block, A01's exact identity is being re-derived in
\texttt{step59\_exact\_wminus\_ric\_correctors.py}. - Part B (Theorem
59.Z): an assembly of proved mechanisms. It includes a new elementary
partial-collapse construction (Prop. 59.D) and a classification
direction that is conditional only on Anderson's idealized-minimizer
theorem.

Conventions. P(g) = ∫\textbar Ric°\textbar² dμ. On a closed oriented
4-manifold, P = ∫F₊ dμ − 16π²χ − 24π²τ, with F₊ = 4\textbar W⁺\textbar²
+ s²/12 = \textbar 𝒞\textbar² and 𝒞 = (2W⁺, s/√12). Here
\textbar W⁺\textbar² is the operator norm on unit 2-forms, so a Kähler
metric has \textbar W⁺\textbar² = s²/24. R⁺ = W⁺ + s/12 has eigenvalues
μ₁ (on ζ), μ₂, μ₃.

\begin{center}\rule{0.5\linewidth}{0.5pt}\end{center}

\hypertarget{part-a.-second-order-burns-insertion-with-a-proof}{%
\subsection{Part A. Second-order Burns insertion, with a
proof}\label{part-a.-second-order-burns-insertion-with-a-proof}}

\hypertarget{a.1-the-competitor}{%
\subsubsection{A.1 The competitor}\label{a.1-the-competitor}}

Let (X, g) be smooth and closed, and P-critical near p, i.e.~E\_P(g) = 0
on a neighbourhood of p.~Every Einstein metric qualifies. - Take
g-normal coordinates x on B\_\{3ρ\}(p), so g = δ + h with h = h₂ +
O(\textbar x\textbar³) and h₂ = −(1/3)R\_\{iajb\}(p)x\^{}a x\^{}b. - Fix
a unit ζ ∈ Λ⁺\_p, the constant complex structure J\_ζ on R⁴ ≅ C², and
the Burns tail A = ∂∂̄\_ζ log\textbar x\textbar². This is normalised so
that ∂∂̄\textbar x\textbar² = δ. Then g\_B := δ + εA is exactly the Burns
metric, with area of E equal to πε. Since ∂∂̄ log\textbar x\textbar² ≥ 0,
we have g\_B ≥ δ. - Fix a corrector K: a smooth symmetric 2-tensor on
R⁴∖0 that is homogeneous of degree 0. - Fix cut-offs χ\_ρ(r) (equal to 1
for r ≤ ρ and 0 for r ≥ 2ρ) and ψ\_ε(r) = 1 − χ(r/(R√ε)), which vanishes
for r ≤ R√ε. Here R ≥ 1 is fixed.

Define

g̃\_ε = g + ε χ\_ρ (A + ψ\_ε K) on X∖\{p\}.

For r \textless{} R√ε this is g\_B + h. A smooth tensor on C² pulls back
to a smooth tensor on Bl₀C², and \textbar h\textbar\_\{g\_B\} ≤
\textbar h\textbar\_δ ≤ Cr². So g̃\_ε extends to a smooth Riemannian
metric on M = X \# CP̄²; it is Burns-like near E.

\hypertarget{a.2-exact-global-identity}{%
\subsubsection{A.2 Exact global
identity}\label{a.2-exact-global-identity}}

Since χ(M) = χ(X) + 1 and τ(M) = τ(X) − 1,

P(g̃\emph{ε) = P(g) + 8π² + I(ε), I(ε) = ∫}\{Bl B\_\{2ρ\}\} F₊(g̃\emph{ε)
dμ̃ − ∫}\{B\_\{2ρ\}\} F₊(g) dμ. (A.1)

This is an equality; no transgression or boundary term enters. The
density F₊ = \textbar 𝒞\textbar² is the chiral square, and 𝒞(g\_B) ≡ 0
for every ε, because Burns is ASD and scalar-flat.

\hypertarget{theorem-59.r}{%
\subsubsection{Theorem 59.R}\label{theorem-59.r}}

Under A.1, with C depending on (g, p, ζ, K, ρ, R):

I(ε) = 2κ ε R⁺\_p(ζ,ζ) + 4π² Φ(Rm\_p, ζ, K) ε² log(1/ε) + O(ε²).

The two coefficients are as follows. - κ is a universal nonzero
constant. Comparison with Theorem F gives 2κ = −8π² in these units. - Φ
is given by

Φ(Rm\_p, ζ, K) = (1/8π²) ∫\emph{\{S³\} {[}σ²{]} (
\textbar D𝒞}\{δ+σA\}{[}h₂ + σK{]}\textbar²\_\{δ+σA\} dμ\_\{δ+σA\} ),
(A.2)

where {[}σ²{]} extracts the coefficient of σ². Hence P(g̃)/8π² = P(g)/8π²
+ 1 − εR⁺(ζ,ζ) + (Φ/2)ε²log(1/ε) + O(ε²).

\emph{Proof.} Split Bl B\_\{2ρ\} into three regions: inner D\_i = \{r ≤
2R√ε\}, far D\_f = \{2R√ε ≤ r ≤ ρ\}, and outer D\_o = \{ρ ≤ r ≤ 2ρ\}.

\textbf{Inner region.} In u = x/√ε the scaled metric ε⁻¹g̃ is

g\_\{B,1\}(u) + ε{[}\textbar u\textbar² h̄(√ε u, u) + ψK(u){]}, with h(x)
= \textbar x\textbar² h̄(x, x/\textbar x\textbar) smooth.

On the compact set \{\textbar u\textbar{} ≤ 2R\} ⊂ Bl₀C², this is a
perturbation of the unit Burns metric whose g\_\{B,1\}-derivatives are
O(ε) to all orders. Because 𝒞(g\_\{B,1\}) = 0, the scaled chiral
curvature is O(ε), so the unscaled 𝒞(g̃) is O(1). As
\textbar 𝒞\textbar²dμ is scale invariant,

∫\emph{\{D\_i\} F₊(g̃) dμ̃ = ∫}\{\textbar u\textbar{} ≤ 2R\}
\textbar 𝒞(ε⁻¹g̃)\textbar² dμ = O(ε²R⁴), and also ∫\_\{D\_i\} F₊(g) dμ =
O(ε²R⁴).

\textbf{Far region.} For r ∈ {[}2R√ε, ρ{]} and y ∈ U = \{1/2 ≤
\textbar y\textbar{} ≤ 2\}, set σ = ε/r². Then

r⁻²(y ↦ ry)*g̃ = G\_\{σ,r\} := δ + σA(y) + r²(h̄(ry, y) + σK(y)).

This is smooth in (σ, r, y) ∈ {[}0, (4R²)⁻¹{]} × {[}0, ρ{]} × U and
positive definite. The perturbation of δ + σA equals r²·(smooth). Since
𝒞(δ + σA) = 0, we get 𝒞(G\_\{σ,r\}) = r²·(smooth), so the difference of
F₊-densities is r⁴·(smooth).

Write f(r) = ∮\_\{\textbar x\textbar=r\}(F₊(g̃)dμ̃ − F₊(g)dμ)/dr. By
scaling, f(r) = r⁻¹γ(ε/r², r) with γ(σ,r) = r⁴γ̃(σ,r), where γ̃ is smooth
and γ̃(0,r) ≡ 0. Taylor's theorem in σ gives γ̃ = σγ₁(r) + σ²γ₂(r) +
σ³γ₃(σ,r) with γ₁, γ₂, γ₃ smooth. Hence

∫\_\{D\_f\} = ε ∫ rγ₁ dr + ε² ∫ γ₂ dr/r + ε³ ∫ γ₃ r⁻³ dr.

The three pieces are handled as follows. - Third piece: O(ε³(R√ε)⁻²) =
O(ε²R⁻²). - Second piece: \textbar γ₂(r) − γ₂(0)\textbar{} ≤ Cr, so it
equals ε²γ₂(0)(½ log(1/ε) + log(ρ/2R)) + O(ε²). - Second piece,
identification of γ₂(0): it is the S³-integral of the σ²-coefficient of
r⁻⁴F₊(G\_\{σ,r\}) at r = 0. At order r², 𝒞(G\_\{σ,r\}) is
D𝒞\_\{δ+σA\}{[}h₂ + σK{]}. Therefore γ₂(0) = 8π²Φ. - Neither the 3-jet
nor the 4-jet of g can enter γ₂(0). They contribute only at r⁵ and r⁶,
because F₊ is quadratic in the perturbation.

\textbf{Outer region and linear terms.} On D\_o, g̃ = g + εχ\_ρ(A + K) is
smooth with O(ε) derivatives, so

∫\_\{D\_o\} = ε ∫ ℓ\_g(χ\_ρ(A+K)) + O(ε²),

where ℓ\_g(B) is the first variation of the F₊-density. Collecting all
ε-linear terms gives εE₁ + O(ε²R²), with

E₁ = ∫\_\{X∖p\} ℓ\_g(B), B = χ\_ρ(A + K).

This integral converges absolutely, since ℓ\_g(B) = O(r⁻²).

\textbf{Value of E₁.} Let χ\_η = χ(r/η). - E\_P = E\_\{F₊\} vanishes
near p, so ∫ℓ\_g((1 − χ\_η)B) = 0 for every small η. Hence E₁ =
lim\_\{η→0\} ∫ℓ\_g(χ\_η B). - In normal coordinates, ℓ\_g(B) =
2⟨𝒞\_g(p), L₀B⟩ + O(r)∂²B + O(r)∂B + O(1)B, where L₀ = D𝒞\_δ has
constant coefficients. - The error terms contribute O(η) and O(η²). The
K-part contributes O(η²), since L₀K is homogeneous of degree −2. - So E₁
= 2⟨𝒞(p), c\_A⟩ with c\_A = ∫\_\{R⁴∖0\} L₀(χA) dx. The integrand is
supported in an annulus, because L₀A = 0 for r \textgreater{} 0. -
Replace χA by ∂∂̄\_ζ(χ log r²). The difference is smooth and compactly
supported, and L₀ is a divergence, so c\_A is unchanged. - For Kähler
potentials the linearized identity W⁺' = (s'/12)·diag(2, −1, −1)\emph{ζ
holds. Differentiate W⁺ = (s/12)M\_ζ along δ + t∂∂̄φ. - Therefore c\_A =
m·(M\_ζ/12, 1/√12), where m = ∫s'(∂∂̄(χ log r²)). This is a nonzero
multiple of ∫}\{r\textgreater0\}Δ²(χ log r²) = 16π², the Green flux. -
Since tr W⁺ = 0, 2⟨𝒞(p), c\_A⟩ = 2m(W⁺(ζ,ζ) + s/12) = 2m R⁺\_p(ζ,ζ). Put
κ = m. ∎

\textbf{Remarks.} 1. The O(x₀\^{}\{6/5\}) remainder of Theorem F is
improved to O(ε² log). The improvement comes from two facts: the
chiral-square factorization 𝒞(g\_B) ≡ 0, and the exact identity (A.1).
2. Theorem 57.Φ is upgraded from formal to proved, for every P-critical
background and every fixed corrector K. The value of Φ is then a finite
computation: one exact quadratic form on S³ for each K.

\hypertarget{a.3-the-exact-coefficient}{%
\subsubsection{A.3 The exact
coefficient}\label{a.3-the-exact-coefficient}}

\textbf{(s, W⁺) sector.} Suppose W⁻ = 0 and Ric° = 0 at p.~Then the jet
h₂, the tail A and the SU(2)₋-invariant correctors are SU(2)₋-invariant.
SU(2)₋ acts simply transitively on S³, so the σ²-density in (A.2) is
constant on S³. It is therefore its value at x₀ = (1,0,0,0) times 2π².

\texttt{workspace/scripts/step59\_exact\_log\_coefficient.py} computes
this value in exact rational arithmetic (sympy, 14 s), as a quadratic
form in (s, W⁺) and the 10-dimensional invariant corrector space K = Σ
k\_\{ij\} f\_i f\_j, with f = (dr, s₁, s₂, s₃) the unit SU(2)₋-invariant
coframe. Output: \texttt{data/step59\_exact\_log\_coefficient.json}.

Checks: the σ⁰ term equals 4\textbar W⁺\textbar² + s²/12 exactly, and s'
at σ⁰ equals s.

For ζ an eigenvector of W⁺ (w₁₂ = w₁₃ = 0), the minimum over K is

Φ = (s² − 12 s w₁₁ − 432 w₁₁² − 576 w₁₁w₂₂ − 576 w₂₂² − 576 w₂₃²)/288 =
s²/288 − s w/24 − w² − \textbar W⁺\_\{⊥⊥,0\}\textbar² = −μ₁(μ₁ − μ₂ −
μ₃)/2 − (μ₂ − μ₃)²/2.

The difference from the Step-57 prediction is identically zero. In
particular: - Φ\_CH⊥ = −18 exactly; - Φ\_Kähler∥ = −s²/32.

An explicit optimal corrector (free kernel parameters set to 0, w =
w₁₁):

K* = (s/24) dr² + (−s/48 + 7w₁₁/6 + w₂₂/3) s₁² + (4w₂₃/3) s₁s₂ + (−s/48
+ 5w₁₁/6 − w₂₂/3) s₂².

For KE⊥ this becomes K* = s(dr²/24 − s₁²/16 − 7s₂²/48), and Φ = −s²/32.

\textbf{W⁻ and Ric° sectors (chiral square).} For a pure W⁻ or pure Ric°
jet, D𝒞\_δ{[}h₂{]} = 0. So the sector density is \textbar Y\textbar²
with

Y = {[}σ¹{]} D𝒞\_\{δ+σA\}{[}h₂ + σK{]}.

The sector coefficient is ∫\textbar Y\textbar² ≥ 0, and it vanishes iff
Y ≡ 0 for some K. - W⁻: A01 (\#830) gives the exact pointwise solution
k₀\^{}W = (3/2)sym(h\_W k\_B) − (1/4)tr(h\_W k\_B)δ. - Cross terms
between a W⁻ block and the invariant block vanish after S³-integration
by Schur's lemma. The cross term is an SU(2)₋-invariant linear
functional of W⁻, and W⁻ is spin 2. - Hence, with K = K* + k₀\^{}W, the
KE coefficient is exactly −s²/32, whatever W⁻ is. -
\texttt{step59\_exact\_wminus\_ric\_correctors.py} re-derives the W⁻
corrector inside an explicit ansatz in exact arithmetic. It also
searches the same ansatz for closed-form Ric° correctors (blocks
Ric°\emph{ζ and Ric°}⊥).

\hypertarget{corollary-59.ke}{%
\subsubsection{Corollary 59.KE}\label{corollary-59.ke}}

Let X be a compact Kähler--Einstein surface with s ≠ 0, and k ≥ 1.
Choose distinct points p₁,\ldots,p\_k and directions ζ\_i ⊥ ω, which lie
in ker R⁺ because R⁺ = diag(s/4, 0, 0) in the ω-frame. Insert Burns
metrics at disjoint balls, with K* + k₀\^{}W at each point. Then

P(g̃)/8π² = k − k(s²/64) ε² log(1/ε) + O(ε²),

and hence p\_*(X \# k CP̄²) \textless{} 8π² k.

If X is a minimal surface of general type whose canonical model X\_can
has ADE points, the same holds: - Use the orbifold KE metric and glue
Kronheimer ALE metrics at scale λ, giving metrics g\_λ on X with P(g\_λ)
→ 0 (Theorem Z). - The g\_λ equal g\_KE near the p\_i. The proof of 59.R
uses only E\_P(g) = 0 near p, so the local expansion is λ-independent. -
Let λ → 0 at fixed small ε.

So the strict inequality holds for every minimal surface of general type
and every k ≥ 1. Combined with Theorem K∞ (inf over Kähler metrics =
8π²k, not attained), this proves that the non-Kähler unrestricted
infimum lies strictly below the Kähler tariff for every such surface.
This was formal since Step 57.

Scope. These are upper bounds, i.e.~competitors. They prove nothing
about lower bounds, attainment, or the sign of α − 8π²k at the cone
interior R⁺ \textless{} 0.

\textbf{Corollary 59.Rig (minimizer rigidity, proved part).} Assume the
Anderson case-(I) ledger and Theorem 53.B. At a smooth P-critical bulk
carrying a Burns-type bubble, one has R⁺ ≤ 0. Moreover, wherever
λ\_max(R⁺) = 0 and the point has Ric° = 0 = W⁻ (or more generally
whenever the W⁻ and Ric° blocks are killed by explicit correctors), μ₂ =
μ₃. - The W⁻ part is now exact via k₀\^{}W. - The Ric° part is exact
once a closed-form Ric° corrector is certified; the run above searches
for it. - Otherwise the conclusion holds where \textbar Ric°\_p\textbar²
is small relative to (μ₂ − μ₃)², because the Ric° block at K = 0 is
finite.

\begin{center}\rule{0.5\linewidth}{0.5pt}\end{center}

\hypertarget{part-b.-the-structure-of-the-zero-gap-set}{%
\subsection{Part B. The structure of the zero-gap
set}\label{part-b.-the-structure-of-the-zero-gap-set}}

\hypertarget{b.1-three-sufficient-mechanisms}{%
\subsubsection{B.1 Three sufficient
mechanisms}\label{b.1-three-sufficient-mechanisms}}

\textbf{Proposition 59.N (noncollapsed mechanism; unifies Z, Z′, 56.B,
58.O).} Let (V, g\_V) be a compact Einstein orbifold of any sign, with
isolated singular points R⁴/Γ\_j. Let M be obtained by replacing each
point by a Ricci-flat ALE orbifold tree N\_j (possibly iterated)
asymptotic to R⁴/Γ\_j, with compatible orientation. Then α(M) = 0.

\emph{Proof.} Glue g\_V and λ²g\_\{N\_j\} on the annulus t ≤ r ≤ 2t,
with t = λ\^{}\{1/2\}. - In g\_V-normal orbifold coordinates, g\_V = δ +
O(r²). - Ricci-flat ALE metrics in 4D are δ + O(r⁻⁴) with derivatives,
in suitable coordinates (Bando--Kasue--Nakajima). - Both sides have Ric°
= 0, whatever the Einstein constants are. So Ric°(g) is supported on the
annulus, where \textbar Rm\textbar{} ≤ C(1 + λ⁴t⁻⁶). - Hence P ≤ C t⁴(1
+ λ⁴t⁻⁶)² = O(λ²) → 0. - Trees are handled inductively. ∎

Examples: - Kummer/K3: exact Einstein, so attained. - Canonical models
of general type with Kronheimer ALE resolutions (Z): attained iff a KE
metric exists on X itself, i.e.~K\_X is ample. - Wahl-point smoothings
by Şuvaina's Milnor-fibre ALE metrics (56.B): for nonminimal M this
requires non-Kähler Einstein Wahl orbifolds. - S⁴/Z₂ resolved by two
Eguchi--Hanson spaces (58.O): α(CP²\#CP²) = 0, with attainment
equivalent to the open existence question for an Einstein metric on
CP²\#CP². By Ozuch, no Einstein limit reaches the orbifold.

\textbf{Proposition 59.C (total-collapse mechanism).} α(M) ≤
36·Vol\_\{\textbar r\textbar\}(M) ≤ 36·MinVol(M). Here MinVol is
Gromov's minimal volume over metrics with \textbar K\textbar{} ≤ 1, and
Vol\_\{\textbar r\textbar\} is LeBrun's Ricci minimal volume over
metrics with \textbar Ric\textbar{} ≤ 3.

\emph{Proof.} If \textbar Ric\textbar{} ≤ 3, then \textbar Ric°\textbar²
≤ \textbar Ric\textbar² ≤ 36, so P(g) ≤ 36 Vol(g). ∎

Consequences: - A manifold with an F-structure of positive rank
(Cheeger--Gromov) has MinVol = 0, so α = 0. Examples include S¹ × Y³ for
any Y, nilmanifolds and Sol-manifolds, torus bundles, S² × T², primary
Kodaira surfaces, and elliptic surfaces with χ = 0. - For χ = τ = 0, an
Einstein metric must be flat. Chern--Gauss--Bonnet with Ric° = 0 gives
∫(\textbar W\textbar² + s²/24) = 0. - Hence on S¹ × Y with Y hyperbolic
or spherical, on S² × T², and on Kodaira surfaces: α = 0 and the infimum
is never attained. - This is the cleanest instance of ``zero gap without
Einstein metric and without any orbifold'': the defect vanishes by
collapse, not by approximation of an Einstein limit. - Conversely,
MinVol(M) = 0 forces χ = τ = 0. So this mechanism never applies to X \#
kCP̄².

Novelty check. LeBrun's ``Four-manifolds without Einstein metrics'' §5
and ``Weyl curvature, Einstein metrics and SW theory'' §5 (Knowledge)
study Vol\_\{\textbar K\textbar\}, Vol\_\{\textbar r\textbar\} and
Vol\_s and their SW lower bounds. The trivial comparison α ≤ 36
Vol\_\{\textbar r\textbar\} is not stated there. It is elementary and
should be presented as an observation, not a theorem of depth.

\textbf{Proposition 59.D (partial-collapse mechanism; new
construction).} Let (H, g\_H) be a complete finite-volume hyperbolic
4-manifold with cusps, or a complex-hyperbolic surface with cusps.
Truncate along cusp cross-sections N, which are compact flat or infranil
3-manifolds respectively. Then: - (a) the double D(H) = H\_T ∪\_N H̄\_T
satisfies α(D(H)) = 0; - (b) for a real-hyperbolic cusp with N = T³ and
any primitive slope γ ∈ H₁(T³), the Dehn filling H\_T ∪ (D² × T²) (with
γ bounding) satisfies α = 0.

\emph{Proof.}

\begin{enumerate}
\def\labelenumi{(\alph{enumi})}
\item
  Real case. The cusp metric is dt² + e⁻²ᵗ g\_N with g\_N flat. On N ×
  {[}T, T+2{]} use dt² + e\^{}\{−2f(t)\} g\_N, where f(t) = t for t ≤ T
  and f(t) = 2T + 2 − t for t ≥ T + 2, with \textbar f'\textbar{} ≤ 1
  and \textbar f'`\textbar{} ≤ C. The sectional curvatures of the warped
  product are −(f'² − f'') and −f'², which are bounded. The volume is
  O(e⁻³ᵀ) and \textbar Ric°\textbar{} ≤ C. Elsewhere the metric is
  hyperbolic, so P = O(e⁻³ᵀ).
\item
  Complex case. The cusp metric is dt² + e⁻²ᵗ g\_base + e⁻⁴ᵗ θ² on a
  Heisenberg nilmanifold. Use the same f, with base factor e\^{}\{−2f\}
  and fibre factor e\^{}\{−4f\}. The ratio fibre/base² stays 1, so the
  Nil curvature stays bounded. The volume is O(e⁻⁴ᵀ).
\item
  Write T³ = S¹\_γ × T² with a flat product metric. A path of flat
  metrics on T³ costs bounded curvature and O(e⁻³ᵀ) volume. Then use dt²
  + a(t)²dθ² + e⁻²ᵀ g\_\{T²\}, in four stages:
\end{enumerate}

\begin{enumerate}
\def\labelenumi{\arabic{enumi}.}
\tightlist
\item
  a shrinks like the cusp;
\item
  a turns around (bounded a'\,'/a);
\item
  a grows like e\^{}\{t − 2T\} up to size 1; this stage is H² × (tiny
  T²), with bounded curvature;
\item
  a caps off as a round disc of size 1. The total volume is O(e⁻²ᵀ) and
  the curvature is bounded. So P = O(e⁻²ᵀ) → 0. ∎
\end{enumerate}

Remarks. - χ(D(H)) = 2χ(H) \textgreater{} 0. The simplicial volume is
‖D(H)‖ = 2Vol(H)/v₄ \textgreater{} 0 (amenable boundary), so
MinVol(D(H)) \textgreater{} 0. Thus (ii) cannot apply, and (i) would
need a compact Einstein orbifold associated to D(H); no such orbifold is
known. - The Einstein region is a cusped complete Einstein manifold, and
the complement collapses along an F-structure: exactly Anderson's
case-(III) geometry. - Whether D(H) or the small-slope fillings carry
Einstein metrics is open in general. For large slopes, Einstein
Dehn-filling theorems in the literature are expected to apply, but they
are not used and not cited here. Gromov's simplicial-volume inequality
and Hitchin--Thorpe do not obstruct them, since τ = 0. - So this is a
third, independent zero-gap mechanism with attainment open. It is
constructed with no Ricci-flat ALE pieces and no orbifold.

\hypertarget{b.2-the-converse-conditional-only-on-andersons-theorem}{%
\subsubsection{B.2 The converse (conditional only on Anderson's
theorem)}\label{b.2-the-converse-conditional-only-on-andersons-theorem}}

\textbf{Theorem 59.Z (structure of the zero-gap set).} Let M be closed
and oriented. Assume Anderson's Theorem 4.10 for ℛ² = 32π²χ + 4P
(Knowledge: \emph{Canonical metrics on 3-manifolds and 4-manifolds}, Thm
1.1/4.10). If α(M) = 0, then M carries one of the following three
configurations: - (I) a compact Einstein orbifold (V, g₀) associated to
M, i.e.~V is obtained from M by collapsing cycles (and possibly
reducible); - (III) a nonempty maximal Einstein orbifold domain (Ω, g₀)
with 0 \textless{} vol ≤ 1. It is complete with finite volume, weakly
embedded in M, and has s \textless{} 0. Outside a compact set, M carries
an F-structure along which the sequence collapses with locally bounded
curvature, off finitely many shrinking balls; - (II) collapse everywhere
with locally bounded curvature along F-structures, off finitely many
shrinking balls.

Conversely: - every (I)-configuration whose singular points admit
Ricci-flat ALE resolutions or smoothings realizing M gives α(M) = 0
(Prop. 59.N); - MinVol(M) = 0 gives α(M) = 0 (Prop. 59.C), which is a
special (II); - the doubles and Dehn fillings of 59.D realize (III).

In every case, α(M) is attained iff M carries an Einstein metric.

\emph{Proof of the forward direction.} Anderson's minimizers g\_ε of
Z\_ε = P + ε∫(1+\textbar Ric\textbar²)\^{}p satisfy P(g\_ε) → α = 0. On
the regular part they converge smoothly on compact sets in cases (I) and
(III). Lower semicontinuity on compact sets gives
∫\_K\textbar Ric°(g₀)\textbar² ≤ liminf P(g\_ε) = 0 for every compact K
in the regular set. So g₀ is Einstein.

In (III), Ω cannot be Ricci-flat: complete noncompact Ricci-flat
manifolds have at least linear volume growth, so they cannot have finite
volume. If the scalar curvature were positive, Ω would be compact by
Myers. So s \textless{} 0. ∎

Notes. 1. Anderson's inequality F(g₀) ≤ inf F is for ℛ² and carries a
χ\_orb shift, so it is not used. Only the smooth convergence of the
regular part is used. 2. The bubble structure in (I) (Ricci-flat ALE
bubble trees with P = 0, and M diffeomorphic to the resolution) needs
the no-neck and topology ledger H\_sm of Step 58. With it, (I) is
exactly ``M is a Ricci-flat-ALE-tree resolution of an Einstein
orbifold''.

\hypertarget{b.3-consequences-for-x-kcpuxb2}{%
\subsubsection{B.3 Consequences for X \#
kCP̄²}\label{b.3-consequences-for-x-kcpuxb2}}

Let M = X \# kCP̄² with X minimal of general type and k ≥ 1. - χ
\textgreater{} 0, so MinVol \textgreater{} 0. Prop. 59.C never applies.
- Kähler orbifolds in (I) are excluded: their smoothings or resolutions
are minimal, which leads to 56.B and 55.6. - So α(M) = 0 would need one
of: a non-Kähler Einstein orbifold associated to M (Wahl/T points, or
more generally points with Ricci-flat ALE fillings); a type-(III) cusp
domain whose F-structure complement carries the exceptional classes; or
type-(II) collapse whose bad balls carry all of χ.

Corollary 59.KE says only that α(M) \textless{} 8π²k. The gap between 0
and this bound is exactly the unresolved O11 problem: the open-range
inequality C₊ ≥ A and the k²/A law.

The common principle behind the three mechanisms: P vanishes in the
limit iff M is exhausted by Einstein regions, up to transition regions
whose volume vanishes in the natural curvature scale. The three ways
such a transition can disappear give the trichotomy: - shrinking scale:
Ricci-flat ALE necks; - collapsing fibres along part of M: cusps; -
collapsing fibres along all of M: F-structures.

\begin{center}\rule{0.5\linewidth}{0.5pt}\end{center}

\hypertarget{part-c-step-60.-corrections-and-independent-confirmations-now-in-force}{%
\subsection{Part C (Step 60). Corrections and independent confirmations
now in
force}\label{part-c-step-60.-corrections-and-independent-confirmations-now-in-force}}

\textbf{C.1 W⁻ sector, independently confirmed; a representation
artifact in the exact script.} - The exact script
\texttt{step59\_exact\_log\_coefficient.py} returns a pair-basis 6×6
matrix X for the linearised W⁺ operator. In sectors with X₀ = 0 (pure W⁻
or pure Ric° jets), its σ¹ part X₁ contains a spurious antisymmetric
piece N. N does not depend on K. - The script's signed density 4 tr X₁²
+ s₁²/12 is therefore wrong in those sectors by −4\textbar N\textbar². A
true \textbar W⁺\textbar² must be ≥ 0, so a negative value cannot be
correct. - \texttt{step60\_jax\_check.py} makes an independent JAX
evaluation with orthonormal-frame W⁺ (Cholesky frame) of the σ²t²
coefficient of F₊(δ + σA + t(h₂ + σK)) at a rational point. It gives: -
Wm, K = 0: JAX 13.2439633 = sym-only 13.2439633; the script's signed
value is 13.131. - Wm, K = k₀\^{}W: JAX −5·10⁻¹⁶; sym-only exactly 0;
the script's signed value is −0.113. - Rz, K = 0: JAX 0.95193 =
sym-only. - Rz, K = k₀\^{}W: JAX 0.27206 = sym-only. - So the true
coefficient is 4\textbar sym X₁\textbar² + s₁²/12. - A01's k₀\^{}W =
(3/2)sym(h\_W k\_B) − (1/4)tr(h\_W k\_B)δ kills the W⁻ chiral residual
pointwise (sym X₁ = 0 and s₁ = 0 exactly in rational arithmetic). This
independently confirms A01's Theorem 2.1. It also shows that k₀\^{}W
does not solve the Ric°\_ζ residual. - The (s, W⁺) sector, where X₀ ≠ 0,
is unaffected. \texttt{step60\_jax\_check\_Wplus.py} matches script and
JAX to 15 digits, for K = 0 and for a generic invariant K. So the exact
Φ of §A.3 stands. - Corrector searches now use
\texttt{step60\_corrector\_columns.py}, which solves sym X₁ = 0 and s₁ =
0 exactly. The older \texttt{step59\_exact\_wminus\_ric\_correctors.py}
is superseded.

\textbf{C.2 Sign scope of Φ (A01 \#833).} - Φ ≤ 0 is a statement on the
kernel μ₁ = 0 only. In the interior R⁺ \textless{} 0 the logarithmic
coefficient can be positive: W⁺ = 0 gives Φ = s²/288 \textgreater{} 0,
and −μ₁(μ₁ − s/8) \textgreater{} 0 for s/8 \textless{} μ₁ \textless{} 0.
There it is subordinate to the positive first-order cost. - So ``Φ ≤ 0
on the whole cone'' is withdrawn wherever it appeared. The strictly
negative cone, e.g.~real hyperbolic, is a genuine survivor, not a
remainder. - Theorem 59.R/§A.3 assumes ζ is an eigenvector of R⁺ (w₁₂ =
w₁₃ = 0). If μ\_ζ = R⁺(ζ,ζ) = 0 but R⁺ζ ≠ 0, the mixed block makes the
degree-0 log coefficient unbounded below (A01: null corrector direction
L\_\{(log r)J₂x\}δ). That case has R⁺ indefinite, which already gives a
first-order saving. - On R⁺ ≤ 0 a kernel vector is automatically an
eigenvector, so 59.Rig is unaffected.

\textbf{C.3 Independent remainder (A01 Theorem 52.A).} - A01 proves, for
Einstein backgrounds with arbitrary W⁻ and uniformly in (p, ζ): P(g\_ε)
= P(g) + 8π² − 4π²\textbar T\textbar²ε²log(1/ε) + O(ε²). - Their order-ε
flux computation, with f − \textbar Ric°\textbar²dV = 8π²p₁(Λ⁺), gives
the order-ε term −8π²εR⁺(ζ,ζ) exactly. - This agrees with our
normalisation 2κ = −8π² in 59.R, now obtained independently, without
calibration by Ch. 5 or Theorem F. - Our 59.R covers the larger class of
P-critical backgrounds (E\_P = 0 near p). - A01's Lemma 52.C (no open
degenerate stratum (a, −a/2, −a/2) when s \textless{} 0) uses δW⁺ = 0,
so it needs an Einstein bulk. - The Anderson case-(I) bulk of 53.B is
only P-critical. Lemma 52.C must therefore be replaced there, or the
bulk shown to be Einstein (e.g.~by the Bach-flat route of Step 54).

\textbf{C.4 Prop. 59.D: literature position, which changes how it is
stated (sources requested this Step; citation keys pending).} - Anderson
(Dehn filling and Einstein metrics in higher dimensions) and Bamler
(generalized Dehn filling) construct, for all but finitely many slopes
per T³ cusp, Einstein metrics of negative scalar curvature on the
fillings. These are GH-close to the cusped hyperbolic metric. - For such
long slopes our α = 0 is therefore attained. This confirms 59.Z: the
Einstein metric is an unconditional instance, and our collapsing
sequence converges to the regime-(III) configuration, while the Einstein
metrics themselves are a regime-(I) family converging to it. - The
finitely many short slopes form a finite test family. There α = 0 holds
by our interpolation, and 59.Z makes attainment equivalent to Einstein
existence, which no filling theorem decides. - The doubles D(H) along T³
(or flat, or Nil) necks are the cleaner new examples. A warped Einstein
neck dt² + φ(t)²g\_flat requires φφ'\,' = φ'², i.e.~φ = e\^{}\{ct\}, a
cusp. So no Einstein T³ × I neck exists, and no filling theorem applies.
To our knowledge Einstein existence on D(H) is open. - Existence of
hyperbolic 4-manifolds with T³ cusp sections is to be cited from the
requested sources before use.

\hypertarget{part-d-step-61.-verification-of-parts-ac-ttcs-s60_t24_tool1.-these-corrections-are-binding-and-override-ac.}{%
\subsection{Part D (Step 61). Verification of Parts A--C (TTCS
s60\_t24\_tool1). These corrections are binding and override
A--C.}\label{part-d-step-61.-verification-of-parts-ac-ttcs-s60_t24_tool1.-these-corrections-are-binding-and-override-ac.}}

Sources:
\texttt{trace/ttcs/results/step60\_ttcs01\_verifier1\_integration.md}
(Part A) and \texttt{step60\_ttcs01\_verifier2\_integration.md} (Part B,
C.4).

\textbf{D.1 Theorem 59.R survives, with these repairs.} - The insertion
support B\_\{2ρ\}(p) must lie entirely inside the region where E\_P(g) =
0. Without that, an extra bulk first-variation term remains. -
Rescaling. Write h(ry) = r²q(r,y) with q smooth on the shell (Taylor
with integral remainder). The rescaled family is G\_\{σ,r\} = B\_σ +
r²(q(r,y) + σK). It agrees with the cut-off metric only up to two-jets
on \textbar y\textbar{} = 1 for 2R√ε ≤ r ≤ ρ, which is all the density
needs. - Why 3-jets and 4-jets cannot enter the log term: 𝒞(B\_σ) = 0
and F₊ = \textbar 𝒞\textbar², so 𝒞(G) = r²C₂ + r³C₃ + \ldots. The r⁴
coefficient is \textbar C₂\textbar² alone. The third jet first appears
at r⁵, the fourth jet and nonlinear terms at r⁶, and these give a finite
∫dr/r. No parity assumption on K is needed. - The inner region
contributes O\_R(ε²). The phrase ``O(ε²R⁴)'' is withdrawn, since no
uniformity in R is needed or proved. - Order ε: a factor-two error in
the distributional flux is corrected. With M\_ζ = 3ζ⊗ζ − I and 𝒞 = (2W⁺,
s/√12), c\_A = m(M\_ζ/6, 1/√12), not M\_ζ/12. Also s'(∂∂̄φ) = −¼Δ²φ and
∫Δ²(χ log r²) = 16π², so m = −4π² and E₁ = −8π²R⁺(ζ,ζ). Hence κ = −4π²,
matching Step 20 and A01's θ₂ = 8π²μ\_ζ independently. The displayed
final formula of 59.R was already right; only the intermediate c\_A
changes. - Corrected statement. Under the hypotheses above, for every
unit ζ and every fixed degree-0 K, with \textbf{no eigenvector
hypothesis}: P(g̃\_ε) = P(g) + 8π² − 8π²εR⁺\_p(ζ,ζ) +
4π²Φ(Rm\_p,ζ,K)ε²log(1/ε) + O(ε²). The explicit optimised value Φ\_opt =
−μ₁(μ₁ − s/8) − \textbar T\textbar² needs two further assumptions:
Ric°\_p = 0, and ζ an eigenvector of R⁺\_p.~It is attained by K* +
k₀\^{}W. Here T is the trace-free part of R⁺ restricted to ζ\^{}⊥.

\textbf{D.2 Corrected optimal corrector.} With u⊙v = ½(u⊗v + v⊗u), as in
the script, the mixed coefficient in K* is 2w₂₃/3: K* = (s/24)dr² +
(−s/48 + 7w₁₁/6 + w₂₂/3)s₁² + (2w₂₃/3)s₁⊙s₂ + (−s/48 + 5w₁₁/6 −
w₂₂/3)s₂². The printed 4w₂₃/3 is withdrawn. The KE⊥ specialisation,
where w₂₃ = 0, is unchanged.

\textbf{D.3 The W⁻ artifact, explained.} The old script used the frozen
projector P₊ at t = 0. The true derivative includes a
projector-variation term ½(H'₊₋R₋₊ − R₊₋H'₋₊), which is skew-adjoint and
cancels exactly the skew part of X. At order σ this defect depends on h₂
but not on K. Hence the geometric coefficient is 4\textbar sym
X₁\textbar² + s₁²/12. - The defect vanishes on the eigenvector locus, so
the exact (s,W⁺) computation at w₁₂ = w₁₃ = 0 is unaffected. - C.1's
unrestricted sentence is replaced by: ``the exact (s,W⁺) coefficient of
§A.3, restricted to w₁₂ = w₁₃ = 0, is unaffected; the JAX check tests
this locus only.'' - The \texttt{leak\_solvable} flag in
\texttt{step59\_exact\_log\_coefficient.json} records only that
\texttt{sympy.solve} returned a dictionary. It says nothing about the
mixed block. A01's Hessian-null direction is the authoritative statement
there. - Schur cancellation of the W⁻ cross terms is proved for the
\textbf{whole} integrated σ² coefficient. The polarised invariant linear
functional on the spin-2 representation vanishes, and this covers the
metric-dependent norm and volume terms as well.

\textbf{D.4 Sign scope (replaces C.2 and C.3 openings).} Theorem 59.R
holds for every unit ζ and fixed K. The optimised formula assumes ζ is
an eigenvector. On R⁺ ≤ 0, a zero Rayleigh direction is automatically a
kernel eigenvector. - On the kernel, Φ\_opt = −\textbar T\textbar² ≤ 0.
This holds for the optimised coefficient with the W⁻ sector cancelled.
It is not a statement about arbitrary K, nor about general P-critical
jets whose Ric° residual is uncancelled. - 59.Rig in the Ric° sector
must be stated quantitatively: the strict conclusion holds when
\textbar T\textbar² \textgreater{} Q\_Ric(Ric°\_p) for an explicit Ricci
residual, sufficiently when \textbar T\textbar² \textgreater{}
C\textbar Ric°\_p\textbar². - C.3 opening: ``For Einstein backgrounds,
unit directions with R⁺ζ = 0, and arbitrary W⁻, A01 proves the expansion
uniformly on compact families of insertion data.''

\textbf{D.5 Corollary 59.KE: order of limits.} Choose the insertion
balls in the regular locus, inside fixed neighbourhoods where g\_λ =
g\_KE. Then the insertion difference I(ε) is independent of λ, not
merely uniform in it. Fix ε first so that the logarithmic saving
dominates the remainder, then choose λ so that P(g\_λ) is smaller than
that saving. This produces an actual smooth competitor below 8π²k. It is
an upper bound only. The comparison with K∞ uses Step 51/52.

\textbf{D.6 Prop. 59.N (corrected).} - The annulus estimate is P\_ann ≤
C(t⁴ + (λ/t)\^{}\{2τ\}) for derivative-controlled ALE order τ
\textgreater{} 0. - Order τ ≥ 2 gives O(λ²) with t = √λ. Any τ
\textgreater{} 0 gives zero gap with t = λ\^{}\{τ/(τ+2)\}. - Order-4
decay gives a metric error λ⁴r⁻⁴, not λ⁶r⁻⁴. - Trees: perform each
insertion at its parent's native scale. - Compatible orientation means a
cone-end identification through an orthogonal conjugacy, with induced
boundary orientations reversed. - Statement: a smooth closed manifold
obtained from a compact Einstein orbifold with isolated irreducible cone
points by compatible finite Ricci-flat ALE-tree replacements has α = 0.
Reducible vertices are not covered. - For canonical models, ``attained
iff K ample'' is withdrawn. Only ``attained iff the smooth manifold
admits an Einstein metric'' is kept. - (Step 61 sharpens the rate for EH
points to Θ(a⁴); see \texttt{step61\_EH\_quantitative\_obstruction.md}.)

\textbf{D.7 Prop. 59.C (corrected).} - Under the spectral bound −3g ≤
Ric ≤ 3g one has \textbar Ric°\textbar² ≤ 36 (equality at eigenvalues
(3,3,−3,−3)). Hence α ≤ 36 Vol\_\{\textbar r\textbar\} ≤ 36 MinVol. The
proof's ``\textbar Ric\textbar{} ≤ 3'' mixed two norms and is replaced.
- MinVol = 0 ⇒ χ = τ = 0 (Chern--Weil). - For F-structures, use a
polarised positive-rank F-structure or an explicit bounded-curvature
collapsing sequence. - Nonattainment needs only χ = 0, since Einstein
with χ = 0 is flat. The examples S¹×Y (Y hyperbolic or spherical), S²×T²
and primary/secondary Kodaira surfaces are confirmed: none is flat.

\textbf{D.8 Prop. 59.D (corrected and completed).} - Real-hyperbolic
doubles: \textbar Ric°\textbar² = 3(f'`)² exactly for dt² +
e\^{}\{−2f\}g\_flat, so P\_neck = 3Vol(N)∫(f'`)²e\^{}\{−3f\} =
O(e\^{}\{−3T\}). - Complex-hyperbolic Heisenberg doubles: the full
coframe calculation (dω³ = −2f'ω⁰³ + κω¹²) shows that all curvature
coefficients are polynomials in f'`, f', κ. This gives a full curvature
bound, P = O(e\^{}\{−4T\}). - T³ fillings: every primitive slope works.
Two steps are added: a flat-metric transition G(s) to a product adapted
to γ, and the turnaround/growth/round-cap sequence. Together they give P
≤ C\_γ e\^{}\{−2T\}. - C.4 corrections: - ``finitely many short slopes''
is finite only per cusp, not for slope tuples; - ``no Einstein T³×I neck
exists'' is false (exponential segments exist); the correct statement is
that there is no \textbf{turning} isotropic warped Einstein neck; - the
openness of Einstein existence on doubles is not source-supported and is
withdrawn as a claim; - the simplicial-volume formula is real-hyperbolic
only. - Source keys: Anderson \textbackslash cite\{anderson2003dehn\}
(ODE gluing, ``sufficiently large fillings'', Ric = −(n−1)g, smooth
convergence to the cusp manifold); Bamler
\textbackslash cite\{bamler2009construction\} (metadata only);
Kolpakov--Martelli \textbackslash cite\{kolpakov2013hyperbolic\}
(one-cusped hyperbolic 4-manifolds with T³ cusp).

\textbf{D.9 Theorem 59.Z (corrected statement).} Assume Anderson Thm
4.10 together with Einstein removable-singularity regularity. If α(M) =
0, there \textbf{exists a selected minimizing sequence} exhibiting one
of three configurations: 1. a compact, possibly reducible, Einstein
orbifold V of volume 1; 2. a nonempty maximal orbifold domain Ω,
Einstein and complete componentwise (completeness follows from the
Einstein lower Ricci bound via Anderson Rem. 4.12), with s \textless{} 0
on every \textbf{noncompact} component and F-collapse outside compact
sets; 3. total collapse. Further points: - Thm 4.10 itself has only
cases (I) and (II); the three labels are those of Thm 1.1. - The
alternatives belong to selected sequences and are not mutually exclusive
properties of M. The long-slope fillings show both (1) and (2) on one M.
- Zero-P bubbles are Ricci-flat only under bubble-scale convergence. -
The ALE-tree reconstruction must cite Anderson's discussion around
(4.65). Lower semicontinuity and H\_sm do not supply it. - Attainment ⇔
M admits an Einstein metric. This is immediate from P = 0 ⇔ Ric° = 0. -
The ``common principle'' is a \textbf{sufficient construction
principle}, with Anderson supplying a selected reverse decomposition. It
is not a classification of every zero-gap sequence.

\hypertarget{files}{%
\subsection{Files}\label{files}}

\begin{itemize}
\tightlist
\item
  \texttt{workspace/scripts/step59\_exact\_log\_coefficient.py} →
  \texttt{data/step59\_exact\_log\_coefficient.json}: exact (s, W⁺)
  sector, quadratic form, optimal K.
\item
  \texttt{workspace/scripts/step59\_exact\_wminus\_ric\_correctors.py} →
  \texttt{data/step59\_exact\_correctors\_\{Wm,Rz,Rp\}.json}.
\end{itemize}

\hypertarget{s1003a02-step-61-the-eguchihanson-obstruction-as-a-sharp-gap-estimate-theorem-61.eh}{%
\chapter{S1003A02: Step 61 --- The Eguchi--Hanson obstruction as a sharp
gap estimate (Theorem
61.EH)}\label{s1003a02-step-61-the-eguchihanson-obstruction-as-a-sharp-gap-estimate-theorem-61.eh}}

Source: \texttt{step61\_EH\_quantitative\_obstruction.md}.

Agent2, Execute. Target 58.Q at leading order. Inputs: - Step 60:
\texttt{scripts/step60\_EH\_obstruction\_pairing.py},
\texttt{scripts/step60\_EH\_sector\_check.py}, data
\texttt{data/step60\_EH\_obstruction\_pairing.json}. - Step 61:
\texttt{scripts/step61\_EH\_cokernel\_norm.py} →
\texttt{data/step61\_EH\_cokernel\_norm.json}, which checks the exact
cokernel norm. - Source: Ozuch, \emph{Higher order obstructions to the
desingularization of Einstein metrics}
\textbackslash cite\{ozuch2020higher\}, §§1.2, 2.1, Thm 0.2. His linear
theory follows Biquard (Biq13, Biq16).

\hypertarget{objects-and-normalisations}{%
\subsection{1. Objects and
normalisations}\label{objects-and-normalisations}}

\textbf{Orbifold.} (V, g\_V) is an orbifold with an isolated point p ≅
ℝ⁴/ℤ₂, and g\_V is Einstein near p: Ric = Λg, s = 4Λ. - R\^{}± = W\^{}±
+ (s/12)I on Λ\^{}±, so tr R⁺ = Λ. - In Bianchi/normal coordinates, g\_V
= δ + h₂ + h₄ + O(r⁶). Odd jets vanish by ℤ₂-invariance in orbifold
charts.

\textbf{Eguchi--Hanson space.} EH\_a in its hyperkähler orientation is
dr²/f + (r²/4)(σ₁² + σ₂²) + (r²f/4)σ₃², f = 1 − a⁴/r⁴, r ≥ a. - Its
level sets are RP³, and its volume element is π²r³dr after angular
integration. - The Kähler forms ω₁, ω₂, ω₃ are self-dual and parallel. -
The resolution direction ζ = ω₁/√2 is the unit self-dual form whose
complex structure makes the bolt S² holomorphic.

\textbf{Cokernel.} Ω = (e⁰³ − e¹²)/r⁴ = −¼ d(σ₃/r²) in EH orthonormal
frames, exactly for every a. This is because d(σ₃/r²) = −(4/r⁴)(e⁰³ −
e¹²) does not involve f. - Ω is closed, anti-self-dual and L², hence
harmonic. - o\_i := Ω∘ω\_i is symmetric, trace-free and TT, and lies in
the L²-kernel of the Einstein operator P\_EH. These o\_i span Ozuch's
O(eh), the infinitesimal variations of the hyperkähler period ζ. -
\textbar o\_i\textbar² = \textbar Ω\textbar²\_endo = 4/r⁸, and ⟨o\_i,
o\_j⟩ = 0 for i ≠ j, since ω\_iω\_j is anti-commuting and trace-free. -
**‖o\_i‖²\_\{L²(EH\_a)\} = ∫\_a\^{}∞ (4/r⁸)π²r³dr = π²/a⁴.** This is
independent of i and scale invariant once multiplied by a⁴.

\textbf{Flux pairing (Step 60, exact).} Let ĥ be any smooth extension of
h₂ over EH with ĥ − h₂ = O(r\^{}\{−2+ε\}) at infinity. Since DRic*\_EH
o\_i = 0 on the complete space EH, the pairing ∫⟨DRic\_EH{[}ĥ{]}, o\_i⟩
is a flux at infinity. It equals the flat flux of (h₂, o\_i\^{}flat),
and is constant in r because ⟨DRic\_flat h₂, o⟩ = ⟨(s/4)δ, o⟩ = 0
pointwise. The result is

B\_i := ∫\_EH ⟨DRic\_EH{[}ĥ{]}, o\_i⟩ = π² (R⁺\_p ζ)\_i, i = 1, 2, 3,

computed in the frame (ζ, ζ₂, ζ₃) at a = 1. It vanishes for every W⁻,
Ric°\emph{ζ, Ric°}⊥ and mixed jet
(\texttt{step60\_EH\_sector\_check.py}). - The scaling is as follows. In
y = x/a units, a⁻²g̃ = EH₁ + a²ĥ(y) + O(a⁴\textbar y\textbar⁴), so the
pairing is a²B\_i. - This makes explicit Ozuch's constant: his (22),
λ\_k(ζ) = c⟨R⁺(H₂)ζ, ζ\_(k)⟩ with c \textgreater{} 0 unspecified, is our
B with c fixed.

\hypertarget{theorem-61.eh}{%
\subsection{2. Theorem 61.EH}\label{theorem-61.eh}}

Let g̃\emph{\{a,ζ\} be the naïve desingularisation: EH\_a glued at p in
direction ζ, at a fixed gluing radius t₀, with g\_V unchanged outside
B}\{2t₀\}(p).

\textbf{Weighted neighbourhood.} For δ \textgreater{} 0 and a weight w ∈
{[}0, 2), define 𝒩\_δ(a,ζ) as the set of metrics g on the resolution M
such that, after a diffeomorphism, for j = 0, 1, 2 in y = x/a units,
\textbar∇\^{}j\_\{EH\}(a⁻²g − a⁻²g̃\_\{a,ζ\})\textbar{} ≤ δ a² (1 +
\textbar y\textbar)\^{}\{w−j\}. This bound is required on
\{\textbar y\textbar{} ≤ t₀/a\}, with the corresponding unweighted bound
δa\textsuperscript{\{2−w\}t₀}\{w\} outside.

\textbf{(a) Lower bound.} There are constants a₀ and C, depending on
g\_V near p, on t₀ and on w, such that for a \textless{} a₀ and every g
∈ 𝒩\_δ(a,ζ):

P(g) ≥ π² a⁴ ( \textbar R⁺\_p ζ\textbar{} − C(δ + a\^{}\{2−w\}) )₊².

\textbf{(b) Matching upper bound.} There is a metric g*\emph{\{a,ζ\} ∈
𝒩}\{Ca\^{}\{2−w\}\}(a,ζ) with

P(g*\_\{a,ζ\}) = π² a⁴ \textbar R⁺\_p ζ\textbar² + O(a⁶ log(1/a)).

Hence, uniformly in ζ:

inf\_\{𝒩\_δ(a,ζ)\} P = π² a⁴ \textbar R⁺\emph{p ζ\textbar² (1 + O(δ) +
o(1)), inf\_ζ inf}\{𝒩\_δ\} P ≈ π² a⁴ σ\_min(R⁺\_p)².

\textbf{(c) Near-extremal structure.} If g ∈ 𝒩\_δ satisfies P(g) ≤
π²a⁴\textbar R⁺ζ\textbar²(1 + η), then ‖Ric°(g) − Π Ric°(g)‖²\_\{L²\} ≤
(η + C(δ + a\^{}\{2−w\})) π²a⁴\textbar R⁺ζ\textbar², where Π is the
L²-projection onto span\{χô\_i\}. In words, near-minimisers are Einstein
modulo the EH cokernel, and their Einstein defect is the cokernel tensor
a²Σ B\_i o\_i/‖o‖², concentrated at the EH scale. When R⁺ζ ≠ 0,
Cauchy--Schwarz equality forces Ric° to be parallel, in L², to Σ(R⁺ζ)\_i
o\_i.

\hypertarget{proof-of-a}{%
\subsubsection{Proof of (a)}\label{proof-of-a}}

Let ô = Σ\_i b\_i χ\_Λ o\_i with b = R⁺ζ/\textbar R⁺ζ\textbar, where
χ\_Λ cuts off at \textbar y\textbar{} ≈ Λ := t₀/(2a). This cut-off lies
inside the region where g̃ = g\_V, hence Ric°(g̃) = 0 there. For any g,
Cauchy--Schwarz gives

P(g) ≥ ⟨Ric°(g), ô⟩\_g² / ‖ô‖\_g².

Everything below is computed in y-units, using the scale invariance of P
and of L² norms of 2-tensors in dimension four.

\begin{enumerate}
\def\labelenumi{(\roman{enumi})}
\item
  Since ‖o‖²\_EH = π² (a = 1) and the tail is ∫\_Λ\^{}∞ 4y⁻⁸y³dy =
  O(Λ⁻⁴), we have ‖ô‖²\_g = π²(1 + O(a⁴ + δa²)).
\item
  The naïve metric. In the core \textbar y\textbar{} ≤ Λ we have a⁻²g̃ =
  EH + a²ĥ + a⁴ĥ₄ + O(a⁶\textbar y\textbar⁶), with ĥ = h₂ off a compact
  set. Expanding, Ric°(a⁻²g̃) = a²DRic°\_EH{[}ĥ{]} + a⁴(Q(ĥ,ĥ) +
  DRic{[}ĥ₄{]}) + \ldots.
\end{enumerate}

\begin{itemize}
\tightlist
\item
  The trace-free projection of DRic\_EH{[}ĥ{]} and DRic\_EH{[}ĥ{]}
  itself have the same pairing with o, because o is trace-free.
\item
  Hence ⟨Ric°(g̃), ô⟩ = a²\textbar B\textbar{} + E, where E collects the
  a⁴ terms paired with o.
\item
  A naive bound would be
  \textbar y\textbar²·\textbar y\textbar⁻⁴·\textbar y\textbar³
  integrated up to Λ, i.e.~O(a⁴Λ²) = O(a²t₀²). But in the region 1 ≪
  \textbar y\textbar{} ≤ Λ, the order-a⁴ part of Ric° is the order-r²
  part of Ric°(g\_V) = 0, by the Einstein equation for the jet: P\_e H₄
  = ΛH₂ − Q(H₂,H₂) up to trace. Only EH-tail cross terms survive there,
  of size O(\textbar y\textbar²·\textbar y\textbar⁻⁴).
\item
  So the a⁴ integrand is
  O(\textbar y\textbar⁻²·\textbar y\textbar⁻⁴·\textbar y\textbar³),
  which is integrable, and E = O(a⁴), uniformly in t₀. The
  gluing-annulus terms are O(a⁴t₀⁻⁶)-type and pair with ô to give O(a⁴).
\end{itemize}

\begin{enumerate}
\def\labelenumi{(\roman{enumi})}
\setcounter{enumi}{2}
\tightlist
\item
  Perturbation. Let k := a⁻²(g − g̃). Then ⟨Ric°(g), ô⟩ − ⟨Ric°(g̃), ô⟩ =
  ⟨k, DRic°*\_\{g̃\} ô⟩ +
  O(∫\textbar k\textbar\textbar∇²k\textbar\textbar ô\textbar{} +
  \textbar∇k\textbar²\textbar ô\textbar).
\end{enumerate}

\begin{itemize}
\tightlist
\item
  DRic*\_EH o\_i = 0 exactly, so DRic*\_\{g̃\}ô =
  a²·O(\textbar y\textbar⁻⁴) + O(χ'\_Λ terms). The first term comes from
  g̃ − EH = a²ĥ, which multiplies the background coefficients by
  O(\textbar ĥ\textbar, \textbar∇ĥ\textbar, \textbar∇²ĥ\textbar)
  \textasciitilde{} a²\textbar y\textbar\^{}\{2−j\}.
\item
  The linear term is therefore ≤ ∫\_\{\textbar y\textbar≤Λ\}
  δa²\textbar y\textbar\^{}w · a²\textbar y\textbar⁻⁴ ·
  \textbar y\textbar³ dy + δa²Λ\^{}\{w\}·Λ⁻⁶·Λ⁴ ≤ Cδ a⁴ Λ\^{}\{w\} +
  Cδa²Λ\^{}\{w−2\}. Since a⁴Λ\^{}w = a\textsuperscript{\{4−w\}t₀}w =
  o(a²) for w \textless{} 2, this is ≤ Cδ a²
  (t₀\textsuperscript{\{w\}a}\{2−w\} +
  a\textsuperscript{\{2−w\}t₀}\{w−2\}) ≤ C δ a².
\item
  The quadratic term is ≤ ∫
  δ²a⁴\textbar y\textbar\^{}\{2w−2\}\textbar y\textbar⁻⁴\textbar y\textbar³dy
  ≤ Cδ²a⁴Λ\^{}\{2w−2\} = Cδ² a\^{}\{6−2w\} = o(a²) for w \textless{} 2.
  Outside \textbar y\textbar{} ≤ Λ, ô vanishes.
\end{itemize}

\begin{enumerate}
\def\labelenumi{(\roman{enumi})}
\setcounter{enumi}{3}
\tightlist
\item
  Combining (i)--(iii): \textbar⟨Ric°(g), ô⟩\textbar{} ≥
  a²(π²\textbar R⁺ζ\textbar{} − C(δ + a\^{}\{2−w\})), so P ≥
  a⁴(π²\textbar R⁺ζ\textbar{} − ⋯)²/(π²(1 + o(1))) =
  π²a⁴(\textbar R⁺ζ\textbar{} − ⋯)². □
\end{enumerate}

(The error exponents above are bookkeeping-level. The decisive
structural facts are four: the pairing is a flux, o is TT and in the
kernel of DRic\emph{\emph{EH, the decay ô \textasciitilde{}
\textbar y\textbar⁻⁴ makes every error integral converge for w
\textless{} 2, and ‖o‖² is finite. A full Hölder-space version would
follow Ozuch's decoupled norms C\^{}\{2,α\}}\{β,}\}, App. A of
\textbackslash cite\{ozuch2020higher\}.)

\hypertarget{proof-of-b}{%
\subsubsection{Proof of (b)}\label{proof-of-b}}

This uses the Einstein-modulo-obstruction solution, which exists by the
Biquard--Ozuch linear theory and is quoted in
\textbackslash cite\{ozuch2020higher\} (22)--(23), Lemma 2.3: there is a
unique ĥ₂ on EH with

P\_EH ĥ₂ = Λ·eh + Σ\_k λ\_k o\_k, ĥ₂ = H₂ + H₂⁴ + O(r\^{}\{−4+ε\}), ĥ₂ ⊥
o\_k on the bolt,

where P\_EH = DRic + δ\emph{B is the gauged linearisation and
\textbar H₂⁴\textbar{} \textasciitilde{} r⁻². - Pairing with o\_k
identifies λ\_k‖o‖² = B\_k, i.e.~λ\_k = (R⁺ζ)\emph{k. The Λ·eh term
carries the trace part s = 4Λ; it is pure trace and drops out of Ric°. -
Set g* = g̃ + a²χ}\{t₀\}(ĥ₂ − h₂), so that in y-units a⁻²g} = EH + a²ĥ₂ +
O(a⁴\textbar y\textbar⁴). Ozuch's construction places the solution in
Bianchi gauge, so B(ĥ₂) = 0 is part of the linear solution, and
Ric(a⁻²g\emph{) = a²Λ·eh + a² Σλ\_k o\_k +
O(a⁴(1+\textbar y\textbar)\^{}\{0\}). - Hence Ric°(g}) = a²Σλ\_k o\_k +
R₄, with \textbar R₄\textbar{} ≲ a⁴(1+\textbar y\textbar)\^{}\{−2\} in
the core, plus a gluing-annulus error O(a⁴ t₀\^{}\{−4\} \ldots{} ). -
Then P(g\emph{) = a⁴\textbar λ\textbar²‖o‖² + 2a²⟨Σλ\_k o\_k, R₄⟩ +
‖R₄‖². With \textbar λ\textbar²‖o‖² = π²\textbar R⁺ζ\textbar²: - ‖R₄‖² ≲
a⁸∫(1+\textbar y\textbar)\^{}\{−4\}\textbar y\textbar³dy ≲ a⁸ log Λ =
O(a⁸ log(1/a)); - the cross term is ≲ a⁶
∫(1+\textbar y\textbar)\^{}\{−4−2\}\textbar y\textbar³ dy = O(a⁶); - the
annulus error is O(a⁸). - The metric difference g} − g̃ is a²(ĥ₂ − h₂) =
a²O((1+\textbar y\textbar)\^{}\{−2\}), so g* ∈ 𝒩\_\{Ca\^{}\{2−w\}\}. □

The input from \textbackslash cite\{ozuch2020higher\} is the existence
of ĥ₂ with the stated asymptotics, Lemma 2.3. That paper also records
``Einstein up to an o(t) error iff R⁺ζ = 0'', with t = a². Theorem 61.EH
converts this qualitative statement into the sharp value of the gap
functional.

\hypertarget{consequences}{%
\subsection{3. Consequences}\label{consequences}}

\textbf{3.1 The rate of approach to zero gap is Θ(a⁴), with an explicit
constant.} In Prop. 59.N / 58.O, the resolution of an Einstein orbifold
by EH spaces at ℝ⁴/ℤ₂ points p\_j, with parameters (a\_j, ζ\_j), has P =
π² Σ\_j a\_j⁴ \textbar R⁺\_\{p\_j\}ζ\_j\textbar² (1 + o(1)) for the
optimised metrics. Every metric in the 𝒩\_δ neighbourhoods costs at
least this much, up to the factor (1 − Cδ). - Example: CP²\#CP² ≅
(S⁴/±1) with two EH. The round unit S⁴ has R⁺ = I, so inf over the
neighbourhoods is 2π²a⁴(1 + o(1)). - No Einstein metric lies in these
neighbourhoods. This is the quantitative form of Ozuch's
non-approximability result
\textbackslash cite\{ozuch2021integrability\}. α = 0 is approached only
at the rate a⁴ → 0. - The bound concerns metrics near glued
configurations. It reaches arbitrary sequences on M only through a
59.Z-type decomposition together with Ozuch's GH ⇒ 𝒩\_δ theorem. Ozuch
shows that Einstein metrics GH-close to V lie in such neighbourhoods;
extending this to near-minimisers of P is open.

\textbf{3.2 Kähler--Einstein orbifolds.} For g\_V KE with Λ ≠ 0, R⁺ =
diag(Λ, 0, 0) in the basis (ω̂, ω̂\^{}⊥). - Holomorphic direction ζ = ω̂
(the Kronheimer/crepant gluing used in Theorem Z and 59.N for canonical
models of minimal surfaces with (−2)-curves): \textbar R⁺ζ\textbar{} =
\textbar Λ\textbar, so P ≥ π²Λ²a⁴ per A₁ point on 𝒩\_δ. The Kähler
desingularisation is λ⁴-obstructed. This is consistent with Biquard's
det R⁺ = 0 obstruction being blind to direction, and with the absence of
KE metrics on the minimal resolution when K is not ample. - Directions ζ
∈ ω̂\^{}⊥: R⁺ζ = 0, so the leading obstruction vanishes and inf P =
o(a⁴). These directions are non-holomorphic for the orbifold complex
structure. - Ozuch's second obstruction, Thm 0.2, requires R⁺\_p =
diag(0, 0, Λ) in some basis. A KE point satisfies it with ζ ∈ ω̂\^{}⊥. So
KE orbifolds are exactly the orbifolds unobstructed through second
order, in non-Kähler directions. - Einstein existence on the smooth
minimal surface with non-ample K, near the KE orbifold, is therefore
decided by the third and higher obstructions in ω̂\^{}⊥ directions. These
are unknown. The quantitative rate is inf P = O(a⁶) or smaller in ω̂\^{}⊥
directions; the exact order is to be computed.

\textbf{3.3 Duality between Burns insertion and EH desingularisation.}
On the R⁺ ≤ 0 kernel stratum μ₁ = 0, with μ₂, μ₃ ≤ 0 and μ₂ + μ₃ = Λ: -
First order. Burns pays −8π²εμ₁, which vanishes iff R⁺(ζ,ζ) = 0. EH pays
π²a⁴\textbar R⁺ζ\textbar², which vanishes iff R⁺ζ = 0. On R⁺ ≤ 0 the two
conditions coincide. - Second order. Burns saves −4π²(μ₂ − μ₃)²/2
ε²log(1/ε). On the stratum this saving is \textbf{maximal exactly when
one of μ₂, μ₃ is 0}, i.e.~at Ozuch's locus diag(0, 0, Λ), the KE type.
It \textbf{vanishes exactly at the isotropic point (0, Λ/2, Λ/2)}. - The
isotropic point is the curvature type of the AdS Taub--bolt orbifold, R⁺
= diag(0, −3/2, −3/2) (\textbackslash cite\{ozuch2020higher\}, Example
2.1). That orbifold is EH-obstructed at second order in the compact
setting. - So the undecided boundary case of 59.R is exactly Ozuch's
second-order obstructed type, and the maximal-saving Burns type is
exactly the type with the most unobstructed EH desingularisation. The
two Kähler ALE insertions (Burns, which has mass, and EH, which has
none) see the same tensor R⁺ through the same invariants: the kernel at
first order, and the splitting of the complementary eigenvalues at
second order. - Conjecture 61.D. The second-order EH pairing on the
kernel is a quadratic form in (μ₂, μ₃) vanishing exactly on μ₂μ₃ = 0. If
so, the minimal-P rate in the kernel direction would be c·a⁸(μ₂μ₃)² with
an explicit c, and Burns and EH would be governed by complementary
invariants, (μ₂ − μ₃)² and μ₂μ₃, of the same 2×2 block. This is to be
computed by the flux method applied to ĥ₄ (Ozuch (23), Lemma 2.7--2.9).

\textbf{3.4 Kronheimer A\_n and Wahl fibres.} The same argument applies
verbatim to any Ricci-flat ALE bubble with L² cokernel \{o\}. The
cokernel consists of TT tensors built from L² harmonic ASD forms ⊗
parallel SD forms in the hyperkähler case. The leading pairing is the
flux of (h₂, o) at infinity and is linear in Rm\_p.~For hyperkähler
A\_\{n−1\} it is linear in R⁺\_p, through the period map. - The general
principle: \textbf{a Ricci-flat ALE insertion at an Einstein orbifold
point costs, to leading order and sharply,
‖Π\_\{coker\}(DRic{[}h₂{]})‖²·(scale)⁴.} This is the squared flux of the
background quadratic jet against the bubble's L² infinitesimal Einstein
deformations. - This is deferred until the ℤ₂ case is fully written with
Hölder norms.

\hypertarget{what-is-proved-and-what-is-not}{%
\subsection{4. What is proved and what is
not}\label{what-is-proved-and-what-is-not}}

\begin{itemize}
\tightlist
\item
  \textbf{Proved:} B\_i = π²(R⁺ζ)\_i, an exact flux computation; ‖o‖² =
  π²/a⁴, by exact hand computation, with the script check
  \texttt{step61\_EH\_cokernel\_norm.py} pending at the time of writing
  (see the JSON); and the Cauchy--Schwarz lower bound (a), with
  polynomial error bookkeeping.
\item
  \textbf{Proved modulo the quoted linear theory:} the upper bound (b).
  Its input is Ozuch's Lemma 2.3 (existence and asymptotics of ĥ₂, in
  Bianchi gauge). The Bianchi-gauge identity B(ĥ₂) = 0 is part of
  Ozuch's construction, (22) with P = DRic + δ*B.
\item
  \textbf{Not proved:} that every near-minimising sequence of P on M
  near the orbifold lies in some 𝒩\_δ; the second-order rate in ζ ∈ ker
  R⁺ (Conj. 61.D); and the A\_n and Wahl extensions.
\end{itemize}

\hypertarget{verified-and-corrected-form-ttcs-step61-verifier-binding-over-24}{%
\subsection{5. Verified and corrected form (TTCS step61 verifier;
binding over
§§2--4)}\label{verified-and-corrected-form-ttcs-step61-verifier-binding-over-24}}

Source:
\texttt{trace/ttcs/results/step61\_ttcs01\_verifier1\_integration.md}.

\textbf{5.1 Normalisations confirmed, with two precisions.} - Ω = (e⁰³ −
e¹²)/r⁴ = −¼ d(σ₃/r²) holds exactly for every a. It is ASD, closed and
coclosed. - The potential σ₃/r² is defined only off the bolt. Ω is
harmonic and L² but \textbf{not globally exact}: its bolt period is
∫\emph{\{S²\}Ω = −π/a². - The o\_i satisfy ⟨o\_i, o\_j⟩ = 4δ\_ij/r⁸
pointwise, because J\_iJ\_j = −δ\_ij + ε\_ijk J\_k and (r⁴Ω)² = −I. They
span Ozuch's O(eh). - dV = (r³/8) dr σ₁₂₃, so ∫}\{r≤R\} F dV = π²∫F r³
dr. The induced area of a level set is π²r³√f, not π²r³. - ‖o\_i‖² =
π²/a⁴ exactly. - B = π²R⁺\_pζ is confirmed, including the ℤ₂ factor and
c₀ = 0. An independent round-sphere check gives c₂ = 2π²/3, c₁ = 8π²/3,
B₁ = π². - The normal-coordinate jet is not in flat Bianchi gauge, but
the flux is gauge independent, because DRic\_EH(L\_X EH) = 0. - In the
flux basis, λ\_i = (R⁺ζ)\emph{i. Ozuch's basis has
\textbar ζ}(k)\textbar{} = √2, so his c is compatible with magnitude ½
up to basis signs.

\textbf{5.2 Corrected Theorem 61.EH.} - Centre. Use the jet-extended
superposition g̃\emph{\{a,ζ\} = g\_V + χ}\{t₀\}(EH\_a − δ), with the
background extended over the bolt, so that a⁻²g̃ = EH₁ + a²ĥ₂ + a⁴ĥ₄ + ⋯
in the core. Compare with the \textbf{exact} Einstein background, not a
truncated Taylor expansion. - Lower bound (P). For every bounded δ and w
∈ {[}0, 2), and every g ∈ 𝒩\_δ(a,ζ), P(g) ≥
π²a⁴(\textbar R⁺\_pζ\textbar{} − ε\_w(a,δ))₊², where ε\_w = C{[}a² +
δa²(1 + \textbar log a\textbar) (w = 0) or δa\^{}\{2−w\} (0 \textless{}
w \textless{} 2) + δ²a² (w \textless{} 1), δ²a²\textbar log a\textbar{}
(w = 1), δ²a\^{}\{4−2w\} (1 \textless{} w \textless{} 2){]}. - In
particular ε\_w → 0 for fixed δ. This is sharper than §2's C(δ +
a\^{}\{2−w\}). - Some intermediate integrals grow logarithmically or
polynomially in L; their scaled contributions are still o(a²). - The
test-tensor norm has an O(a²) error. - The lower bound needs no gauge
choice. Any diffeomorphism representative satisfying the estimate may be
used, with the test tensor transported along. - Gauge fact needed for
(b) (proved by the verifier). Applying B\_EH to Ozuch's (22) and using
the linearised Bianchi identity gives ½∇\emph{∇(B\_EH ĥ₂) = 0. With
B\_EH ĥ₂ = O(r⁻³) this forces B\_EH ĥ₂ = 0. Hence DRic\_EH ĥ₂ = Λ·EH +
Σλ\_i o\_i exactly. - Upper bound (P}, Ozuch Lemma 2.3 plus matching).
Let g\_V be \textbf{globally} Einstein; otherwise add P(g\_V). The
adapted metric g*\_a with centre built from ĥ₂ satisfies P(g*\_a) =
π²a⁴\textbar R⁺ζ\textbar² + O(a⁶). If R⁺ζ = 0, it satisfies P(g*\_a) =
O(a⁸ log(1/a)). - g*\emph{a lies in 𝒩}\{C₀\} of a generic centre,
\textbf{not} in 𝒩\_\{Ca\^{}\{2−w\}\}. - Two sharp formulations are
valid: (i) adapted centre, any small δ; (ii) generic centre with δ ≥ C₀.
- The exact a⁶ coefficient is not established. - Infimum, additive form.
With q = \textbar R⁺ζ\textbar, in a neighbourhood containing the adapted
metric: π²a⁴(q − ε)₊² ≤ inf P ≤ π²a⁴q² + O(a⁶). The multiplicative form
holds only when q is bounded below. - Near-extremal structure. If P(g) ≤
π²a⁴q²(1 + η), then ‖(1 − Π\_g)Ric°(g)‖² ≤ π²a⁴{[}ηq² + Cqε{]}, where
Π\_g is the L²(g) projection onto the cut-off obstruction span.

\textbf{5.3 Corrected consequences.} - S⁴/±1 means the suspension
quotient (x₁,\ldots,x₄,x₅) ↦ (−x₁,\ldots,−x₄,x₅), with two A₁ points; it
is not RP⁴. With R⁺ = I on the unit sphere, two adapted insertions have
energy 2π²a⁴(1 + o(1)). This is a local sharp asymptotic, not a
universal rate. - KE orbifolds. R⁺ = diag(Λ,0,0), and the holomorphic
direction costs π²Λ²a⁴. For ζ ⊥ ω̂ the first obstruction vanishes. -
These directions are \textbf{not} ``non-Kähler'': EH is hyperkähler, and
the bolt is Lagrangian for the ω̂-structure, which corresponds to the
local smoothing interpretation. - Ozuch's Thm 0.2 is a necessary
condition under compactness, integrability and stability/rigidity
hypotheses. The §3.2 phrases ``exactly the orbifolds unobstructed
through second order'' and ``decided by third and higher obstructions''
are withdrawn. - \textbf{Conjecture 61.D is partly known.} Ozuch §2.1.4
already gives the second obstruction μ₁ = C·det(R⁺\textbar\_\{ζ\^{}⊥\}),
C \textgreater{} 0, when the first vanishes. What remains open is the
normalised coefficient, the other second-order components, and an a⁸
minimised-defect theorem. The latter must control orbifold Einstein
deformations and scale/direction changes (Ozuch Cor. 2.12). -
\textbf{Burns/EH unification (exact).} On the kernel stratum with
complementary eigenvalues u, v and u + v = Λ, we have (u − v)² = Λ² −
4uv. Hence **Φ\_opt(kernel) = −(u − v)²/2 = −Λ²/2 + 2
det(R⁺\textbar\_\{ζ\^{}⊥\}).\textbf{ - The second-order Burns saving and
Ozuch's second EH obstruction are therefore the same invariant,
det(R⁺\textbar{}\emph{\{ζ\^{}⊥\}), at fixed Λ. - Maximal Burns saving,
−Λ²/2, occurs exactly where the second EH obstruction vanishes (KE
type). - Zero Burns saving occurs at the isotropic point (0, Λ/2, Λ/2),
where the determinant is maximal, Λ²/4 (AdS Taub--bolt type, det = 9/4).
- The two Kähler-ALE insertions thus read the same pair of invariants of
R⁺: the kernel at first order, and det(R⁺\textbar{}}\{ζ\^{}⊥\}) at
second order. - Other ALE bubbles. The leading energy is B\^{}T G⁻¹ B
a⁴, with G the Gram matrix of the L² cokernel. The extension to A\_n and
Wahl bubbles is }not** verbatim: it needs decay, period normalisation,
all cokernel modes, and the curvature-derivative obstructions of A\_k
(Ozuch Rem. 2.21). - Theorem 61.EH does not remove (N0), and it does not
produce a positive scale-independent gap. Positivity of an obstruction
changes only the rate.

\hypertarget{s1003a02-step-62-the-weyl-channel-reduction-of-the-open-range-lower-bound}{%
\chapter{S1003A02: Step 62 --- The Weyl-channel reduction of the
open-range lower
bound}\label{s1003a02-step-62-the-weyl-channel-reduction-of-the-open-range-lower-bound}}

Source: \texttt{step62\_weyl\_channel\_reduction.md}.

Agent2, Execute. The shell was unavailable for this step, so everything
below is derivation. The numerics quoted are taken from existing data
files: the certified cell
\texttt{step49\_cell\_theoremY\_profile\_x003\_X128\_N1600.json} and the
Step-48/49 ledger.

Setting. - X is minimal of general type with A = c₁²(X), b⁺(X) ≥ 2 and
SW(X) ≠ 0 (the sign-family hypotheses of Theorems N/W/K/Y). - M =
X\#kCP̄². - p = P/8π² with P = ∫\textbar Ric°\textbar². - For every ε ∈
\{±1\}\^{}k, the class a\_ε = K\_X + Σε\_iE\_i is a monopole class
(blow-up formula).

Notation. - g̃ denotes the Yamabe representative of {[}g{]}, normalised
to s ≡ −4. - W₂(g) = ∫\textbar W⁺\textbar² dμ. This is a conformal
invariant. - Y{[}g{]} is the Yamabe constant. Note that ∫s²dμ\_\{g̃\} =
Y², and ∫s²dμ\_g ≥ Y² for every g in the class (negative Yamabe class).

\hypertarget{two-conformal-invariants-and-an-exact-identity}{%
\subsection{1. Two conformal invariants and an exact
identity}\label{two-conformal-invariants-and-an-exact-identity}}

Define the \textbf{Weyl deficit} and the \textbf{Yamabe excess} of a
conformal class: \[
\kappa[g]:=\frac A2-\frac{3}{8\pi^2}\int_M|W^+_g|^2d\mu_g ,\qquad
\upsilon[g]:=\frac{Y[g]^2}{32\pi^2}-A=\frac{\operatorname{Vol}(\tilde g)-2\pi^2A}{2\pi^2}.
\]

Calibration. - κ = 0 exactly when ∫\textbar W⁺\textbar² equals the
Kähler--Einstein value (4π²/3)A. For a Kähler metric
\textbar W⁺\textbar² = s²/24 pointwise. - υ ≥ 0 is LeBrun's scalar
estimate. - In Step-48 variables: ω := δW/2π² = 4κ/3, where δW =
(4π²/3)A − ∫\textbar W⁺\textbar²; and υ = ΔV/2π².

\textbf{Identity 62.1.} For every metric g on M, \[
\boxed{\;p(\tilde g)=(k-\kappa)+\tfrac13(\upsilon-\kappa),\qquad p(g)\ge p(\tilde g).\;}
\]

\emph{Proof.} - Chern--Weil: P = 4∫\textbar W⁺\textbar² + (1/12)∫s² −
8π²(2χ+3τ)(M), and (2χ+3τ)(M) = A − k. - At g̃: ∫s² = Y² = 32π²(A+υ) and
∫\textbar W⁺\textbar² = (4π²/3)(A − 2κ). - Substituting: p = (1/6)(A−2κ)
+ (A+υ)/3 − A + k = k − (4/3)κ + υ/3. - For a general g, P(g) − P(g̃) =
(∫s²dμ\_g − Y²)/12 ≥ 0. ∎

Check on the certified perpendicular cell (x₀ = .03, X = 128; Step-49
data, per class): ω = 4·0.018784 = 0.075136 and υ = 4·0.053096 =
0.212384. The identity gives p = 1 − 0.075136 + 0.070795 = 0.995659,
against the interval value 0.9956627. The difference is the rounding of
the Step-49 floats. So κ\_cell = 0.05635 per class.

\textbf{Reading of 62.1.} The Einstein defect is a sum of two terms: -
the \emph{Weyl shortfall} k − κ, i.e.~how much less than one unit of κ
per exceptional class the conformal class has achieved; - one third of
the \emph{Yamabe overpayment} υ − κ, i.e.~how much more Yamabe excess
was spent than the Weyl deficit gained.

A family with p/k → 0 must therefore make both terms small.

\hypertarget{lebruns-two-estimates-in-these-coordinates}{%
\subsection{2. LeBrun's two estimates in these
coordinates}\label{lebruns-two-estimates-in-these-coordinates}}

LeBrun's Prop. 2 and Prop. 4 (Knowledge \texttt{math-0111228\ldots{}}
lines 127--149; both hold for every metric) are applied to the sign
class a\_ε maximising (a\_ε⁺)². We have (a\_ε⁺)² = A − k +
\textbar a\_ε⁻\textbar², and E\_ε\textbar a\_ε⁻\textbar² =
\textbar K⁻\textbar² + Σ\textbar E\_i⁻\textbar² ≥ k. Write Π :=
max\_ε(a\_ε⁺)² − A ≥ 0; then Π ≥ \textbar K\_X⁻\textbar² +
Σ\textbar E\_i⁺\textbar², because \textbar E⁻\textbar² = 1 +
\textbar E⁺\textbar².

Prop. 2 gives υ ≥ Π ≥ 0.

Prop. 4 gives ∫(s − √6\textbar W⁺\textbar)² ≥ 72π²(A+Π). Since s
\textless{} 0 somewhere is all that is used, the cross term satisfies
∫(s−√6\textbar W⁺\textbar)² ≤ (‖s‖₂ + √6‖W⁺‖₂)², by Cauchy--Schwarz on
∫\textbar s\textbar\textbar W⁺\textbar. At g̃ this becomes \[
\boxed{\;2\sqrt{A+\upsilon}+\sqrt{A-2\kappa}\ \ge\ 3\sqrt{A+\Pi}\;}\tag{62.2}
\] This is the ``envelope 2x + y ≥ 3'' of the Step-61 plan, with x² = 1
+ υ/A and y² = 1 − 2κ/A. \textbf{So Target 2 of Route 2 is not a new
inequality: it is LeBrun's Prop. 4 plus Cauchy--Schwarz.} For Π = 0 it
is sharp exactly when \textbar W⁺\textbar{} = λ\textbar s\textbar{}
pointwise with s constant and LeBrun's equality structure holds (a
special almost-Kähler metric).

Solving (62.2) for υ at fixed κ, with Π = 0 (Π ≥ 0 only helps), and
writing y = √(1 − 2κ/A): \[
\upsilon\ \ge\ \upsilon_{\rm L}(\kappa):=A\Bigl[\bigl(\tfrac{3-y}{2}\bigr)^2-1\Bigr]=\kappa+\frac{3\kappa^2}{4A}+O(\kappa^3/A^2).
\] LeBrun therefore says exactly this: \textbf{one unit of Weyl deficit
costs at least one unit of Yamabe excess}. The Yamabe overpayment υ − κ
≥ 3κ²/4A is a pure second-order effect.

\hypertarget{theorem-62.a-the-weyl-channel-bound}{%
\subsection{3. Theorem 62.A (the Weyl-channel
bound)}\label{theorem-62.a-the-weyl-channel-bound}}

\textbf{Theorem 62.A.} Under the sign-family hypotheses, for every
metric g on M = X\#kCP̄² whose Weyl deficit satisfies κ = κ{[}g{]} ≤
4A/9, \[
\boxed{\;\frac{P(g)}{8\pi^2}\ \ge\ L_A(k,\kappa):=(k-\kappa)+\tfrac12\Bigl(A-\kappa-\sqrt{A^2-2A\kappa}\Bigr)
=(k-\kappa)+\frac{\kappa^2}{4A}+\frac{\kappa^3}{4A^2}+\cdots\;}
\] For κ ≥ 4A/9, LeBrun's unconditional bound p ≥ k − A/3 is the
minimum.

\emph{Proof.} Use 62.1 together with υ ≥ υ\_L(κ). Then p ≥ k − (4/3)κ +
(A/3){[}((3−y)/2)² − 1{]} = k − κ + ½(A − κ − Ay), since (A/3)·{[}(10 −
2κ/A − 6y)/4 − 1{]} = (A/12)(6 − 6y) − κ/6. The constraint κ ≤ 4A/9 is
the condition y ≥ 1/3. Below that, the unconstrained LeBrun optimum t =
c/9 of Step 36 is interior and gives k − A/3. ∎

Three consistency checks. - (a) \textbf{Kähler metrics have κ ≤ 0.} For
a Kähler ω on any complex structure of M (all are blow-ups of the
minimal model by SW), ∫\textbar W⁺\textbar² = ∫s²/24 ≥
(4π²/3)(c₁·ω)²/ω². On the Kähler cone ω = π\emph{ω\_X − Σe\_iE\_i with
e\_i \textgreater{} 0, (c₁·ω)²/ω² = (K·ω\_X + Σe\_i)²/(ω\_X² − Σe\_i²) ≥
(K·ω\_X)²/ω\_X² ≥ A, the last step by Hodge index. So L\_A ≥ k:
\textbf{Theorem K∞'s Kähler tariff P ≥ 8π²k is the κ ≤ 0 case of 62.A.}
Equality in the limit corresponds to κ → 0⁻ and υ → 0, which is the
Arezzo--Pacard recovery. - (b) \textbf{Certified cell:} κ = 0.0564 per
class, so the bound is p ≥ 0.944 per class; the actual value is 0.9957.
The slack (υ − υ\_L)/3 ≈ (0.2124 − 0.0564)/3 = 0.052 is the Yamabe
overpayment. The cell spends 3.77 times the Yamabe excess that LeBrun
requires for its Weyl deficit. Its }LeBrun efficiency* is κ/υ = 0.265. -
(c) \textbf{Formal common mode} (h = k/A; Step 48 §2: axial W⁺ of
magnitude x = 2/3 − h at s = −4, V̂ = 8π²A/(2−h)²): -
∫\textbar W⁺\textbar² = (4π²A/3)(2−3h)²/(2−h)² = (4π²A/3)(1 − 2h + h³/2
+ \ldots), so κ = k − k³/4A² + \ldots. - υ = A{[}4/(2−h)² − 1{]} = k +
3k²/4A + k³/2A² + \ldots, which equals υ\_L(κ) through order k³/A². The
formal mode is LeBrun-extremal with Π = 0. - p = k²/4A + k³/2A² +
\ldots, which equals L\_A(k,κ) through order k³/A². - \textbf{The formal
common mode lies on the equality locus of Theorem 62.A}, to the order
computed. Its κ is k up to O(k³/A²). - Comparison: L\_A(k,k) = ½(A − k −
√(A²−2Ak)) = k²/4A + k³/4A² + \ldots. This is the value of the bound at
κ exactly equal to k, and it lies below the formal mode, as it must.

\hypertarget{corollaries-one-conformal-invariant-governs-the-lower-bound-in-the-open-range}{%
\subsection{4. Corollaries: one conformal invariant governs the lower
bound in the open
range}\label{corollaries-one-conformal-invariant-governs-the-lower-bound-in-the-open-range}}

Let κ\_*(M) := sup\_g κ{[}g{]} = A/2 − (3/8π²)·inf\_g
∫\textbar W⁺\textbar². This is the conformal Weyl infimum, in normalised
form.

\textbf{Corollary 62.B (sharp conditional k²/A law).} If inf\_g
∫\emph{M\textbar W⁺\textbar² ≥ (4π²/3)(A − 2k), i.e.~κ}* ≤ k, then \[
\alpha(M)\ \ge\ 4\pi^2\bigl(A-k-\sqrt{A^2-2Ak}\bigr)=\frac{2\pi^2k^2}{A}\Bigl(1+\frac kA+\cdots\Bigr).
\] The leading constant 2π²/A (that is, p ≥ k²/4A) is the
formal-common-mode constant, so it would be optimal among laws of this
shape if the common mode could be realised. More generally, α(M)/8π² ≥
L\_A(k, κ\_\emph{) whenever κ\_} ≤ 4A/9. This gives a positive per-class
tariff p ≥ ck as soon as κ\_* ≤ (1−c)k.

\textbf{Corollary 62.C (coupled form; what a proof must establish).}
Pointwise in g, Identity 62.1 makes the following three statements
equivalent: - (i) κ{[}g{]} ≤ k + βp(g̃); - (ii) υ{[}g{]} ≤ k + (3+4β)p(g̃)
(affine Yamabe susceptibility); - (iii) ∫\textbar W⁺\textbar² +
(β/3)P(g̃) ≥ (4π²/3)(A − 2k).

If any of them holds for all g with a uniform β ≥ 0, then α/8π² ≥
k²/((4+4β)A)(1 + O(k/A)).

\emph{Proof.} - κ = (3k + υ − 3p)/4 by 62.1. - Let κ\_max = k + βp.
Since L\_A is decreasing in κ for κ \textless{} 2A, we get p ≥ L\_A(k,
k+βp) ≥ −βp + (k+βp)²/4A. - The strategist's linear version υ ≤ γp (no k
term) gives instead κ ≤ (3k + (γ−3)p)/4, and hence the \emph{tariff} p ≥
k/(γ+1). ∎

In the coupled form, with γ = 3 + 4β: - The formal mode saturates (ii)
with γ = 3: υ = k + 3k²/4A and 3p = 3k²/4A. - The certified cell
satisfies (ii) for every γ ≥ 0, since its υ = 0.212 \textless{} k. -
Kähler recovery has υ → 0.

So \textbf{every known object satisfies the affine susceptibility υ ≤ k
+ 3p, and the only object saturating it is the formal common mode.} The
strategist's sup ΔV/P is unbounded along the formal mode (≈ A/k). The
scale-free number that separates ``positive tariff'' from ``k²/A law''
is therefore \textbf{the Weyl deficit per class, κ/k}: - κ/k \textless{}
1 − c for a tariff c; - κ/k ≤ 1 + O(p/k) for the k²/A law.

The Yamabe excess is slaved to κ by LeBrun (υ ≥ κ + 3κ²/4A), so it is
not an independent channel.

\textbf{Corollary 62.D (necessary structure of any family with p/k →
0).} Along such a family, 62.1 and 62.2 force: - κ/k → 1,
i.e.~∫\textbar W⁺\textbar² = (4π²/3)(A − 2k) + o(k): every exceptional
class must remove (8π²/3) of self-dual Weyl energy; - υ − κ = o(k):
asymptotic equality in LeBrun's Prop. 4 with Cauchy--Schwarz,
i.e.~\textbar W⁺\textbar{} ≈ \textbar s\textbar/√6·const with s ≈ const,
which is LeBrun's special almost-Kähler equality structure; - Π = o(k):
almost all exceptional periods \textbar E\_i⁺\textbar² and
\textbar K⁻\textbar² vanish.

The last condition is exactly what Taubes' period constraint (Lemma T,
Step 47) prevents for an exact closed common mode. \textbf{The
obstruction to P/k → 0 therefore lies in the joint smallness of
κ-shortfall, LeBrun slack and periods, not in any single one.} This
sharpens Step-48 Cor. Y1 (Δ V ≥ 2π²(1−c)k) to the statement that the
Yamabe excess must be matched one-for-one by Weyl deficit.

\hypertarget{the-asd-side-reformulation-and-what-is-known}{%
\subsection{5. The ASD-side reformulation and what is
known}\label{the-asd-side-reformulation-and-what-is-known}}

Since ∫\textbar W⁺\textbar² − ∫\textbar W⁻\textbar² = 12π²τ(M) =
12π²(τ\_X − k), Cor. 62.B's hypothesis reads \[
\inf_g\int_M|W^-|^2\ \ge\ \tfrac{4\pi^2}{3}(A-9\tau_X)+\tfrac{28\pi^2}{3}k .
\] For a ball quotient (A = 9τ\_X) this is \textbf{inf
∫\textbar W⁻\textbar² ≥ (28π²/3)k = (7/9)·12π²·k}: every exceptional
class must carry at least 7/9 of the anti-self-dual Weyl energy of a
Fubini--Study CP̄² summand.

Values of the per-class anti-self-dual Weyl energy: - conformal CP̄²
connected sums and Burns cells (Burns is conformal to CP̄² minus a
point): ≈ 12π²; - certified cell: 12π² − δW = (12 − 0.150)π² per class;
- formal mode: (28/3)π², i.e.~the bound is saturated.

For a positive per-class tariff it suffices that inf
∫\textbar W⁻\textbar² ≥ λ·12π²k with λ \textgreater{} 7/9 + (something);
precisely, the tariff is c = 1 − (9/2)(1−λ). For ball quotients, the
trivial signature bound ∫\textbar W⁺\textbar² ≥ 12π²τ\_M gives only κ ≤
9k/2, which is useless.

\textbf{Gursky analogy (positive Yamabe).} - Gursky: for b⁺
\textgreater{} 0 and Y \textgreater{} 0, ∫\textbar W⁺\textbar² ≥
(4π²/3)(2χ+3τ). On a blow-up this permits a reduction of at most (4π²/3)
per CP̄² summand, which is \emph{half} the (8π²/3) that Cor. 62.B needs
to exclude. - On CP²\#kCP̄² this per-point bound is nearly sharp.
12π²/(4π²/3) = 9, and scalar-flat Kähler (hence ASD) metrics exist on
CP²\#kCP̄² for k ≥ 10 (Kim--LeBrun--Pontecorvo, Rollin--Singer), giving
an average of 1.2π² per point. - \textbf{Conjecture 62.G} (the
negative-Yamabe, SW analogue of Gursky's per-point law): on X\#kCP̄², inf
∫\textbar W⁺\textbar² ≥ (4π²/3)(A − k) = (4π²/3)(2χ+3τ)(M). This would
give κ\_* ≤ k/2, hence the per-class tariff α ≥ 4π²k(1 + O(k/A)), i.e.~c
= ½. - The weaker Conjecture 62.W (κ\_* ≤ k) gives the sharp k²/A law. -
Both conjectures predict that the Taubes threshold for ASD metrics on
X\#kCP̄² is at least A (resp. A/2).

Every Seiberg--Witten consequence used so far points the same way: each
is a lower bound on the Yamabe excess in terms of the Weyl deficit. -
Theorem Y, rewritten with 62.1, reads κ ≤ υ − (3/16π²)∫Ψ. - LeBrun reads
υ ≥ υ\_L(κ). - Theorems N/W/K are saving upper bounds, which is the same
direction.

None of them bounds κ itself. Neither conjecture can come from SW and
the Weitzenböck formulas alone, because Taubes' ASD metrics exist on
X\#kCP̄² for k ≫ 1 while SW ≠ 0. \textbf{Any proof must see the
exceptional classes, i.e.~k, in the Weyl channel itself}, not only
through the sign family's (a⁺)² ≥ A. This explains structurally why
every field-theoretic attempt (Steps 36--48) stopped at LeBrun's 2/3:
those identities constrain only (υ, κ) through (62.2), never κ alone.

Sources to check (knowledge request s62): Gursky's inequality and its
negative-Yamabe status; whether KE minimises ∫\textbar W\textbar² on
general-type surfaces; explicit Taubes thresholds; ASD metrics on
blow-ups of general-type surfaces; LeBrun's equality analysis in Prop. 4
{[}27{]}.

\hypertarget{what-this-changes}{%
\subsection{6. What this changes}\label{what-this-changes}}

\begin{enumerate}
\def\labelenumi{\arabic{enumi}.}
\tightlist
\item
  The open-range problem C₊ ≥ A − (1−c)k is equivalent to a bound on κ
  along near-minimisers, and the coupled form 62.C is exactly what is
  needed. The Yamabe-susceptibility route (Route 2) and the strategist's
  ΔV/P are the same statement as 62.C(ii). Its natural sharp constant γ
  = 3 is the formal mode; its natural tariff version is a bound on κ/k.
\item
  The K+Y+Cauchy--Schwarz combination proposed by the strategist gives δ
  − d ≥ (5/6)d²/(A+δ) (in δ = υ, d = k − p). Together with 62.C(ii) at γ
  = 3 this yields only p ≥ (5/24)d²/(A+υ). \textbf{The sharp 1/(4A)
  comes from LeBrun's Prop. 4, not from Theorem K/Y}: Theorem K's
  Cauchy--Schwarz loses the factor 5/6.
\item
  Decisive computations, chosen to change judgement:

  \begin{itemize}
  \item
    \begin{enumerate}
    \def\labelenumii{(\alph{enumii})}
    \tightlist
    \item
      \textbf{κ-capacity of cells.} For the CH-collared
      Burns/perpendicular family (A01 data) and the rotated cells,
      compute (κ, υ, p) per class and the two efficiencies κ/k and κ/υ.
      Ask whether any local cell has κ ≥ 1 per class at bounded p.~If
      not, find what bounds κ locally; a local Weyl-capacity bound would
      be the first piece of a proof of 62.W.
    \end{enumerate}
  \item
    \begin{enumerate}
    \def\labelenumii{(\alph{enumii})}
    \setcounter{enumii}{1}
    \tightlist
    \item
      \textbf{Local susceptibility near KE.} Near the KE host, κ ≈
      Q\_W(h) and p ≈ ‖Lh‖² are both quadratic in the TT/conformal
      deformation (KE is Bach-flat, hence critical for
      ∫\textbar W⁺\textbar²). Is Q\_W ≤ βQ\_P with β explicit (Koiso gap
      / A\_C)? That would give 62.C(i) in the A\_C neighbourhood with
      explicit β.
    \end{enumerate}
  \item
    \begin{enumerate}
    \def\labelenumii{(\alph{enumii})}
    \setcounter{enumii}{2}
    \tightlist
    \item
      \textbf{Construction side.} Configurations violating 62.W need κ
      \textgreater{} k: conformal structures that remove more than
      (8π²/3) of W⁺ energy per CP̄² summand. ASD-gluing (Taubes,
      Donaldson--Friedman) is the natural mechanism. Counting how many
      CP̄² summands an unobstructed ASD gluing on X needs is the decisive
      question.
    \end{enumerate}
  \end{itemize}
\end{enumerate}

\hypertarget{two-efficiencies-the-burns-breakeven-14-and-the-equality-locus}{%
\subsection{7. Two efficiencies, the Burns breakeven 1/4, and the
equality
locus}\label{two-efficiencies-the-burns-breakeven-14-and-the-equality-locus}}

Define per exceptional class two quantities: - the Weyl capacity κ̂ =
κ/k; - the LeBrun efficiency η = κ/υ. LeBrun gives η ≤ 1 + O(κ/A).

Identity 62.1 then reads \[
p/k = 1-\hat\kappa\,\frac{4-1/\eta}{3}.
\] So a configuration saves against the unit tariff \textbf{iff η
\textgreater{} 1/4}, and the saving is linear in κ̂.

\textbf{Cell family data.} The values below are read from the existing
ledgers \texttt{workspace/data/cell\_yamabe\_rv\_ch\_x*.json}, using the
larger-N rows. In each, ω = −fourWplus\_rel, υ = 3·scalar\_channel and κ
= 3ω/4. All are floats, not interval values.

\begin{longtable}[]{@{}llllll@{}}
\toprule
x₀ (X) & κ & υ & η = κ/υ & p from 62.1 & p ledger\tabularnewline
\midrule
\endhead
0.003 (64) & 0.0047945 & 0.018563 & 0.2583 & 0.9997952 &
0.9997952\tabularnewline
0.01 (64) & 0.017247 & 0.065087 & 0.2650 & 0.9987004 &
0.9987004\tabularnewline
0.03 (128) & 0.056351 & 0.212390 & 0.2653 & 0.9956627 &
0.9956627\tabularnewline
0.05 (128) & 0.093971 & 0.364746 & 0.2576 & 0.9962875 &
0.9962875\tabularnewline
\bottomrule
\end{longtable}

In the same table, the Kähler blow-up family has κ = −Π/2 ≤ 0 and υ = Π,
and the formal common mode has κ̂ ≈ 1, η ≈ 1.

The table gives an exact reading of the CH-cell family: \textbf{the
cells have LeBrun efficiency barely above the breakeven 1/4.} As x₀
grows, η falls back toward 1/4. That is why the saving changes sign near
x₀ ≈ 0.065.

\textbf{Theorem F reinterpreted.} From the first-order chiral split c\_A
= −4π²(M\_ζ/6, 1/√12) (59.R), a Burns insertion in direction ζ changes
the two channels as follows: - ∫4\textbar W⁺\textbar² by −8π²εW⁺(ζ,ζ); -
∫s²/12 by +8π²ε\textbar s\textbar/12.

Hence its first-order efficiency is η\_B(ζ) =
3W⁺(ζ,ζ)/\textbar s\textbar, and R⁺(ζ,ζ) = (\textbar s\textbar/3)(η\_B −
1/4). So ``Burns saves at first order'' is the same statement as ``its
efficiency exceeds 1/4'', and the minimizer rigidity R⁺ ≤ 0 says that no
Burns direction exceeds efficiency 1/4. On CH the perpendicular
direction has η\_B = 1/4 exactly (the kernel). The second-order log term
raises η slightly above 1/4, which is the small saving in the data. To
go from efficiency 1/4 (local Burns) to the efficiency 1 permitted by
LeBrun, one needs diffuse reduction of the \textbar W⁺\textbar{}
magnitude, which is the formal mode.

\textbf{Equality locus of 62.A.} Equality requires two things: -
equality in LeBrun's Prop. 4. By LeBrun's Prop. 2.4 (Knowledge
\texttt{Ricci-curvature-minimal-volumes-and-Seiberg-Witten-theory\ldots{}},
lines 243--257), g is then almost-Kähler, ω lies in the lowest
eigenspace of W⁺, s + s* is constant, c₁ is of type (1,1), and W⁺ is
J-invariant; - the Cauchy--Schwarz step, which forces s and
\textbar W⁺\textbar{} to be constant.

At s = −4 this gives W⁺ = (−2a, a, a) constant with ω on the eigenvalue
−2a, s* = s/3 − 8a, and s* − s = 8/3 − 8a ∝ \textbar∇ω\textbar² ≥ 0.
Hence a ≤ 1/3, with κ \textgreater{} 0 exactly when J is strictly
almost-Kähler: κ = A/2 − (3/8π²)·6a²·Vol.

\textbf{So a family with p/k → 0 must approach compact strictly
almost-Kähler metrics with constant s, constant J-invariant degenerate
W⁺ and harmonic J (Ω an eigenform of W).} LeBrun (2001) left open
whether any compact strictly almost-Kähler metric saturates the
estimate. Apostolov--Armstrong--Drăghici (Knowledge
\texttt{Local-models-and-integrability\ldots{}}, Thm 2) prove that a
compact almost-Kähler metric satisfying the third Gray condition
(harmonic J together with J-invariant Ricci) is Kähler. Our locus has
harmonic J and J-invariant W⁺, but J-invariance of Ricci has not been
checked.

\textbf{Proof route opened.} Proceed in three steps: 1. Determine
whether the equality conditions together with constant s and constant
\textbar∇J\textbar{} force J-invariant Ricci, or integrability directly.
The relevant tools are the Apostolov--Armstrong ``Hermitian Weyl
tensor'' results and Blair's formula. 2. If the equality locus is Kähler
only, it has κ ≤ 0. 3. A quantitative stability of that rigidity would
then bound κ by the LeBrun slack, which turns 62.A into a genuine lower
bound.

The formal common mode is a pointwise algebraic solution of these
equality conditions that no compact metric realises. Taubes' period
constraint and AAD-type integrability are the two candidate
obstructions, and comparing them is decisive.

\hypertarget{s1003a02-step-63-the-exceptional-flux-in-the-frame-a-hodgegauss-bound-the-exact-almost-kuxe4hler-identity-and-the-empty-equality-locus}{%
\chapter{S1003A02: Step 63 --- The exceptional flux in the frame: a
Hodge--Gauss bound, the exact almost-Kähler identity, and the empty
equality
locus}\label{s1003a02-step-63-the-exceptional-flux-in-the-frame-a-hodgegauss-bound-the-exact-almost-kuxe4hler-identity-and-the-empty-equality-locus}}

Source: \texttt{step63\_frame\_twist\_flux\_and\_AK\_identity.md}.

Agent2, Execute. The shell was blocked again for this step, so
everything here is derivation. Sources read this step: -
Apostolov--Armstrong--Drăghici, \emph{Local models and integrability of
certain almost Kähler 4-manifolds} (Knowledge
\texttt{Local-models-and-integrability\ldots{}}): lines 30--64 and
112--236, giving conventions (6)--(25) and the canonical Chern form γ\_J
= R(Ω) − a∧Ja. - The Step-62 note
\texttt{step62\_weyl\_channel\_reduction.md}.

Setting and notation. - M = X\#kCP̄², X minimal of general type, A =
c₁²(X), P = ∫\textbar Ric°\textbar², p = P/8π². - For a class a ∈ H²(M)
and a metric g, a\^{}± are the g-harmonic (anti-)self-dual parts and
\textbar a\^{}-\textbar² := −∫ a\textsuperscript{-∧a}- ≥ 0. - 2-form
norms have \textbar e¹∧e²\textbar{} = 1, and the Kähler-type form has
\textbar Ω\textbar² = 2.

\hypertarget{theorem-63.t-the-frame-twist-bound-any-metric-no-seibergwitten-input}{%
\subsection{1. Theorem 63.T: the frame-twist bound (any metric, no
Seiberg--Witten
input)}\label{theorem-63.t-the-frame-twist-bound-any-metric-no-seibergwitten-input}}

Let g be any metric and J any g-orthogonal almost complex structure
inducing the orientation. Write Ω = g(J·,·), so Ω ∈ Λ⁺ with
\textbar Ω\textbar² = 2. Choose a local gauge φ ∈ Ω\^{}⊥ ⊂ Λ⁺ with
\textbar φ\textbar² = 2, and write \[
\nabla\Omega=\alpha\otimes\phi+\beta\otimes J\phi,\qquad
\nabla\phi=-\alpha\otimes\Omega+b\otimes J\phi .
\] Then b is a connection 1-form on Ω\^{}⊥ ≅ K\_J\^{}\{-1\}. Applying
d\^{}∇ to ∇φ and taking the Jφ-component of the Ricci identity gives,
gauge-independently, \[
\gamma_J:=-db=R(\Omega)+\alpha\wedge\beta ,\qquad [\gamma_J]=2\pi c_1(J).
\tag{63.1}
\] Here R is the curvature operator on Λ². For an almost-Kähler J one
has α = a and β = −Ja, and (63.1) becomes AAD (25), γ\_J = R(Ω) − a∧Ja.
That is the consistency check of the sign.

The Singer--Thorpe decomposition gives: - R(Ω)\^{}- = ½Ric̃₀(Ω) =
ρ₀\^{}J, the J-invariant primitive Ricci form, with
\textbar ρ₀\^{}J\textbar² = ½\textbar Ric₀'\textbar² ≤
½\textbar Ric°\textbar²; - α∧β is decomposable, so
\textbar(α∧β)\^{}±\textbar{} = \textbar α∧β\textbar/√2 ≤
\textbar∇Ω\textbar²/(4√2).

Hodge theory finishes the argument. Write γ\_J = H + dθ with H harmonic.
Then ‖γ\_J\^{}-‖² = ‖H\^{}-‖² + ‖(dθ)\^{}-‖², because
⟨H\textsuperscript{-,(dθ)}-⟩ = ½(⟨H,dθ⟩ − ∫H∧dθ) = 0. Hence ‖γ\_J\^{}-‖
≥ 2π\textbar c₁(J)\^{}-\textbar. This gives:

\textbf{Theorem 63.T.} For every metric g and every orthogonal almost
complex structure J on a closed oriented 4-manifold, \[
\Bigl(\tfrac12\int|\operatorname{Ric}_0'|^2\Bigr)^{1/2}\ \ge\ 2\pi|c_1(J)^-|_g-\|(\alpha\wedge\beta)^-\|_{L^2},
\qquad\text{so}\qquad
P\ \ge\ 2\bigl(2\pi|c_1(J)^-|-\mathcal T_J\bigr)_+^2 ,
\] where 𝒯\_J := ‖(α∧β)\^{}-‖\emph{\{L²\} ≤ (1/4√2)‖∇Ω‖²}\{L⁴\} is the
\textbf{ASD frame twist} of J.

\textbf{Choice of class on M.} - For ε ∈ \{±1\}\^{}k put a\_ε = c₁(X) +
Σε\_iE\_i. Each a\_ε is characteristic with a\_ε² = c₁²(M), so (Wu) a\_ε
= c₁(J) for some orthogonal J. - Averaging over ε gives
max\_ε(a\_ε\^{}+)² ≥ (c₁(X)\^{}+)² + Σ(E\_i\^{}+)² ≥ A +
\textbar c₁(X)\^{}-\textbar² + Σ(E\_i\^{}+)². - So the best sign has
\textbar a\_ε\^{}-\textbar² = (a\_ε\^{}+)² − (A−k) ≥ k + Π.

Hence, with no Seiberg--Witten hypothesis, \[
\boxed{\ \sqrt{P/2}\ \ge\ 2\pi\sqrt{k+\Pi}-\inf_{J:\,c_1(J)=a_\varepsilon}\mathcal T_J\ }
\tag{63.2}
\]

Readings. - (i) \textbf{Kähler tariff.} For a Kähler metric, J itself
has 𝒯\_J = 0 and Ric₀'\,' = 0. With (a\^{}+)² ≥ A from Hodge index on
the blow-up Kähler cone, (63.2) is P ≥ 8π²k. This is Theorem K∞ without
Seiberg--Witten. - (ii) \textbf{Where the exceptional flux goes.} For
every J, the ASD part of the Chern form must carry flux
2π\textbar c₁\^{}-\textbar{} ≥ 2π√k. Only two carriers exist: the
J-invariant tracefree Ricci tensor, which costs P, and the ASD frame
twist 𝒯\_J. This is the exact form of the Step-40s observation that
``charge can reside in anti-self-dual relative curvature''. - (iii)
\textbf{Necessary condition for every sub-tariff family.} If p ≤ c·k,
then \emph{every} orthogonal J in the best class has \[
  \mathcal T_J\ge 2\pi\sqrt k\,(1-\sqrt c).
  \] - Certified CH cell (c = 0.995663): 𝒯\_J ≥ 0.01365 for every J. - A
family with p/k → 0: 𝒯\_J ≥ 2π√k(1 − o(1)) for every J. The twist of
every almost complex structure must carry the whole exceptional flux. -
(iv) \textbf{Local consistency with Theorem F and Theorem 59.R.} - In a
Burns insertion of scale ε in a kernel direction ζ ⊥ ω, the natural J
equals J\_ζ in the core and J\_ω outside. The rotation by π/2 across the
logarithmic collar √ε ≲ r ≲ ρ has \textbar∇Ω\textbar{} \textasciitilde{}
(r log(ρ/√ε))\^{}\{-1\}, so 𝒯\_J² ≲ (log 1/ε)\^{}\{-3\}. - (63.2)
therefore permits a saving of at most O((log 1/ε)\^{}\{-3/2\}). The
actual savings, ε (first order) and ε²log (second order), lie far inside
this.

\hypertarget{proposition-63.ak-the-exact-almost-kuxe4hler-identity}{%
\subsection{2. Proposition 63.AK: the exact almost-Kähler
identity}\label{proposition-63.ak-the-exact-almost-kuxe4hler-identity}}

Now let (g, J, Ω) be almost-Kähler. AAD conventions: - ∇Ω = a⊗φ − Ja⊗Jφ;
- u := \textbar a\textbar² = \textbar∇Ω\textbar²/4; - s* − s =
\textbar∇Ω\textbar² = 4u; - W⁺ = W₁ + W₂ + W₃, with conformal scalar
curvature κ\_J = 3⟨W⁺Ω, Ω⟩ = s + 6u; - \textbar W₁\textbar² = κ\_J²/24;
- W₂ = ½(ρ\_*'`⊗Ω + Ω⊗ρ\_\emph{'\,'), so \textbar W₂\textbar² =
\textbar ρ\_}''\textbar².

From (63.1), γ\_J\^{}+ = ((s + s\emph{)/8)Ω + ρ\_}'\,', so
\textbar γ\_J\^{}+\textbar² = (s+2u)²/8 + \textbar ρ\_*'\,'\textbar²;
and γ\_J\^{}- = ρ₀ − (a∧Ja)\^{}-.

Combine three exact relations: - Chern--Weil: P = 4∫\textbar W⁺\textbar²
+ ∫s²/12 − 8π²c₁²(M); - ∫\textbar γ\_J\^{}+\textbar² −
∫\textbar γ\_J\^{}-\textbar² = 4π²c₁(J)² = 4π²c₁²(M); - the pointwise
identity (s+6u)²/6 + s²/12 − (s+2u)²/4 = su + 5u².

They give:

\textbf{Proposition 63.AK.} For every almost-Kähler metric on a closed
4-manifold, \[
\boxed{\ P=2\|\gamma_J^-\|^2+\int u\,(5u+s)\,d\mu+2\int|\rho_*''|^2+4\int|W_3^+|^2\ }
\tag{63.3}
\] with ‖γ\_J\^{}-‖² ≥ 4π²\textbar c₁(J)\^{}-\textbar². Hence \[
P\ \ge\ 8\pi^2|c_1(J)^-|^2-\int u\,(|s|-5u)_+ .
\]

Checks. - Kähler (u = 0): (63.3) is P = ∫\textbar Ric₀'\textbar². This
is correct because Ric'\,' = 0. - \textbf{Sekigawa's theorem is (63.3)
at P = 0, s ≥ 0.} Every term is then nonnegative, so u ≡ 0. - AK
Einstein with s \textless{} 0. Here γ\^{}- = −(a∧Ja)\^{}- and ‖γ\^{}-‖²
= ½∫u². Then (63.3) gives \[
  \int u\,(|s|-6u)=2\int|\rho_*''|^2+4\int|W_3^+|^2,\qquad
  \tfrac12\int u^2\ge4\pi^2|c_1(J)^-|^2 .
  \] A strictly AK Einstein metric therefore needs ‖∇Ω‖⁴\_\{L⁴\} = 16∫u²
≥ 128π²\textbar c₁\^{}-\textbar², and also u ≤ \textbar s\textbar/6 on
average in the weighted sense above. These are two integral constraints
for A01's Goldberg problem. Novelty against the literature (Armstrong's
and Drăghici's integral formulas) is not yet checked. - Saving against
8π²\textbar c₁\^{}-\textbar² requires torsion in the window 0
\textless{} u \textless{} \textbar s\textbar/5. The pointwise maximum is
s²/20. Optimizing (63.3) and Chern--Weil pointwise over (s, u)
reproduces exactly LeBrun's k − A/3, with multiplier λ = 4/3. So
\textbf{(63.3) by itself sits at the LeBrun barrier.} New information
comes only from the ASD-side split γ\^{}- = ρ₀ − t\^{}- (Theorem 63.T)
together with the moments of u.

\hypertarget{the-equality-locus-of-theorem-62.a-is-empty-for-every-k-1}{%
\subsection{3. The equality locus of Theorem 62.A is empty for every k ≥
1}\label{the-equality-locus-of-theorem-62.a-is-empty-for-every-k-1}}

Equality in 62.A requires three things: - LeBrun equality (Prop. 2.4):
AK at the Yamabe metric, Ω on the lowest eigenspace of W⁺, W₂ = W₃ = 0;
- the Cauchy--Schwarz step: s and \textbar W⁺\textbar{} constant; - Π =
0.

Then κ\_J = s + 6u is constant, so \textbf{u = \textbar∇Ω\textbar²/4 is
constant}. Two independent arguments exclude this unless k = 0 and the
metric is Kähler.

\begin{itemize}
\tightlist
\item
  \textbf{(a) Chern classes (Armstrong; LeBrun; AAD Lemma 2).} If u ≡
  const \textgreater{} 0, the Nijenhuis tensor never vanishes. Then
  c₁(I) = 3c₁(J), which forces 5χ + 6τ = 0. But N is a section of
  Λ\textsuperscript{\{0,1\}⊗Λ}\{0,2\} ≅
  T\textsuperscript{\{1,0\}⊗K}\{-1\}, whose c₂ is c₂ + 2c₁². On M this
  is \[
  5\chi+6\tau=c_2(X)+2A-k>0
  \] for every k \textless{} c₂(X) + 2A, so in particular in the whole
  open range. If u ≡ 0, J is Kähler and κ ≤ 0 (Step 62 §3a).
\item
  \textbf{(b) Periods (Taubes; Li--Liu).} Π = 0 forces E\_i\^{}+ = 0,
  hence E\_i·{[}Ω{]} = 0 because {[}Ω{]} ∈ H⁺\_g. But Ω is symplectic,
  and ±E\_i is represented by a symplectic (−1)-sphere, which has
  positive area.
\end{itemize}

Answer to the scope question on AAD. Their Theorem 2 is stated for the
\emph{third} Gray condition, i.e.~Ric'\,' = 0 together with W(Ω) =
(κ/6)Ω (lines 54--62). The 62.A locus has W(Ω) = (κ/6)Ω and W₃ = 0 with
constants, but not Ric'\,' = 0. Argument (a) makes AAD unnecessary.
\textbf{So rigidity on the locus holds and is cheap. The whole content
of the open-range problem is stability.}

\hypertarget{what-the-combined-system-allows-the-saturating-asd-patch-configuration}{%
\subsection{4. What the combined system allows: the saturating
``ASD-patch''
configuration}\label{what-the-combined-system-allows-the-saturating-asd-patch-configuration}}

The integral system available so far consists of: - Chern--Weil (Euler
and signature); - LeBrun's Prop. 4 for the best sign; - Yamabe; - the
flux bound 63.T.

All four are evaluated in the AK model at s ≡ −4, with W₂ = W₃ = 0 and
the torsion density u a function. Write the volume as V = 2π²(A + υ),
and set h = k/A.

\textbf{(A) Constant torsion (u ≡ ū).} Here 63.AK and 63.T, with ρ₀ = 0
allowed, give exactly \[
(A+\upsilon)\,\bar u(4-3\bar u)\ge4(A-k),\qquad
p\ \ge\ (A-k)\frac{(2-3\bar u)^2}{\bar u(4-3\bar u)} .
\] - The bound vanishes only at ū = 2/3 = \textbar s\textbar/6, which is
the ASD point W⁺ = 0. Its value there is κ = A/2, the Taubes regime. -
In the formal-common-mode regime ū ≈ 2k/3A it gives p ≳ 3A²/2k ≫ k. -
\textbf{The formal common mode, which is LeBrun-extremal with nearly
constant torsion, is therefore excluded as a p ≪ k state by the flux
bound 63.T}, independently of Taubes' periods. Its 1/(4A) coefficient
came from an algebra that never imposed the ASD flux
2π\textbar c₁\^{}-\textbar{} ≥ 2π√k. - Directly: with uniform torsion
and κ ≈ k, 𝒯² = 4π²κ²/(9A). So 63.T gives p ≥ k(1 − κ/(3√(kA)))² ≈ k.

\textbf{(B) Variable torsion.} - \emph{Variance enters the 62.A slack.}
The variance of u enters the Cauchy--Schwarz slack of 62.A with an
explicit constant. To leading order the slack in p is σ = ¼A·Var(u),
i.e. \[
  \int(u-\bar u)^2=8\pi^2\sigma .
  \] - \emph{What flux transport by twist requires.} Carrying the flux
by twist needs ½∫u² ≈ 4π²k when p ≈ 0. That forces σ ≈ k, then κ ≈ 2k
and υ ≈ 5k. - \emph{The extreme distribution saturates everything.} The
extreme distribution is u = 2/3 (W⁺ = 0) on patches of total volume
18π²k and u = 0 elsewhere. It saturates \textbf{all} constraints
simultaneously: - LeBrun Prop. 4 gives υ = 5k exactly, since 36V −
20V\_p = 72π²A; - the flux bound holds with ½∫u² = 4π²k; - and p = υ −
5k = 0. - \emph{Per-patch data.} Each patch is an Einstein, s = −4, W⁺ =
0 region of volume 18π², with 2χ + 3τ = Vol/6π² = 3. These are
\textbf{exactly the numbers of an orientation-reversed fake projective
plane}: χ = 3, τ = −1, ASD Einstein, Vol = 2π²c₁² = 18π², and rational
intersection form ⟨−1⟩, the same as CP̄².

\textbf{Interpretation (changes a global judgement).} - Every integral
constraint used in Steps 36--63 --- Chern--Weil, the SW sign family,
Hodge flux, Yamabe and Weitzenböck --- depends only on rational
cohomology and Chern numbers. None can distinguish a CP̄² summand from a
rational-homology CP̄² such as \textbackslash overline\{FPP\}. - We
expect the SW blow-up formula to extend to negative-definite summands
with b₁ = 0. This needs a source. - Within that information, the
saturating configuration is X ⊔ k·\textbackslash overline\{FPP\} with p
= 0. What it ignores is the \textbf{joining}. A two-ended AE neck
(χ(S³×ℝ) = 0) has P = 2∫\textbar W\textbar² + ∫s²/12 + 32π² ≥ 32π²,
i.e.~4 units, by Gauss--Bonnet with two Euclidean ends. - So a positive
tariff on X\#kCP̄², or even the k²/A law, cannot come from such
inequalities alone. It must price one of two things: - \textbf{the
transition between twist-carrying (W⁺ ≈ 0) and W⁺-carrying (Kähler-type)
regions}, which in the degenerate limit costs at least a neck; - or the
smooth/homotopy type of the summand: CP̄²∖B is simply connected, whereas
an ASD Einstein s \textless{} 0 piece with χ = 3, τ = −1 is a
ball-quotient type. - This is the precise version of Step 62's remark
that ``k must enter the Weyl channel itself''. Conjecture 62.W/62.G now
has a concrete competitor geometry to defeat.

\hypertarget{consequences-for-the-target-choice-strategist-78}{%
\subsection{5. Consequences for the target choice (strategist
\#78)}\label{consequences-for-the-target-choice-strategist-78}}

\begin{itemize}
\tightlist
\item
  \textbf{62.W and 62.G.} The flux bound excludes the formal-mode
  mechanism (diffuse torsion, κ ≈ k) for p ≪ k. What remains for p ≪ k
  is the patch mechanism, which needs κ ≈ 2k \textgreater{} k. So the
  surviving low-cost scenario violates 62.W (κ ≤ k) itself. Either 62.W
  fails or the patch mechanism is too expensive.
\item
  \textbf{Plausible lower-bound mechanism.} A Kähler-type ↔ ASD-type
  transition cost, i.e.~a quantitative ``domain-wall'' estimate for the
  frame twist, plausibly gives a positive per-class tariff (62.G-type).
  It would also explain why every known object has p/k ≥ 0.9957.
\item
  \textbf{Next decisive questions.}

  \begin{itemize}
  \tightlist
  \item
    (Q1) A local wall estimate. In a region where the twist 𝒯 carries
    flux 2π, must the transition from u ≈ 0 to u ≈ \textbar s\textbar/6,
    i.e.~from \textbar W⁺\textbar{} ≈ \textbar s\textbar/√24 to W⁺ ≈ 0,
    cost P ≳ c \textgreater{} 0, uniformly in the shape? The relevant
    quantity is the Euler/signature balance on the wall: the
    Gauss--Bonnet neck bound is the degenerate limit.
  \item
    (Q2) A test of the competitor. Is
    α(X\#k\textbackslash overline\{FPP\}) \textless{} 8π²k?
    Equivalently, can negative Einstein pieces be joined more cheaply
    than a 4-unit neck through a fat transition? If yes, the
    CP̄²-versus-\textbackslash overline\{FPP\} distinction is essential
    to any tariff on X\#kCP̄².
  \item
    (Q3) Make 63.T effective for non-AK near-minimisers. Use J :=
    σ(ψ)/\textbar σ(ψ)\textbar{} from the SW spinor of the best class,
    modified near zeros of ψ. Control 𝒯\_J by the SW slack in L⁴ of
    ∇ψ/\textbar ψ\textbar, as in the AK model where the slack controls
    ∫(u − ū)².
  \end{itemize}
\end{itemize}

\hypertarget{status-1}{%
\subsection{Status}\label{status-1}}

\begin{longtable}[]{@{}ll@{}}
\toprule
\begin{minipage}[b]{0.47\columnwidth}\raggedright
Item\strut
\end{minipage} & \begin{minipage}[b]{0.47\columnwidth}\raggedright
Status\strut
\end{minipage}\tabularnewline
\midrule
\endhead
\begin{minipage}[t]{0.47\columnwidth}\raggedright
63.T\strut
\end{minipage} & \begin{minipage}[t]{0.47\columnwidth}\raggedright
Proved: Hodge theory plus the Gauss formula (63.1). Sign cross-checked
against AAD (25); a symbolic check is still owed.\strut
\end{minipage}\tabularnewline
\begin{minipage}[t]{0.47\columnwidth}\raggedright
63.AK\strut
\end{minipage} & \begin{minipage}[t]{0.47\columnwidth}\raggedright
Proved, given AAD's formulas, by exact algebra. Sekigawa is recovered.
Independent verification pending.\strut
\end{minipage}\tabularnewline
\begin{minipage}[t]{0.47\columnwidth}\raggedright
§3 emptiness\strut
\end{minipage} & \begin{minipage}[t]{0.47\columnwidth}\raggedright
Proved, modulo LeBrun Prop. 2.4's exact equality statement and the
Taubes/Li--Liu (−1)-sphere theorem (cite).\strut
\end{minipage}\tabularnewline
\begin{minipage}[t]{0.47\columnwidth}\raggedright
§4\strut
\end{minipage} & \begin{minipage}[t]{0.47\columnwidth}\raggedright
Model computations: AK, s constant, W₂ = W₃ = 0, linearized slack. The
\textbackslash overline\{FPP\} identification is exact numerology. The
neck bound is Gauss--Bonnet. Nothing here is a theorem about general
metrics.\strut
\end{minipage}\tabularnewline
\bottomrule
\end{longtable}

\hypertarget{verification-ttcs-step63_ttcs01-verifier-generator-binding-corrections}{%
\subsection{6. Verification (TTCS step63\_ttcs01, verifier + generator)
--- BINDING
corrections}\label{verification-ttcs-step63_ttcs01-verifier-generator-binding-corrections}}

Source:
\texttt{trace/ttcs/results/step63\_ttcs01\_verifier1\_integration.md}
and \texttt{\ldots{}generator1\_integration.md}.

\textbf{Confirmed in full:} - 63.T, including (63.1) γ\_J = R(Ω) + α∧β.
The sign, the absence of any factor 2, and gauge invariance are all
correct, and (63.1) holds without dΩ = 0. R(Ω)\^{}- = ρ₀ with
\textbar ρ₀\textbar² = ½\textbar Ric₀'\textbar². The Hodge step and the
Wu existence step hold. - Sharpening: for the maximising sign,
\textbar a\_ε\^{}-\textbar² = k + Π holds with \emph{equality}. - Scope:
the Kähler reduction must be applied to the actual parallel J. - 63.AK:
every coefficient is correct. - The Sekigawa check holds. - The
AK-Einstein identities hold: ∫u(\textbar s\textbar{} − 6u) =
2∫\textbar ρ\_\emph{'\,'\textbar² + 4∫\textbar W₃\textbar² and
∫\textbar∇Ω\textbar⁴ ≥ 128π²\textbar c₁(J)\^{}-\textbar². The first is a
weighted integral identity, }not* a pointwise bound u ≤
\textbar s\textbar/6. - The algebraic LeBrun barrier is confirmed
(minimum k − A/3 at u = 1/2). - §3, emptiness of the equality locus for
k ≥ 1. - The period argument is already contained in LeBrun Prop 3.2
(Taubes), so Li--Liu is not needed. - The Chern-class argument needs
only 5χ + 6τ ≠ 0, i.e.~k ≠ c₂(X) + 2A. - The constant-positive-torsion
bound p ≥ (A−k)(2−3u)²/(u(4−3u)) is valid for 0 \textless{} u ≤ 2/3. -
The two-ended AE neck bound P ≥ 32π² is correct for complete two-ended
AE limits. It is sharp, attained by (1 + a/r²)²g\_\{ℝ⁴\}.

\textbf{Corrected:} 1. \textbf{Step 62.} - The intermediate coefficient
is (2/3)(A − 2κ), not (1/6)(A − 2κ). The final identity 62.1 is
unchanged. - 62.A holds for every 0 ≤ κ ≤ A/2. The function L\_A
\emph{decreases} only up to 4A/9 and then increases; its minimum is k −
A/3. The phrase ``decreasing for κ \textless{} 2A'' is wrong. - For κ
\textless{} 0, use p ≥ k − (4/3)κ. 2. \textbf{Corollary 62.D is
withdrawn as stated.} The term k − κ can be negative and cancel (υ −
κ)/3. So p/k → 0 implies only κ ≥ k − p and υ = 4κ − 3k + 3p, together
with (62.2). It does not imply κ/k → 1, nor Π = o(k), nor convergence to
the equality locus. The patch densities (κ, υ) = (2k, 5k) with p = 0 are
exactly such a cancellation. 3. \textbf{Formal-mode torsion is u = h =
k/A, not 2k/(3A).} It follows from the eigenvalues (−x, x/2, x/2) with x
= 2/3 − h, so κ\_J = −4 + 6h. Consequently: - the constant-torsion bound
gives p ≳ A²/k; - the best-class twist bound gives p ≥ k(1 − ½√(k/A) +
\ldots)². Both exclude the formal mode \emph{as an exact AK model with
these hypotheses}. Neither excludes non-AK near-minimisers. 4.
\textbf{Variance and slack.} The relation ∫(u − ū)² = 8π²σ is only
leading order, valid as m → 0, B/A → 1, Π = 0 and with saturated LeBrun.
LeBrun slack contributes to σ independently. 5. \textbf{The patch model
is a formal density only.} - It needs k ≤ A/4, since the residual Kähler
volume is 2π²(A − 4k). - It does \emph{not} saturate the Cauchy--Schwarz
step. - With p = 0, W₂ = W₃ = 0 and Ric'\,' = 0, AAD (37) gives d\^{}J u
= 0. So u is constant, and no smooth AK-Einstein two-phase metric
exists; any transition must activate W₃, ρ\_*'\,' or Ric'`. The
governing identity is 2d\^{}Ju = 2⟨δW₃,Ω⟩ + 2δρ\_\emph{'\,' − ⟨ρ\_}'',
∇\_JΩ⟩ + ⟨A\^{}+, Ω⟩, with A\_Z = (d\^{}∇Ric'`)\emph{Z + ι}\{JZ\}dρ. -
\textbackslash overline\{FPP\} has b⁺ = 0 and no orientation-compatible
ACS. It is an ASD-Einstein competitor, not an AK patch. - For the full
pieces: X ⊔ k\textbackslash overline\{FPP\} has 2χ + 3τ = A + 3k, and
with full-volume pieces κ = 0 and υ = 9k. The connected-sum identity
then gives p = 4k: four units per join, consistent with the neck bound.
The patch model has silently removed four units of host volume per
class. - \textbf{The ``blindness to CP̄² versus
\textbackslash overline\{FPP\}'' sentence of §4 is downgraded.} The cap
estimates below do not distinguish the two summands, since both
punctures have (χ, τ) = (2, −1). But no low-cost competitor has been
exhibited, and the SW extension to negative-definite summands is
unsourced. 6. \textbf{The §1(iv) collar-twist estimate is withdrawn.} A
one-parameter rotation of Ω in a parallel frame has α∧β = 0. Twist and
frame rotation are different: the latter can cost nothing, with
∫\textbar∇Ω\textbar⁴ = 8π²Δθ⁴/L³ → 0. 7. \textbf{Arbitrary-J twist
coercivity ‖t‖ ≤ √(CP) + r is false.} On Σ×Σ' with equal Einstein
constants, P = 0, yet a J in a characteristic class with c₁\^{}- ≠ 0
carries all its flux by twist. Any coercivity must use a
\emph{canonically selected} J, e.g.~from the SW spinor; the obstacle
there is the unweighted ∫\textbar∇(ψ/\textbar ψ\textbar)\textbar⁴ near
the zeros of ψ.

\textbf{New rigorous wall estimate (generator; proved by localized
Chern--Weil).} Let U ≅ (\#\_m CP̄²)∖B⁴ be a compact region with boundary.
Define 𝔟(∂U) as the Gauss--Bonnet plus signature (APS-corrected)
boundary terms and Q(U) = ∫\_U(4\textbar W⁺\textbar² + s²/12). Then \[
P(U)=Q(U)+8\pi^2(m-2)+2\mathfrak b(\partial U).
\] Suppose U replaces a ball B of a reference metric h with boundary
data matched up to \textbar 𝔟\_g − 𝔟\_h\textbar{} ≤ δ. Then \[
P_g(U)\ \ge\ 8\pi^2m-Q_h(B)-2\delta,\qquad
\text{and for a KE reference with } s=-4:\ \ P_g(U)\ge8\pi^2m-4\operatorname{Vol}_h(B)-2\delta .
\] - This holds for arbitrarily fat interiors. - With exact matching (δ
= 0) and Vol\_h(B) ≤ π², it gives a tariff of ½ per cap. - A one-ended
AE exceptional bubble has P ≥ 8π²m, with equality for Burns. The same
holds for a punctured \textbackslash overline\{FPP\}. - It is the region
form of the Step-54 Lemma 54.1′ charge bound, now with an explicit
replaced-bulk term. - \textbf{The only escape in this formulation is a
large replaced host volume or a large boundary-transgression error.} The
quantitative target that would close the open range is a
transgression--coercivity inequality (8π² − 𝔟(∂U))₊ ≤ q·P(U ∪ collar) +
ε, which gives P ≥ (8π²m − 2ε)/(1 + 2q). - Cheap joinings exist only
through collapsed interfaces (CH-cusp collars give P = O(e\^{}\{−4T\})).
That is the familiar collapse channel, not an exceptional S³ cap.

\textbf{Revised reading of the step.} The frame-twist theorem and the AK
identity are solid tools, and the equality locus is empty. The
quantitative content now sits in two places: - (i) a canonically
selected J, which arbitrary J cannot provide; - (ii) the
transgression/replaced-volume ledger of exceptional caps, which needs a
coercivity estimate for 𝔟(∂U) against P on the collar.

Item (ii) is the cleanest formulation of the open-range problem obtained
so far. Its Bianchi-IX restricted version, an exceptional bolt with
CH-type outer boundary data, is computable with existing diagonal tools:
compute I\_U(h, L) = 8π²(m−2) + 2𝔟(h, L) + inf Q.

\hypertarget{s1003a02-step-64-the-diffuse-frame-twist-target-what-the-aad-transport-identity-controls-and-why-the-sw-bloch-field-is-singular}{%
\chapter{S1003A02: Step 64 --- The diffuse frame-twist target, what the
AAD transport identity controls, and why the SW Bloch field is
singular}\label{s1003a02-step-64-the-diffuse-frame-twist-target-what-the-aad-transport-identity-controls-and-why-the-sw-bloch-field-is-singular}}

Source: \texttt{step64\_diffuse\_twist\_AAD\_and\_SW\_zeros.md}.

Agent3, Execute. The shell remained unavailable, so the exact local
check was derived and staged rather than executed. This note corrects
the route judgment after Step 63: neither a positive per-class tariff
nor a \(k^2/A\) law is presently favored by the proved inequalities.

\hypertarget{the-actual-low-defect-construction-target-supplied-by-theorem-63.t}{%
\subsection{1. The actual low-defect construction target supplied by
Theorem
63.T}\label{the-actual-low-defect-construction-target-supplied-by-theorem-63.t}}

Let \(M=X\#k\overline{\mathbb {CP}}^{\,2}\), let \(a=a_\varepsilon\) be
a sign class maximizing \((a^+)^2\), and let \(h_a^-\) denote the
harmonic ASD representative normalized so that
\([2\pi h_a^-]=2\pi a^-\). For a smooth orthogonal almost-complex
structure \(J\) with \(c_1(J)=a\), write \[
 \gamma_J^-=\rho_0^J+t_J^-,\qquad t_J^-=(\alpha\wedge\beta)^-,
\] with the sign understood as in (63.1); in the almost-Kähler
convention \(t_J^-=-(a\wedge Ja)^-\). If \(P=o(k)\), Theorem 63.T and
Hodge orthogonality require \[
 \|\rho_0^J\|_2=o(\sqrt{k}),\qquad
 \|t_J^-\|_2=2\pi\sqrt{k+\Pi}+o(\sqrt{k}),
\tag{64.1}
\] and the efficient version is stronger: \[
 t_J^-=-2\pi h_a^-+o_{L^2}(\sqrt{k}),
\tag{64.2}
\] up to the convention sign. Thus the exceptional classes should be
nearly ASD (\(\Pi\simeq0\)), their harmonic representatives should be
diffuse, and the rotation of the \((2,0)\)-plane should have an ASD
quadratic twist nearly equal to the harmonic flux.

For an almost-Kähler structure put \(u=|\nabla\Omega|^2/4\). Then \[
 |t_J^-|^2=\frac{u^2}{2},
\] so (64.1) forces \[
 \frac12\int_Mu^2\ge4\pi^2(k+\Pi)-o(k).
\tag{64.3}
\] At the negative-KE volume scale \(V\simeq2\pi^2A\), the smallest
diffuse RMS size is \[
 u_{\rm rms}\ge 2\sqrt{\frac{k+\Pi}{A}}+o(\sqrt{k/A}),
 \qquad |\nabla\Omega|_{\rm rms}\asymp (k/A)^{1/4}.
\tag{64.4}
\] This corrects the old common mode: its torsion \(u=k/A\) is too small
by a square root to carry the ASD Chern flux. But (64.4) does
\textbf{not} imply a tariff. For \(k/A\ll1\), this torsion lies in the
negative window of the exact identity \[
 P=2\|\gamma_J^-\|^2+\int u(5u+s)+2\|\rho_*''\|^2+4\|W_3^+\|^2.
\tag{64.5}
\] The positive \(W_2/W_3\) energy needed to balance the negative linear
torsion term can be financed by the simultaneous deletion of the
Kähler-type \(W_1\) component; (64.5) is exactly the negative Goldberg
balance. Hence neither 62.W nor 62.G follows from 63.T.

\textbf{Construction problem 64.C.} Construct, or obstruct, a sequence
\((g_j,J_j)\) with the following properties simultaneously: 1. \(g_j\)
is the scalar \(-4\) Yamabe representative and \(P(g_j)=o(k)\) (or has a
prescribed smaller scale); 2. \(c_1(J_j)=a_{\varepsilon_j}\),
\(\Pi_j=o(k)\), and (64.2) holds; 3. \(u_j\) is diffuse with RMS size
\(2\sqrt{k/A}\), while its unavoidable zeros/transition cores occupy
vanishing volume; 4. the actual Levi-Civita Gauss/Ricci/Bianchi
equations hold, not only the pointwise algebraic identities.

This is the precise cheap-state template to test. It is more informative
than a collection of nearly unit-cost local cells, which never enters
the diffuse scale.

\hypertarget{a-quantitative-weak-form-of-the-aad-transition-identity}{%
\subsection{2. A quantitative weak form of the AAD transition
identity}\label{a-quantitative-weak-form-of-the-aad-transition-identity}}

For every almost-Kähler four-manifold, AAD Corollary 1 gives, since
\(\kappa-s=6u\), \[
\boxed{
2d^Ju=
2\langle\delta W_3^+,\Omega\rangle
+2\delta\rho_*''
-\langle\rho_*'',\nabla_J\Omega\rangle
+\langle A^+,\Omega\rangle,
}
\tag{64.6}
\] where \[
 A_Z=(d^\nabla\operatorname{Ric}'')_Z+\iota_{JZ}d\rho,
 \qquad d\rho=*\bigl(J\delta\operatorname{Ric}''\bigr).
\] The last identity is AAD (26), so the derivatives in \(A\) involve
the anti-Hermitian Ricci tensor, not a new curvature component.

Integrate (64.6) against a smooth one-form \(\eta\). Integrating each
divergence term by parts and using \(|\nabla\Omega|=2\sqrt u\) gives a
universal algebraic constant \(C\) such that \[
\begin{split}
2|\langle d^Ju,\eta\rangle|
\le C\bigl(&\|W_3^+\|_2+\|\rho_*''\|_2+\|\operatorname{Ric}''\|_2\bigr)\\
&\times\bigl(\|\nabla\eta\|_2+\|\sqrt u\,\eta\|_2\bigr).
\end{split}
\tag{64.7}
\] No derivative of curvature remains in (64.7). Thus, for the weighted
negative norm \[
 \|\xi\|_{\mathcal H_u^{-1}}
 :=\sup_{\eta\ne0}\frac{|\langle\xi,\eta\rangle|}
 {\|\nabla\eta\|_2+\|\sqrt u\,\eta\|_2},
\] one has \[
 \|d^Ju\|_{\mathcal H_u^{-1}}
 \le C\bigl(\|W_3^+\|_2+\|\rho_*''\|_2+\sqrt P\bigr).
\tag{64.8}
\] Using (64.5), a weaker closed estimate is \[
 \|d^Ju\|_{\mathcal H_u^{-1}}^2
 \le C\left[P+\int u(|s|-5u)_+\right].
\tag{64.9}
\] Equations (64.7)--(64.9) are the rigorous scale at which AAD
transports torsion amplitude into curvature components. They control
only the nonconstant part of \(u\). They do not control the mean diffuse
mode, the orientation of \(t_J^-\), or its harmonic projection.

\hypertarget{why-64.8-does-not-yet-give-a-gap}{%
\subsubsection{Why (64.8) does not yet give a
gap}\label{why-64.8-does-not-yet-give-a-gap}}

A positive constant value of \(u\) is topologically impossible in the
open range because the Nijenhuis section has \[
 c_2(T^{1,0}\otimes K^{-1})=5\chi+6\tau=c_2(X)+2A-k\ne0.
\] But topology does not impose a quantitative size on the zero cores. A
profile equal to \(u_0\) away from finitely many balls of radius
\(\epsilon\), and vanishing inside them, has \[
 \|u-u_0\|_2^2=O(u_0^2\epsilon^4)
\] per isolated core. Its derivative therefore tends to zero in the
natural exact \(H^{-1}\) scale. A degree-one phase of the Nijenhuis
section on the linking \(S^3\) has critical boundary energy, but the
bulk \(L^2\) energy of its first derivative scales like
\(u_0\epsilon^2\) and can vanish. The four-dimensional \(L^2\) curvature
scale is therefore supercritical for pricing isolated zeros solely
through their degree.

Consequently, the desired estimate \[
 \|t_J^-\|_2\le C\bigl(\|W_3\|_2+\|\rho_*''\|_2+\sqrt P\bigr)
\tag{64.10}
\] does not follow from (64.6): (64.6) annihilates the constant mode and
sees no quantitative radius for the topologically required cores. A
theorem of the form (64.10) needs additional input that couples the
\textbf{harmonic projection and orientation} of \(t_J^-\) to the second
AAD Bianchi identity, or an elliptic/noncollapse statement that prevents
the Nijenhuis zeros from concentrating. This is exactly the global
information missing from arbitrary near-minimizers.

The second AAD identity does contain the orientation information, \[
0=(d u\wedge Z^\flat)''+\frac\kappa4\nabla_{JZ}\Omega-\nabla_{\rho_0(Z)}\Omega
+2(\delta W_3^+)_Z''+\nabla_{\rho_*''(Z)}\Omega
+(\nabla_{JZ}\rho_*'')''+(A_Z^+)''.
\tag{64.11}
\] However, pairing it with \(\nabla\Omega\) is the negative-scalar
Goldberg problem: the zeroth-order term has the noncoercive sign that
allows strict almost-Kähler negative Einstein geometry. Steps 43 and 63
already identified the resulting derivative loss. Thus (64.11) is a
genuine route only with a new spectral/noncollapse input, not by
Cauchy--Schwarz alone.

\hypertarget{the-sw-selected-bloch-field-does-not-directly-enter-theorem-63.t}{%
\subsection{3. The SW-selected Bloch field does not directly enter
Theorem
63.T}\label{the-sw-selected-bloch-field-does-not-directly-enter-theorem-63.t}}

Let \(\Phi\) be a positive spinor, \(\rho=|\Phi|^2\), and on
\(\{\Phi\ne0\}\) let \[
 n=\frac{q(\Phi)}{|q(\Phi)|}:M\setminus Z(\Phi)\to S(\Lambda^+)
\] be its unit Bloch direction. Theorem N controls the weighted quantity
\[
 \int\rho^2|\nabla n|_{\mathcal W_2}^2.
\] Theorem 63.T instead requires a smooth nonvanishing orthogonal
almost-complex form on all of \(M\), and controls the unweighted
quadratic twist of its moving frame. These are not the same object near
a spinor zero.

\hypertarget{proposition-64.h-the-transverse-zero-model}{%
\subsubsection{Proposition 64.H (the transverse-zero
model)}\label{proposition-64.h-the-transverse-zero-model}}

Suppose \(\Phi:\mathbb R^4\to\mathbb C^2\) has a nondegenerate zero at
the origin. After bounded linear changes, its leading unit spinor is
\(x/r:S^3\to S^3\), and its Bloch field is the radial Hopf map \[
 n(x)=\frac1{r^2}
 \bigl(2(x_1x_3+x_2x_4),\ 2(x_2x_3-x_1x_4),\ x_1^2+x_2^2-x_3^2-x_4^2\bigr).
\tag{64.12}
\] Then \[
 |\nabla n|^2=\frac8{r^2},\qquad
 |n^*\omega_{S^2}|^2=\frac{16}{r^4}.
\tag{64.13}
\] For \(\Omega=\sqrt2\,n\) and a moving orthonormal frame of
\(\Omega^\perp\), \[
 |(\alpha\wedge\beta)^-|^2=\frac8{r^4}.
\tag{64.14}
\] Therefore on an annulus \(\epsilon<r<R\), \[
 \int|\nabla n|^4=128\pi^2\log(R/\epsilon),\qquad
 \int|(\alpha\wedge\beta)^-|^2=16\pi^2\log(R/\epsilon).
\tag{64.15}
\] By contrast, because \(\rho\asymp r^2\), \[
 \int_{B_R}\rho^2|\nabla n|^2=O(R^6),
 \qquad \int_{B_R}\rho|\nabla n|^2=O(R^4).
\tag{64.16}
\] The boundary map \(n|_{S^3}\) is the nontrivial Hopf class in
\(\pi_3(S^2)\), so it cannot be filled by a smooth unit self-dual form
inside the ball while it is kept fixed outside.

\emph{Proof.} The Hopf map on the unit \(S^3\) has two horizontal
singular values equal to \(2\), giving (64.13). In a tangent frame of
\(S^2\), \(\alpha\) and \(\beta\) are orthogonal of squared norm
\(4/r^2\); a decomposable two-form has equal SD and ASD norms, giving
(64.14). Since \(|S^3|=2\pi^2\), radial integration gives (64.15). The
weighted estimates are immediate. The homotopy statement is the defining
Hopf class. The exact algebra has been staged in
\texttt{workspace/scripts/step64\_hopf\_spinor\_twist.py}.

\textbf{Consequences.} 1. The Bloch field of a generic isolated SW
spinor zero is not a global orthogonal almost-complex structure. It
cannot be repaired only inside a small ball without altering the
exterior homotopy data. 2. Even on the punctured ball its unweighted
\(L^4\) rotation and \(L^2\) frame twist diverge logarithmically. Thus
no estimate of the 63.T twist by Theorem N's weighted rotation can hold
without explicit zero-set boundary terms. 3. Theorem N is intentionally
blind to this singularity: the amplitude weight makes the local
contribution tend to zero as \(R^6\). 4. A usable canonical-field
theorem must retain the determinant connection and the zero indices,
rather than normalize the spinor and then invoke 63.T. Equivalently, one
needs a punctured Gauss formula with boundary Hopf/zero charges, or a
different smooth field supplier.

This is a structural obstruction to Q3 of Step 63, not merely a missing
estimate. The external literature retrieved this step supplies frequency
and weighted control for Dirac/SW zero sets, but no four-dimensional
theorem controlling the unweighted normalized-spinor twist; this absence
is consistent with the local logarithmic model.

\hypertarget{what-remains-mathematically-live}{%
\subsection{4. What remains mathematically
live}\label{what-remains-mathematically-live}}

The two serious alternatives are now sharper.

\hypertarget{lower-bound-route}{%
\subsubsection{Lower-bound route}\label{lower-bound-route}}

Prove a harmonic-projection estimate for a \textbf{smooth} best-class
\(J\): \[
 \|\mathsf H^-(t_J^-)\|_2^2
 \le C\,\mathcal E_{\rm geom}(g,J),
\tag{64.17}
\] where \(\mathcal E_{\rm geom}\) is controlled by \(P\) together with
already priced Weyl/Yamabe terms, and where the quantifiers exclude
collapsed interfaces only through a proved global alternative. AAD
(64.6) controls amplitude variation, while (64.11) is the necessary
orientation equation. Neither alone proves (64.17).

\hypertarget{construction-route}{%
\subsubsection{Construction route}\label{construction-route}}

Realize (64.2)--(64.4) with tiny topological zero cores. The twist
should be nearly harmonic and diffuse, while \(W_2/W_3\) implement the
negative Goldberg balance. A construction would refute every positive
per-class tariff and determine whether the remaining scale is \(k^2/A\),
smaller, or zero. Local cells do not test this regime.

\hypertarget{conditional-theorem-that-is-already-available}{%
\subsubsection{Conditional theorem that is already
available}\label{conditional-theorem-that-is-already-available}}

If an Anderson-type selected near-minimizer decomposition supplies a
noncollapsed smooth bulk and a uniform lower harmonic radius at every
Nijenhuis/field-transition core, then the AAD identities and local trace
estimates can price the finitely many cores. Without such an entrance,
shrinking cores and collapsed interfaces remain genuine escape
mechanisms. This is why the unrestricted lower bound still cannot be
inferred from the local EH/Burns theory.

\hypertarget{status-and-route-judgment}{%
\subsection{5. Status and route
judgment}\label{status-and-route-judgment}}

\begin{itemize}
\tightlist
\item
  The diffuse-twist scale (64.3)--(64.4) is a direct consequence of the
  verified Theorem 63.T.
\item
  The weak transport estimate (64.7) is obtained directly from AAD
  Corollary 1 by integration by parts; it has no curvature-derivative
  assumption.
\item
  Proposition 64.H is an elementary exact local calculation; its script
  is staged but not executed because the shell lease is still blocked.
\item
  The conclusion about canonical SW normalization is exact for a
  transverse isolated zero. General SW zero sets may be more singular,
  not better for the unweighted critical norm.
\item
  No positive tariff, \(k^2/A\) bound, or cheap closed construction
  follows yet. The proved work identifies the correct construction and
  the precise analytic obstruction: the harmonic projection of quadratic
  frame twist, including zero/core and collapse data.
\end{itemize}

\hypertarget{step-chronology}{%
\chapter{Step chronology}\label{step-chronology}}

The synthesis gives the mathematical trajectory for every Step. Repeated
handoff summaries are omitted here because their proof content is
already retained in the working notes.

\hypertarget{s1003a01-step-48}{%
\section{S1003A01, Step 48}\label{s1003a01-step-48}}

Started 10-07 21:53:22 CST; status at export: completed; turns 18, tools
18.

\hypertarget{s1003a01-step-49}{%
\section{S1003A01, Step 49}\label{s1003a01-step-49}}

Started 10-07 22:46:31 CST; status at export: completed; turns 19, tools
18.

\hypertarget{s1003a01-step-50}{%
\section{S1003A01, Step 50}\label{s1003a01-step-50}}

Started 10-08 00:49:07 CST; status at export: completed; turns 20, tools
24.

\hypertarget{s1003a01-step-51}{%
\section{S1003A01, Step 51}\label{s1003a01-step-51}}

Started 10-08 01:48:10 CST; status at export: completed; turns 17, tools
19.

\hypertarget{s1003a01-step-52}{%
\section{S1003A01, Step 52}\label{s1003a01-step-52}}

Started 10-08 03:00:43 CST; status at export: completed; turns 28, tools
35.

\hypertarget{s1003a01-step-53}{%
\section{S1003A01, Step 53}\label{s1003a01-step-53}}

Started 10-08 04:27:04 CST; status at export: completed; turns 40, tools
40.

\hypertarget{s1003a01-step-54}{%
\section{S1003A01, Step 54}\label{s1003a01-step-54}}

Started 10-08 06:20:18 CST; status at export: completed; turns 15, tools
32.

\hypertarget{s1003a01-step-55}{%
\section{S1003A01, Step 55}\label{s1003a01-step-55}}

Started 10-08 08:17:43 CST; status at export: running; turns 16, tools
22.

\hypertarget{s1003a02-step-53}{%
\section{S1003A02, Step 53}\label{s1003a02-step-53}}

Started 10-07 21:53:26 CST; status at export: completed; turns 17, tools
22.

\hypertarget{s1003a02-step-54}{%
\section{S1003A02, Step 54}\label{s1003a02-step-54}}

Started 10-07 22:21:02 CST; status at export: completed; turns 15, tools
18.

\hypertarget{s1003a02-step-55}{%
\section{S1003A02, Step 55}\label{s1003a02-step-55}}

Started 10-07 23:03:14 CST; status at export: completed; turns 16, tools
21.

\hypertarget{s1003a02-step-56}{%
\section{S1003A02, Step 56}\label{s1003a02-step-56}}

Started 10-07 23:51:19 CST; status at export: completed; turns 16, tools
25.

\hypertarget{s1003a02-step-57}{%
\section{S1003A02, Step 57}\label{s1003a02-step-57}}

Started 10-08 00:39:08 CST; status at export: completed; turns 27, tools
37.

\hypertarget{s1003a02-step-58}{%
\section{S1003A02, Step 58}\label{s1003a02-step-58}}

Started 10-08 01:25:58 CST; status at export: completed; turns 41, tools
46.

\hypertarget{s1003a02-step-59}{%
\section{S1003A02, Step 59}\label{s1003a02-step-59}}

Started 10-08 02:59:28 CST; status at export: completed; turns 16, tools
21.

\hypertarget{s1003a02-step-60}{%
\section{S1003A02, Step 60}\label{s1003a02-step-60}}

Started 10-08 03:50:27 CST; status at export: completed; turns 33, tools
42.

\hypertarget{s1003a02-step-61}{%
\section{S1003A02, Step 61}\label{s1003a02-step-61}}

Started 10-08 05:26:14 CST; status at export: completed; turns 30, tools
34.

\hypertarget{s1003a02-step-62}{%
\section{S1003A02, Step 62}\label{s1003a02-step-62}}

Started 10-08 06:58:16 CST; status at export: completed; turns 21, tools
26.

\hypertarget{s1003a02-step-63}{%
\section{S1003A02, Step 63}\label{s1003a02-step-63}}

Started 10-08 07:33:27 CST; status at export: completed; turns 12, tools
16.

\hypertarget{s1003a02-step-64}{%
\section{S1003A02, Step 64}\label{s1003a02-step-64}}

Started 10-08 08:54:34 CST; status at export: running; turns 8, tools
15.

\hypertarget{delivered-mathematical-collaboration}{%
\chapter{Delivered mathematical
collaboration}\label{delivered-mathematical-collaboration}}

\hypertarget{message-825-s1003a01-step-48}{%
\section{Message 825, S1003A01 Step
48}\label{message-825-s1003a01-step-48}}

Sent 10-07 22:13:58 CST.

A01 Step48 reply to \#824. (1) Your S\^{}4/Z2 cell has an explicit
geometric model that fixes its baseline: it is the iota-quotient of a
connected-sum neck S\^{}3x{[}-1,1{]} joining two CH collars, with
iota(x,t)=(sigma x,-t), sigma=diag(1,-1,-1,-1) on S\^{}3 (the only way
to get isolated fixed points: an SO(4) element fixing a line fixes a
2-plane, so sigma must reverse S\^{}3 orientation; d iota=-Id at the two
fixed points, hence two A1 points; C\^{}2/\{+-1\} carries both
orientations so EH can be glued with the -2 spheres E1+-E2 in M's
orientation). For a shrinking neck, Chern-Weil gives Delta P
-\textgreater{} -8pi\^{}2 Delta T\_+ = 32pi\^{}2 for the two-ended neck
(Delta chi=-2, Delta tau=0), so the quotient cell costs 16pi\^{}2 =
8pi\^{}2 per absorbed class: exactly standard, consistent with your
charge-2 statement. Any saving therefore must come from the CH germ,
exactly as for Burns. (2) Construction warning:
sigma=q-\textgreater qbar maps left-invariant Berger/triaxial metrics to
RIGHT-invariant ones, so an iota-invariant neck with exact CH collars at
both ends is NOT in our SU(2)-diagonal Bianchi-IX class (the CH collar
is left-Berger; its iota image at -t is right-Berger). The
cross-sections must interpolate left- to right-invariant structure
through the middle; the natural symmetric reduction is the U(1)xU(1)
(toric/T\^{}2) one, not SU(2). So the cell cannot be tested by my
perp\_rigorous machinery without a new T\^{}2-cohomogeneity-two solver;
I do not plan to build it now. If you want a cheap first test, a
U(2)-invariant neck that is CH at one end and the round-cone C\^{}2/Z2
model at the other (orbifold side handled by your Kronheimer EH) is in
the Bianchi-IX class after passing to the double cover; I can supply the
frame/Cartan code (workspace/scripts/cell\_ricci.py, perp\_lm.py) and
the bolt-regularity substitutions if you build it. (3) My focus this
loop: Problems 1 and 4. For Problem 1 I confirmed no reverse-orientation
actual class exists in general (Mbar = Xbar \# k CP2 is a sum of two
b+\textgreater=1 pieces when b-(X)\textgreater=1, so SW vanishes;
fake-projective-plane bases with k=1 only give a b+=1 chamber input), so
the two-orientation cross-pairing has no supply there; the
same-orientation 2\^{}k sign family only adds the average/max
information already in LeBrun and N1 (consistent with your Theorem B).
For Problem 4 I am mapping exactly what the full AAD Bianchi system
(their Lemma 4 / Cor 1 eqs (37)-(38)) plus the W+ Bochner identity
force; I will send the precise integral identity and what it does/does
not close. Question for you: does your Step-44 near-CH monopole branch
control give anything for an ARBITRARY constant-length SD harmonic form
(not the monopole n\_g)? Near a negative KE metric I can show any such
form converges to +-omega0 (holomorphic 2-form components must have
constant length, impossible on ample K), so Goldberg near KE reduces to
Koiso-type local rigidity of Einstein metrics; I would like to know if
your slice theorem already gives that Einstein metrics near g\_Gamma are
KE.

\hypertarget{message-826-s1003a02-step-53}{%
\section{Message 826, S1003A02 Step
53}\label{message-826-s1003a02-step-53}}

Sent 10-07 22:14:16 CST.

A02 Step53 (after reading the 2026-10-07 supplements; verification
pending, task s53\_t14\_tool2). Note:
Sessions/S1003A02/workspace/notes/step53\_anderson\_entrance\_bubble\_exclusion\_attainment.md.

\begin{enumerate}
\def\labelenumi{(\arabic{enumi})}
\item
  Entrance. Anderson, Asian J. Math 10 (2006), Thm 4.10, uses
  P+eps\emph{int(1+\textbar Ric\textbar{}\textsuperscript{2)}p on unit
  volume and gives a specific minimizing sequence: case (I) orbifold
  bulk + ALE trees with energy equality, or case (II) collapse. Our
  Lemma 53.0: the case-(I) equality is exactly p\_}(M)=P(V,g0)+sum
  P(N\_k), since the orbifold Euler terms telescope. So the step-0 issue
  for arbitrary sequences is off the critical path.
\item
  Thm 53.F': Theorem F generalizes to any ASD scalar-flat AE bubble with
  homogeneous r\^{}-2 tail: P=P(g)+P(N)-8 lambda\^{}2
  \textless R\^{}+\_p, U Phi\^{}\{++\}(N)
  U\textsuperscript{T\textgreater+O(lambda}\{12/5\}). Optimizing over U
  gives the von Neumann sum.
\item
  Thm 53.B: a smooth-point bubble tree can be moved to any point/frame
  without changing the diffeotype. So in any case-(I) minimizing
  configuration with a Burns-type bubble, R\^{}+\_\{g0\}\textless=0 on
  every bulk component. Consequences: s0\textless=0 constant; s0=0
  forces ASD; s0\textless0 forces
  \textbar W\textsuperscript{+\textbar{}}2\textless=s\^{}2/24. With
  your/web Chapter-5 CH cell, no bulk component contains an exact CH
  region.
\item
  Thm 53.D: if there are no Gamma!=e points, every bulk component
  attains p\_\emph{(V\_a) and p\_}(M)=sum p\_*(V\_a)+sum P(N\_k). Smooth
  bubbling transfers attainment rather than destroying it.
\item
  Thm 53.E: if Y(M)\textless0, some such component has Y\textless0.
  Proof: Green-function positivity of scalar-flat AE compactifications
  plus Kobayashi. Problem 3 (non-Bach-flat smooth optimum) thus reduces
  to excluding (a) collapse, (b) nontrivial-Gamma absorption points, (c)
  Einstein bulk.
\end{enumerate}

Testable question for you, given your construction strength: the open
borderline is p\_\emph{(V\#CPbar2)\textless p\_}(V)+8pi\^{}2 at a
minimizer with R\^{}+\textless=0 and no CH region. That means the x\^{}2
log(1/x) kernel-direction coefficient in a general real-analytic
background, beyond CH. By U(2)\_zeta invariance the tail-curvature terms
see only s and W\^{}+(zeta,zeta); the diagonal squash modes c may couple
to W\^{}- diagonal entries. Can your Chapter-5 machinery give the
coefficient as min\_c{[}Q0(c)+L\_\{Rm\_p\}(c){]} for a KE (non-CH)
background, i.e.~does it stay negative for all W\^{}-?

\hypertarget{message-827-s1003a02-step-54}{%
\section{Message 827, S1003A02 Step
54}\label{message-827-s1003a02-step-54}}

Sent 10-07 22:42:46 CST.

A02 Step54. Note
workspace/notes/step54\_k1\_ball\_quotient\_attainment\_theorem.md;
checks in scripts/step54\_charge\_and\_entropy\_checks.py and
data/step54\_charge\_and\_entropy\_checks.json. Verification is running.
Treat everything below as conditional where marked.

\begin{enumerate}
\def\labelenumi{(\arabic{enumi})}
\tightlist
\item
  Testable theorem 54.A. Setting: X a ball quotient that is not a fake
  projective plane, so b+ \textgreater= 3; M = X\#CPbar2.
\end{enumerate}

\begin{itemize}
\tightlist
\item
  Unconditional: p\_*(M) \textless{} 8pi\^{}2, by the Ch.5 analytic
  CH-collar cell.
\item
  Conditional hypotheses: an Anderson case-(I) configuration with only
  trivial-group cone points; lower semicontinuity; bubbles scalar-flat
  AE with derivative-controlled decay of some order tau \textgreater{}
  0.
\item
  Then every smooth bubble is a one-ended homotopy-sphere bubble. The
  bulk V is connected and homeomorphic to M, so E stays in the bulk.
\item
  Adding energy equality and removability: g0 attains p\_*(V), and there
  are no bubbles unless some bubble compactifies to an exotic S\^{}4. g0
  has constant s \textless{} 0. So g0 is either Einstein (Problem 1 at
  k=1, inside LeBrun's open window k \textless{} c2) or a smooth
  non-Bach-flat global minimizer (Problem 3).
\item
  Proof ingredients:

  \begin{enumerate}
  \def\labelenumi{(\alph{enumi})}
  \tightlist
  \item
    Two-orientation bubble cost: P(N) \textgreater= 8pi\^{}2 max(c(N),
    c(Nbar), 0) with c = 4(e-1) + 4b1 + b\^{}- - 5b\^{}+. A two-ended
    neck costs 32pi\^{}2, matching your 8pi\^{}2-per-class S4/Z2
    baseline. This excludes the X ⊔ (CPbar2, FS) neck-pinch.
  \item
    Indefinite bubbles: hat-N is PSC (Green function), and a PSC summand
    with b+ \textgreater= 1 kills SW when b+(M) \textgreater= 2.
    Neck-stretching works even when b+(rest) = 0; tell me if you know a
    clean citation or a counterexample.
  \item
    pi1: Gamma is freely indecomposable, and Hopf gives b2(hat N)
    \textgreater= b2(X) if pi1(hat N) = Gamma. So each hat N is a
    homotopy sphere.
  \end{enumerate}
\item
  General k: sum{[}4(e-1) + 4b1 + beta{]} \textless{} k, so not all
  exceptional classes can escape into smooth bubbles.
\item
  Open gates: collapse; orbifold points; bubble-scale P-criticality of
  Anderson's perturbed sequence. The rescaled coefficient
  eps*ell\^{}\{4-4p\} is only bounded by naive eps-regularity. Partial
  collapse result: a free-circle-action region is spin, so it cannot
  carry E.
\end{itemize}

\begin{enumerate}
\def\labelenumi{(\arabic{enumi})}
\setcounter{enumi}{1}
\item
  Entropy test. Can BCG + Bishop exclude Einstein metrics on X\#CPbar2
  for complex-hyperbolic X, as Sambusetti's argument does for
  real-hyperbolic X? No, not by itself. With Ric = -6, h(CH2) = 4
  against the Bishop bound 3sqrt2, BCG gives only Vol \textgreater=
  (64/81)Vol(X). LeBrun's int s\^{}2 \textgreater= 32pi\^{}2 c1\^{}2(X)
  already gives Vol \textgreater= Vol(X), which is stronger. A
  by-product: an Einstein metric on M cannot have Kahler-type W+ (that
  would force Vol = (1 - 1/3c2)Vol(X)). If you see a sharper
  entropy/calibration input (e.g.~Kahler-hyperbolic or Toledo-type), it
  would decide the dichotomy.
\item
  Your \#825 question, local KE rigidity near g\_Gamma. Use the Koiso
  route, not A\_C. On CH2 (Ric = -6), Rcirc on S\^{}2\_0 has eigenvalues
  -2 (J-invariant) and +4 (J-anti) (our step44 JSON). This gives
  \textless(nabla*nabla - 2Rcirc)h, h\textgreater{} \textgreater=
  2\textbar h\textbar\^{}2 on TT tensors. So there are no infinitesimal
  Einstein deformations, and Koiso local rigidity applies: Einstein
  metrics C\^{}\{2,alpha\}-close (or H\^{}s, s \textgreater{} 4, in
  Bianchi gauge) to g\_Gamma are isometric to it up to scale. The gap is
  pointwise, and Lemma M gives quotient-uniform inverses, so the
  threshold is uniform over finite covers in uniformly local Holder
  norms. H\^{}2-closeness is not enough. For general negative KE
  surfaces, Koiso gives Einstein deformations contained in complex
  deformations, and Einstein metrics near KE are KE when the Kuranishi
  space is smooth. So your Goldberg-near-KE reduction closes for ball
  quotients in C\^{}\{2,alpha\}.
\item
  Coordination. Your Ch.5 machinery is what is needed for
  p\_\emph{(M\_k) \textless{} p\_}(M\_\{k-m\}) + 8pi\^{}2 m at bulk
  minimizers with R+ \textless= 0. That is the input converting the
  general-k inequality into attainment. I will not duplicate it unless
  you say you are not doing it.
\end{enumerate}

\hypertarget{message-828-s1003a02-step-54}{%
\section{Message 828, S1003A02 Step
54}\label{message-828-s1003a02-step-54}}

Sent 10-07 22:58:14 CST.

A02 Step54 follow-up: independent verification of 54.A is done
(trace/ttcs/results/step54\_ttcs01\_verifier1\_integration.md).
Corrections to my previous message. (a) Proof order: SW must remove b+
of every bubble cap FIRST. Ordinary connected-sum vanishing does this if
b+(rest)\textgreater=1; otherwise the PSC cap plus the one-sided neck
argument does. Only then does the charge give P(N) \textgreater=
8pi\^{}2{[}4(e-1)+4b1+b2{]}. Charge \textless= 0 alone does NOT force
e=1 or b1=0: two-ended scalar-flat AE S2xS2 (Green functions) has c =
cbar = 0. (b) (U) p\_* \textless{} 8pi\^{}2 holds for EVERY ball
quotient, fake planes included, since Ch.5 builds the cell for every R;
fake planes are excluded only by the SW step. (c) New for general k: the
blowdown map is classifying and Gamma is freely indecomposable, so
exactly one vertex carries Gamma, maps with degree one to X, and has b+
\textgreater= 3. Hence it is bulk, every bubble cap is simply connected,
and bubble H2 lies in span\{E\_i\}. With p\_\emph{(M\_k) \textless{}
8pi\^{}2 k: total exceptional rank in smooth bubbles \textless{} k; all
bubbles are one-ended for k \textless= 4. (d) Bubble P-criticality comes
from compact replacement once smooth bubble convergence holds (no kappa
-\textgreater{} 0 needed). Also eps}Q\_p(g\_i) -\textgreater{} 0 on
genuine near-minimizers. (e) The Einstein/W+ side remark had the wrong
direction: Einstein metrics on X\#CPbar2 have
int\textbar W+\textbar\^{}2 \textless= int s\^{}2/24 - 2pi\^{}2,
strictly below the Kahler-type value. (f) The Koiso gap 2 is an
integrated TT Codazzi-Bochner estimate, not pointwise; cover uniformity
uses Lemma M (step45/46 notes). (g) Orbifold parity (absorbing E needs
a\^{}2 = -1) is withdrawn: even m is lattice-allowed. Remaining
geometries for k=1: collapse, nontrivial-group orbifold points,
exotic-homotopy-sphere bubbles.

\hypertarget{message-829-s1003a02-step-56}{%
\section{Message 829, S1003A02 Step
56}\label{message-829-s1003a02-step-56}}

Sent 10-08 00:35:38 CST.

A02 Step56 → A01 (answering \#826 and continuing \#825). Note:
workspace/notes/step56\_orbifold\_attainment\_zero\_energy\_bound\_Tpoint\_interaction.md.
Step-55 note with verifier repairs:
workspace/notes/step55\_k1\_orbifold\_gate\_charges\_lattices\_collapse.md
plus step56 §0.

\begin{enumerate}
\def\labelenumi{\arabic{enumi}.}
\tightlist
\item
  Conservation law Let M = V° ∪\emph{L N, where N is a bubble with end
  S³/Γ. With c}±(N) = 2/\textbar Γ\textbar{} − 2χ ∓ 3τ ∓ 3η (η(n,q) =
  −4s(q,n)):
\end{enumerate}

\begin{itemize}
\tightlist
\item
  T\_+\^{}orb(V) = T\_+(M) + c\_+(N);
\item
  (P(V)+P(N))/8π² = C\_+(V) + w(N) − (A−1), where w(N) =
  (1/2π²)∫\_N\textbar W⁺\textbar². So bubbling never changes the
  Kähler-type baseline; it only moves W⁺-charge. Consequences:
\item
  A sub-8π² configuration needs C\_+(V) + Σw(N\_j) \textless{} A,
  together with the additive budget Σ max(c\_+, c\_−, 0) \textless{} 1.
\item
  Charge alone excludes a channel only when it saturates this budget.
\item
  Excluding orbifold channels is the same open-range inequality C\_+ ≥ A
  on a non-Kähler orbifold. LeBrun's 2A/3 fails there for the same
  reason it fails on M. Exact checks:
\item
  LeBrun O(−n): P = 8π²(n−2)²/n.~So O(−4) costs exactly 8π², and E or 2E
  cannot leave through a simply connected cyclic-end bubble.
\item
  Warning: for non-simply-connected regions this fails. B̄\_\{t,a\} ♮
  (CP̄² minus a ball) has charges (−3+4/t², −5).
\end{itemize}

\begin{enumerate}
\def\labelenumi{\arabic{enumi}.}
\setcounter{enumi}{1}
\tightlist
\item
  Parity law at k=1 (Wahl ball B\_\{n,a\}, end 1/n²(1,na−1), charges (0,
  −4(1−1/n²)) for all n) If such a Ricci-flat bubble occurs in the
  complex orientation:
\end{enumerate}

\begin{itemize}
\tightlist
\item
  the resolution Y has π₁ = Γ and Y → X has degree one;
\item
  the chain embeds primitively in a negative-definite unimodular lattice
  of rank r+1 (H₁-torsion argument). Such an embedding exists iff 1−na
  is a unit square mod n², iff n is odd or 8 \textbar{} n.~Consequences:
\item
  EH/Z₂ (n=2) is impossible in the zero-charge orientation. The standard
  B₂ ⊂ CP̄² sits in the reversed orientation, with charge −3.
\item
  The first case is n=3, chain {[}5,2{]}, realized by 2E₁+E₂ and E₂+E₃
  in \#3CP̄². The complement is simply connected, so the blowdown Z is
  homeomorphic to CP̄². A diffeomorphism is unproved.
\item
  Kähler+Siu shows the bulk cannot be a complex orbifold at that point:
  the degree-one map is birational, and blow-up counting gives 3r+2 ≤
  3r+1, a contradiction.
\end{itemize}

\begin{enumerate}
\def\labelenumi{\arabic{enumi}.}
\setcounter{enumi}{2}
\tightlist
\item
  Zero-energy bound (Prop 56.U, proved) For any smooth orbifold metric g
  on an orbifold V with Wahl points, p\_*(Milnor smoothing) ≤ P(V,g),
  with gluing error O(λ²), using the Ricci-flat ALE metrics of Şuvaina.
\end{enumerate}

\begin{itemize}
\tightlist
\item
  Corollary: an Einstein orbifold with Wahl points whose smoothing is M
  gives α(M) = 0.
\item
  Conditional Theorem 56.A: p\_*(M) is attained by a P-critical metric
  with s\textless0, either on M or on such an orbifold. The minimizer is
  Einstein (then α=0) or non-Bach-flat. Hypotheses: case (I), SW\_L, and
  all nontrivial-group bubbles Ricci-flat.
\end{itemize}

\begin{enumerate}
\def\labelenumi{\arabic{enumi}.}
\setcounter{enumi}{3}
\tightlist
\item
  Local interaction at a Wahl point (negative result for the cheap
  route)
\end{enumerate}

\begin{itemize}
\tightlist
\item
  Γ-invariant curvature at a Wahl point with n≥3 is 4-dimensional:
  type-D W\^{}± along the fixed ω, ω̄, and Ric° ∝ ω∘ω̄.
\item
  The bubble's quadrupole tail is phase-independent, so scale is the
  only effective modulus; no relocation or frame freedom.
\item
  Lemma 56.L: for a P-critical bulk, the term linear in a homogeneous
  degree −4 tail has no λ⁴log(1/λ) part, because ∇P(g0)=0 and divergence
  fluxes give only powers. I checked numerically that the S³-averaged
  coefficient vanishes, for generic curvature and for invariant
  curvature, and with both Ricci-flat and non-Ricci-flat tails. The
  reason: a (2,0) quadrupole has zero average in Sym²T*.
\item
  So there is no Biquard--Ozuch-type pointwise condition at log order.
  The first interaction is a λ⁴ local form: a bulk-flux pairing of the
  quadrupole with ∇Rm\_p, plus the bubble's P-polarizability (quadratic
  in R\_p). Script:
  workspace/scripts/step56\_Tpoint\_log\_interaction.py and the controls
  script.
\end{itemize}

Questions and testable requests: (a) Your U(2) CH → C²/Z₂ test is now
known to be the n=2 case, which is lattice-forbidden at k=1 in the
zero-charge orientation. The live cell is CH-collared O\_\{9,2\}: a
1/9(1,2) point, T²-symmetric. Can cell\_ricci.py be extended to
cohomogeneity two (T²)? Alternatively, use the reversed B₂ channel
X\#P̄(1,1,4), which needs C\_+ \textless{} A−3 on the bulk (charge −3,
w=3). Inequality to test: P(CH-collared O\_\{9,2\} cell) \textless{}
8π²·0.99566\ldots{} (the Burns cell value). If it holds, 56.U makes it
an upper bound for X\#Z, and gives a nonattainment mechanism. (b) If you
have tools for the P-polarizability of a Ricci-flat ALE, meaning the
renormalized min of δ²P on B\_\{n,a\} with r² growth R\_p, its sign
structure on the 4-dimensional invariant space decides whether a
pointwise condition exists at Wahl points. (c) Collapse remains
topological: does any finite-isotropy F-structure region carry an
odd-square class? A free circle action gives spin, so no. Cyclic
π₁-image only puts the class in ZE, which does not exclude E.

\hypertarget{message-830-s1003a01-step-50}{%
\section{Message 830, S1003A01 Step
50}\label{message-830-s1003a01-step-50}}

Sent 10-08 01:38:39 CST.

A01 Step50 -\textgreater{} A02 (answers \#826/\#827(4)/\#829; note
Sessions/S1003A01/workspace/notes/step50\_wminus\_completed\_square\_and\_general\_KE\_construction.md;
verifier trace/ttcs/results/step50\_ttcs01\_verifier1\_integration.md).

\begin{enumerate}
\def\labelenumi{(\arabic{enumi})}
\tightlist
\item
  Ch.5 log saving extends to every compact negative KE surface, at every
  point, perpendicular zeta (ker R+\_p). The W\^{}- part of the
  Burns-background interaction is an exact completed square.
\end{enumerate}

\begin{itemize}
\tightlist
\item
  CGB reduction: \textbar Ric0\textbar\^{}2 dV =
  (4\textbar W+\textbar\^{}2 + s\^{}2/12) dV - 16pi\^{}2 e - 24pi\^{}2 l
  (l = signature form). The S\^{}3-averaged transgression is analytic in
  (y=r\^{}2, z=eps/r\^{}2), and r d/dr kills the diagonal coefficients
  z\^{}a y\^{}a. So the eps\^{}2 log corner coefficient F can be
  computed from 4\textbar W+\textbar\^{}2 + s\^{}2/12.
\item
  Burns is exactly ASD and scalar-flat, so F depends ONLY on Rm\_p, with
  no 4-jet. W\^{}- enters only as min over k0 of
  int{[}4\textbar W+\_1\textsuperscript{W\textbar{}}2 +
  (s\_1\textsuperscript{W)}2/12{]} \textgreater= 0. W\^{}- can never
  help.
\item
  Explicit corrector (exact sympy check: all 10 components vanish
  identically as rational functions; extended to all W\^{}- by
  SU(2)\emph{- irreducibility): k0\^{}W = (3/2) sym(h\_W kB) - (1/4)
  tr(h\_W kB) delta, h\_W = -(1/3) W\^{}-}\{ikjl\} x\^{}k x\^{}l, kB =
  sigma1\^{}2 + sigma2\^{}2. It solves W+'`(h\_W,kB) + W+'(k0\^{}W) = 0
  and s'`(h\_W,kB) + s'(k0\^{}W) = 0 at degree -2.
\item
  Independent Ric0/Laurent check on H2xH2: F(c* + t k0\^{}W) = -18 +
  5.6(1-t)\^{}2. The degree-6 run (1305 correctors) stays at -18.000,
  not below, as kappa \textgreater= 0 predicts.
\item
  So F\_KE = F\_CH(c*) = -18 and P/8pi\^{}2 = 1 - 9 eps\^{}2 log(1/eps)
  + O\_\{X,p,T\}(eps\^{}2). In Ric=-g units: 1 - (1/4) x\^{}2 log(1/x) +
  O(x\^{}2), x = Area(E)/pi.
\item
  Order-eps term: a Burns first-variation flux gives ell =
  -R+\_p(zeta,zeta) in P/8pi\^{}2 units. It is independent of every
  bounded degree-0 corrector and of the blending, and vanishes in the
  kernel orientation. This matches your -8 lambda\^{}2 \textless R+\_p,
  U Phi U\^{}T\textgreater.
\item
  Remaining exact input: Ch.5's symbolic F\_CH(c*) = -18 (our numerics
  agree to 1e-12). The analytic two-scale remainder follows Ch.5 5.3-5.5
  with S\^{}3-averaged analytic densities; parity kills odd jets.
\end{itemize}

\begin{enumerate}
\def\labelenumi{(\arabic{enumi})}
\setcounter{enumi}{1}
\tightlist
\item
  Consequences (all upper constructions, not lower bounds):
\end{enumerate}

\begin{itemize}
\tightlist
\item
  p\_\emph{(X\#CPbar2) \textless{} 8pi\^{}2 = p\_}(X) + 8pi\^{}2 for
  every negative KE X; no CH region is needed.
\item
  Closed bounds via the Yamabe shortcut: P(ghat) - P(g) = (S(ghat) -
  S(g))/12 and Y\textless0 give E(ghat) \textless= 4P(g)/(P(g) +
  8pi\^{}2 A), given LeBrun supply S \textgreater= 32pi\^{}2 A. Hence
  limsup alpha/h \textless= 4(1 - delta\_X0 theta) on unramified covers
  of ANY fixed negative KE base. Order of limits: T, then eps, then j.
\item
  For your chain: an energy-equality case-(I) configuration with KE bulk
  (X, g\_KE) plus one Burns bubble cannot occur. The bubble can be moved
  to a perpendicular frame (your 53.B), and the log term then strictly
  lowers P.
\item
  Not shown: exclusion of non-KE Einstein bulk, attainment, or
  p\_\emph{(M\_k) \textless{} p\_}(M\_\{k-m\}) + 8pi\^{}2 m. A
  non-Einstein bulk also carries a first-order bulk variation eps
  DP\_\{g0\}{[}chi(kB+k0){]}; this reduces to ell = -R+(zeta,zeta) only
  if the bulk is P-critical, which your minimizers are.
\end{itemize}

\begin{enumerate}
\def\labelenumi{(\arabic{enumi})}
\setcounter{enumi}{2}
\tightlist
\item
  Next on my side, offered to you as a precise target. By the CGB form,
  F(Rm\_p, zeta) for ANY background is a U(2)\_zeta-invariant quadratic
  form in Rm\_p alone. The invariants:
\end{enumerate}

\begin{itemize}
\tightlist
\item
  5 in (s, W+): s\^{}2, s W+\emph{zz, W+\emph{zz\^{}2,
  \textbar W+}\{z,perp\}\textbar\^{}2,
  \textbar W+}\{perp,perp,0\}\textbar\^{}2;
\item
  2 in Ric0: \textbar Ric0(zeta)\textbar\^{}2 and
  \textbar Ric0\_perp\textbar\^{}2, each a nonnegative minimized square
  with no cross terms;
\item
  W\^{}-: coefficient 0. I will compute these 7 coefficients next. That
  decides the sign of the log term at R+ \textless= 0 non-KE and
  non-Einstein bulk points, which is what your general-k comparison
  needs. Tell me which R+ \textless= 0 shapes and normalization you want
  tabulated first.
\end{itemize}

\begin{enumerate}
\def\labelenumi{(\arabic{enumi})}
\setcounter{enumi}{3}
\tightlist
\item
  On \#829(a)/(b): a CH-collared O\_\{9,2\} cell needs T\^{}2
  cohomogeneity two. My cell\_ricci machinery is Bianchi-IX only, so it
  is not quick. The P-polarizability of a Wahl ALE at quadratic R\_p
  growth is the same kind of `corner coefficient' computation, and the
  CGB form may reduce it to Rm\_p data in the same way. I can try that
  after (3) if you supply the B\_\{n,a\} metric representation you are
  using.
\end{enumerate}

\hypertarget{message-831-s1003a02-step-58}{%
\section{Message 831, S1003A02 Step
58}\label{message-831-s1003a02-step-58}}

Sent 10-08 01:56:39 CST.

A02 Step58 → A01 (response to \#830; Theorem 57.Φ confirmed, extended,
and coordinated).

\begin{enumerate}
\def\labelenumi{(\arabic{enumi})}
\tightlist
\item
  Agreement on KE backgrounds. Your \#830 derivation of F\_KE = −18 via
  the CGB reduction and explicit W⁻ corrector k0\^{}W = (3/2)sym(h\_W
  kB) − (1/4)tr(h\_W kB)δ is confirmed by our independent computation.
  Our general formula:
\end{enumerate}

Φ = −μ1(μ1 − μ2 − μ3)/2 − (μ2 − μ3)²/2,

where μ1,μ2,μ3 are eigenvalues of R⁺ = W⁺ + s/12 on Λ⁺ (operator norm
convention, sum = s/4). In the kernel μ1 = 0: Φ = −(μ2 − μ3)²/2 ≤ 0. KE
gives (s/4, 0, 0) so Φ = −s²/32 = −18 at s = −24.

W⁻ coefficient: exactly 0, confirmed at D = 4 and D = 6 (independent
product backgrounds with large W⁻ and Ric° at D = 6 give errors ≤
3e−15). Your completed-square argument is the algebraic proof; our
numerical convergence check is the computational verification.

Ric° coefficient: exactly 0 (D = 4 and D = 6). Your ``nonnegative
minimized squares'' prediction confirmed.

\begin{enumerate}
\def\labelenumi{(\arabic{enumi})}
\setcounter{enumi}{1}
\item
  Product-background cross-check. Eight backgrounds Σ₁(K₁)×Σ₂(K₂) in
  conformal coordinates (nonzero quartic jet C, Ric° ≠ 0, W⁻ ≠ 0), both
  layouts, K₁K₂ ∈ \{(−1,−1),(−1,−3),(−1,1),(−2,0.5),(1,2)\}: all match
  −s²/32 at D = 6 to machine precision. Lesson: convergence in corrector
  space depends on the gauge of the background jet; conformal
  coordinates need D = 6 where normal coordinates need D = 4.
\item
  Consequences we carry forward from 57.Φ:
\end{enumerate}

\begin{itemize}
\tightlist
\item
  Every KE surface saves at every point: p\_*(X\#kCPbar2) \textless{}
  8π²k for all minimal general-type (formal; the rigorous remainder is
  the next target).
\item
  Minimizer rigidity (conditional on the Anderson ledger and 53.B): at a
  P-critical s \textless{} 0 bulk with Burns bubbles, R⁺ ≤ 0 and the
  boundary μ1 = 0 forces μ2 = μ3. This excludes KE-type and CH-type
  points and gives \textbar W⁺\textbar² \textless{} s²/24 strictly
  everywhere.
\item
  Unification with Theorem H: R⁺ \textless{} 0 is the cost regime
  (hyperbolic), R⁺\textbar\_boundary is the saving regime (KE). Burns
  bubbling is not excluded for R⁺ \textless{} 0 because those metrics
  genuinely pay; the two known families are the interior and boundary of
  the cone.
\end{itemize}

\begin{enumerate}
\def\labelenumi{(\arabic{enumi})}
\setcounter{enumi}{3}
\tightlist
\item
  Division of labour.
\end{enumerate}

\begin{itemize}
\tightlist
\item
  We own: the formal coefficient formula, the corrector data (D = 4 and
  D = 6 M\_eff matrices and K bases), the product-background tests, the
  minimizer rigidity theorem, and the consolidation (one hypothesis
  ledger at
  workspace/notes/step58\_consolidated\_hypothesis\_ledger.md).
\item
  You own: the CGB reduction proof and the explicit W⁻ corrector (sympy
  verified); the Ch. 5 two-scale factorization code.
\item
  Critical path: the rigorous remainder. For KE backgrounds (convergent
  jets), the averaged CGB form gives an analytic expansion; the Ch. 5
  5.3--5.5 matching extends with S³-averaged densities. Do you take
  this? If yes, I concentrate on the balanced-family matching law
  (Problem 2: per-cell saving λ²ε²log/4 at Ric = λg; batch interactions;
  optimization of ε at fixed density). If not, we both try
  independently.
\end{itemize}

\begin{enumerate}
\def\labelenumi{(\arabic{enumi})}
\setcounter{enumi}{4}
\tightlist
\item
  Specific requests/offers.
\end{enumerate}

\begin{itemize}
\item
  \begin{enumerate}
  \def\labelenumi{(\alph{enumi})}
  \tightlist
  \item
    For your 7-coefficient tabulation (\#830(3)): the full cone scan is
    Φ(μ₁,μ₂,μ₃) = −μ₁(μ₁−μ₂−μ₃)/2 − (μ₂−μ₃)²/2 = −(μ₁² − μ₁(μ₂+μ₃))/2 −
    (μ₂−μ₃)²/2. At R⁺ ≤ 0 with μ₁ = s/4+ε: Φ ≈ −(s/4)(s/4−s/4)/2
    +\ldots{} The sign structure on the cone \{μ ≤ 0, Σμ = s/4\} is: Φ ≤
    0 always (the interior has first-order cost so no log term; the
    boundary has second-order saving). The ``non-KE non-Einstein'' shape
    μ = (0, a, −a+s/4), a ≠ 0, s \textless{} 0, gives Φ = −(2a − s/4)²/2
    \textless{} 0. So the log term is always ≤ 0 or absent; this is the
    ``negative on the whole cone'' confirmation.
  \end{enumerate}
\item
  \begin{enumerate}
  \def\labelenumi{(\alph{enumi})}
  \setcounter{enumi}{1}
  \tightlist
  \item
    The CH-collared O\_\{9,2\} T²-cell (\#829(a)): understood that it
    needs cohomogeneity-two extension. If you supply an explicit
    T²-metric representation of B\_\{3,1\} in your cell\_ricci format,
    we can set up the ODE system. We are using the GH moment-map
    coordinates from step56\_Tpoint\_local\_structure.py.
  \end{enumerate}
\item
  \begin{enumerate}
  \def\labelenumi{(\alph{enumi})}
  \setcounter{enumi}{2}
  \tightlist
  \item
    The polarizability at a Wahl point (\#830(4)): it is the same kind
    of corner coefficient in ν = Rm\_p, with the (2,0) quadrupole tail
    replacing the Burns tail. Lemma 56.L shows no λ⁴log term at a
    P-critical bulk. The first nontrivial interaction is O(λ⁴) --- a
    finite computation but expensive. The CGB form may reduce it as you
    suggest; we defer to the balanced-family law first unless you take
    it up.
  \end{enumerate}
\end{itemize}

Files:
workspace/notes/step57\_kernel\_burns\_coefficient\_general\_background.md
(§§1--6 updated);
workspace/notes/step58\_consolidated\_hypothesis\_ledger.md;
workspace/data/step58\_products/ (8 JSON outputs);
workspace/data/step57\_general\_form\_D\{2,4,6\}\_q8\_24.json and
*\_invfit.json.

\hypertarget{message-832-s1003a02-step-58}{%
\section{Message 832, S1003A02 Step
58}\label{message-832-s1003a02-step-58}}

Sent 10-08 02:51:18 CST.

A02 Step58 addendum → A01. (1) After our TTCS verifier: the Chern--Weil
pointwise argument for W⁻/Ric° ≥ 0 is invalid; replaced by the chiral
square 𝒞=(2W⁺,s/√12), Burns 𝒞≡0, sector coefficient \textbar 𝒜\_B +
D𝒞\_δ{[}K{]}\textbar² ≥ 0 --- your explicit k₀\^{}W is exactly the
solution of D𝒞{[}K{]}=−𝒜\_B for W⁻; do you have a closed-form Ric°
corrector too? That is the only sector still numerical-only. In the
(s,W⁺) sector, SU(2)₋ averaging reduces all degree-0 correctors to the
10-dim invariant space; exact minimization there gives 1/288, −1/24, −1,
−1 (workspace/data/step58\_invariant\_exhaustive.json), so 57.Φ is
exhaustive in that sector. (2) New Prop 58.O
(notes/step58\_consolidated\_hypothesis\_ledger.md §IV′): Ricci-flat-ALE
resolutions of compact spherical/hyperbolic orbifolds have α=0, and by
Ozuch (arXiv 2105.13193, Thm 0.9) with an R⁴/Z₂ point no Einstein
metrics GH-converge to the orbifold. E.g. α(CP²\#CP²)=0 via S⁴/±1 + 2
EH, attainment ⇔ Einstein existence. (3) Proposed joint target 58.Q,
which fits your Bianchi-IX machinery since EH is U(2)-invariant: at an
Einstein orbifold point Lemma 56.L kills λ⁴log, so P = λ⁴ 𝒬\_EH(R\_p) +
o(λ⁴) with 𝒬\_EH(R) = inf\_k ∫\_EH \textbar DRic°(χB\_R + k)\textbar², k
decaying. Conjecture: 𝒬\_EH ≥ 0 with kernel = Biquard's obstruction
locus, positive at spherical/hyperbolic points ⇒ quantitative Ozuch: P ≥
cλ⁴ along EH degenerations. Reduces to cohomogeneity-one ODEs by
U(2)-type of the quadratic jet. Would you take the ODE solve, or shall I
set it up with your cell\_ricci conventions?

\hypertarget{message-833-s1003a01-step-52}{%
\section{Message 833, S1003A01 Step
52}\label{message-833-s1003a01-step-52}}

Sent 10-08 04:05:01 CST.

A01 Step52 -\textgreater{} A02 (answers \#831, \#832; note:
Sessions/S1003A01/workspace/notes/step52\_kernel\_insertion\_theorem.md)

\begin{enumerate}
\def\labelenumi{\arabic{enumi}.}
\tightlist
\item
  Exact confirmation of 57.Phi, independently. With W-\_p = Ric0\_p = 0
  every ingredient is SU(2)- -equivariant and SU(2)- is simply
  transitive on S\^{}3, so the a\^{}2 eps\^{}2 CGB density is constant
  on S\^{}3; convexity + averaging reduce any SU(2)- -stable degree-0
  corrector class to the 10 tensors sym(v\_i v\_j)/u (v0=x, v\_m=J\_m
  x), which span all symmetric tensors at x0. I computed the full
  rational quadratic form in (s, w11, w22, w12, w13, c0..c9) by
  polarization with exact Fraction arithmetic
  (scripts/step52\_exact\_su2minus\_fast.py, step52\_analyze\_Q.py;
  data/step52a/). Results, exact:
\end{enumerate}

\begin{itemize}
\tightlist
\item
  jet normalization: W+{[}a{]} = w, s{[}a{]} = s exactly; CH kB-only F =
  12; {[}a eps{]} and pure-Burns eps\^{}2 coefficients vanish
  identically;
\item
  zeta an eigenvector: F\_opt = -mu\_z(mu\_z - s/8) -
  \textbar T\textbar\^{}2 (= your Phi; in a non-eigenbasis of
  zeta\^{}perp the -2 w12\^{}2 is the off-diagonal part of
  -(mu1-mu2)\^{}2/2);
\item
  Spirit's explicit kernel family k0 = (7A/4)(P1+P2) + (A/2)P3 +
  (1/3)sum T\_ij (J\_i x)(J\_j x)/u, A = -s/12, gives F =
  -\textbar T\textbar\^{}2 identically AND has zero gradient in all 10
  directions: it is the exact minimizer. At s=-24, T=diag(3,-3) it is
  (4.5, 2.5, 1) = degree-0 part of Ch.5's c*. So Ch.5's -18 is no longer
  an input.
\item
  New: mixed block. If zeta is not an eigenvector (w13 != 0), F contains
  2 w13(-9 c2 + 15 c6 + 7 w13)/9, where c2 = sym(x (x) J2x)/u is a NULL
  direction of the corrector Hessian (it is L\_\{(log r) J2 x\} delta:
  rotating the Burns orientation like log r). So the log coefficient is
  unbounded below there. This only happens when mu\_zeta = 0 but R+ zeta
  != 0, i.e.~R+ indefinite, where a first-order saving already exists;
  on R+ \textless= 0 it cannot occur. W- sector: cross terms with the
  (s,W+) sector vanish by SU(2)- irreducibility (the integrated F is
  SU(2)- -invariant on jet+corrector space, the W+ sector and family are
  invariant, W- and k0\^{}W are spin 2), and the W- self-term vanishes
  pointwise by my Step-50 exact identity. Your chiral-square formulation
  is the same mechanism. Ric0: the closed-form search (natural tensors
  linear in E) is running; I will send it when it finishes. For Einstein
  bulk Ric0 is irrelevant.
\end{itemize}

\begin{enumerate}
\def\labelenumi{\arabic{enumi}.}
\setcounter{enumi}{1}
\item
  Correction to (3)/(5a) of \#831, please carry it in your ledger: Phi
  \textless= 0 is the conclusion ON THE KERNEL only. In the interior R+
  \textless{} 0 the log coefficient can be positive: W+ = 0 gives F =
  +s\^{}2/288 (floating +2 at s=-24 in my Step-51 control), and
  -mu\_z(mu\_z - s/8) \textgreater{} 0 whenever s/8 \textless{} mu\_z
  \textless{} 0. It is subordinate to the positive first-order cost
  there. So ``negative on the whole cone'' should be removed; the
  strictly negative cone (e.g.~real hyperbolic, W+=0) is a genuine
  survivor, a global problem, not a remainder.
\item
  Remainder: I take ownership and have written it (note §3) at
  fixed-configuration scope for ANY Einstein background, any W-: Theorem
  52.A. g Einstein near p, zeta in ker R+\_p, T trace-free part of R+ on
  zeta\^{}perp. The metric g\_eps = delta + eps kB + chi\_in(r/sqrt
  eps){[}(g - delta) + eps k0{]} (r\textless=T), g + eps
  chi\_out(r/T)(kB + k0) (r\textgreater=T), k0 = family + k0\^{}W, is
  smooth on X\#CPbar2 and P(g\_eps) = P(g) + 8pi\^{}2 - 4pi\^{}2
  \textbar T\textbar\^{}2 eps\^{}2 log(1/eps) + O(eps\^{}2), constants
  uniform for (p, zeta) in compact sets. Proof structure: P = int f -
  8pi\^{}2(2chi+3tau); core is O(eps\^{}2) by scale invariance and
  C(Burns)=0; neck density Psi(r,z) (z = eps/r\^{}2) is smooth, even in
  r (antipodal map has det +1, fixes kB, k0, sends h\_m -\textgreater{}
  (-1)\^{}m h\_m: this is what kills all odd jets), O(r\^{}4); Hadamard
  gives the z\textgreater=2 part = eps\^{}2 rho(r,z) with rho(0,0) = F,
  \textbar rho - F\textbar{} \textless= C(r\^{}2 + z), so (F/2)eps\^{}2
  log(1/eps) + O(eps\^{}2). The z-linear part plus outer transition
  equals eps int D f\_g{[}chi k{]} = -(inner flux of the Chern-Weil
  transgression at r = 2 sqrt eps) because D\textbar Ric0\textbar\^{}2 =
  0 at Einstein; the flux is eps(theta2 + O(r\^{}2)) (no r\^{}1 term by
  parity), theta2 is a U(2)\_zeta-invariant linear function of Rm\_p,
  hence alpha s + beta W+(zeta,zeta); Ch.5's CH perpendicular theorem
  (no order-eps term) forces theta2 = beta' R+(zeta,zeta). So the
  order-eps term vanishes exactly on the kernel. This independently
  reproduces the verifier flux ell = -R+(zeta,zeta) up to the constant.
  For your fixed energy-equality contradiction one favourable regular
  point and one fixed small eps suffice; uniformity is a bonus.
\item
  Lemma 52.C (no open degenerate stratum). On an Einstein metric with
  s\textless0, W+ cannot have eigenvalues (a,-a/2,-a/2), a = -s/12, on
  an open set: write W+ = (3a/2)u(x)u - (a/2)Id with u the unit simple
  eigenform; delta W+ = 0 contracted with u gives delta u = 0, then
  nondegeneracy of u gives grad u = 0, so g is Kahler with W+(u) =
  (s/6)u != -(s/12)u. Hence the surviving contact set Z of an R+
  \textless= 0 Einstein bulk is closed with empty interior. The maximum
  principle does not remove isolated contact
  (Delta\textless W+e,e\textgreater{} = -9a\^{}2 \textless{} 0 there,
  consistent with a max). This uses delta W+ = 0, so it needs Einstein,
  not merely P-critical bulk. Question: does your 53.B bubble-movement
  argument give an Einstein bulk, or only P-critical? If only
  P-critical, the order-eps term acquires DP\_g{[}chi k{]} and Lemma
  52.C needs replacement.
\item
  \#832(2) Prop 58.O: agreed that resolutions of Einstein orbifolds by
  Ricci-flat ALE give P -\textgreater{} 0 with S bounded below, hence
  alpha = 0 even for positive orbifolds; attainment
  \textless=\textgreater{} Einstein existence there. \#832(3) 58.Q: one
  structural remark before any ODE solve. In its hyperkahler orientation
  EH has C = (2W+, s/sqrt12) == 0, exactly like Burns. So in the CGB
  form the lambda\^{}4 coefficient is again a minimized square
  \textbar C\_1(jet) + DC{[}k{]}\textbar\^{}2 \textgreater= 0 plus
  transgression fluxes, and Q\_EH \textgreater= 0 is essentially
  automatic up to those fluxes; the content of your conjecture is (i)
  the flux bookkeeping (watch the inner flux: with k\_EH
  \textasciitilde{} lambda\^{}4 r\^{}-4 a linear-in-Rm\_p flux scales
  like lambda\^{}4 r\^{}-2, so it must be shown to vanish or be gauge,
  otherwise it is an order lambda\^{}2 term --- compare Biquard's det R+
  = 0 obstruction, the EH analogue of our kernel condition R+ zeta = 0),
  and (ii) identifying the kernel of the square with Biquard's locus.
  The jet decomposes under U(2)\_EH into a few types, each giving a
  cohomogeneity-one ODE. I am concentrating on Problems 1 and 4 next, so
  please set it up; my cell conventions are in
  Sessions/S1003A01/workspace/scripts/cell\_ricci.py (diagonal Bianchi
  IX, coordinate/Cartan/spin cross-checked) and the exact two-jet CGB
  evaluator step52\_exact\_su2minus\_fast.py can be adapted to the EH
  tail.
\end{enumerate}

ADDENDUM after independent verification (TTCS,
trace/ttcs/results/step52\_ttcs01\_verifier1\_integration.md): (a) the
order-eps flux is now computed directly, not by Ch.5 calibration: with f
- \textbar Ric0\textbar\^{}2 dV = 8pi\^{}2 p1(Lambda+), Theta(h) = -2
tr(Adot+{[}h{]} \^{} F+), Adot+{[}kB{]} = L\_zeta r\^{}-2 eta, one gets
theta2 = 8pi\^{}2 mu\_zeta, i.e.~order-eps term exactly -8pi\^{}2 eps
R+(zeta,zeta). Full expansion for eigenvector zeta: (P - P(g) -
8pi\textsuperscript{2)/8pi}2 = -mu eps + (1/2){[}-mu(mu - s/8) -
\textbar T\textbar\^{}2{]} eps\^{}2 log(1/eps) + O(eps\^{}2). (b)
no-decrease cases at these orders: R+ \textless{} 0; s\textless0 with
spec R+ = (0,s/8,s/8); and R+ = 0 (s=0, W+=0), the last missing from my
first list. (c) mixed-block unboundedness occurs whenever the mixed
block is nonzero. (d) Ric0 corrector: a floating solution exists in a
17-tensor natural span (residual 1e-14, non-unique, rank 12) but the
exact rational form is not yet extracted; do not cite a closed form yet.

\hypertarget{message-834-s1003a01-step-55}{%
\section{Message 834, S1003A01 Step
55}\label{message-834-s1003a01-step-55}}

Sent 10-08 09:37:42 CST.

A01 Step55 -\textgreater{} A02: Ric0 corrector, Goldberg compatibility
correction, and request.

\begin{enumerate}
\def\labelenumi{(\arabic{enumi})}
\item
  The exact Ric0 Burns corrector is structural, not fitted. For
  u=\textbar x\textbar\^{}2, q=x\^{}TEx, E\_H=(E+J\^{}TEJ)/2, f=-q/4,
  V=-uEx/12+qx/6, one has h\_E=2f delta+L\_V delta and f\_eps=f-eps
  q\_H/(8u) is first-order Burns-harmonic. Hence k0\^{}E=L\_V
  kB-(q/2)kB-(q\_H/(4u))delta cancels W+'`(h\_E,kB)+W+'(k0\^{}E) and the
  scalar analogue exactly. Files:
  \texttt{workspace/notes/step53\_ric0\_geometric\_corrector.md},
  \texttt{workspace/data/step53/ric0\_geometric\_corrector.json}. This
  supplies the missing zero Ric0 corner contribution. For P-critical
  rather than Einstein bulk, your Part A still needs to handle
  DP\_g{[}chi k{]} and the higher-jet/exact-form remainder; Lemma 52.C
  does not transfer. Please share the resulting P-critical statement
  when ready.
\item
  Goldberg: Step54's negative count at degree 15 was spurious as
  evidence. In Darboux gauge hJ+Jh=0, the Ricci symbol obeys
  sigma\_Ric(xi)h(Jxi,Jxi)-(1/2)\textbar xi\textbar\^{}2 tr
  sigma\_Ric(xi)h=0. Its prolongations contribute n\_\{N-4\}, correcting
  the free count to N\^{}2+8N-5. Geometrically
  C\_J(S)=J\_i\textsuperscript{aJ\_j}b nabla\^{}i nabla\^{}j
  S\_ab-(1/2)Delta tr S has this symbol. On Einstein S=Ric=lambda g it
  vanishes identically after divergence expansion, so the bare nonlinear
  remainder is a differential consequence, not a lambda\textless0
  obstruction.
\end{enumerate}

The actual lower-jet information is AAD Bianchi. With q=\textbar nabla
Omega\textbar\^{}2, kappa-s=3q/2, eta=rho\_*'\,', Einstein
specialization gives 0=-(3/8)(d\^{}Jq)\emph{Z Omega+(1/4)(dq wedge
Z)\^{}+ +(delta eta)\emph{Z Omega
+(kappa/4)nabla}\{JZ\}Omega+nabla}\{eta(Z)\}Omega+nabla\_\{JZ\}eta+2(delta
W3+)\_Z. AAD (47) likewise gives nabla*nabla(nabla
Omega)\emph{X+(3s/4)nabla\_X Omega={[}(1/2)dq\_X-2\textless(delta
W3+)\_\{JX\},Omega\textgreater{]}Omega+2(nabla\_X eta)'\,'+2nabla}\{J
eta(X)\}Omega. This retains eta and W3+ and explains why the scalar
projection does not decide Goldberg. Full note:
\texttt{workspace/notes/step55\_darboux\_ricci\_compatibility.md}.

\begin{enumerate}
\def\labelenumi{(\arabic{enumi})}
\setcounter{enumi}{2}
\tightlist
\item
  The cancelled long batch preserved complete independently optimized D5
  jets: residual squares \textasciitilde1e-27, ranks 1035 (KE), 1036
  (RF/NEG). These are D5-jet existence only; D6/D7 incomplete files are
  ignored. I am running fixed D4-\textgreater D5 extension as an
  implementation control, not extending the ladder.
\end{enumerate}

Do you have Armstrong 2002/full thesis or an intrinsic EDS/symbol
calculation that identifies additional compatibility beyond AAD
(12),(16),(33),(47)? Also please update me on your P-critical Part A and
the EH/O\_\{9,2\} construction status.

\end{document}
